\documentclass[11pt,reqno]{amsart}

\usepackage[T1]{fontenc}
\usepackage[utf8]{inputenc}
\usepackage{lmodern}
\usepackage{microtype}
\usepackage{amsmath,amssymb,mathtools}
\usepackage{booktabs}
\usepackage{enumitem}
\usepackage[margin=1.05in]{geometry}
\usepackage{xcolor}
\usepackage[colorlinks=true,linkcolor=blue!55!black,citecolor=blue!55!black,urlcolor=blue!65!black]{hyperref}
\hypersetup{pdftitle={Kummer Division, Polynomial Secants, and Trace Rigidity},pdfauthor={}}
\ifdefined\pdftrailerid
  \pdftrailerid{}
\fi

\newtheorem{theorem}{Theorem}[section]
\newtheorem{proposition}[theorem]{Proposition}
\newtheorem{lemma}[theorem]{Lemma}
\newtheorem{corollary}[theorem]{Corollary}
\newtheorem{criterion}[theorem]{Criterion}
\theoremstyle{definition}
\newtheorem{definition}[theorem]{Definition}
\newtheorem{problem}[theorem]{Problem}
\newtheorem{conjecture}[theorem]{Conjecture}
\theoremstyle{remark}
\newtheorem{remark}[theorem]{Remark}
\newtheorem{warning}[theorem]{Warning}

\newcommand{\R}{\mathbb R}
\newcommand{\C}{\mathbb C}
\newcommand{\Q}{\mathbb Q}
\newcommand{\Z}{\mathbb Z}
\newcommand{\F}{\mathbb F}
\newcommand{\Gm}{\mathbb G_m}
\newcommand{\Ga}{\mathbb G_a}
\newcommand{\cO}{\mathcal O}
\newcommand{\cR}{\mathcal R}
\newcommand{\cT}{\mathcal T}
\newcommand{\cJ}{\mathcal J}
\newcommand{\Tr}{\operatorname{tr}}
\newcommand{\Trd}{\operatorname{Trd}}
\newcommand{\Nrd}{\operatorname{Nrd}}
\newcommand{\Gap}{\operatorname{Gap}}
\newcommand{\Comm}{\operatorname{Comm}}
\newcommand{\Span}{\operatorname{span}}
\newcommand{\Gal}{\operatorname{Gal}}
\newcommand{\SL}{\operatorname{SL}}
\newcommand{\GL}{\operatorname{GL}}
\newcommand{\PSL}{\operatorname{PSL}}
\newcommand{\PGL}{\operatorname{PGL}}
\newcommand{\SU}{\operatorname{SU}}
\newcommand{\Sp}{\operatorname{Sp}}
\newcommand{\SO}{\operatorname{SO}}
\newcommand{\PGU}{\operatorname{PGU}}
\newcommand{\Ad}{\operatorname{Ad}}
\newcommand{\End}{\operatorname{End}}
\newcommand{\Spec}{\operatorname{Spec}}
\newcommand{\diag}{\operatorname{diag}}
\newcommand{\rs}{\mathrm{rs}}
\newcommand{\ssc}{\mathrm{sc}}
\newcommand{\aad}{\mathrm{ad}}
\newcommand{\BC}{\mathrm{BC}}
\newcommand{\dash}{\mathbin{\dashrightarrow}}
\newcommand{\sslash}{\mathbin{\! /\! / \!}}

\title[Kummer division and trace rigidity]{Kummer Division, Polynomial Secants, and Trace Rigidity}
\author{}
\date{Working paper, 26 September 2026}

\begin{document}

\begin{abstract}
We develop an algebraic and arithmetic theory of division correspondences and
polynomial secants in semisimple groups.  On the regular semisimple locus,
division by \(m\) is a finite Kummer torsor under the \(m\)-torsion of the
universal centralizer.  In type \(A_1\), Hilbert 90 turns the projective
square-root correspondence into the rational formula \([1+g]\).  We prove
that this phenomenon is exceptional: a rational section of the generic
projective \(m\)-th-root correspondence exists only when \(m=2\) and the
group is a product of groups of type \(A_1\).  In higher type \(A\),
projective polynomial secants \([f(g)]\) take the place of roots.  Their
fields of definition detect the Galois action on the Dynkin diagram and
thereby distinguish inner from outer forms.

Combining secants with double-coset dynamics and arithmetic descent gives
uniform gaps for signed trace sumsets, a trace-gap classification for lattices
in \(\SL_d(\C)\), and a reduced-trace classification for lattices in
\(\SL_m(\mathbb H)\).  For compact Fuchsian trace sparsity we identify the
canonical role of the subgroup generated by squares and prove that, on word
balls of radius \(R\), every fiber of two trace coordinates has size
\(\exp(o(R))\) over every finitely generated characteristic-zero field.  We
then deduce the sharp universal bounds
\(D_H(X)\ge X^{1-o(1)}\) and
\(N_H^*(L)\ge e^{(1/2-o(1))L}\) for every closed hyperbolic surface.
The same \(X^{1-o(1)}\) lower bound holds in every component indexed by
every finite quotient.  Apart from the inverse identity
\(A+A^{-1}=(\Tr A)I\), the additive representation multiplicities inside
an exact trace fiber are \(\exp(o(R))\) on a word ball of radius \(R\), once
three fixed trace coordinates are used to recover \(A\).  Consequently, the
usual small-doubling arguments do not by themselves prove arithmeticity.

To prove arithmeticity, we fix an index \(N\) in the Chebyshev recurrence
and vary the group element to which the recurrence is applied.  A
reduced-divisor form of Roth's theorem gives a lower bound with coefficient
\(2N-4-o(1)\).  Expansion modulo \(p^2\) is enough to choose elements
\(b^mh_m\) for which every prime occurs to uniformly bounded order in the
product of two consecutive recurrence values.  A uniform fiber estimate for
two trace coordinates, followed by an explicit congruence counting
argument, gives an upper bound with coefficient \((2N-1)/2\).  Letting
\(N\) tend to infinity shows that the global recurrence height attached to
every element of the subgroup generated by squares is at most half its
original translation length.  At an algebraic character, the distinguished
real place already contributes exactly this half-length.  All other
nonnegative local terms must therefore vanish.  It follows that the traces
of the subgroup generated by squares are algebraic integers, their trace
field is totally real, and every nondistinguished conjugate lies in
\([-2,2]\).  Takeuchi's criterion \cite{Takeuchi} then gives arithmeticity.
If the original character were transcendental, a
nearby algebraic Fuchsian specialization in the same component would satisfy
\(\ell_y(\gamma)\le\ell_x(\gamma)\) for every \(\gamma\ne1\).
Thurston's asymmetric metric then forces the two characters to be equal,
contradicting the choice of an algebraic specialization of a transcendental
character.  Thus a cocompact Fuchsian group with linearly growing signed
trace set is arithmetic.  Hence the compact Schmutz linear-growth and compact
Sarnak bounded-clustering conjectures both hold.

The proof does not settle several stronger statements about the recurrence
of a single fixed element.  For every algebraic Fuchsian character satisfying
the linear trace-growth hypothesis and every hyperbolic \(a\), expansion
modulo rational squares gives the one-term estimate
\[
 \limsup_{n\to\infty}\frac1n
 \log N\operatorname{rad}_S(q_n(\Tr a))\le\frac{\ell(a)}2.
\]
The matching universal lower bound is implied by number-field \(abc\), but
at \(s=3/\sqrt2\) it is equivalent to full exponential growth of
\(\operatorname{rad}(2^n-1)\), and implies infinitely many base-\(2\)
non-Wieferich primes.
Thus the universal pointwise assertion remains open.  So does
superapproximation for arbitrary ideals in general twisted
restriction-of-scalars groups.  Neither statement is used in the proof of
the compact conjectures.  The paper also proves a Lucas--Wieferich
localization theorem and an exact entropy--pole-degree inequality for
algebraic curves on which prescribed trace equalities persist.
\end{abstract}

\maketitle
\tableofcontents

\section{Introduction and main results}

Power maps on reductive groups are naturally correspondences rather than
single-valued rational operations.  Their branches are controlled by regular
centralizer tori, their fields of definition by Kummer cohomology, and their
monodromy by Weyl and Galois actions.  Polynomial secants \([f(g)]\) are
rational maps obtained from these correspondences without choosing a root.
When paired with matrix traces, they express certain products in terms of
finitely many additive trace differences.  The same construction therefore
applies to rational root formulas, descent questions for algebraic group
forms, and arithmetic rigidity of sparse trace sets.

The first results concern algebraic groups and their representations.  We
then apply polynomial secants and commensurator rigidity to trace gaps.  In
the rank-one application, let
\(\Gamma<\PSL_2(\F)\), \(\F=\R\) or \(\C\), be a lattice, let
\(\widetilde\Gamma\) be its full inverse image in \(\SL_2(\F)\), and write
\[
 T(\Gamma)=\{\Tr\gamma:\gamma\in\widetilde\Gamma\}.
\]
Thus \(T(\Gamma)=-T(\Gamma)\); it is a set of values, so repeated traces are
counted only once.  For a subset \(S\) of \(\R\) or \(\C\), define
\[
 \Gap(S)=\inf\{|x-y|:x,y\in S,\ x\ne y\}.
\]
The set is \emph{uniformly discrete} when \(\Gap(S)>0\).  We say that
\(S\) has \emph{bounded clustering}, and write \(\BC(S)\), if
\[
 \sup_{z}\#\bigl(S\cap\{w:|w-z|\le1\}\bigr)<\infty.
\]
The radius \(1\) is immaterial: the same property results from any fixed
positive radius, because a ball of one fixed radius can be covered by
finitely many balls of another.  For \(S\subset\R\), \emph{linear growth}
means
\[
 \#(S\cap[-X,X])=O(X)\qquad(X\to\infty).
\]
Uniform discreteness implies bounded clustering, and bounded clustering
implies linear growth, but neither converse holds for arbitrary subsets of
\(\R\).

We use the following arithmetic conventions throughout.  Two subgroups are
\emph{commensurable} if their intersection has finite index in both, and
\[
 \Comm_G(\Lambda)=
 \{g\in G:g\Lambda g^{-1}\text{ is commensurable with }\Lambda\}.
\]
Here a \emph{lattice} in a locally compact group means a discrete subgroup
whose quotient has finite invariant volume.
For a linear algebraic \(k\)-group \(\mathbf G\), an arithmetic subgroup is
one commensurable with the integral points of a fixed \(\cO_k\)-model of
\(\mathbf G\); changing the faithful representation or integral model does
not change this commensurability class.

In rank one, an \emph{admissible quaternion algebra} \(B/k\) is one whose
chosen archimedean realization is \(M_2(\F)\) and whose other archimedean
factors are compact.  Thus, in the Fuchsian case, \(k\) is totally real and
\(B\) is ramified at every real place except the chosen one; in the Kleinian
case, \(k\) has one complex place up to conjugation and \(B\) is ramified at
every real place.  An order \(\cO\subset B\) is a subring containing \(1\)
which is a full \(\cO_k\)-lattice, and
\(\cO^1=\{b\in\cO:\Nrd(b)=1\}\).  A lattice is \emph{arithmetic} if,
after conjugation and projectivization, it is commensurable with the image of
some \(\cO^1\); it is \emph{derived from the order} if it is actually
contained in that image (and hence has finite index there).  Whenever a
normal subgroup inside the integral norm-one group is needed, we take the
normal core of its intersection with the arithmetic lattice.  This produces
a finite-index normal subgroup, denoted \(\Lambda_0\) or \(\Gamma_0\), that
is contained in the relevant integral norm-one group.

We also fix references for several standard results used repeatedly below.
Northcott's finiteness theorem is
\cite[Theorem~1.6.8]{BombieriGubler}.  Borel density is used in the form
\cite{BorelDensity}; the finite-covolume theorem for arithmetic groups is
\cite{BorelHarishChandra}; and Skolem--Noether is used in its usual
central-simple-algebra form \cite[Section~12.6]{Pierce}.  For a non-elementary
convex-cocompact Fuchsian group \(H\), orbital counting gives
\[
 \#\{h\in H:d(o,ho)\le R\}=\exp((\delta_H+o(1))R),
\]
and primitive closed geodesics have the corresponding prime-geodesic
asymptotic; see \cite{Roblin} (and \cite[Chapter~9]{Buser} in the compact
surface case).  These citations cover later invocations of ``orbit
counting'' and the ``prime geodesic theorem.''

With these conventions, the positive trace-gap classification
\cite{Bogachev}, which is used in Sections~5 and 9, is
\[
 \Gamma\text{ is derived from an admissible quaternion algebra},\qquad
 \Gap T(\Gamma)>0,
 \qquad
 \overline{T(\Gamma)-T(\Gamma)}\ne\F,
\tag{1.1}
\]
where the three conditions are equivalent.  The closure in the last
condition is essential.

The principal results are as follows.

\begin{enumerate}[label=\textup{(\roman*)}]
\item The projective square root \([1+g]\) is Hilbert--90 division by two on
the regular centralizer torus.  The obstruction to an actual rational root is
a Kummer class.
\item A rational generic projective \(m\)-th-root section exists for a
semisimple simply connected group only in the square map on products of
\(\SL_2\).
\item For type \(A_{d-1}\), \(d\ge3\), a polynomial secant with at least two
terms has target descent to the adjoint group precisely for an inner form.  In
an outer form its unitary norm is generically noncentral.
\item If \(\Gamma\) is derived, then its entire difference set, and more
generally every fixed integral signed trace sumset, is uniformly discrete with
an explicit norm bound.  Otherwise the first difference set is dense.
\item For \(d\ge3\), non-density of trace differences for a lattice in
\(\SL_d(\C)\) is equivalent to containment with finite index in the norm-one
group of an order in a degree-\(d\) central simple algebra over an imaginary
quadratic field.
\item For \(m\ge2\), the analogous reduced-trace classification holds for
lattices in \(\SL_m(\mathbb H)\); the defining field is \(\Q\).
\item The compact Sarnak and Schmutz conjectures are equivalent to a
trace-gap implication for the subgroup \(\Gamma^{(2)}\) generated by
squares.  For a fixed regular semisimple matrix \(B\), every fiber of the
coordinate pair \(A\mapsto(\Tr A,\Tr(AB))\) has size \(\exp(o(R))\) on a
word ball of radius \(R\), over every finitely generated
characteristic-zero field.
Consequently every closed hyperbolic surface satisfies
\(D_H(X)\ge X^{1-o(1)}\), or equivalently has at least
\(e^{(1/2-o(1))L}\) distinct primitive lengths up to \(L\); the exponent is
sharp.
\item The lower bound \(X^{1-o(1)}\) holds separately in every component
indexed by a finite quotient, and linear trace growth forces mixed energy
\(X^{3-o(1)}\) between all pairs of components in odd-order quotients.  On
the other hand, in an exact trace fiber the
only matrix sum with large multiplicity is the forced inverse identity
\(A+A^{-1}=(\Tr A)I\); every other sum has representation multiplicity
\(\exp(o(R))\) on a word ball of radius \(R\).
\item Invariant lattices at all prime divisors of a proper normal model
descend an absolutely irreducible representation to a central simple algebra
over its constant field; the vertical divisors produce an order.  Within an
admissible form, arithmeticity also follows from a non-elementary normal
subgroup in an order, one bounded spanning adjoint orbit, or nonexpanding
adjoint directions on growing word balls.  At any divisor where the adjoint
image is unbounded, every direction instead expands exponentially.
\item On a character locus preserving all trace collisions at a Fuchsian
point, two evaluations of \(A=\rho(a)\) give trace fibers of size
\(\exp(o(R))\) in a word ball of radius \(R\), as soon as their relative
matrix has distinct eigenlines.  Hence restriction to
each of Hao's complementary subgroups is birational on every
positive-dimensional irreducible locus through the point on which these
collisions persist, and has injective differential on every reduced
subvariety of such a locus.  At an algebraic point of the exceptional
trace-collision locus, the full adjoint trace field is already generated
by squared traces on each such subgroup.
Finite covers strengthen this to an exponential family of globally distinct
trace functions that coincide on the rational closure of the point.
\item At a divisorial valuation, powers of a loxodromic element
\(k\in\ker\phi\) separate the two spectral coefficients of each trace
equality that holds identically on the locus.  This gives
\[
 h_{\Tr}(x)\ge
 \frac{\delta_H}
 {1+\ell_x(k)\mathfrak c_v(a,k;H)/\ell_v(k)}.
\]
Every unbounded divisorial realization contains such an element \(k\).
Thus all divisorial degenerations with sufficiently large local translation
are excluded; the remaining case is the explicit opposite inequality
between local translation and cancellation.
\item Linear trace growth produces a set of
\(e^{(1/2-o(1))R}\) traces on which all Chebyshev recurrence dilates have
simultaneously controlled sumsets.  On a function field of a curve, the
pairwise coprime Chebyshev zero divisors satisfy a Chinese-remainder
uncertainty theorem, yielding the exact entropy--pole-degree inequality
(13.10bh10).
\item An explicit injective map converts uniform trace counts modulo
\((q_{n-1}(s)q_n(s))\) into a comparison between global recurrence height
and geometric translation length.  Fixing the recurrence index, Roth's
theorem gives a lower bound for almost the entire reduced divisor of two
consecutive terms.  A sieve over varying group elements, using expansion
only modulo \(p^2\), chooses \(a_m=b^mh_m\) for which this divisor has
bounded prime exponents.  Superapproximation for moduli whose prime
exponents are uniformly bounded, together with the uniform estimate
\[
 \#\{g:d_x(o,go)\le Cm,
       \Tr_y g=\xi,\ \Tr_y(a_mg)=\zeta\}=\exp(o_C(m)),
\]
then gives
\[
 \mathcal J_{k_y}(\Tr_yb)\le\frac{\ell_x(b)}2
\]
for every algebraic faithful cocompact Fuchsian character \(y\) in the
rational closure of \(x\).
\item At an algebraic Fuchsian point the distinguished real place gives
the reverse inequality, and Takeuchi's criterion \cite{Takeuchi} gives
arithmeticity.
If the original point were transcendental, its rational closure would
contain a nearby algebraic Fuchsian point in the same component.  The
preceding inequality implies that every marked length at the specialization
is at most the corresponding original marked length.  Thurston's asymmetric
metric \cite{ThurstonStretch} forces the two marked hyperbolic structures to
coincide, contradicting transcendence.  Thus linear trace growth forces arithmeticity
unconditionally, proving both compact Schmutz and compact Sarnak.
\item Independently, superapproximation for moduli \(m^E\), with \(m\)
square-free and \(E\) fixed, proves the stronger
one-term bound
\[
 \limsup_n\frac1n
 \log N\operatorname{rad}_S(q_n(s))\le\frac{\ell(a)}2.
\]
It also shows that any failure of the radical bound for a fixed recurrence
must occur in the initial Lucas--Wieferich multiplicities.
\item The stronger universal one-term Chebyshev radical hypothesis
\[
 \limsup_n n^{-1}\log N\operatorname{rad}_S(q_n(s))
 =\mathcal J_k(s)
\]
is implied by number-field \(abc\).  At \(s=3/\sqrt2\), however, it is
exactly the assertion that
\(\operatorname{rad}(2^n-1)\) has exponential rate \(2^n\), which implies
infinitely many base-\(2\) non-Wieferich primes.  This explains why the
unconditional proof varies the group element rather than proving the stronger
pointwise statement.
\end{enumerate}

The main unconditional conclusion is
Theorem~\ref{thm:unconditional-compact-trace-sparsity}.  The earlier
arbitrary-ideal and number-field \(abc\) implications are retained because
they prove stronger statements about a single recurrence and show why that
approach encounters familiar open problems.  Proposition
\ref{prop:crt-diagonal-obstruction} gives a finite-ring example in which
uniform estimates modulo separate prime powers do not combine into a
uniform estimate modulo their product.  The proof using varying group
elements does not require this invalid inference.  Section~16 records the
logical dependencies among these results, and Section~17 compares them with
the literature.

\section{The \texorpdfstring{$A_1$}{A1} square root is Hilbert 90}

\subsection{Quadratic centralizer algebra}

Recall that a finite-dimensional associative algebra \(D/k\) is
\emph{central simple} if its center is \(k\) and it has no nonzero proper
two-sided ideals.  Its degree is
\(\deg D=\sqrt{\dim_kD}\), and it is \emph{split} over an extension \(K/k\)
when \(D\otimes_kK\simeq M_{\deg D}(K)\).  The reduced characteristic
polynomial, trace \(\Trd\), and norm \(\Nrd\) are the ordinary characteristic
polynomial, trace, and determinant after such a splitting; Galois descent
makes them independent of the splitting.  A quaternion algebra is exactly a
central simple algebra of degree two.

An \emph{involution} on a central simple algebra is an anti-automorphism
\(*\) of order two: \((xy)^*=y^*x^*\) and \((x^*)^*=x\).
It is of the \emph{first kind} if it fixes the center pointwise.  It is of the
\emph{second kind}, or \emph{unitary}, if the center is a quadratic
extension of the fixed field and \(*\) induces the nontrivial automorphism of
that extension.  Thus an involution of the second kind is not linear over the
center; it is linear only over its fixed subfield.
If \(k\) is a number field, an \emph{order} in a central simple \(k\)-algebra
\(D\) is a subring containing \(1\) which is also a full
\(\cO_k\)-lattice in
\(D\), namely a finitely generated \(\cO_k\)-submodule spanning \(D\) over
\(k\).

Let \(k\) have characteristic different from two, let \(B/k\) be a quaternion
algebra, and let \(g\in\SL_1(B)\) be regular semisimple; equivalently,
\(\Trd(g)^2-4\ne0\).  Indeed, the reduced characteristic polynomial is
\[
 X^2-\Trd(g)X+1,
\]
because \(\Nrd(g)=1\).  Its discriminant is
\(\Trd(g)^2-4\), so the nonzero-discriminant hypothesis makes
\[
 E=k[g]
\]
a quadratic \'etale \(k\)-algebra.  Its nontrivial \(k\)-involution sends
\(g\) to the other root,
\(\bar g=\Trd(g)-g=g^{-1}\).  Regularity excludes
\(\Trd(g)=-2\), and
\[
 \Nrd(1+g)=(1+g)(1+g^{-1})=2+\Trd(g)\ne0.
\]
Thus \(u=1+g\) is invertible.  Then
\[
 \bar u=1+g^{-1}=g^{-1}(1+g),
 \qquad
 \frac{u}{\bar u}=g,
\tag{2.1}
\]
and
\[
 u\bar u=2+g+g^{-1}=2+\Trd(g).
\tag{2.2}
\]
Equation (2.1) is Hilbert 90 in its most elementary form.  After adjoining a
square root of (2.2), normalization gives the claimed root.  Explicitly,
\(u/\bar u=g\) implies \(u^2=(u\bar u)g\); division by a chosen square root
of the scalar \(u\bar u\) therefore gives
\[
 R(g)=\frac{1+g}{\sqrt{2+\Trd(g)}},
 \qquad R(g)^2=g.
\tag{2.3}
\]
The two choices of the scalar square root differ by a sign, so the normalized
object is a two-sheeted algebraic correspondence, or equivalently a map on
the quadratic Kummer cover on which that square root has been chosen.
Projectivization removes exactly that scalar ambiguity, making the following
class rational over \(k\):
\[
 R_{\aad}(g)=[1+g]\in\PGL_1(B),
 \qquad R_{\aad}(g)^2=[g].
\tag{2.4}
\]

\begin{proposition}[Kummer obstruction in type \(A_1\)]
For \(g\in\SL_1(B)(k)\) with \(2+\Trd(g)\ne0\),
\[
 g\text{ has a square root in }\SL_1(B)(k)
 \quad\Longleftrightarrow\quad
 2+\Trd(g)\in k^{\times2}.
\]
Thus the obstruction is the class of \(2+\Trd(g)\) in
\(H^1(k,\mu_2)=k^\times/k^{\times2}\).
\end{proposition}

\begin{proof}
If \(2+\Trd(g)\) is a square, formula (2.3) gives a \(k\)-rational
root.  Conversely, suppose \(h^2=g\).
Cayley--Hamilton for \(h\) gives
\[
 h^2-\Trd(h)h+1=0,
\]
so \(g+1=\Trd(h)h\).  Taking reduced norms yields
\[
 \Nrd(g+1)=2+\Trd(g)
 =\Nrd(\Trd(h)h)=\Trd(h)^2\Nrd(h)=\Trd(h)^2.
\]
\end{proof}

The reduced Cayley--Hamilton equation
\(g^2-\Trd(g)g+1=0\) also gives the secant identity
\[
 (g-1)^2=g^2-2g+1=(\Trd(g)-2)g.
\tag{2.5}
\]
Thus \([1+g]\), when \(2+\Trd(g)\ne0\), and \([g-1]\), when
\(\Trd(g)-2\ne0\), are projective square roots.  The
advantage of the second is its linear trace derivative.  In any split matrix
realization, for every fixed matrix \(\gamma\), trace linearity gives
\[
 \Tr(\gamma(g-1))=\Tr(\gamma g)-\Tr(\gamma).
\tag{2.6}
\]

\subsection{Geometric meaning}

For a loxodromic element of \(\PSL_2(\C)\), write its eigenvalues as
\(e^{\pm\lambda/2}\), where \(\lambda=\ell+i\theta\) is complex translation
length.  A branch of \(R(g)\) has eigenvalues \(e^{\pm\lambda/4}\).  It fixes
the same oriented axis because it is a rational function of \(g\), hence
commutes with \(g\) and has the same eigenspaces.  It translates by
\(\ell/2\) and rotates by \(\theta/2\); changing the square-root branch
multiplies the lift by \(-I\), which is invisible in \(\PSL_2\).  In
\(\PSL_2(\R)\) this is the half-translation along the same geodesic.
Equation (2.6) gives the arithmetic link: when \(\gamma\) and
\(\gamma g\) are lattice elements, \(\Tr(\gamma(g-1))\) is literally the
difference of two lattice trace values.  A density theorem for the relevant
double coset can therefore be read as a density theorem for additive trace
differences.

\section{Universal division torsors and rational sections}

\subsection{Construction from the universal centralizer}

Let \(G\) be connected reductive, let \(G^{\rs}\) be its regular semisimple
locus, and let
\[
 \cJ\longrightarrow G^{\rs},\qquad \cJ_g=Z_G(g),
\]
be the regular centralizer group scheme.  Equivalently, it is the pullback
to \(G^{\rs}\) of the universal centralizer on the regular conjugacy stack.
Concretely,
\[
 \cJ=\{(g,h)\in G^{\rs}\times G:gh=hg\},
\]
with group law in the second coordinate.  The tautological section is
\(\iota(g)=(g,g)\).
It is a torus of relative dimension \(\operatorname{rank}G\).  If the
characteristic does not divide \(m\), define
\[
 \cR_m=[m]^{-1}(\iota)\longrightarrow G^{\rs}.
\tag{3.1}
\]
Thus, more concretely,
\[
 \cR_m=\{(g,h)\in G^{\rs}\times G:h^m=g\};
\]
the commutation relation is automatic because every element commutes with
its own powers.

\begin{proposition}
The morphism \(\cR_m\to G^{\rs}\) is a finite \(\cJ[m]\)-torsor of degree
\(m^{\operatorname{rank}G}\).  At \(g\in G^{\rs}(k)\), the obstruction to a
\(k\)-rational root is the Kummer class
\[
 \delta_m(g)\in H^1(k,\cJ_g[m]).
\]
\'Etale-locally on \(G^{\rs}\), (3.1) is ordinary division by \(m\) on a
split torus.  When \(G\) is split, one may use the usual cameral cover; in
general one first makes a splitting base change.  Weyl and Galois monodromy
permute the branches.
\end{proposition}

\begin{proof}
Multiplication by \(m\) on a torus is finite, faithfully flat, and, under the
characteristic hypothesis, finite \`etale.  Its kernel is \(\cJ[m]\).  Pulling it
back along \(\iota\) proves the torsor statement.  After a finite
\'etale splitting base change, \(\cJ_g\simeq\Gm^r\), where
\(r=\operatorname{rank}G\), and its \(m\)-torsion is
\(\mu_m^r\).  This explains the degree \(m^r\).

At a \(k\)-point \(g\), the exact Kummer sequence
\[
 1\longrightarrow\cJ_g[m]\longrightarrow\cJ_g
   \mathop{\longrightarrow}^{[m]}\cJ_g\longrightarrow1
\]
has a connecting map
\(\cJ_g(k)\to H^1(k,\cJ_g[m])\).  The image of the element \(g\) is
\(\delta_m(g)\).  By exactness, this class vanishes precisely when
\(g=[m](h)=h^m\) for some \(h\in\cJ_g(k)\).  For split \(G\), the cameral
cover is the finite cover which chooses and labels the maximal torus
containing a regular element---the higher-rank analogue of choosing its
eigenvalues.  It removes the Weyl-group ambiguity and splits the preceding
torus locally.
\end{proof}

Thus roots do generalize geometrically, but normally as a finite
correspondence rather than a rational single-valued map.

The correspondence has an intrinsic torus presentation.  Let \(T\) be a
maximal torus after a splitting base change, let \(N=N_G(T)\), and put
\[
 \widetilde{\cR}_m(T)=
 \{(t,u)\in T^{\rs}\times T:u^m=t\}.
\]

\begin{proposition}[torus and quotient-stack model]\label{prop:root-stack}
After the chosen splitting base change there is a canonical equivalence
\[
 [\cR_m/G]\simeq[\widetilde{\cR}_m(T)/N].
\tag{3.2}
\]
In particular, the geometric branches form a torsor under \(T[m]\), Weyl
monodromy acts through \(N/T\), and over a nonsplit field this action is
combined with the Galois action on the based root datum.
\end{proposition}

\begin{proof}
The quotient stack notation remembers not only a conjugacy class but also
the conjugating elements and stabilizer; this is why the normalizer \(N\),
rather than only the abstract Weyl group \(N/T\), appears.
Every regular semisimple element is conjugate to an element of \(T^{\rs}\),
and two such elements are conjugate precisely through \(N\).  If \(u^m=t\)
and \(t\) is regular, then \(u\) is regular as well: for every root \(\alpha\),
\(\alpha(u)^m=\alpha(t)\ne1\).  Thus a commuting root pair is conjugate to a
pair in \(\widetilde{\cR}_m(T)\): once \(t\) is in \(T\), its centralizer is
exactly \(T\), so the commuting root \(u\) lies in the same torus.  The pair
has exactly the indicated normalizer
ambiguity.  This is the standard regular-semisimple equivalence
\([G^{\rs}/G]\simeq[T^{\rs}/N]\) with the division correspondence pulled
back along it.
\end{proof}

There is also a uniform geometric meaning for every real division branch.

\begin{proposition}[Jordan displacement under division]
Let \(G\) be a connected real semisimple linear group, and let \(g\in G\) be
regular semisimple.  If \(h^m=g\), then the Jordan projections satisfy
\[
 \lambda(h)=\frac1m\lambda(g).
\tag{3.3}
\]
For the associated symmetric space, the translation lengths therefore obey
\(\ell(h)=\ell(g)/m\).
\end{proposition}

\begin{proof}
Recall the definition.  Every element \(x\) of a real linear semisimple
group has a unique multiplicative Jordan decomposition
\[
 x=x_{\mathrm e}x_{\mathrm h}x_{\mathrm u}
\]
into commuting elliptic, hyperbolic, and unipotent factors.  The hyperbolic
factor is conjugate to a unique element \(\exp H\) with \(H\) in a fixed
closed positive Weyl chamber \(\mathfrak a^+\).  This element
\(H\) is the Jordan projection \(\lambda(x)\).

Apply uniqueness to \(h^m=g\).  If
\(h_{\mathrm u}=\exp N\), then the unipotent part of \(g\) is
\(h_{\mathrm u}^m=\exp(mN)\).  The element \(g\) is semisimple, so this part
is the identity; characteristic zero gives \(N=0\).  Thus \(h\) is
semisimple.  Conjugate its hyperbolic part to
\(\exp H\), \(H\in\mathfrak a^+\).  Its \(m\)-th power is
\(\exp(mH)\), and \(mH\) remains in \(\mathfrak a^+\).
The exponential map is injective on the positive split torus
\(A=\exp\mathfrak a\).  More explicitly, choose the simple-root characters
\(\alpha_1,\ldots,\alpha_r\).  The map
\[
 A\longrightarrow\R^r,\qquad
 a\longmapsto
   \bigl(\log|\alpha_1(a)|,\ldots,\log|\alpha_r(a)|\bigr)
\]
is injective on the positive split torus, and the chosen closed Weyl chamber
is described by linear inequalities in these logarithmic coordinates.  The
absolute values remove the signs which can occur on other components of a
real torus; on \(A=\exp\mathfrak a\) the character values are already
positive.  For every real character \(\chi\) one has
\[
 \log|\chi((\exp H)^m)|
   =m\log|\chi(\exp H)|.
\]
Thus taking the \(m\)-th power is literally multiplication by \(m\) in
logarithmic character coordinates.  The simple-root characters determine
\(H\), so the unique chamber
representative for the hyperbolic part of \(g\) is \(mH\).  Therefore
\(\lambda(g)=m\lambda(h)\), which is (3.3).

Finally, for the standard \(G\)-invariant metric on the symmetric space,
the minimum displacement of a semisimple element \(x\) is
\(\|\lambda(x)\|\), where the norm comes from the invariant inner product
on \(\mathfrak a\).  Homogeneity of that norm gives
\(\ell(g)=m\ell(h)\).
\end{proof}

\subsection{Classification of rational projective root sections}

The next result classifies when the finite division correspondence can
collapse to a rational projective branch.  Notice that no
conjugation-equivariance hypothesis is needed.

\begin{theorem}[classification of rational projective root sections]\label{thm:root-no-go}
Let \(G_{\ssc}\) be a nontrivial semisimple simply connected group over an
algebraically closed field of characteristic zero, and let
\(\pi:G_{\ssc}\to G_{\aad}\) be the adjoint quotient.  Suppose that, for some
\(m\ge2\), a rational map on a nonempty open set satisfies
\[
 s:G_{\ssc}\dashrightarrow G_{\aad},
 \qquad s(g)^m=\pi(g).
\tag{3.4}
\]
Then \(m=2\) and every simple factor of \(G_{\ssc}\) is of type \(A_1\).
Conversely, for a product of type-\(A_1\) factors the componentwise map
\(g\mapsto[1+g]\) satisfies (3.4) for \(m=2\).
\end{theorem}

\begin{proof}
Choose a regular semisimple point \(g\) in the domain and put
\(T_{\ssc}=Z_{G_{\ssc}}(g)\).  The domain of the rational map is open, so its
intersection with \(T_{\ssc}\) contains a nonempty open neighbourhood of a
suitable such \(g\).  For generic \(t\) in this intersection, equation (3.4)
implies that \(s(t)\) commutes with \(s(t)^m=\pi(t)\).  Since \(\pi(t)\) is
regular, its centralizer is \(T_{\aad}\).  Hence restriction gives a rational
map \(T_{\ssc}\dashrightarrow T_{\aad}\).

Write
\[
 P=X^*(T_{\ssc}),\qquad Q=X^*(T_{\aad});
\]
these are the weight and root lattices.  Pullback by the central isogeny
\(\pi\) identifies \(Q\) with the root sublattice of \(P\).  For every
\(\alpha\in Q\),
\[
 (s^*\alpha)^m=\pi^*\alpha=\chi^\alpha
\quad\text{in }k(T_{\ssc}).
\tag{3.5}
\]
In the fraction field of the Laurent UFD \(k[P]\), a monomial
\(\chi^\alpha\) is an \(m\)-th power only if \(\alpha\in mP\).  To see this,
write a putative root as a coprime quotient \(a/b\).  Valuation at every
irreducible divisor shows that all exponents in \(a^m\) and \(b^m\) must
cancel; hence both \(a\) and \(b\) are Laurent monomials times scalars.
Comparing their character exponents then gives \(\alpha=m\beta\) for some
\(\beta\in P\).  Therefore
\[
 Q\subset mP.
\tag{3.6}
\]

Write a simple root in the fundamental-weight basis as
\[
 \alpha_i=\sum_j a_{ji}\omega_j,
\]
where \((a_{ji})\) is the Cartan matrix.  In every irreducible root system
of rank at least two, some off-diagonal entry is \(-1\).  The corresponding
simple root therefore has coefficient \(-1\) in this basis and cannot be
divisible by any \(m\ge2\) in \(P\).  Apply this argument to each simple
factor.  In type \(A_1\), by contrast,
\(Q=2\Z\omega\subset P=\Z\omega\), and (3.6) forces \(m=2\).  Formula
(2.4) proves the converse.
\end{proof}

\begin{remark}
The theorem concerns a map \(G_{\ssc}\dashrightarrow G_{\aad}\).  It does not
give a rational square-root map on \(\PGL_2\) itself.  It is also distinct from
surjectivity of power maps, rational sections of the conjugacy quotient, and
Cayley birationality; see Steinberg~\cite{SteinbergPower},
Popov~\cite{PopovCrossSections}, and Lemire--Popov--Reichstein~\cite{LPRCayley}.
\end{remark}

\section{Projective secants and form detection}

Let \(\rho:G\to\GL(V)\) be a linear representation and let
\(0\ne f\in k[z,z^{-1}]\).  Laurent evaluation is legitimate because
\(\rho(g)\) is invertible.  On the open set where \(f(\rho(g))\) is
invertible,
define the projective secant
\[
 \Sigma_{f,\rho}(g)=[f(\rho(g))]\in\PGL(V).
\tag{4.1}
\]
This construction commutes with base change and conjugation of the
representation.  Its target need not lie in the projective image of \(G\).
That failure records the tensors which cut out \(\rho(G)\) inside
\(\GL(V)\).

For self-dual classical representations this target condition is explicit.
Let \(\dagger\) be the adjoint involution of a nondegenerate bilinear or
Hermitian form and assume \(\rho(g)^\dagger=\rho(g)^{-1}\).  The coefficients
of \(f\) lie in the fixed field, so
\[
 f(\rho(g))^\dagger=f(\rho(g)^{-1}).
\]
An invertible operator \(a\) is a similitude of the form exactly when
\(a^\dagger a\) is a scalar.  Put
\[
 q_f(z)=f(z^{-1})f(z).
\tag{4.2}
\]
Then \([f(\rho(g))]\) is a projective similitude precisely when
\[
 f(\rho(g))^\dagger f(\rho(g))=q_f(\rho(g))
\quad\text{is scalar}.
\tag{4.3}
\]

\begin{theorem}[classical secant obstruction]\label{thm:classical-secant}
Let \(k\) have characteristic zero and let \(f\ne0\).  Consider either the
standard representation, over an algebraically closed field, of
\[
 \Sp_{2r}\ (r\ge2),\qquad
 \SO_N\ (N\ge3),
\]
or the standard representation of an outer unitary \(k\)-form of type
\(A_{d-1}\), \(d\ge3\), tested after base change to a splitting field.
The generic projective secant \(\Sigma_{f,\rho}\) lies in the corresponding
projective similitude group---and in the unitary case has target descent to
the adjoint \(k\)-group---if and only if \(f\) is a Laurent monomial.
\end{theorem}

\begin{proof}
If the smallest and largest exponents occurring in \(f\) are \(a<b\), the
coefficient of \(z^{b-a}\) in \(q_f\) is the nonzero product of the two
extreme coefficients.  This term comes only from pairing exponent \(a\) in
\(f(z^{-1})\) with exponent \(b\) in \(f(z)\), so it cannot cancel.  Thus
\(q_f\) is constant exactly when \(f\) is a Laurent monomial
\(cz^a\), \(c\ne0\).

On a standard symplectic or even orthogonal torus one may write
\[
 \diag(z_1,\ldots,z_r,z_1^{-1},\ldots,z_r^{-1}).
\]
Equation (4.3) then has the eigenvalues
\(q_f(z_i)=q_f(z_i^{-1})\).  In rank at least two, scalarity forces the
nonconstant Laurent polynomial \(q_f\) to take the same value on two
independent variables, which is impossible.  An odd orthogonal torus has in
addition the eigenvalue \(1\); already \(q_f(z_1)=q_f(1)\) forces
\(q_f\) to be constant.  After splitting a unitary form, a diagonal torus
has eigenvalues \(z_1,\ldots,z_d\), subject only to
\(\prod z_i=1\).  When \(d\ge3\), \(z_1,z_2\) vary independently and the
remaining coordinates absorb the determinant relation, with the same
conclusion.  Conversely, for a
monomial \(f(z)=cz^a\), equation (4.3) is scalar.
\end{proof}

The excluded symplectic rank-one case is
\(\Sp_2=\SL_2\), whose projective similitude group is all of \(\PGL_2\).
This tensor-theoretic exception is the same \(A_1\) exception found in the
rational root-section theorem.

\subsection{Polynomial secants in type \texorpdfstring{\(A\)}{A}}

Let \(k\) have characteristic zero and let \(G\) be a simply connected
\(k\)-form of type \(A_{d-1}\).  For a
nonzero Laurent polynomial \(f\in k[z,z^{-1}]\), consider the projective polynomial
secant
\[
 s_f(g)=[f(g)]
\tag{4.4}
\]
on the open set where \(f(g)\) is invertible.

Here ``descent'' means \emph{target descent}: after scalar extension, the
projective class lands in the canonical adjoint \(k\)-subgroup.  It does not
mean that the map factors through the central quotient of its source; indeed
\([f(\zeta g)]\) need not equal \([f(g)]\).

Here a \(k\)-form of \(\SL_d\) means a \(k\)-group which becomes isomorphic
to \(\SL_d\) over an algebraic closure.  It is \emph{inner} when its Galois
descent maps use only conjugations, and \emph{outer} when the
transpose--inverse diagram automorphism occurs.  Equivalently, the image of
its descent cocycle in the order-two automorphism group of the
\(A_{d-1}\) Dynkin diagram is trivial in the inner case and nontrivial in the
outer case.

\subsection{Inner forms}

The simply connected inner forms of type \(A_{d-1}\) are precisely
\(\SL_1(D)\) for degree-\(d\) central simple \(k\)-algebras \(D\), and their
adjoint groups are \(\PGL_1(D)\).  Every positive or negative Laurent power
of \(g\) lies in \(D\).  Hence \(f(g)\in D\), and invertibility cuts out a
Zariski-open set on which (4.4) is automatically the \(k\)-rational map
\[
 s_f:G\dashrightarrow\PGL_1(D)=G_{\aad}.
\tag{4.5}
\]
This is the higher-rank remnant of the rationality of \([1+g]\).
In this inner case, ``the defining algebra'' is merely shorthand for the
central simple algebra \(D/k\) from which the group \(\SL_1(D)\) is
constructed; the word ``defining'' imposes no additional condition on
\(D\).

\subsection{Outer forms}

Suppose instead that the diagram action is nontrivial and
\[
 G=\SU(E,*).
\]
The \emph{defining unitary datum} \((L/k,E,*)\) consists of the quadratic
field \(L/k\), a degree-\(d\) central simple \(L\)-algebra \(E\), and an
\emph{involution of the second kind} \(*\).  Explicitly, \(*\) is a
\(k\)-linear anti-automorphism of \(E\) of order two and
\(\ell^*=\bar\ell\) for \(\ell\in L\), where the bar is the nontrivial
element of \(\Gal(L/k)\).  The special unitary group is the \(k\)-group
whose \(k\)-points are
\[
 \SU(E,*)(k)=
 \{g\in E^\times:g^*g=1,\ \Nrd_{E/L}(g)=1\}.
\]
For example, if \(E=\End_L(V)\) and \(*\) is adjoint with respect to a
nondegenerate Hermitian form \(h\) on \(V\), this is the usual special
unitary group \(\SU(h)\).
Thus \(G=\SU(E,*)\) is the simply connected outer form attached to this
datum.  The restriction of scalars
\(\operatorname{Res}_{L/k}\PGL_1(E)\) is the \(k\)-group whose \(k\)-points
are \(\PGL_1(E)(L)\).
In this manuscript, the occasional shorthand ``the defining algebra'' of an
outer form always means this \emph{full unitary datum} \((L/k,E,*)\),
not the central simple \(L\)-algebra \(E\) by itself.  Indeed, \(E\) alone
does not record either the descent field \(k\) or the semilinear involution
that defines the \(k\)-form.
Inside \(\operatorname{Res}_{L/k}\PGL_1(E)\), a projective class \([a]\)
lies in the projective unitary group precisely when \(a^*a\in L^\times\)
(then this scalar is automatically fixed by \(*\), because
\((a^*a)^*=a^*a\), and hence lies in \(k^\times\)).  On \(G\) one has
\(g^*=g^{-1}\), so
\[
 f(g)^*f(g)=f(g^{-1})f(g).
\tag{4.6}
\]

We also fix the meaning of a distinguished split factor.  If \(v\) is the
archimedean place used to realize the lattice and
\(L\otimes_k k_v\simeq k_v\times k_v\), the two factors correspond to the
two places of \(L\) above \(v\), and \(*\) exchanges them.  Choosing one
place \(w\mid v\) selects the factor \(E_w\).  When
\(E_w\simeq M_d(\F)\), projection to that factor identifies the local
special unitary group with \(\SL_d(\F)\); this chosen projection is the
\emph{distinguished realization}.  The arithmetic datum is assumed compact
at every other archimedean factor whenever the resulting image is called a
lattice in a single factor.

\begin{theorem}[detection of inner and outer forms by polynomial secants]\label{thm:secant-detector}
Let \(d\ge3\), and let \(f\in k[z,z^{-1}]\) have at least two nonzero
coefficients.  The generic
projective secant \([f(g)]\) has target descent to the adjoint \(k\)-group of
a type-\(A_{d-1}\) form if and only if that form is inner.
\end{theorem}

\begin{proof}
The inner assertion is (4.5).  In the outer case put
\[
 q_f(z)=f(z^{-1})f(z).
\]
If the smallest and largest exponents occurring in \(f\) are \(a<b\), then
the coefficient of \(z^{b-a}\) in \(q_f\) is the product of the two extreme
coefficients and is nonzero.  Hence \(q_f\) is nonconstant; equivalently,
\(q_f\) is constant only when \(f\) is a unit of \(k[z,z^{-1}]\), that is, a
monomial.

After base change to an algebraic closure, the involution exchanges the
standard module with its dual; on a split diagonal torus this sends
\(\diag(z_1,\ldots,z_d)\) to
\(\diag(z_1^{-1},\ldots,z_d^{-1})\).  Take a generic such element with
\(\prod z_i=1\).  Equation (4.6) is diagonal with
entries \(q_f(z_i)\).  For \(d\ge3\), two eigenvalues can be varied
independently subject to the determinant relation, so these entries are not
generically equal when \(q_f\) is nonconstant.  Thus (4.6) is generically
noncentral, and \([f(g)]\) does not descend to the projective unitary group.
\end{proof}

For the basic secant \(f(z)=z-\lambda\), \(\lambda\in k^\times\), the scalar
\(\lambda\) is fixed by \(*\), and direct multiplication gives
\[
 (g-\lambda)^*(g-\lambda)
 =(g^{-1}-\lambda)(g-\lambda)
 =(1+\lambda^2)I-\lambda(g+g^{-1}).
\tag{4.7}
\]
In particular, \([g-I]\) descends in an outer form exactly on the proper
locus where \(g+g^{-1}\) is scalar.  When \(d=2\), Cayley--Hamilton says
\(g^2-\Tr(g)g+I=0\).  Multiplying by \(g^{-1}\) gives
\(g+g^{-1}=\Tr(g)I\), so the obstruction disappears; indeed every symmetric
Laurent polynomial in \(g\) is scalar.  This is the same exceptional rank in
which (2.5) makes the secant a projective square root.

\subsection{The Galois action being detected}

For \(d\ge3\),
\[
 \operatorname{Aut}(\SL_d)=\PGL_d\rtimes\langle\iota\rangle,
 \qquad \iota(g)=(g^{-1})^{\mathsf t}.
\]
A \(k\)-form is inner when its descent cocycle has trivial image in the
diagram automorphism group.  More explicitly, composing the descent cocycle
with
\(\operatorname{Aut}(\SL_d)\to\langle\iota\rangle\simeq\Z/2\)
gives the diagram character.  A trivial character gives an inner form
\(\SL_1(D)\).  A nontrivial character cuts out a quadratic extension \(L/k\);
over \(L\), Galois conjugation identifies the standard representation with
its dual and produces the unitary datum \((L/k,E,*)\).  Formula (4.7)
detects precisely that exchange.

This should not be confused with the distinction between an arithmetic
lattice and a lattice \emph{derived} from an order.  Type \(A_1\) has no
nontrivial diagram automorphism, so all of its forms are inner; nevertheless
an arithmetic \(A_1\)-lattice can fail to be derived.  In rank one, the
square-root argument detects the latter Kummer and integrality obstruction.  Thus the
diagram character is an invariant of the algebraic \(k\)-form, whereas
being derived from an order is an integral and commensurability property of
a particular lattice.

We record the two dynamical facts used in the next corollary.  First, if an
invertible matrix commensurates the rational subgroup of an outer arithmetic
lattice, its distinguished realization is projectively a global unitary
similitude; this is the outer case of Lemma~\ref{lem:normalizer}.  Second, the
double-coset trace theorem~\cite{Bogachev} says that, for an invertible
matrix outside that commensurator, the values
\(\{\Tr(\gamma A):\gamma\in\Lambda_0\}\) are dense in the distinguished
field, component by component.

The finite-index rational subgroup used here always exists.  Indeed, an
outer arithmetic lattice is, by definition, commensurable with the integral
points of a unitary \(\cO_k\)-model.  Intersect the two groups and, if a
normal subgroup is wanted, take the intersection of its finitely many
conjugates in the lattice.  The resulting subgroup is still finite index and
consists of genuine unitary norm-one points.  This is the same normal-core
construction used in the introduction.

\begin{corollary}[trace density for outer forms]\label{cor:outer-witness}
Let \(\Lambda_0\) be the finite-index subgroup consisting of rational
unitary norm-one points of an outer arithmetic lattice.  Assume that at the
unique noncompact archimedean place \(v\), its unitary datum has a split
factor \(E_w\simeq M_d(\F)\), for a chosen \(w\mid v\), giving the
distinguished realization \(\Lambda_0<\SL_d(\F)\), \(d\ge3\).
Let \(f(z)=\sum_jc_jz^j\) have at least two nonzero coefficients.  For generic
\(g\in\Lambda_0\), \([f(g)]\) does not commensurate the lattice, and hence
\[
 \left\{\sum_jc_j\Tr(\gamma g^j):\gamma\in\Lambda_0\right\}
\tag{4.8}
\]
is dense in \(\F\).
\end{corollary}

Here ``generic'' means every lattice point in a fixed nonempty Zariski-open
set on which \(f(g)\) is invertible and
\(f(g^{-1})f(g)\) is noncentral.  Borel density
\cite{BorelDensity} guarantees lattice points in that open set.

\begin{proof}
If \([f(g)]\) commensurates, the unitary version of
Lemma~\ref{lem:normalizer} writes its distinguished realization as \(s b\),
where \(b\in E^\times\) and \(b^*b\in L^\times\).  Both \(f(g)\) and \(b\)
are global elements of \(E\); since their quotient realizes the scalar
\(sI\), faithfulness of the chosen realization makes that quotient central
in \(E\).  Hence \([f(g)]=[b]\) is a projective unitary similitude, contrary
to Theorem~\ref{thm:secant-detector} for generic \(g\).  The double-coset
trace theorem~\cite{Bogachev} applies to arbitrary invertible matrices in
\(\GL_d(\F)\), on every nonzero-determinant component, and therefore makes
\(\{\Tr(\gamma f(g))\}\) dense, and trace linearity gives (4.8).
\end{proof}

\begin{warning}
The element \(g\) in this fixed-element statement is initially chosen in the
rational unitary subgroup, where \(g^*=g^{-1}\) is justified.  For an arbitrary
element of a larger arithmetic overgroup, a determinant-one lift may only be a
unitary similitude.  Borel density in the finite-index rational subgroup is
enough for all form-exclusion arguments.
\end{warning}

\section{Exact additive consequences of the positive-gap theorem}

Suppose \(\Gamma\) is derived from an admissible quaternion algebra over a
number field \(k\) of degree \(n\).  Every trace is an algebraic integer; at
each nondistinguished archimedean embedding its absolute value is at most two.
The same product-formula argument used for \(T\) applies to much larger fixed
sumsets.

For an integral vector \(c=(c_1,\ldots,c_q)\), set
\[
 S_c(T)=\left\{\sum_{i=1}^q c_it_i:t_i\in T(\Gamma)\right\},
 \qquad C(c)=\sum_{i=1}^q|c_i|.
\]

\begin{theorem}[uniform gaps for fixed trace sumsets]\label{thm:sumset-gap}
If \(\Gamma\) is derived over \(k\), then \(S_c(T)\) is uniformly discrete.
More precisely,
\[
 \Gap S_c(T)\ge
 \begin{cases}
 (4C(c))^{1-n},&\F=\R,\\[2mm]
 (4C(c))^{1-n/2},&\F=\C.
 \end{cases}
\tag{5.1}
\]
For the zero vector the assertion is interpreted trivially.
\end{theorem}

\begin{proof}
Let \(x\ne y\) belong to \(S_c(T)\).  Their difference is a nonzero algebraic
integer
\[
 z=\sum_i c_i(t_i-u_i).
\]
At any nondistinguished archimedean embedding,
\(|\sigma z|\le4C(c)\).  In the Fuchsian case the norm inequality is
\[
 1\le |N_{k/\Q}(z)|
 \le |z|(4C(c))^{n-1}.
\]
In the Kleinian case the distinguished conjugate pair contributes \(|z|^2\)
and the other \(n-2\) embeddings contribute at most
\((4C(c))^{n-2}\).  Solving the two inequalities for \(|z|\) proves (5.1).
\end{proof}

Taking \(c=(1,-1)\) gives the promised sharpening.

\begin{corollary}[exact difference-set dichotomy]\label{cor:difference}
Put \(D=T(\Gamma)-T(\Gamma)\).  If \(\Gamma\) is not derived, then \(D\) is
dense in \(\F\).  If \(\Gamma\) is derived over a degree-\(n\) field, then
\[
 \Gap D\ge
 \begin{cases}
 8^{1-n},&\F=\R,\\[1mm]
 8^{1-n/2},&\F=\C.
 \end{cases}
\]
\end{corollary}

The first assertion is Bogachev's theorem; the second is Theorem
\ref{thm:sumset-gap}.  Thus the main theorem is not merely a dichotomy between
positive and zero gap of \(T\): it is a dichotomy between a uniformly discrete
first additive derivative and a dense one.

\section{A complex higher-rank trace-gap theorem}

We now give the complex analogue of the corresponding classification for
\(\SL_d(\R)\).  We use Margulis arithmeticity
\cite[Chapter~IX]{Margulis}, the standard classification of arithmetic
type-\(A\) forms \cite{PlatonovRapinchuk}, and the following
commensurator-normalizer fact.  The
Skolem--Noether step in that fact is the classical theorem for central
simple algebras \cite[Section~12.6]{Pierce}.

\begin{lemma}[projective normalizer]\label{lem:normalizer}
Let \(\F=\R\) or \(\C\), let \(K\) be a number field with a chosen
archimedean embedding into \(\F\), and let \(D/K\) be a degree-\(d\) central
simple algebra satisfying \(D\otimes_K\F\simeq M_d(\F)\).  Let
\(\Lambda_0<D^\times\cap\SL_d(\F)\) be Zariski dense.  If
\(A\in\GL_d(\F)\) commensurates \(\Lambda_0\), then
\[
 A=s b\qquad(s\in\F^\times,\ b\in D^\times).
\tag{6.1}
\]
For an outer unitary form \((E/L,*)\), where \(L/K\) is quadratic, let
\(\Lambda_0<\SU(E,*)(K)\) be Zariski dense and use a chosen embedding, or a
chosen factor when \(L\otimes_K\R\) is split, to identify the distinguished
component of \(E\) with \(M_d(\F)\).  If \(A\) commensurates its image, then
\[
 A=sb\qquad(s\in\F^\times,\ b\in E^\times,\ b^*b\in L^\times).
\]
\end{lemma}

\begin{proof}
Put \(H=\Lambda_0\cap A^{-1}\Lambda_0A\).  It is finite index and Zariski
dense, and \(AHA^{-1}\subset\Lambda_0\).  Choose \(d^2\) elements in \(H\)
which are \(\F\)-linearly
independent.  They form a \(K\)-basis of \(D\).  Conjugation by \(A\) therefore
preserves \(D\) and fixes its center.  By Skolem--Noether
\cite[Section~12.6]{Pierce} it agrees on \(D\)
with conjugation by some \(b\in D^\times\).  The matrix \(b^{-1}A\) centralizes
\(D\otimes_K\F=M_d(\F)\), and is scalar.  In the outer case Zariski density
implies that the \(L\)-span of \(H\) is \(E\); the same basis argument and
Skolem--Noether over \(L\) give \(A=sb\), \(b\in E^\times\).  Applying the
unitary equation to a Zariski-dense finite-index subgroup shows that
\(b^*b\) centralizes \(E\), hence lies in \(L^\times\).
\end{proof}

\begin{theorem}[complex higher-rank classification]\label{thm:complex}
Let \(d\ge3\), let \(\Lambda<\SL_d(\C)\) be a lattice, and put
\(T_\Lambda=\{\Tr g:g\in\Lambda\}\).  The following are equivalent.
\begin{enumerate}[label=\textup{(\arabic*)}]
\item After conjugation, \(\Lambda\) has finite index in \(\cO^1\), where
\(\cO\) is an order in a degree-\(d\) central simple algebra \(D/K\) and
\(K\) is imaginary quadratic.
\item For some imaginary quadratic field \(K\subset\C\),
\(T_\Lambda\subset\cO_K\).
\item \(\Gap T_\Lambda>0\) in the Euclidean metric on \(\C\).
\item \(\overline{T_\Lambda-T_\Lambda}\ne\C\).
\end{enumerate}
When these conditions hold, \(\Gap T_\Lambda\ge1\).
\end{theorem}

\begin{proof}
The implications (1)\(\Rightarrow\)(2)\(\Rightarrow\)(3)\(\Rightarrow\)(4)
are immediate.  Indeed, reduced traces of order elements lie in \(\cO_K\),
and a nonzero \(z\in\cO_K\) satisfies
\(|z|^2=|N_{K/\Q}z|\ge1\).

Assume (4).  If \(g\in\Lambda\) and \(\det(g-I)\ne0\), then
\[
 \{\Tr(\gamma(g-I)):\gamma\in\Lambda\}
 \subset T_\Lambda-T_\Lambda.
\]
Bogachev's double-coset trace theorem applies to every invertible
\(A\in\GL_d(\C)\): scalar normalization places \(A\) in \(\SL_d(\C)\), and
trace maps the corresponding nonzero-determinant hypersurface onto \(\C\).
It therefore forces \(g-I\) to commensurate \(\Lambda\).  Margulis
arithmeticity \cite[Chapter~IX]{Margulis} applies to \(\SL_d(\C)\), viewed as a
real group of real rank \(d-1\ge2\).

We first exclude an outer arithmetic form.  Pass to its finite-index rational
unitary subgroup and choose, by Borel density \cite{BorelDensity}, an element \(g\) satisfying
\[
 \det(g-I)\ne0,
 \qquad g+g^{-1}\text{ is not scalar}.
\]
The secant \(A=g-I\) commensurates.  Lemma~\ref{lem:normalizer} writes
\(A=sb\), with \(b\in E^\times\) and \(b^*b\in L^\times\).  Since also
\(A\in E\), injectivity of the distinguished realization shows that
\(Ab^{-1}=sI\) is central in \(E\), and hence \(s\in L\).  Consequently
\[
 (g-I)^*(g-I)=2I-g-g^{-1}
\]
would be central, a contradiction.  Hence the arithmetic form is inner,
say \(\SL_1(D)\) over a number field \(K\).

There is only one noncompact archimedean factor.  Every complex place of an
inner type-\(A_{d-1}\) group gives \(\SL_d(\C)\).  At a real place the local
group is either \(\SL_d(\R)\), or, when \(d\) is even,
\(\SL_{d/2}(\mathbb H)\); both are noncompact for \(d\ge3\).  Thus \(K\) has
no real places and exactly one complex place.  It is imaginary quadratic.

It remains to show that the entire lattice, not merely a finite-index subgroup,
lies in \(D^1\).  Every \(g\in\Lambda\) commensurates the initial arithmetic
subgroup, so Lemma~\ref{lem:normalizer} writes
\(g=s b\) with \(s\in\C^\times\), \(b\in D^\times\).  If, in addition,
\(g\) is nonscalar and \(\det(g-I)\ne0\), write similarly
\(g-I=t c\).  Then
\[
 c=\frac{s}{t}b-\frac1t I\in D.
\]
The elements \(I,b\) are \(K\)-linearly independent and remain
\(\C\)-linearly independent after scalar extension.  Comparing coordinates in
a \(K\)-basis of \(D\) gives \(s/t,1/t\in K\), hence \(s,t\in K\), and
\(g\in D\).  Its reduced norm is its determinant at the distinguished
embedding, so \(\Nrd(g)=\det(g)=1\) and \(g\in D^1\).

Let \(U\) be this nonempty good Zariski-open set.  For an arbitrary
\(h\in\Lambda\), the two open conditions \(x\in U\) and \(hx\in U\) have a
common point of \(\Lambda\) by Zariski density.  Both \(x\) and \(hx\) lie in
\(D^1\), and therefore \(h=(hx)x^{-1}\in D^1\).  This two-open-set argument is
needed: the countable subset \(D\subset M_d(\C)\) is not complex-Zariski
closed.

Choose an initial order \(\cO_0\subset D\) whose norm-one group is
commensurable with \(\Lambda\).  For every \(g\in\Lambda\), some positive
power \(g^r\) lies in \(\cO_0^1\).  Each eigenvalue \(\alpha\) of \(g\) then
has \(\alpha^r\) integral, hence is itself integral.  Since \(\Tr g\in K\),
we obtain \(\Tr g\in\cO_K\).

Finally choose \(d^2\) elements \(\gamma_i\in\Lambda\) forming a \(K\)-basis
of \(D\).  The reduced-trace Gram determinant
\[
 m=\det(\Trd(\gamma_i\gamma_j))
\]
is a nonzero element of \(\cO_K\).  Cramer's rule puts all coordinates of all
elements of \(\Lambda\) in the fixed fractional ideal \(m^{-1}\cO_K\).
Therefore \(\cO=\Span_{\cO_K}\Lambda\) is a finite \(\cO_K\)-module; it
contains \(1\), spans \(D\), and is closed under multiplication.  It is an order
with \(\Lambda<\cO^1\).  Since a discrete overgroup of a lattice has finite
index, \([\cO^1:\Lambda]<\infty\).  This proves (1).
\end{proof}

\begin{remark}
We have not found the precise trace-difference criterion in
Theorem~\ref{thm:complex} in the literature.  The proof combines standard
results, and we make no claim here about priority for this particular
formulation.
\end{remark}

\section{Secant accumulation and further local forms}

In absolute degree at least three, repeated secants can prove arithmeticity
without using the special root formula available in type \(A_1\).

\begin{lemma}[secant accumulation]\label{lem:secant-accumulation}
Let \(g\in\GL_N(\C)\), \(N\ge3\), be regular semisimple, with every eigenvalue
off the unit circle.  If a subgroup \(C<\PGL_N(\C)\) contains
\([g^n-I]\) for every \(n\ge1\), then \(C\) is nondiscrete.
\end{lemma}

\begin{proof}
All \(g^n-I\) are invertible.  Put
\[
 S_n=(g^n-I)(g-I)^{-1},\qquad
 Q_n=S_{n+1}S_n^{-1}=(g^{n+1}-I)(g^n-I)^{-1}.
\]
Their projective classes lie in \(C\).  On the \(g\)-eigenline with
eigenvalue \(z\), the operator \(Q_n\) acts by
\[
 q_n(z)=\frac{z^{n+1}-1}{z^n-1}\longrightarrow
 \begin{cases}
 z,&|z|>1,\\
 1,&|z|<1.
 \end{cases}
\tag{7.1}
\]
Thus \(Q_n\) converges to an invertible operator.  Its projective classes are
not eventually constant.  Indeed, if \(q_n(z)=c q_n(w)\) for every large
\(n\), cross-multiplication gives
\[
 (z-cw)(zw)^n+(c-z)z^n+(cw-1)w^n+(1-c)=0.
\]
Linear independence of exponential sequences with distinct nonzero bases
(and \(z,w\ne1\)) shows that the ratio can be constant only for \(z=w\) or
\(zw=1\); in the latter case
\(q_n(z^{-1})=z^{-1}q_n(z)\).  A regular spectrum of cardinality at least three
cannot be a single inversion orbit.  Hence \(C\) contains an infinite
convergent sequence, and is nondiscrete.
\end{proof}

\begin{remark}
For spectrum \(\{z,z^{-1}\}\), the classes \([Q_n]\) are constant.  Thus
this proof fails in absolute degree two, precisely the degree in which the
projective square root is rational.
\end{remark}

Suppose a lattice in an inner local type-\(A\) group has non-dense reduced
trace differences.  Whenever \(g^n-I\) is invertible, the identity
\[
 \Trd(\gamma(g^n-I))=\Trd(\gamma g^n)-\Trd(\gamma)
\]
and the double-coset theorem force \([g^n-I]\) into the commensurator.  A
regular element with the spectrum required in
Lemma~\ref{lem:secant-accumulation} then makes the commensurator nondiscrete.
In higher real rank this recovers Margulis arithmeticity.  For suitable
rank-one inner forms it gives a proof by a different argument.

The same argument gives a quaternionic theorem. In rank one, its proof uses
secant accumulation; in higher rank, the conclusion already follows from
Margulis arithmeticity.

We shall use the following standard consequence of Shah's orbit-closure
theorem \cite{Shah}.

\begin{lemma}[double-coset dichotomy]\label{lem:double-coset}
Let \(G\) be a connected noncompact almost simple real Lie group generated by
one-parameter unipotent subgroups, and let \(\Gamma<G\) be a lattice.  For every
\(a\in G\), either \(a\in\Comm_G(\Gamma)\), or \(\Gamma a\Gamma\) is dense in
\(G\).
\end{lemma}

\begin{proof}
Let the diagonal copy \(H=\Delta G\) act on the finite-volume homogeneous
space
\[
 (G\times G)/(\Gamma\times\Gamma)
\]
through the point \(x=(e,a)(\Gamma\times\Gamma)\).  Shah's theorem makes the
orbit closure homogeneous.  A Lie subalgebra of
\(\mathfrak g\oplus\mathfrak g\) containing \(\Delta\mathfrak g\) is either
\(\Delta\mathfrak g\) or all of \(\mathfrak g\oplus\mathfrak g\): subtracting
a diagonal element from any nondiagonal element produces a nonzero ideal in
one factor.  Connectedness of the orbit closure therefore gives the following
dichotomy.  Either \(Hx\) is closed of finite volume, in which case
\[
 \operatorname{Stab}_H(x)\simeq\Gamma\cap a\Gamma a^{-1}
\]
is a lattice in \(G\), so the two lattices are commensurable; or \(Hx\) is
dense.  In the second case, approximate an arbitrary
\((e\Gamma,y\Gamma)\) by \((h_i\Gamma,h_i a\Gamma)\).  Choosing
\(\gamma_i,\delta_i\in\Gamma\) with \(h_i\gamma_i\to e\) and
\(h_i a\delta_i\to y\) gives
\(\gamma_i^{-1}a\delta_i\to y\), proving density of the double coset.
\end{proof}

\begin{theorem}[quaternionic trace-gap classification]\label{thm:quaternionic}
Let \(m\ge2\), let
\[
 G_m=\SL_m(\mathbb H)=\SL_1(M_m(\mathbb H)),
\]
and use reduced trace.  For a lattice \(\Lambda<G_m\), the following are
equivalent:
\begin{enumerate}[label=\textup{(\arabic*)}]
\item There are a degree-\(2m\) central simple algebra \(D/\Q\), an order
\(\cO\subset D\), and an \(\R\)-algebra isomorphism
\(D\otimes_\Q\R\simeq M_m(\mathbb H)\) such that, after conjugation in
\(G_m\), \(\Lambda\) has finite index in the image of \(\cO^1\).
\item \(\Trd(\Lambda)\subset\Z\);
\item the reduced-trace set has positive gap;
\item its difference set is not dense in \(\R\).
\end{enumerate}
When these conditions hold, the gap is at least one.
\end{theorem}

\begin{proof}
The first three implications are immediate.  We prove (4)\(\Rightarrow\)(1).
First note the quaternionic version of the double-coset trace corollary.  If
\(A\in M_m(\mathbb H)^\times\), multiplication by a positive real scalar
normalizes its reduced norm to one without changing whether it commensurates
\(\Lambda\).  If it does not commensurate, Lemma~\ref{lem:double-coset}
makes \(\Lambda A\Lambda\), after that normalization, dense in \(G_m\).
Reduced trace is cyclic, and it maps \(G_m\) onto \(\R\): for
\(x>0\),
\[
 \diag(x,-x^{-1},1,\ldots,1)\in G_m,
 \qquad
 \Trd=2(x-x^{-1}+m-2).
\tag{7.2}
\]
Consequently \(\{\Trd(\gamma A):\gamma\in\Lambda\}\) is dense in \(\R\).

Assume now that the reduced-trace difference set is not dense.  For every
\(g\in\Lambda\) and every \(n\ge1\) for which \(g^n-I\) is invertible,
\[
 \Trd(\gamma(g^n-I))=\Trd(\gamma g^n)-\Trd(\gamma)
\]
shows that \(g^n-I\) commensurates \(\Lambda\).

For \(m\ge3\), Margulis arithmeticity \cite[Chapter~IX]{Margulis} applies because
\(\operatorname{rank}_\R G_m=m-1\ge2\).  For \(m=2\), pass to a torsion-free
finite-index subgroup.  The complex regular-semisimple (simple-spectrum) locus
in the standard four-dimensional complex realization is a nonempty real
Zariski-open set, so Borel density \cite{BorelDensity} gives such a \(g\) in that subgroup.
A nontrivial semisimple element of a torsion-free discrete rank-one group
cannot be elliptic, so \(g\) is axial.  Its four standard complex eigenvalues
have moduli \(r,r,r^{-1},r^{-1}\) for some \(r\ne1\), and in particular all
lie off the unit circle.  Thus every \([g^n-I]\) lies in the projective
commensurator, and
Lemma~\ref{lem:secant-accumulation} makes that commensurator nondiscrete.
Margulis's commensurator criterion \cite[Chapter~IX]{Margulis} again gives
arithmeticity.

  The standard classification of arithmetic type-\(A_{2m-1}\) forms
  \cite{PlatonovRapinchuk} now gives
either an inner central simple algebra or an outer unitary form.  We first rule
out the latter.  (A quadratic extension can split at the distinguished real
place, so local quaternionic type alone does not do this.)  Pass to the
finite-index rational unitary subgroup.  Borel density
\cite{BorelDensity} gives an element
\(g\) for which \(g-I\) is invertible and \(g+g^{-1}\) is not central.  The
secant \(g-I\) commensurates.  The outer version of the normalizer argument in
Lemma~\ref{lem:normalizer}, with the distinguished simple factor now
\(M_m(\mathbb H)\), writes
\[
 g-I=s b,
 \]
where \(b\) belongs to the central simple \(L\)-algebra \(E\) occurring in
the unitary datum and
\(b^*b\in L^\times\).  Because \(g-I\) and \(b\) are global elements, the
identity \((g-I)b^{-1}=sI\) in the faithful distinguished factor forces
\(s\in L^\times\).  It follows that
\[
 (g-I)^*(g-I)=2I-g-g^{-1}
\]
is central, a contradiction.  The arithmetic form is therefore inner, given
by a central simple algebra \(D/K\).

Every complex completion of
\(D\) gives a noncompact \(\SL_{2m}(\C)\), while at a real completion its
norm-one group is \(\SL_{2m}(\R)\) or \(\SL_m(\mathbb H)\), both noncompact
because \(m\ge2\).  The arithmetic construction has only the distinguished
noncompact factor.  Hence \(K\) has exactly one archimedean place, which is
real, and therefore \(K=\Q\).

It remains to pass from commensurability to actual containment in an order.
We use the local normalizer argument from Lemma~\ref{lem:normalizer}: if
\(A\in M_m(\mathbb H)^\times\) commensurates a Zariski-dense subgroup of
\(D^\times\), then
\[
 A=s b\qquad(s\in\R^\times,\ b\in D^\times).
\tag{7.3}
\]
Indeed, a basis chosen in a finite-index intersection shows that conjugation
by \(A\) preserves \(D\); Skolem--Noether
\cite[Section~12.6]{Pierce} gives \(b\), and the centralizer
of \(D\otimes_\Q\R=M_m(\mathbb H)\) in that algebra is \(\R\).

For a nonscalar \(g\in\Lambda\) with \(g-I\) invertible, apply (7.3) to both
\(g=s b\) and \(g-I=t c\).  From
\[
 c=(s/t)b-(1/t)I\in D
\]
and linear independence of \(I,b\), coordinate comparison gives
\(s/t,1/t\in\Q\).  Thus \(g\in D\), and its reduced norm is one.  The
two-open-set argument used in Theorem~\ref{thm:complex} propagates this from
the good Zariski-open set to every element of \(\Lambda\), so
\(\Lambda<D^1\).

Choose an order \(\cO_0\subset D\) whose norm-one group is commensurable with
\(\Lambda\).  Every \(g\in\Lambda\) has a positive power in \(\cO_0^1\), so
its reduced eigenvalues, and hence \(\Trd(g)\in\Q\), are integral.  Therefore
\(\Trd(\Lambda)\subset\Z\).  Finally choose \((2m)^2\) elements of \(\Lambda\)
forming a \(\Q\)-basis of \(D\).  The nondegenerate reduced-trace pairing and
Cramer's rule put every element of \(\Lambda\) in one fixed fractional
\(\Z\)-lattice.  Hence \(\cO=\Span_\Z\Lambda\) is an order, contains
\(\Lambda\), and is closed under multiplication.  Since \(\cO^1\) is a
discrete overgroup of the lattice \(\Lambda\), the index is finite.  This is
(1), and integer traces give the stated gap.
\end{proof}

\subsection{A quantitative root-commensurator criterion}

The root idea does give an unconditional arithmeticity criterion in every
rank-one group.

\begin{theorem}\label{thm:root-commensurator}
Let \(G\) be a connected center-free simple rank-one isometry group,
\(\Gamma<G\) a lattice, \(h\) its volume entropy, and \(m\ge2\).  Let
\(P_m(L)\) count primitive axial \(\Gamma\)-conjugacy classes \([\gamma]\) of
length at most \(L\) for which some \(r\in\Comm_G(\Gamma)\) satisfies
\(r^m=\gamma\).  If \(\Gamma\) is nonarithmetic, then
\[
 P_m(L)\ll e^{hL/m}.
\tag{7.4}
\]
Consequently, an exponential rate strictly greater than \(h/m\), and in
particular the existence of such a root for every primitive axial element,
forces arithmeticity.
\end{theorem}

\begin{proof}
For nonarithmetic \(\Gamma\), its commensurator \(\Delta\) is discrete and
contains \(\Gamma\) with finite index.  If \(r^m=\gamma\), then
\(\ell(r)=\ell(\gamma)/m\).  A \(\Delta\)-conjugacy class meets \(\Gamma\) in
at most \([\Delta:\Gamma]\) \(\Gamma\)-classes.  The prime-geodesic upper
bound for \(\Delta\) counts \(O(e^{hR})\) axial classes of length at most
\(R\).  Taking \(R=L/m\) gives (7.4).
\end{proof}

Counting axial classes is important in the nonuniform case: parabolic classes
have zero translation length and cannot be included in this estimate.

\section{Decomposition of arithmetic traces by square classes}

An arithmetic lattice need not be derived, and its full trace set need not
have positive gap.  What survives is a finite decomposition indexed by
square classes, compatible with the Galois action.

\begin{proposition}[finite decomposition of an arithmetic trace set]\label{prop:arith-colors}
Let \(\Gamma<\PSL_2(\R)\) be arithmetic.  There is a finite-index normal
subgroup \(H\lhd\Gamma\), derived from an order, such that, for every coset
\(aH\), the signed
trace set
\[
 T(aH)=\{\Tr\widetilde{ah}:h\in H\}
\]
is uniformly discrete.  Consequently \(T(\Gamma)\) is a finite union of
uniformly discrete sets and has bounded clustering.  Differences between
two components may nevertheless be dense.
\end{proposition}

\begin{proof}
Take a normal finite-index subgroup \(H\) contained in the norm-one group of
an order in the defining quaternion algebra \(B/k\).  A lift of a coset
representative commensurates that order group.  Skolem--Noether
\cite[Section~12.6]{Pierce} writes it as
\(s_ab_a\), with \(b_a\in B^\times\), while the determinant condition gives
\(s_a^2\in k^\times\).  Hence
\[
 \Tr(s_ab_ah)=s_a\Trd(b_ah).
\tag{8.1}
\]
The second factor lies in a fixed fractional \(\cO_k\)-ideal.  Put
\(L_a=k(s_a)\), of degree at most two over \(k\).  After one fixed
denominator is cleared, all differences in (8.1) are algebraic integers of
\(L_a\).  At every extension of a nondistinguished archimedean embedding of
\(k\), they are bounded: the corresponding norm-one group is compact, and
the factors \(s_a,b_a\) are fixed.  Over the distinguished real place there
are either one or two embeddings of \(L_a\); in the quadratic case they send
\(s_a\) to \(\pm s_a\), so both absolute values equal the one at the
distinguished place.
The norm inequality in \(L_a\) therefore gives a positive separation
constant for (8.1).  There are only finitely many cosets.
\end{proof}

For the Fricke extension in Bogachev's manuscript,
\[
 T(\Gamma)=2\Z\cup\sqrt2\,\Z.
\tag{8.2}
\]
Each of the two displayed subsets is separated, while Pell approximants
using one value from each subset give zero total gap and
\(2\Z+\sqrt2\Z\subset T-T\), a dense subgroup of \(\R\).

\subsection{Ordinary versus invariant trace field}

Let
\[
 K=\Q(\Tr\widetilde\Gamma),\qquad
k=\Q(\Tr\widetilde\Gamma^{(2)})
   =\Q((\Tr\widetilde\gamma)^2:\gamma\in\Gamma).
\]
This is the ordinary/invariant trace-field convention of
\cite[Chapter~3]{MaclachlanReid}; in particular, the invariant trace field is
unchanged on passage to a finite-index non-elementary subgroup.
Because \(\Gamma\) is finitely generated, the standard finite trace-generation
theorem for \(\SL_2\) characters expresses all word traces as integral
polynomials in finitely many basic word traces~\cite{Horowitz}.  Together with
\[
 (\Tr\widetilde\gamma)^2=\Tr(\widetilde\gamma^2)+2\in k,
\tag{8.3}
\]
this shows that \(K/k\) is a finite multiquadratic extension.  There is a
more precise compatibility with the quotient by the subgroup generated by
squares.

\begin{proposition}[Galois signs are mod-two characters]\label{prop:galois-signs}
Assume \(\Gamma\) is nonelementary.  For every
\(\sigma\in\Gal(K/k)\), there is a character
\[
 \varepsilon_\sigma:\Gamma\longrightarrow\{\pm1\}
\]
such that
\[
 \sigma(\Tr\widetilde\gamma)
 =\varepsilon_\sigma(\gamma)\Tr\widetilde\gamma.
\tag{8.4}
\]
Consequently \(\Gal(K/k)\) embeds in
\(H^1(\Gamma,\F_2)\).  For a fixed class in
\(\Gamma/\Gamma^{(2)}\), all its traces lie in a one-dimensional
\(k\)-eigenspace \(s_vk\subset K\).
\end{proposition}

\begin{proof}
Work first on the full inverse image \(\widetilde\Gamma\), with its standard
representation \(\rho\), viewed over \(\C\).  Extend \(\sigma\) to a field
automorphism of \(\C\) and apply it entrywise (equivalently, realize the
irreducible character over a finite algebraic extension of its trace field and
extend \(\sigma\) there).  After semisimplification this gives a representation
\(\rho^\sigma\) whose character is \(\sigma(\Tr\rho)\).  The adjoint
characters agree, since
\[
 \Tr\Ad(\widetilde\gamma)=(\Tr\widetilde\gamma)^2-1\in k.
\]
Because \(\Gamma\) is nonelementary, its projective image is Zariski dense;
the adjoint representation is irreducible.  Brauer--Nesbitt therefore
identifies the two adjoint representations.  Since every automorphism of
\(\PGL_2\) is inner, after conjugating \(\rho^\sigma\) their projectivizations
agree.  Hence
\[
 \rho^\sigma(\delta)=\widetilde\varepsilon_\sigma(\delta)\rho(\delta)
 \qquad(\delta\in\widetilde\Gamma),
\]
where multiplicativity and determinant one make
\(\widetilde\varepsilon_\sigma:\widetilde\Gamma\to\{\pm1\}\) a character.
It takes the value \(+1\) on the central element \(-I\), because its trace
\(-2\) is fixed by \(\sigma\), and therefore descends to \(\Gamma\).  This
gives (8.4), and it is trivial on squares.

The twisting character is unique.  Indeed, if \(\rho\simeq\rho\otimes\chi\)
for a nontrivial quadratic character \(\chi\), then \(\Tr\rho\) vanishes on
the nontrivial coset of \(\ker\chi\).  That coset is Zariski dense because
\(\rho(\widetilde\Gamma)\) is Zariski dense and \(\ker\chi\) has finite index,
which is impossible for the nonzero trace function.  Applying \(\sigma\) to
the twist for \(\tau\), and then using the twist for \(\sigma\), shows that
\(\varepsilon_\sigma\varepsilon_\tau\) is the twisting character for
\(\sigma\tau\).  Uniqueness therefore makes
\(\sigma\mapsto\varepsilon_\sigma\) a homomorphism, including on trace-zero
elements.  Its kernel fixes the
trace-generated field \(K\), so it is injective.  A multiquadratic extension
decomposes into one-dimensional eigenspaces over \(k\), proving the last
assertion.
\end{proof}

The two Galois actions occur at different levels.  The diagram action
distinguishes inner and outer forms of higher type \(A\), whereas the
characters in (8.4) distinguish the components indexed by square classes
inside one rank-one arithmetic commensurability class.

\section{Reduction to the subgroup generated by squares}

Let
\[
 H=\Gamma^{(2)}
\]
be the subgroup generated by the elements \(\gamma^2\),
\(\gamma\in\Gamma\).  Every automorphism carries a square to a square, so
\(H\) is characteristic.  Every element of \(\Gamma/H\) has order at most
two; a group of exponent two is abelian because
\((xy)^2=1=x^2y^2\) implies \(xy=yx\).  Thus \(\Gamma/H\) is an elementary
abelian two-group.  It is finite because \(\Gamma\) is finitely generated,
and it is the largest quotient of \(\Gamma\) with these properties.  For a
closed surface, it is \(H_1(\Gamma,\F_2)\), so \(H\) is the fundamental
group of the mod-two homology cover.

\begin{theorem}[arithmeticity and the square subgroup]\label{thm:kummer-reduction}
For a Fuchsian lattice \(\Gamma\), the following are equivalent:
\begin{enumerate}[label=\textup{(\arabic*)}]
\item \(\Gamma\) is arithmetic;
\item \(\Gamma^{(2)}\) is derived from a quaternion algebra;
\item \(\Gap T(\Gamma^{(2)})>0\);
\item \(T(\Gamma^{(2)})-T(\Gamma^{(2)})\) is not dense in \(\R\).
\end{enumerate}
Consequently, the compact bounded-clustering conjecture is equivalent to
\[
 \BC(T(\Gamma))\Longrightarrow \Gap T(\Gamma^{(2)})>0,
\tag{9.1}
\]
and the compact linear-growth conjecture is equivalent to
\[
 \#(T(\Gamma)\cap[-X,X])=O(X)
 \Longrightarrow \Gap T(\Gamma^{(2)})>0.
\tag{9.2}
\]
\end{theorem}

\begin{proof}
The classical invariant-trace-field characterization says that \(\Gamma\) is
arithmetic exactly when \(\Gamma^{(2)}\) is derived from the invariant
quaternion algebra; see \cite{Takeuchi} and
\cite[Chapter~8]{MaclachlanReid}.  Bogachev's positive trace-gap theorem applied to the
finite-index lattice \(H\) makes (2), (3), and (4) equivalent.  Bounded
clustering and linear growth do pass from \(T(\Gamma)\) to its subset
\(T(H)\), but this alone does not give a gap.  The exact reformulations follow
instead by replacing, pointwise, the conclusion ``\(\Gamma\) is arithmetic''
in the Sarnak or Schmutz implication by its already proved equivalent
``\(\Gap T(H)>0\).''  Thus each implication in (9.1) and (9.2) entails the
original conjecture, and conversely the corresponding original conjecture,
together with (1)\(\Leftrightarrow\)(3), entails that implication.
\end{proof}

Thus (9.1) is not a weaker substitute for Sarnak's conjecture.  It is an
equivalent formulation in which the finite ambiguity coming from square
classes has been removed.  The reformulation is useful because
\(\Gamma^{(2)}\) is characteristic and \(\Gamma/\Gamma^{(2)}\) is the
maximal elementary abelian two-quotient of \(\Gamma\).

\begin{corollary}[finite-cover formulation]
A Fuchsian lattice is arithmetic if and only if it has a finite-index subgroup
\(H\) whose signed trace set is uniformly discrete.  If it is arithmetic, one
may moreover choose \(H\) normal so that every coset trace set is uniformly
discrete.
\end{corollary}

\begin{proof}
The forward statement is Proposition~\ref{prop:arith-colors}.  Conversely,
positive trace gap makes \(H\) derived by Bogachev's theorem, and arithmeticity
passes to the finite-index overgroup \(\Gamma\).
\end{proof}

\section{The trace algebra graded by square classes}

Put \(V=\Gamma/\Gamma^{(2)}\), written additively.  For \(v\in V\), let
\(T_v\) be the signed traces of lifts of elements in that coset.  The Fricke
identity gives a coherent multiplication law which an arbitrary coloring does
not possess.

\begin{theorem}[trace identities respecting square classes]\label{thm:graded-trace}
For \(a,b\in V\),
\[
 T_aT_b\subset T_{a+b}+T_{a+b}.
\tag{10.1}
\]
For \(t\in T_a\),
\[
 t^2-2\in T_0.
\tag{10.2}
\]
If \(R_\Gamma=\Z[T(\Gamma)]\), then
\[
 R_\Gamma=\Z+\Span_\Z T(\Gamma).
\tag{10.3}
\]
More precisely, with \(M_a=\Span_\Z T_a\),
\[
 M_aM_b\subset M_{a+b}.
\tag{10.4}
\]
The sum of the \(M_a\) need not be direct.
\end{theorem}

\begin{proof}
Choose lifts \(A,B\) with traces \(x\in T_a\), \(y\in T_b\).  Then
\[
 xy=\Tr(AB)+\Tr(AB^{-1}).
\tag{10.5}
\]
Because \(V\) has exponent two, both \(AB\) and \(AB^{-1}\) have class
\(a+b\), proving (10.1).  Cayley--Hamilton gives
\(t^2-2=\Tr(A^2)\), which proves (10.2).  Finally (10.5) reduces every product
of trace generators to an integral sum of traces, inductively proving (10.3)
and (10.4).
\end{proof}

\begin{corollary}[the trace algebra over its neutral component]\label{cor:finite-kummer-algebra}
Put \(A_0=\Z+M_0\) and
\(A=\Z[T(\Gamma)]\).  Then \(A_0\) is a ring and \(A\) is a finite integral
\(A_0\)-algebra generated by elements satisfying quadratic equations over
\(A_0\).
\end{corollary}

\begin{proof}
Equations (10.2) and (10.4) make \(A_0\) a ring.  The trace algebra of a
finitely generated \(\SL_2\)-group is generated by finitely many word traces,
say \(t_1,\ldots,t_r\).  Each satisfies
\[
 X^2-(\Tr(\gamma_i^2)+2)=0
\]
over \(A_0\).  Hence
\(A=A_0[t_1,\ldots,t_r]\) is integral and finite over \(A_0\).
\end{proof}

These identities show why an approximate-ring theorem is relevant.  They also
show why it cannot yet be applied: an approximate subring requires both
\(X+X\) and \(XX\) to be covered by finitely many additive translates of
\(X\).  Formula (10.1) only puts products in a two-fold sumset.

\subsection{An abstract graded-ring counterexample}

\begin{proposition}[a scalar model that does not come from group traces]\label{prop:scalar-countermodel}
Let
\[
 S=2\Z\cup2\sqrt2\,\Z,
 \qquad S_0=S,\quad S_1=2S.
\tag{10.6}
\]
Thus \(S_0\cup S_1=S\) and \(\pm2\in S_0\).  Then:
\begin{enumerate}[label=\textup{(\roman*)}]
\item \(S\) has bounded clustering and linear growth, but \(\Gap S=0\) by
Pell approximation;
\item \(S_iS_j\subset S_{i+j}\) for \(i,j\in\Z/2\), and
\(s^2-2\in S_0\) for every \(s\in S\);
\item every Chebyshev trace polynomial \(P_n\), defined by
\(P_0=2,P_1=x,P_{n+1}=xP_n-P_{n-1}\), respects the grading:
\(P_n(S_i)\subset S_{ni}\);
\item \(\Z[S]=\Z+\Span_\Z(S)=\Z[2\sqrt2]\).
\end{enumerate}
The component declared neutral is \(S_0=S\), and hence has zero gap.  Moreover,
\(S+S\) contains the dense group \(2\Z+2\sqrt2\Z\), while a finite union of
translates of \(S\) is locally finite.  Thus \(S\) is not an approximate
subring.
\end{proposition}

\begin{proof}
Each unit interval meets each of the two lattices in a uniformly bounded
number of points, and their counts in \([-X,X]\) are linear.  Continued
fractions for \(\sqrt2\) give \(|2p-2\sqrt2q|\to0\), proving zero gap.  Put
\(B=\Z\cup\sqrt2\Z\), so \(S_0=2B\), \(S_1=4B\), and \(BB\subset B\).
It follows immediately that
\(S_0S_0\subset S_0\), \(S_0S_1\subset S_1\), and
\(S_1S_1\subset S_0\).  Also \(s^2-2\in2\Z\subset S_0\).

The polynomial \(P_n\) has the parity of \(n\).  If its argument lies in \(2B\), its
even powers and constant term give an element of \(2\Z\), while its odd powers
remain in the same one of the two lattices making up \(2B\).  Thus
\(P_n(S_0)\subset S_0\).  If its argument lies in \(4B\), the same observation gives an
element of \(4B\) for odd \(n\), and of \(2\Z\) for even \(n\).  This proves
the graded Chebyshev assertion.  The ring identity is direct.  Finally,
\(2\Z+2\sqrt2\Z\) is dense in \(\R\),
whereas every finite union of translates of \(S\) is locally finite.  Thus
no finite collection of additive translates of \(S\) covers \(S+S\).
\end{proof}

This model satisfies every purely scalar operation in
Theorem~\ref{thm:graded-trace}, including the displayed two-component
multiplication law, but its neutral component has zero gap and its whole set is not an
approximate subring.  Therefore closure under the Fricke identity,
Chebyshev iteration, finite
generation of the trace ring, and a grading by square classes cannot prove
(9.1) without information from noncommutative multiplication or hyperbolic
geometry.  The arithmetic example (8.2)
shows that the whole-set obstruction is not an artificial pathology.

There is, however, a sharp distinction between assigning degrees to a scalar
model and realizing those degrees by traces of group elements.  The full
Fricke relation forces purity in any direct grading.

\begin{theorem}[Fricke purity for a direct grading]
\label{thm:fricke-purity}
Let \(Q\) be an elementary abelian two-group, and let
\[
 R=\bigoplus_{q\in Q}R_q\subset\R
\]
be a direct sum of additive subgroups such that
\(R_pR_q\subset R_{p+q}\) and \(2\in R_0\).  Let
\(\Delta<\SL_2(\R)\) be a group for which every trace is nonzero and
homogeneous:
\[
 \Tr(\Delta)\subset\bigcup_{q\in Q}R_q.
\]
For \(g\in\Delta\), let \(\varepsilon(g)\) be the unique degree of
\(\Tr g\).  Then
\[
 \varepsilon:\Delta\longrightarrow Q
\]
is a homomorphism.  In particular,
\[
 \Tr(\Delta^{(2)})\subset R_0.
\tag{10.7}
\]
If \(\Delta\) is the full lift of a torsion-free cocompact Fuchsian lattice
and \(R_0\) is uniformly discrete, then the lattice is arithmetic.
\end{theorem}

\begin{proof}
The identity has trace \(2\in R_0\), and
\(\Tr(g^{-1})=\Tr g\), so inversion preserves degree.  Put
\(p=\varepsilon(g)\) and \(q=\varepsilon(h)\).  The Fricke identity says
\[
 \Tr g\,\Tr h=\Tr(gh)+\Tr(gh^{-1}).
\tag{10.8}
\]
Its left side is a nonzero element of \(R_{p+q}\).  Each term on the right
is a nonzero homogeneous element.  Because the sum of the \(R_r\) is direct,
neither right-hand term can have degree different from \(p+q\): a component
in any other degree would have to vanish, and if both terms had the same
wrong degree their sum could not equal a nonzero element of \(R_{p+q}\).
Consequently
\[
 \varepsilon(gh)=\varepsilon(gh^{-1})=p+q.
\]
Thus \(\varepsilon\) is a homomorphism.  Since \(Q\) has exponent two, it
kills every square and hence the subgroup generated by squares, proving
(10.7).

For a torsion-free cocompact Fuchsian group there are no trace-zero elements.
The central element \(-I\) has degree zero, so \(\varepsilon\) descends from
the full lift.  If \(R_0\) is uniformly discrete, (10.7) gives a positive
trace gap on the square subgroup.  Theorem~\ref{thm:kummer-reduction} then
gives arithmeticity.
\end{proof}

For example, fix any integer \(q\ge1\) and put
\[
 R_0=2\Z,\qquad R_1=2q\sqrt2\,\Z.
\]
These groups form a direct \(\Z/2\)-grading: their products have the required
degrees.  Hence, if all traces of a torsion-free cocompact Fuchsian lift lay
in \(R_0\cup R_1\), Theorem~\ref{thm:fricke-purity} would force every trace
on the square subgroup into \(2\Z\), and the group would be arithmetic.  The
scalar example (10.6) obtains its zero neutral gap only by declaring the whole
union to have degree zero; that declaration cannot agree with the canonical
quotient by square classes.

\section{Why the trace identities do not yet give an approximate ring}

A set \(X\) in a ring is \emph{symmetric} if \(0\in X\) and \(X=-X\).
It is an \emph{approximate subring} if there is a finite set \(F\) such that
\[
 X+X\ \cup\ XX\subset F+X.
\]
Two subsets of an additive group are \emph{additively commensurable} if each
is contained in finitely many additive translates of the other.
Krupi\'nski and Machado prove the following 2026 structure theorem
\cite{KrupinskiMachado}.

\begin{theorem}[Krupi\'nski--Machado, specialized form]\label{thm:KM}
Let \(X\) be an infinite uniformly discrete approximate subring spanning a
finite-dimensional simple real algebra \(A\).  Then there are a number field
\(K\), an archimedean place \(v\), a \(K_v\)-algebra structure on \(A\), and a
\(K_v\)-basis \(e_1,\ldots,e_r\) such that \(X\) is additively commensurable
with
\[
 \sum_{i=1}^r\cO_{K,v}e_i,
 \qquad
 \cO_{K,v}=\{a\in K:|a|_w\le1\text{ for every }w\ne v\}.
\tag{11.1}
\]
For \(A=\R\), these are the classical Pisot--Salem model sets.
\end{theorem}

The model set on the right has the three features that occur in Takeuchi's
criterion \cite{Takeuchi}: a number field, integrality at the finite places,
and boundedness at the nondistinguished archimedean places.  Additive
commensurability, however, does not imply that the original set or the
finitely many translating elements lie in that number field.  One must also
relate the model set to the trace field and to the components indexed by
square classes.  In addition, the approximate-ring hypothesis is strictly
stronger than (10.1).

\begin{lemma}[covering close differences already gives gap]\label{lem:additive-cover-gap}
Let \(S\subset\R\) be locally finite.  If
\[
 (S-S)\cap(-\eta,\eta)\subset F+S
\]
for a finite set \(F\) and some \(\eta>0\), then \(\Gap S>0\).
\end{lemma}

\begin{proof}
The finite union \(F+S\) is locally finite, so zero is isolated in it.  Every
nonzero difference smaller than \(\eta\) belongs to this union; all other
differences are already at least \(\eta\).  Hence the nonzero differences are
bounded away from zero.
\end{proof}

Thus a finite-translate cover of the sufficiently small differences in
\(T(\Gamma^{(2)})\) would already prove the required trace gap, without an
application of Theorem~\ref{thm:KM}.  By contrast, finding a
small-doubling subset at selected scales leaves the other trace values
uncontrolled, and those values may still contain arbitrarily close pairs.

\subsection{Finite-place boundedness and arithmeticity}

We shall use the following elementary consequence of the product formula.

\begin{lemma}[separation by the product formula]\label{lem:galois-window}
Let \(k\) be a number field with a distinguished real embedding \(v_0\), and
let \(I\) be a fractional \(\cO_k\)-ideal, that is, a nonzero finitely
generated \(\cO_k\)-submodule of \(k\) whose \(k\)-linear span is \(k\).
Equivalently, some nonzero \(a\in\cO_k\) satisfies
\(aI\subset\cO_k\).
Suppose \(S\subset I\) and,
for every other archimedean embedding \(\sigma\), there is a constant
\(C_\sigma\) such that
\[
 |\sigma(s)|\le C_\sigma\qquad(s\in S).
\tag{11.2}
\]
Then \(v_0(S)\subset\R\) is uniformly discrete.
\end{lemma}

\begin{proof}
Choose \(a\in\cO_k\setminus\{0\}\) with \(aI\subset\cO_k\).  For distinct
\(s,t\in S\), put \(z=a(s-t)\).  Then \(z\) is a nonzero algebraic integer,
so \(N_{k/\Q}(z)\) is a nonzero rational integer.  If complex embeddings are
counted with their usual multiplicities, then
\[
 1\le |N_{k/\Q}(z)|
   =|v_0(z)|\prod_{\sigma\ne v_0}|\sigma(z)|.
\]
For \(\sigma\ne v_0\), the triangle inequality and (11.2) give
\[
 |\sigma(z)|\le 2|\sigma(a)|C_\sigma.
\]
There are only finitely many archimedean embeddings, so the product of these
upper bounds is a fixed constant \(M\).  It follows that
\[
 |v_0(s-t)|=|v_0(z)|/|v_0(a)|
 \ge \bigl(|v_0(a)|M\bigr)^{-1}.
\]
This lower bound is independent of \(s\) and \(t\), which is precisely
uniform discreteness.
\end{proof}

The hypothesis of Lemma~\ref{lem:galois-window} controls an entire scalar
set simultaneously.  For arithmeticity, a weaker matrix-valued condition is
enough: one adjoint orbit which spans the Lie algebra may be controlled
instead.

\begin{theorem}[arithmeticity inherited from a normal subgroup]
\label{thm:normal-skeleton}
Let \(\Gamma<\SL_2(\R)\) be a lattice and let
\(\Delta\triangleleft\Gamma\) be non-elementary.  Suppose that, under a
fixed splitting at its distinguished real place, \(\Delta\) is contained in
the norm-one group of an order in an admissible quaternion algebra \(B/k\).
Then \(\Gamma\) is arithmetic.

It is enough, in particular, to find one noncentral \(g\in\Gamma\) whose
entire \(\Gamma\)-conjugacy orbit is contained in one such order.
\end{theorem}

\begin{proof}
The non-elementary group \(\Delta\) is absolutely irreducible.  Hence, by
Burnside's theorem,
\[
 k[\Delta]\otimes_{k,v_0}\R=M_2(\R),
\]
and therefore \(k[\Delta]=B\).  The algebra
\[
 \cO_\Delta=\cO_k[\Delta]
\]
is an \(\cO_k\)-submodule of the given order.  Since \(\cO_k\) is
Noetherian, it is a finite \(\cO_k\)-module; because it spans \(B\), it is
itself an order.  Normality says that conjugation by
every \(\gamma\in\Gamma\) preserves \(\Delta\), and therefore preserves both
\(B=k[\Delta]\) and \(\cO_\Delta\).  By Skolem--Noether
\cite[Section~12.6]{Pierce}, this conjugation on
\(B\) is induced by an element \(b\in B^\times\).  After passage to the
distinguished split factor, \(b^{-1}\gamma\) centralizes
\(B\otimes_{k,v_0}\R=M_2(\R)\), hence is scalar, and
\(b\cO_\Delta b^{-1}=\cO_\Delta\).  Thus the projective image of \(\Gamma\)
lies in the orientation-preserving projective normalizer of
\(\cO_\Delta\).

For an admissible quaternion algebra, the projective order normalizer is an
arithmetic lattice and is commensurable with the norm-one group of the order
\cite[Chapter~8]{MaclachlanReid}.
It contains the projective image of \(\Gamma\).  Since an inclusion between
two lattices has finite index, \(\Gamma\) is arithmetic.

For the last assertion, take \(\Delta\) to be the normal closure of \(g\).
Every conjugate of \(g\) has reduced norm one and lies in the given order.
Its inverse also lies in the order (for a norm-one quaternion,
\(u^{-1}=\Trd(u)-u\)), and the order is closed under multiplication.
Consequently the entire normal closure is contained in the norm-one group of
that order.  It remains only to check that this normal closure is
non-elementary.
Every elementary normal subgroup of a non-elementary Fuchsian group has
trivial projective image, equivalently it is contained in the center of
\(\SL_2(\R)\).  Since \(g\) is noncentral, \(\Delta\) is non-elementary.
\end{proof}

The preceding criterion assumes that the number-field algebra is already
known.  The next two results construct such an algebra, and then an order in
it, from boundedness at divisorial valuations.  They apply over arbitrary
finitely generated coefficient fields and treat rational descent and
integrality in a single argument.

We first record the codimension-one patching fact used in both results.

\begin{lemma}[patching lattices on a normal scheme]
\label{lem:reflexive-lattice-patching}
Let \(X\) be a normal noetherian integral scheme with function field \(F\),
and let \(V\) be a finite-dimensional \(F\)-vector space.  For every prime
divisor \(D\in X^{(1)}\), prescribe an
\(\mathcal O_{X,D}\)-lattice \(L_D\subset V\).  Suppose that, on each affine
open subset of \(X\), these lattices agree with the localizations of one
fixed free lattice except at finitely many prime divisors.  Then
\[
 \mathcal E(U)=\bigcap_{D\in U^{(1)}}L_D\subset V
\]
defines a coherent reflexive sheaf whose stalk at every prime divisor \(D\)
is exactly \(L_D\).
\end{lemma}

\begin{proof}
The assertion is local, so write \(X=\operatorname{Spec}A\), identify the
fixed lattice with \(A^d\), and let \(S\) be the finite set of exceptional
height-one primes.  Choose \(0\ne f\in A\) so that, for every
\(\mathfrak p\in S\),
\[
 fA_{\mathfrak p}^d\subset L_{\mathfrak p}
       \subset f^{-1}A_{\mathfrak p}^d.
\]
For nonexceptional height-one primes put
\(L_{\mathfrak p}=A_{\mathfrak p}^d\), and set
\[
 M=\bigcap_{\operatorname{ht}\mathfrak p=1}L_{\mathfrak p}\subset F^d.
\]
Then \(fA^d\subset M\subset f^{-1}A^d\).  Hence \(M\) is a finite
\(A\)-module.  We verify that localization recovers the prescribed lattice.
If \(x\in L_{\mathfrak p}\), only finitely many height-one valuations prevent
\(x\) from lying in the other prescribed lattices.  By prime avoidance there
is an element \(s\in A\setminus\mathfrak p\) with
sufficiently large valuation at each of them.  Then \(sx\in M\), so
\(M_{\mathfrak p}=L_{\mathfrak p}\).

A finite torsion-free module over a normal domain is reflexive precisely
when it is the intersection, inside its generic fibre, of its localizations
at the height-one primes.  The displayed definition of \(M\), together with
the equality of its height-one localizations just proved, therefore gives
\(M=M^{**}\).  These modules agree on overlaps because all were defined by
the same valuation conditions, and hence glue to the asserted coherent
reflexive sheaf.
\end{proof}

\begin{theorem}[divisorial descent to the constant field]
\label{thm:constant-field-descent}
Let \(K\) be a field of characteristic zero, let \(F/K\) be a finitely
generated extension in which \(K\) is relatively algebraically closed, and
let \(X/K\) be a normal projective integral variety with function field
\(F\).  Let \(V\) be a \(d\)-dimensional \(F\)-vector space, let
\(\Gamma\) be finitely generated, and let
\(\rho\colon\Gamma\to\SL_F(V)\) be absolutely irreducible.  Suppose that,
for every prime divisor \(D\in X^{(1)}\), there is an
\(\cO_{X,D}\)-lattice \(L_D\subset V\) satisfying
\[
 \rho(\gamma)L_D=L_D\qquad(\gamma\in\Gamma).
\]
Then there are a central simple \(K\)-algebra \(B\) of degree \(d\) and an
\(F\)-algebra isomorphism
\[
 B\otimes_KF\simeq\End_F(V)
\]
under which
\[
 \rho(\Gamma)\subset
 B^1:=\{b\in B^\times:\operatorname{Nrd}_{B/K}(b)=1\}.
\]
\end{theorem}

Here \(K\) is the \emph{constant field} in \(F\): relative algebraic
closedness means that every element of \(F\) algebraic over \(K\) already
lies in \(K\).  A lattice over the discrete valuation ring \(\cO_{X,D}\)
is a free \(\cO_{X,D}\)-submodule of rank \(d\) that spans \(V\) over \(F\).

\begin{proof}
Fix an \(F\)-basis of \(V\).  Away from the finitely many divisors that
support a pole of a matrix entry of a chosen finite generating set or of its
inverse, the standard lattice is invariant.  At each exceptional divisor,
choose one of the invariant lattices in the hypothesis.  The resulting
codimension-one lattices differ from one fixed rational lattice at only
finitely many divisors.  Lemma~\ref{lem:reflexive-lattice-patching} patches
them to the coherent reflexive sheaf
\[
 \mathcal E(U)=
 \{s\in V:s\in L_D\text{ for every }D\in U^{(1)}\};
\]
whose stalk at \(D\) is the prescribed lattice \(L_D\).  Each
\(\rho(\gamma)\) preserves every \(L_D\), and the codimension-one criterion
for reflexive sheaves therefore gives
\[
 \rho(\Gamma)\subset\operatorname{Aut}_X(\mathcal E).
\]

Because \(X\) is proper over \(K\),
\[
 M:=H^0\!\left(X,\mathcal End(\mathcal E)\right)
\]
is a finite-dimensional \(K\)-algebra.  Put
\(B=K[\rho(\Gamma)]\subset M\).  Absolute irreducibility and Burnside's
theorem say that the natural map
\[
 B\otimes_KF\longrightarrow\End_F(V)
\]
is surjective.  We include the descent argument that makes this map an
isomorphism.  The image of
\(\operatorname{rad}(B)\otimes_KF\) is a nilpotent two-sided ideal of the
full matrix algebra, hence is zero.  Since \(B\) itself embeds in
\(\End_F(V)\), this gives \(\operatorname{rad}(B)=0\).  A central idempotent
of \(B\) becomes a central idempotent of the full matrix algebra and is
therefore zero or one; consequently \(B\) has a single simple factor.  Its
center is a finite field extension of \(K\) whose image in
\(\End_F(V)\) consists of scalar matrices, so it embeds in \(F\).
Relative algebraic closedness forces \(Z(B)=K\).  Thus \(B\) is central
simple.  Then \(B\otimes_KF\) is simple, and its displayed nonzero
surjection is an isomorphism.  In particular \(\deg B=d\).  Reduced norm
becomes determinant after extension to \(F\), which proves
\(\rho(\Gamma)\subset B^1\).
\end{proof}

\begin{theorem}[integral descent on a proper model]
\label{thm:proper-model-descent}
Let \(K\) be a number field, let \(F/K\) be finitely generated with \(K\)
relatively algebraically closed in \(F\), and let \(\mathcal X\) be a
normal integral scheme, projective and flat over
\(\operatorname{Spec}\cO_K\), with function field \(F\).  Let \(V\),
\(\Gamma\), and \(\rho\) be as in
Theorem~\ref{thm:constant-field-descent}.  If \(\rho(\Gamma)\) preserves an
\(\cO_{\mathcal X,D}\)-lattice for every prime divisor
\(D\in\mathcal X^{(1)}\), then there are a degree-\(d\) central simple
\(K\)-algebra \(B\), an \(\cO_K\)-order \(\cO\subset B\), and an isomorphism
\[
 B\otimes_KF\simeq\End_F(V)
\]
for which \(\rho(\Gamma)\subset\cO^1\).
\end{theorem}

\begin{proof}
Choose the standard lattice away from the finitely many pole divisors and
patch the chosen invariant lattices to a reflexive sheaf \(\mathcal E\) on
\(\mathcal X\), exactly as above.  Properness now makes
\[
 H^0\!\left(\mathcal X,\mathcal End(\mathcal E)\right)
\]
a finite \(\cO_K\)-module.  Applying
Theorem~\ref{thm:constant-field-descent} on the generic fibre shows that
\(B:=K[\rho(\Gamma)]\) is central simple of degree \(d\) and splits over
\(F\).  The algebra
\[
 \cO:=\cO_K[\rho(\Gamma)]
\]
is a submodule of that finite module of global endomorphisms, hence is
finite over the noetherian ring \(\cO_K\).  It is torsion-free, contains
one, is closed under multiplication, and spans \(B\); it is therefore an
order.  Determinant one again becomes reduced norm one, so
\(\rho(\Gamma)\subset\cO^1\).
\end{proof}

For a finitely generated field \(F/\Q\), the relative algebraic closure of
\(\Q\) in \(F\) is a number field.  Thus
Theorem~\ref{thm:proper-model-descent} does not presuppose a number-field
form.  Boundedness at horizontal divisors descends the generic
representation to its constant field, while boundedness at vertical
divisors produces the order.

\begin{corollary}[Fuchsian divisorial descent]
\label{cor:fuchsian-divisorial-descent}
In Theorem~\ref{thm:proper-model-descent}, let \(d=2\), and suppose that one
real realization of the projective image of \(\Gamma\) is a cocompact
Fuchsian lattice.  Assume moreover that
\(K=\Q(\operatorname{Trd}\Gamma)\).  For each embedding
\(\sigma\colon K\hookrightarrow\C\), form the realization of
\(\Gamma\subset B^1\) in
\((B\otimes_{K,\sigma}\C)^1\).  Assume that its image is relatively compact
whenever \(\sigma\) is not the distinguished real embedding.  Then \(K\) is
totally real,
\(B\) splits at the distinguished real place and ramifies at every other
real place, and \(\Gamma\) is arithmetic.
\end{corollary}

\begin{proof}
The distinguished realization identifies \(B\) with \(M_2(\R)\).  A compact
subgroup of \(\SL_2(\C)\) is conjugate into \(\SU(2)\), so all its traces
are real.  Since the reduced traces of \(\Gamma\) generate \(K\), every
embedding of \(K\) has real image; hence \(K\) is totally real.  At another
real place, if \(B\) split, the compact conjugate of \(\Gamma\) would lie
in a compact subgroup of \(\SL_2(\R)\).  Its real Zariski closure would
then be proper, contrary to Zariski density, which follows by transport
from the distinguished Fuchsian realization.  Hence \(B\) ramifies there.
The norm-one group of the order constructed in the theorem is consequently an
arithmetic Fuchsian lattice containing \(\Gamma\).  An inclusion of
lattices has finite index, so \(\Gamma\) is arithmetic.
\end{proof}

We write \(\mathbb A_{k,f}\) for the finite adele ring of \(k\), namely the
restricted product of the completions \(k_v\) over the nonarchimedean places
with respect to their integer rings \(\cO_{k_v}\).

\begin{theorem}[arithmeticity from one bounded adjoint orbit]
\label{thm:one-orbit}
Let \(k\) be a number field, let \(G/k\) be an absolutely simple adjoint
group, and let \(v_0\) be an archimedean place.  Assume that
\(G(k_v)\) is compact at every archimedean \(v\ne v_0\).  Let
\(\Gamma<G(k)\) be Zariski dense and suppose that its image in
\(G(k_{v_0})\) is a lattice.  Choose \(0\ne X\in\mathfrak g(k)\).

Suppose that the orbit
\[
 \mathcal O_X=\{\Ad(\gamma)X:\gamma\in\Gamma\}
\tag{11.3}
\]
is bounded in \(\mathfrak g(k_v)\) at every finite place \(v\), and lies in
a fixed integral lattice for all but finitely many \(v\).  Then \(\Gamma\)
is arithmetic.  Equivalently, it is enough that \(\mathcal O_X\) have
compact closure in the finite-adelic Lie algebra
\(\mathfrak g(\mathbb A_{k,f})\).
\end{theorem}

\begin{proof}
The \(k\)-linear span of \(\mathcal O_X\) is invariant under \(\Gamma\),
hence under its Zariski closure.  It is therefore a nonzero ideal in the
simple Lie algebra \(\mathfrak g\), and so equals \(\mathfrak g\).  Choose
\(\gamma_1,\ldots,\gamma_N\in\Gamma\) such that
\(X_i=\Ad(\gamma_i)X\) is a \(k\)-basis.

Relative to this basis, the \(i\)-th column of \(\Ad(\gamma)\) is the
coordinate vector of
\[
 \Ad(\gamma)X_i=\Ad(\gamma\gamma_i)X,
\]
which belongs to the same orbit.  Thus \(\Ad(\Gamma)\) is bounded at every
finite place and integral outside a fixed finite set.  The same statement
for \(\gamma^{-1}\) shows that, relative to the basis \(X_i\), both
\(\Ad(\gamma)\) and its inverse are bounded at every finite place and
preserve one fixed global integral lattice outside a finite set \(S\).
For \(v\notin S\), \(\Ad(\Gamma)\) therefore lies in the compact-open
stabilizer of that lattice.  For \(v\in S\), its compact closure preserves
some \(\cO_{k_v}\)-lattice and is contained in a compact-open subgroup
\(K_v<G(k_v)\).  Enlarging \(S\) to accommodate an integral model of
\(G\), the product of these stabilizers is a compact-open subgroup
\(K_f<G(\mathbb A_{k,f})\) containing \(\Gamma\).  Here adjointness matters:
it makes \(\Ad\) a faithful closed immersion.

At every nondistinguished archimedean place the entire local group is compact.
The diagonal intersection
\[
 G(k)\cap\left(G(k_{v_0})\times
      \prod_{v\mid\infty,\,v\ne v_0}G(k_v)\times K_f\right)
\]
is an arithmetic lattice in the full archimedean product, by the
Borel--Harish-Chandra theorem \cite{BorelHarishChandra}.  Since all factors other than \(v_0\) are
compact, its projection is a lattice \(\Lambda<G(k_{v_0})\).  It contains
\(\Gamma\).  An inclusion of one lattice in another has finite index, so
\([\Lambda:\Gamma]<\infty\).
\end{proof}

In type \(A_1\) the orbit criterion can be checked using only three families
of actual word traces, once the rational form is fixed.  More precisely,
assume that the projective image of \(\Gamma\) lies in a fixed admissible
adjoint \(k\)-form \(G(k)\).  Let \(g\in\Gamma<\SL_2(\R)\) be noncentral, put
\(t=\Tr g\), and set
\[
 X=g-\frac t2I.
\]
Assume \(X\in\mathfrak g(k)\), and choose \(h_1,h_2,h_3\in\Gamma\) so that
\(X_i=h_iXh_i^{-1}\in\mathfrak g(k)\) is a basis of \(\mathfrak g\);
Zariski density guarantees such a choice.  The trace pairing is
nondegenerate, and
\[
 \Tr\bigl((\gamma X\gamma^{-1})X_i\bigr)
 =\Tr(\gamma g\gamma^{-1}h_i g h_i^{-1})-\frac{t^2}{2}.
\tag{11.4}
\]
The three explicitly displayed traces on the right are therefore
\(k\)-linear coordinates
of the orbit (11.3).  If they are contained in fixed fractional ideals at
the finite places and are uniformly bounded at every nondistinguished
archimedean place, then Theorem~\ref{thm:one-orbit} proves arithmeticity.
The assumption that the ambient admissible \(k\)-form is already fixed is
essential: membership of three scalar families in a number field does not by
itself construct that rational form.  Subject to this necessary hypothesis,
a finite collection of trace coordinates which faithfully records the matrix
orbit suffices; no simultaneous bound for every scalar trace is logically
necessary.

There is a still weaker criterion which uses only one exact trace fiber.  It
is stated using the coordinate map
\(z(A)=(\Trd(AB_1),\Trd(AB_2),\Trd(AB_3))\).

\begin{theorem}[arithmeticity from coordinate sums in one trace fiber]
\label{thm:one-fiber-window}
Let \(B/k\) be an admissible quaternion algebra with distinguished split
real place \(v_0\), and let \(\Lambda<\SL_1(B)(k)\) have Zariski-dense
projective image which is a lattice in
\(\PGL_1(B)(k_{v_0})\).  Choose a noncentral \(g\in\Lambda\), put
\[
 B^0=\{Y\in B:\Trd(Y)=0\},\qquad
 t=\Trd(g),\qquad X=g-\frac t2 1\in B^0,
\]
and choose \(h_1,h_2,h_3\in\Lambda\) so that
\(X_i=\Ad(h_i)X\) is a \(k\)-basis of \(B^0\).  Set
\[
 B_i=h_i g h_i^{-1},\qquad
 z(A)=\bigl(\Trd(AB_1),\Trd(AB_2),\Trd(AB_3)\bigr),
\]
and
\[
 Z_g^{\rm orb}=\{z(\gamma g\gamma^{-1}):\gamma\in\Lambda\}.
\]
If the following translated set of coordinatewise trace sums
\[
 \mathcal W_g=z(g)+Z_g^{\rm orb}-(t^2,t^2,t^2)
\]
has compact closure in \(\mathbb A_{k,f}^3\), then \(\Lambda\) is
arithmetic.  Equivalently, it is enough that the three coordinates of
\(\mathcal W_g\) be contained in fixed fractional ideals of \(k\).  This is
a criterion inside the specified quaternion algebra \(B/k\); it does not
construct \(k\) or \(B\) from the real trace set.
\end{theorem}

\begin{proof}
For \(C=\gamma g\gamma^{-1}\), define
\[
 P_\gamma=g+C-t1=X+\Ad(\gamma)X\in B^0.
\]
For context, Cayley--Hamilton gives the following intertwining identities:
\[
 (A+C-t1)A=C(A+C-t1),\qquad
 (A+C-t1)C=A(A+C-t1)
\]
whenever \(A,C\in\SL_1(B)\) have reduced trace \(t\).  Thus
\(P_\gamma\), when invertible, conjugates \(g\) to \(C\) projectively.  The
proof below only needs its linear expression
\(P_\gamma=X+\Ad(\gamma)X\), so no invertibility assumption is required.
Since \(B_i=X_i+\frac t2 1\) and
\(\Trd(P_\gamma)=0\),
\[
 \begin{split}
 \Trd(P_\gamma X_i)
 &=\Trd(P_\gamma B_i)\\
 &=\Trd(gB_i)+\Trd(CB_i)-t\Trd(B_i)\\
 &=z_i(g)+z_i(C)-t^2.
 \end{split}
\]
Consequently \(\mathcal W_g\) is exactly the image of
\(\{P_\gamma:\gamma\in\Lambda\}\) under the \(k\)-linear isomorphism
\[
 B^0\longrightarrow k^3,\qquad
 P\longmapsto
 \bigl(\Trd(PX_1),\Trd(PX_2),\Trd(PX_3)\bigr).
\]
This is an isomorphism because the reduced-trace pairing on \(B^0\) is
nondegenerate and the \(X_i\) form a basis.  Hence the hypothesis is
equivalent to finite-adelic precompactness of the set of matrices
\(\{P_\gamma\}\).  To justify the equivalent formulation in the statement,
recall that a diagonally embedded subset of \(k^3\) has compact closure in
\(\mathbb A_{k,f}^3\) exactly when it is bounded at every finite place and
integral at all but finitely many places.  Clearing the finitely many
remaining denominators puts each coordinate in one fractional
\(\cO_k\)-ideal.
Translation by the fixed rational vector \(X\) preserves precompactness, so
\[
 \{\Ad(\gamma)X:\gamma\in\Lambda\}
 =\{P_\gamma-X:\gamma\in\Lambda\}
\]
has compact closure in \(B^0(\mathbb A_{k,f})\).  Apply
Theorem~\ref{thm:one-orbit} to \(G=\PGL_1(B)\).  Admissibility gives the
compact nondistinguished archimedean factors, and the theorem gives
arithmeticity.
\end{proof}

The full fiber of the coordinate map \(z\),
\(Z_t=\{z(A):A\in\Lambda,\ \Trd(A)=t\}\) contains
\(Z_g^{\rm orb}\).  Finite-place boundedness of
\(z(g)+Z_t-t^2(1,1,1)\) is therefore sufficient, but unnecessarily strong:
the theorem needs only the subset arising from the conjugacy orbit of \(g\).
This makes the remaining extraction problem precise.  It is enough to bound,
at every finite place, the matrices
\(g+\gamma g\gamma^{-1}-t1\); one need not first construct an integral model
for the entire scalar trace set.

For the high-multiplicity application even the basepoint need not be fixed.
The following asymptotic version is strictly weaker than a bounded full orbit.

\begin{theorem}[arithmeticity from nonexpanding adjoint directions]
\label{thm:asymptotic-orbit}
Let \(G/k\) be absolutely simple and adjoint, with
\(G(k_v)\) compact at every archimedean place \(v\ne v_0\).  Let
\(\Gamma<G(k)\) be finitely generated and Zariski dense, and suppose that
its image at \(v_0\) is a lattice.  Fix a word metric and a finite set \(S\)
of finite places such that
\(\Gamma\subset\mathcal G(\cO_{k,S})\) for an integral model
\(\mathcal G\).  Suppose that for every \(n\) there is
\(0\ne X_n\in\mathfrak g(k)\) such that, for every \(v\in S\),
\[
 \max_{|\gamma|\le n}\|\Ad(\gamma)X_n\|_v
 \le C_v\|X_n\|_v,
\]
where \(C_v\) is independent of \(n\).  Then \(\Gamma\) is arithmetic.
As in Theorem~\ref{thm:one-orbit}, the hypotheses already include a fixed
\(k\)-form of \(G\); the conclusion is not a method for recovering that form
from a representation over \(\R\).
\end{theorem}

\begin{proof}
Fix \(v\in S\).  We regard the same rational vector \(X_n\) in
\(\mathfrak g(k_v)\), but we are free to multiply it by a scalar in
\(k_v^\times\), since the displayed ratio of norms is homogeneous in
\(X_n\).  Normalize it to lie in a fixed compact annulus which is bounded
away from zero.  Such an annulus is compact over a nonarchimedean local
field.  Since \(S\) is finite, successive extraction gives one subsequence
which converges at every \(v\in S\), say to
\(0\ne X_v\in\mathfrak g(k_v)\).  Fix \(\gamma\in\Gamma\).  For every
sufficiently large index in this subsequence one has \(|\gamma|\le n\);
continuity of \(\Ad(\gamma)\) and the assumed inequality then give
\[
 \|\Ad(\gamma)X_v\|_v\le C_v\|X_v\|_v.
\]
The constant is independent of \(\gamma\), so the full orbit
\(\Ad(\Gamma)X_v\) is bounded.

Its \(k_v\)-span is a nonzero \(\Ad(\Gamma)\)-invariant subspace.  Zariski
density makes it invariant under \(G_{k_v}\), and absolute simplicity makes
the adjoint representation irreducible; the span is therefore all of
\(\mathfrak g(k_v)\).  Choose
\(\delta_1,\ldots,\delta_N\in\Gamma\) such that
\(\Ad(\delta_i)X_v\) is a basis.  The \(i\)-th column of
\(\Ad(\gamma)\) in this basis is the coordinate vector of
\(\Ad(\gamma\delta_i)X_v\), which belongs to the bounded orbit.  Thus all
matrix entries of \(\Ad(\Gamma)\) are bounded.  The same statement applies
to inverses because \(\Gamma\) is a group.  Hence
\(\Ad(\Gamma)\) has compact closure in \(\GL(\mathfrak g(k_v))\).
Since \(G\) is adjoint, \(\Ad:G\to\GL(\mathfrak g)\) is a faithful closed
immersion, and \(\Gamma\) is relatively compact in \(G(k_v)\).

This holds for each \(v\in S\).  Outside \(S\), the assumed integral model
places \(\Gamma\) in the compact-open subgroup
\(\mathcal G(\cO_{k_v})\).  At each \(v\in S\), the compact closure of
\(\Gamma\) is contained in some compact-open subgroup of \(G(k_v)\).
Taking their restricted product gives a compact-open
\(K_f<G(\mathbb A_{k,f})\) containing \(\Gamma\).  The
Borel--Harish-Chandra argument in Theorem~\ref{thm:one-orbit} now proves
arithmeticity.
\end{proof}

In type \(A_1\), take varying elements \(g_n\), put
\(t_n=\Trd(g_n)\), \(X_n=g_n-\frac{t_n}{2}1\), and define
\[
 P_{n,\gamma}=g_n+\gamma g_n\gamma^{-1}-t_n1
              =X_n+\Ad(\gamma)X_n.
\]
The hypothesis of Theorem~\ref{thm:asymptotic-orbit} is equivalent, up to
fixed constants, to uniform bounds
\(\|P_{n,\gamma}\|_v\ll_v\|X_n\|_v\) on the growing conjugacy balls.  Choose
three fixed traceless matrices \(B_i\) whose trace-pairing functionals form
a basis of the dual space.  The preceding matrix bounds are then equivalent
to bounds for the three scalar quantities
\[
 \Trd(g_nB_i)+\Trd(\gamma g_n\gamma^{-1}B_i)
       -t_n\Trd(B_i).
\]
Thus the trace fiber of high multiplicity may vary with the scale.  What is
needed is a nonzero direction in each such fiber whose adjoint orbit, up to
the corresponding word radius, remains bounded by a fixed multiple of the
size of that direction at every relevant finite place.

This condition has an exact quantitative converse.

\begin{lemma}[uniform expansion over archimedean and divisorial fields]
\label{lem:projective-escape}
Let \(K\) be \(\R\), \(\C\), or a complete discretely valued field, and let
\(\rho:\Gamma\to\GL(V)\) be an irreducible finite-dimensional
representation of a finitely generated group.  Fix a norm on \(V\); in the
discretely valued case take the norm defined by a lattice.  Put
\[
 m_n=\inf_{0\ne x\in V}\ \max_{|\gamma|\le n}
       \frac{\|\rho(\gamma)x\|}{\|x\|}.
\]
Then \(\rho(\Gamma)\) is bounded if and only if
\(\sup_nm_n<\infty\).  If it is unbounded, there are
\(c>0\) and \(\lambda>1\) such that \(m_n\ge c\lambda^n\) for every \(n\).
Over \(\R\), \(\C\), or a nonarchimedean local field, boundedness is
equivalent to relative compactness.  Over an arbitrary complete discretely
valued field, boundedness is equivalent to preservation of a lattice.
\end{lemma}

\begin{proof}
The word balls give the supermultiplicativity
\(m_{n+r}\ge m_nm_r\).  Indeed, for a fixed \(x\ne0\), choose a word
\(\delta\) of length at most \(r\) which realizes the maximum over the
\(r\)-ball.  Applied to \(\rho(\delta)x\), the definition of \(m_n\)
provides a word \(\gamma\) of length at most \(n\) with
\[
 \|\rho(\gamma\delta)x\|
 \ge m_n\|\rho(\delta)x\|
 \ge m_nm_r\|x\|.
\]
Taking first the maximum and then the infimum proves the inequality.

Suppose first that \(K=\R\) or \(\C\) and that
\(\sup_n m_n=M<\infty\).  Choose unit vectors \(x_n\) for which the maximum
over the \(n\)-ball is at most \(M+1/n\).  Compactness of the unit sphere
gives a subsequence converging to a unit vector \(x\).  For fixed
\(\gamma\in\Gamma\), passage to the limit after \(n\ge|\gamma|\) gives
\(\|\rho(\gamma)x\|\le M\).  Thus the full orbit of \(x\) is bounded.

Now let \(K\) be complete discretely valued, with valuation ring \(R\) and
uniformizer \(\pi\), and identify the chosen lattice with \(R^d\).  After
normalizing \(x\) to be primitive, every ratio occurring in the finite
maximum defining \(m_n\) lies in the discrete value group.  The identity is
in every word ball, so these ratios are at least one.  Consequently the
infimum is attained by a primitive vector.  Choose an integer \(a\) with
\(M\le|\pi|^{-a}\), and define the descending \(R\)-submodules
\[
 L_n=\{x\in R^d:\rho(\gamma)x\in\pi^{-a}R^d
                     \text{ for every }|\gamma|\le n\}.
\]
Every \(L_n\) contains a primitive vector.  For \(r\ge1\), let
\[
 I_{n,r}=(L_n+\pi^rR^d)/\pi^rR^d.
\]
For fixed \(r\), this descending chain stabilizes because
\((R/\pi^rR)^d\) has finite module length; finiteness of the residue field is
not needed.  Write its stable value as \(N_r\).  We have \(N_1\ne0\), and
the reduction maps \(N_{r+1}\to N_r\) are surjective: choose \(n\) beyond
the two stabilization indices, and use the surjection from the image of
\(L_n\) modulo \(\pi^{r+1}\) onto its image modulo \(\pi^r\).  Starting
with a nonzero element of \(N_1\), choose compatible lifts through all
\(N_r\).  Completeness identifies their inverse limit with a primitive
vector \(x\in R^d\).  For fixed \(n\), the vector \(x\) belongs to
\(L_n+\pi^rR^d\) for every \(r\).  The module \(L_n\), being a finite
intersection of inverse images of lattices, is closed, so \(x\in L_n\).
Thus \(x\) again has a bounded full orbit.

In either case, irreducibility gives orbit vectors
\(\rho(\delta_1)x,\ldots,\rho(\delta_{\dim V})x\) which form a basis.
In that basis every column of \(\rho(\gamma)\) is the coordinate vector of
another point of the bounded orbit.  Hence \(\rho(\Gamma)\) is bounded; the
same assertion for inverses is automatic because \(\Gamma\) is a group.
Conversely, a bounded group plainly makes the \(m_n\) uniformly bounded.
Over a complete discrete valuation ring, a bounded group preserves the
lattice \(\sum_{\gamma\in\Gamma}\rho(\gamma)R^d\); boundedness sandwiches
this module between two homothetic lattices, and noetherianity makes it a
lattice.  The converse is immediate.  Over a locally compact field, bounded
matrices together with bounded inverses have compact closure.

If \(\rho(\Gamma)\) is unbounded, the preceding equivalence says that
\(m_n\) is unbounded, so
\(m_N>1\) for some \(N\).  Write \(n=qN+r\), with \(0\le r<N\).
Since the identity belongs to every word ball, \(m_r\ge1\), and
supermultiplicativity gives
\[
 m_n\ge m_N^q\ge m_N^{-1}\bigl(m_N^{1/N}\bigr)^n.
\]
Taking \(c=m_N^{-1}\) and \(\lambda=m_N^{1/N}>1\) proves the final claim.
\end{proof}

\begin{corollary}[descent from uniformly nonexpanding vectors in word balls]
\label{cor:valuative-nonexpansion}
Use the notation of Theorem~\ref{thm:proper-model-descent}.  For a prime
divisor \(D\in\mathcal X^{(1)}\), let \(F_D\) be the completion of \(F\) at
its discrete valuation.  Suppose that, for every \(D\) and every \(n\ge1\),
there is a nonzero \(x_{D,n}\in V\otimes_FF_D\) such that
\[
 \max_{|\gamma|\le n}\|\rho(\gamma)x_{D,n}\|_D
       \le C_D\|x_{D,n}\|_D,
 \tag{11.5}
\]
where \(C_D\) is independent of \(n\).  Then \(\rho(\Gamma)\) descends to
the norm-one group of an order in a degree-\(d\) central simple algebra over
\(K\), as in Theorem~\ref{thm:proper-model-descent}.

If \(d=2\), the distinguished real projective image is a cocompact Fuchsian
lattice, the reduced traces generate \(K\), and every nondistinguished
archimedean realization is relatively compact, then \(\Gamma\) is
arithmetic.
\end{corollary}

\begin{proof}
Global absolute irreducibility remains absolute after extension to \(F_D\).
Lemma~\ref{lem:projective-escape} applied over \(F_D\) turns (11.5) into
boundedness of \(\rho(\Gamma)\) at \(D\), hence into an invariant
\(\mathcal O_{\mathcal X,D}\)-lattice.  Theorem
\ref{thm:proper-model-descent} applies simultaneously to all divisors.  The
last assertion is Corollary~\ref{cor:fuchsian-divisorial-descent}.
\end{proof}

The same conclusion follows from nonexpansion in the adjoint representation.
Indeed, boundedness in \(\PGL_d(F_D)\) lifts to boundedness in
\(\SL_d(F_D)\): in Cartan coordinates the adjoint bound controls all
differences of exponents, while determinant one says that their sum is zero.

At any divisorial or finite place where a proposed nonarithmetic group has
unbounded adjoint image, Lemma~\ref{lem:projective-escape} says that
\emph{every} choice of
\(X_n\) has some \(|\gamma|\le n\) with
\[
 \|\Ad(\gamma)X_n\|_v\ge c_v\lambda_v^n\|X_n\|_v.
\]
For large \(n\), the factor \(c_v\lambda_v^n\) is greater than one.  The two
summands in
\(P_{n,\gamma}=X_n+\Ad(\gamma)X_n\) then have different norms, so the
ultrametric equality gives
\(\|P_{n,\gamma}\|_v=\|\Ad(\gamma)X_n\|_v\).  Since the three chosen
trace-pairing coordinates form a fixed linear coordinate system, at least
one of them
also grows exponentially, up to constants depending only on that system.
This gives a valuative sufficient condition for the compact conjectures that
does not assume a number-field form in advance.  It is enough to deduce
(11.5) for the adjoint representation from sparse real traces, at every
horizontal and vertical divisor of one proper model, and then to prove that
the nondistinguished archimedean realizations are compact.  Corollary
\ref{cor:valuative-nonexpansion} then implies arithmeticity.  If (11.5)
fails at a divisor, the failure is quantitative: every nonzero vector is
expanded by at least \(c_v\lambda_v^n\) somewhere in the word ball of
radius \(n\).

The following stronger scalar statement would also settle the conjectures.
It combines two logically separate tasks: constructing a common number
field from the trace set and proving integrality and boundedness at all of
its nondistinguished places.

\begin{conjecture}[adelic trace decomposition by square classes]\label{conj:window}
Let \(\Gamma\) be a cocompact Fuchsian lattice whose trace set has bounded
clustering, or merely linear growth, and put
\(V=\Gamma/\Gamma^{(2)}\).  There exist a number field \(k\) with a
distinguished real place, Kummer representatives \(s_v\), fractional ideals
\(I_v\), and constants \(C_{v,\sigma}\) such that every trace in the
component indexed by \(v\)
has the form
\[
 \Tr\widetilde\gamma=s_v a_\gamma,
 \qquad a_\gamma\in I_v,
 \qquad |\sigma(a_\gamma)|\le C_{v,\sigma}
\tag{11.6}
\]
for every nondistinguished archimedean \(\sigma\).  For the identity
component one
may take \(s_0\in k^\times\).
\end{conjecture}

For the trace gap on \(\Gamma^{(2)}\), (11.6) is stronger than logically
necessary: it would suffice to place all differences in the neutral
component in one fractional ideal with uniformly bounded nondistinguished
conjugates.  The component-by-component form is stated because it also
recovers the canonical decomposition of arithmetic traces by square classes.

\begin{proposition}
Conjecture~\ref{conj:window} implies both the compact Sarnak and compact
Schmutz conjectures.
\end{proposition}

\begin{proof}
Apply Lemma~\ref{lem:galois-window} to the identity component
\(T(\Gamma^{(2)})\), after absorbing \(s_0\) into its fractional ideal.  This
gives positive gap.  Theorem~\ref{thm:kummer-reduction} gives arithmeticity.
\end{proof}

The content of (11.6) is precisely what arbitrary graph coloring omits.  Any
bounded-clustering subset of \(\R\) can be finitely colored into separated
sets: connect two values at distance less than one and greedily color the
resulting bounded-degree graph.  Such colors need not factor through
\(\Gamma/\Gamma^{(2)}\), respect (10.1), lie in one number field, or have
bounded conjugates.

\subsection{Why algebraicity alone is insufficient}

Even a fixed number field and algebraic integrality do not give a gap at the
distinguished embedding.
The two-color set (10.6) already proves this.  A one-point-per-unit-interval
variant, with \(\alpha=\sqrt2\), is the set
\[
 S=\{n+\{\alpha n^2\}:n\in\Z\}\subset\Q(\sqrt2).
\]
This set has one point in each unit interval, and therefore has bounded
clustering.
Weyl equidistribution of
\((\{\alpha n^2\},\{\alpha(n+1)^2\})\) gives adjacent pairs whose
differences tend to zero, while the second embedding is unbounded.  Variants
over \(\Q\) with denominators
tending through primes show the independent finite-place obstruction.  The
trace ring can also be misleading: under one real embedding, \(\cO_k\) is
dense whenever \([k:\Q]>1\).  Derived trace sets are separated not because
their additive span is discrete, but because their images under every
nondistinguished embedding remain uniformly bounded.

\section{Trace multiplicity, fixed trace coordinates, and additive energy}

\subsection{A finite-type divisor theorem}

We first prove a divisor bound over an arbitrary finitely generated field of
characteristic zero.  This generality is necessary because the matrix
entries of a finitely generated group need not be algebraic numbers.

Let \(L/\Q\) be finitely generated, and let \(K\) be the algebraic closure of
\(\Q\) inside \(L\); then \(K\) is a number field.  Choose a transcendence
basis \(t=(t_1,\ldots,t_r)\), so \(L/K(t)\) is finite.  After enlarging a
finite set \(S\) of places of \(K\) and inverting one nonzero polynomial
\(a\in\cO_{K,S}[t]\), one may choose a normal domain \(R\subset L\), with
fraction field \(L\), which is finite free over
\[
 A=\cO_{K,S}[t_1,\ldots,t_r,a^{-1}].
\]
Here \(S\) contains all archimedean places and
\[
 \cO_{K,S}=\{x\in K:v_{\mathfrak p}(x)\ge0
       \text{ for every finite place }\mathfrak p\notin S\}
\]
is the ring of \(S\)-integers.  Thus its elements may have denominators only
at the finitely many finite places belonging to \(S\).
Fix an \(A\)-basis \(e_1,\ldots,e_D\) of \(R\).  For \(C>0\), let
\(R_C[n]\) consist of the elements which can be written
\[
 \sum_{j=1}^D a^{-m}P_j(t)e_j,
 \qquad
 0\le m\le Cn,\quad
 \deg P_j\le Cn,\quad h_{\rm aff}(P_j)\le Cn,
\tag{12.1}
\]
where \(P_j\in\cO_{K,S}[t]\) and
\[
 h_{\rm aff}(P_j)=h_K\bigl([1:(\operatorname{coefficients of }P_j)]\bigr)
 \le Cn.
\]
Here \(h_K\) is the normalized logarithmic projective Weil height: it is the
sum, over all places of \(K\), of the logarithm of the largest absolute value
of the homogeneous coordinates, with the usual local-degree weights.
Thus the height is \emph{affine}, not merely the projective height of the
coefficient vector: in particular, the constant polynomial \(N\in\Z\) has
height \(\log N\), rather than zero.  This normalization is essential for
controlling vertical prime content.  Changing the chosen model, basis, or
finite generating set only changes \(C\).

\begin{theorem}[finite-type divisor bound]\label{thm:finite-type-divisor}
For every fixed \(C\),
\[
 \sup_{0\ne q\in R_C[n]}
 \#\{(b,c)\in R_C[n]^2:bc=q\}=\exp(o_C(n)).
\tag{12.2}
\]
The \(o_C(n)\) is uniform in \(q\).
\end{theorem}

\begin{proof}
We first count divisors in the polynomial ring before passing to its finite
normal extension.  Enlarge \(S\) so that \(\cO_{K,S}\) is a principal ideal
domain; this is possible because the ideal class group of \(K\) is finite.
Thus \(\cO_{K,S}[t]\) is a unique factorization domain.

We put a positive weight on each of its prime elements.  A vertical prime
is a prime \(\mathfrak p\notin S\) of \(\cO_K\); give it weight
\(\log N\mathfrak p/\log 2\), so every vertical weight is at least one.
A horizontal prime is represented, up to
multiplication by \(K^\times\), by a primitive irreducible polynomial
\(f\in K[t]\).  At a nonarchimedean place \(v\), let \(M_v(f)\) be the
Gauss norm, the maximum \(v\)-absolute value of its coefficients.  At an
archimedean place put
\[
 \log M_v(f)=
 \int_{(S^1)^r}\log|f(z_1,\ldots,z_r)|_v\,d\mu.
\]
The global multivariable Mahler height is the normalized sum
\[
 m_K(f)=\frac1{[K:\Q]}\sum_v [K_v:\Q_v]\log M_v(f).
\]
The product formula makes \(m_K\) unchanged when \(f\) is multiplied by a
nonzero scalar.  Gauss norms are multiplicative, and logarithms turn the
archimedean identity \(M_v(fg)=M_v(f)M_v(g)\) into addition.  Hence
\(m_K(fg)=m_K(f)+m_K(g)\).

The standard Mahler--coefficient inequalities compare \(m_K(f)\) with the
projective height of the primitive coefficient vector up to a constant times
\(\deg f\).  Choose a constant \(C_r\), depending only on the number of
variables, so that
\[
 w(f)=C_r\deg f+m_K(f)
\]
is at least one for every nonconstant irreducible \(f\), and controls both
\(\deg f\) and its projective coefficient height.  This weight is additive
under multiplication.  Moreover, only finitely many horizontal primes have
\(w(f)\le T\): their degrees and coefficient heights are bounded, so this is
Northcott's finiteness theorem
\cite[Theorem~1.6.8]{BombieriGubler}.  A vertical prime is indexed by a
prime ideal of the fixed ring of integers, so only finitely many have norm,
and hence weight, at most a given bound.

We use the following elementary weighted divisor observation.  Suppose prime
weights satisfy \(w(p)\ge1\), have finite sublevel sets, and
\[
 \sum_p e_p w(p)\le C_0n.
\]
Then
\[
 \log\prod_p(e_p+1)=o(n).
\tag{12.3}
\]
Indeed, fix \(T\).  The finitely many primes of weight at most \(T\)
contribute \(O_T(\log n)=o(n)\).  For the remaining primes,
\(\log(e+1)\le e\log2\), and hence their contribution is at most
\(C_0n\log2/T\).  First let \(n\to\infty\), and then let \(T\to\infty\).

If \(Q\in\cO_{K,S}[t]\) has total degree and affine coefficient height
\(O(n)\), the sum of the weights of its horizontal prime factors, with
multiplicity, is \(O(n)\) by additivity and the coefficient comparison.
The weighted multiplicity of its vertical content is also \(O(n)\): after
choosing a primitive polynomial, the norm of the removed coefficient content
is bounded by the affine height.  Equation (12.3) therefore says
that \(Q\) has only \(\exp(o(n))\) divisors.  Inverting the fixed polynomial
\(a\) does not change this estimate: its finitely many prime factors become
units, and clearing a power \(a^{O(n)}\) adds only polynomially many choices.

Now take \(0\ne q\in R_C[n]\).  Multiplication by \(q\) on the fixed free
\(A\)-module \(R\) has a \(D\times D\) matrix whose entries, after clearing
\(a^{O(n)}\), have degree and coefficient height \(O(n)\).  Its determinant
is the field norm
\[
 Q=N_{L/K(t)}(q).
\]
It has the same complexity bounds.  Let \(P\) be a height-one prime of the
normal domain \(R\), away from the already inverted boundary, and put
\(p=P\cap A\).  The norm valuation formula is
\[
 v_p(Q)=\sum_{P\mid p}f(P/p)v_P(q),
\tag{12.4}
\]
where the residue degrees \(f(P/p)\) are positive and the number of primes
above \(p\) is at most \(D\).

If \(bc=q\) in \(R\), normality makes every height-one localization
\(R_P\) a discrete valuation ring.  Since \(b,c\in R\), it follows that
\[
 0\le v_P(b)\le v_P(q)
\]
at every such \(P\).  Formula (12.4) shows that the possible collections of
all valuations \(v_P(b)\) are bounded by a fixed power, depending only on
\(D\), of the number of divisors of \(Q\).  There are therefore
\(\exp(o(n))\) possible divisor patterns.

It remains to count elements having one fixed pattern.  If \(b,b'\) have the
same valuations at every height-one prime, then \(u=b'/b\) and \(u^{-1}\)
belong to every localization \(R_P\).  Normality gives the Krull intersection
identity
\[
 R=\bigcap_{\operatorname{ht}P=1}R_P,
\]
so \(u\in R^\times\).  Take a fixed normal projective compactification of
the generic fiber
\(\operatorname{Spec}(R\otimes_{\cO_{K,S}}K)\).  Orders along its finitely
many boundary divisors map \(R^\times\), modulo constant units, into a fixed
lattice \(\Z^s\); this is the Rosenlicht--Samuel unit theorem
\cite{Rosenlicht,SamuelUnits}.  The fixed
filtration bounds the \emph{absolute} valuation of \(b\) and \(b'\) by
\(O(n)\) along every boundary divisor.  Indeed, their displayed coordinate
expressions give the lower bound.  Applying the field norm gives a rational
function of degree and affine coefficient height \(O(n)\), hence of boundary
valuation \(O(n)\); the norm-valuation formula and the lower bounds at all
other divisors above the same base divisor give the corresponding upper
bound.  Thus every boundary order of \(u=b'/b\) is \(O(n)\), and there are
only \(n^{O(1)}\) possible order vectors.  Choose once and for all generators
for their image lattice.  A representative \(u_0\) made from these
generators with exponents \(O(n)\) has logarithmic height \(O(n)\).  In
\(b'=k u_0b\), clear the common denominator \(a^{O(n)}\) and compare one
nonzero coefficient in the fixed basis; the affine coefficient bounds give
\(h(k)=O(n)\).  The residual constant \(k\) is an \(S'\)-unit of \(K\), for
one fixed finite \(S'\).
Dirichlet's \(S'\)-unit theorem
\cite[Chapter~I, Sections~7 and~11]{Neukirch} gives only
\(n^{O(1)}\) such constants.
Thus each divisor pattern contributes only polynomially many factors, which
proves (12.2).

For completeness, the model used in the statement always exists.  Noether
normalization makes \(L\) finite over \(K(t)\).  Start with the integral
closure of \(\cO_{K,S}[t]\) in \(L\); after initially enlarging \(S\), this
normalization is finite.  Invert the finitely many denominators needed to
contain all prescribed entries and their inverses.  Generic freeness then
provides one further polynomial localization after which the resulting
normal algebra is finite free over \(A\).  After one final finite enlargement
of \(S\), the ring \(\cO_{K,S}\) is a principal ideal domain, as required
above.
\end{proof}

\begin{theorem}[subexponential fibers for two trace coordinates]\label{thm:finite-type-fiber}
Let \(F/\Q\) be finitely generated, let
\(\Delta<\SL_2(F)\) be finitely generated, and let \(B\in\Delta\) be regular
semisimple.  For any word metric,
\[
 \sup_{x,z}
 \#\{A\in\Delta:|A|\le n,\ \Tr A=x,\ \Tr(AB)=z,\
       \langle A,B\rangle\text{ absolutely irreducible}\}
 =\exp(o(n)).
\tag{12.5}
\]
\end{theorem}

\begin{proof}
Pass to one finite extension of \(F\) which splits \(B\), and conjugate
\[
 B=\diag(\lambda,\lambda^{-1}),\qquad\lambda\ne\lambda^{-1}.
\]
Choose the normal arithmetic model \(R\) above to contain the entries of the
conjugated finite generating set and of its inverses.  Repeated multiplication
by these fixed matrices shows that every entry of a word of length at most
\(n\) belongs to \(R_C[n]\) for one fixed \(C\): total degrees, denominator
orders, and logarithmic coefficient heights all grow at most linearly with
the number of factors.

Write \(A=(\begin{smallmatrix}a&b\\c&d\end{smallmatrix})\).  The two
specified values \(\Tr A=x\) and \(\Tr(AB)=z\) determine the diagonal
entries:
\[
 a=\frac{z-\lambda^{-1}x}{\lambda-\lambda^{-1}},
 \qquad
 d=\frac{\lambda x-z}{\lambda-\lambda^{-1}}.
\tag{12.6}
\]
The determinant equation then fixes
\[
 bc=q:=ad-1.
\tag{12.7}
\]
The pair \(\langle A,B\rangle\) is absolutely reducible exactly when it has
a common \(B\)-eigenline, which in these coordinates is equivalent to
\(bc=0\).  Thus absolute irreducibility is precisely \(q\ne0\).

If a fiber is nonempty, choose one matrix in it.  Its entries show that the
fixed \(a,d,q\) have complexity \(O(n)\).  Theorem
\ref{thm:finite-type-divisor} gives only \(\exp(o(n))\) pairs \((b,c)\) of
the required complexity, after enlarging the fixed filtration constant
\(C\), satisfying (12.7).  Once \(a,b,c,d\) are fixed, so
is \(A\).  The estimate is uniform in \(x,z\).
\end{proof}

The same divisor argument applies uniformly to a family of quadratic norm
equations.  We can therefore study addition inside an exact trace fiber
without separately diagonalizing over each varying quadratic field.

\begin{lemma}[a uniform divisor bound for quadratic norm equations]
\label{lem:quadratic-norm-divisor}
Use the arithmetic model and filtration \(R_C[n]\) of (12.1).  For every
fixed \(C>0\), uniformly in \(d,N\in R_C[n]\) with \(dN\ne0\),
\[
 \#\{(y,v)\in R_C[n]^2:y^2-dv^2=N\}=\exp(o_C(n)).
\]
The conclusion remains true if the four filtration constants are different
but fixed.
\end{lemma}

\begin{proof}
Put \(E_d=\operatorname{Frac}(R)[z]/(z^2-d)\), a quadratic \'etale
algebra which may be a product, and let \(R_d\) be the normalization of
\(R\) in \(E_d\).  A solution determines \(\alpha=y+vz\in R_d\) with
\[
 N_{E_d/\operatorname{Frac}(R)}(\alpha)=N.
\]
The map \((y,v)\mapsto\alpha\) is injective because \(1,z\) is a basis of
the rank-two algebra.

The divisor-pattern argument in Theorem~\ref{thm:finite-type-divisor} is
uniform in this varying quadratic algebra.  Indeed, \(R_d\) has degree at
most \(2D\) over \(A\).  Applying the norm-valuation formula down to \(A\)
bounds the possible height-one valuations of \(\alpha\), by a fixed power
depending only on \(2D\), by the divisor count of
\[
 N_{\operatorname{Frac}(R)/K(t)}(N).
\]
The latter has degree and affine coefficient height \(O_C(n)\).  Thus there
are \(\exp(o_C(n))\) possible divisor patterns, uniformly in \(d,N\).

It remains to count associates uniformly as \(d\) varies.  Let
\(X=\operatorname{Spec}(R\otimes_{\cO_{K,S}}K)\), fix a normal projective
compactification \(\overline X\), and write its boundary as
\(D_1\cup\cdots\cup D_s\).  Let \(\overline Y_d\) be the normalization of
\(\overline X\) in \(E_d\).  Above each \(D_i\) there are at most two prime
divisors.  Their orders embed \(R_d^\times\), modulo its constant units, in
\(\Z^{2s}\).  For quotients of two admissible \(\alpha\)'s every such order
is \(O_C(n)\).  Indeed,
\[
 \alpha+\bar\alpha=2y,\qquad \alpha\bar\alpha=N.
\]
At a boundary valuation the orders of \(y\) and \(N\) are \(O_C(n)\).  If
the two root valuations differ, their minimum is the valuation of \(2y\);
if they agree, each is half the valuation of \(N\).  Thus both are
\(O_C(n)\), giving only \(n^{O(1)}\) boundary vectors.

For one boundary vector, the quotient of two units is a constant unit in
the constant \'etale algebra of \(E_d\), of degree at most
\(2[K:\Q]\) over \(\Q\).  Its height is \(O_C(n)\): use
\(\alpha^{-1}=\bar\alpha/N\) and express its trace and norm as rational
functions of degree and affine coefficient height \(O_C(n)\).  These
constants are \(S'\)-units for one fixed finite set of rational places.  In
the full archimedean-plus-\(S'\) logarithmic embedding their rank is bounded
in the fixed degree.  Northcott's theorem
\cite[Theorem~1.6.8]{BombieriGubler} gives a uniform positive lower
height for a non-torsion algebraic number of that degree, and Minkowski's
second theorem then shows that an \(O(n)\) box contains only \(n^{O(1)}\)
such units.  The split case is identical componentwise.  There are
\(\exp(o(n))\) possible divisor patterns, so summing over them proves the
lemma.
\end{proof}

\begin{theorem}[additive representations inside a fixed trace fiber]
\label{thm:same-trace-sum-fiber}
Let \(F/\Q\) be finitely generated and let
\(\Delta<\SL_2(F)\) be finitely generated.  Assume that the identity is the
only element of \(\Delta\) having trace \(2\); this holds, in particular,
for the full lift of a torsion-free cocompact Fuchsian lattice.  For a word
metric put
\[
 \mathcal F_t(n)=\{A\in\Delta:|A|\le n,\ \Tr A=t\}.
\]
Uniformly in \(t\in F\) and \(S\in M_2(F)\) with \(S\ne tI\),
\[
 \#\{(A,C)\in\mathcal F_t(n)^2:A+C=S\}=\exp(o(n)).
\]
For the excluded sum \(S=tI\), the pairs are exactly
\((A,A^{-1})\), so the fiber has size \(|\mathcal F_t(n)|\).
\end{theorem}

\begin{proof}
For a nonempty fiber choose one pair.  In one fixed normal arithmetic model,
all entries of \(A,C\), and hence \(t,S\), have complexity \(O(n)\).  Put
\[
 P=S-tI=A+C-tI
  =\begin{pmatrix}r&s\\u&-r\end{pmatrix},
 \qquad d=r^2+su=-\det P,
\]
and write
\[
 A=\frac t2I+\begin{pmatrix}x&b\\c&-x\end{pmatrix},
 \qquad q_0=\frac{t^2}{4}-1.
\]
After inverting \(2\) in the fixed model, the equations
\(\det A=\det C=1\) are exactly
\[
 x^2+bc=q_0,\qquad 2rx+sc+ub=d.
\]

First suppose \(d=0\).  Cayley--Hamilton gives
\(P=A-C^{-1}\), and
\[
 \det(A-C^{-1})=2-\Tr(AC).
\]
Thus \(\Tr(AC)=2\).  The hypothesis forces \(AC=I\), whence
\(P=A+A^{-1}-tI=0\), contrary to \(S\ne tI\).  Hence this case is empty.

Assume \(d\ne0\) and set \(q=q_0-d/4\).  Polarizing the determinant in the
preceding system gives \(\det(A-C)=-4q\).  If \(q=0\), then
\(2-\Tr(AC^{-1})=\det(A-C)=0\), so again the trace-two hypothesis gives
\(A=C=S/2\).  There is at most one pair.

It remains to consider \(q\ne0\).  If \(s\ne0\), put
\[
 y=2(sx-rb),\qquad v=2b-s.
\]
A direct calculation from the two determinant equations yields
\[
 y^2-dv^2=s^2(4q_0-d)=4s^2q\ne0.
\]
The map \((x,b,c)\mapsto(y,v)\) is injective: recover \(b\) from \(v\),
then \(x\) from \(y\), and finally \(c\) from the second determinant
equation.  Lemma~\ref{lem:quadratic-norm-divisor} gives \(\exp(o(n))\)
possibilities.  If \(s=0\) but \(u\ne0\), transpose the calculation, using
\[
 y=2(ux-rc),\qquad v=2c-u.
\]
Finally, if \(s=u=0\), then \(r\ne0\), the second determinant equation gives
\(x=r/2\), and
\[
 bc=q_0-r^2/4=q\ne0.
\]
Theorem~\ref{thm:finite-type-divisor} gives \(\exp(o(n))\) possibilities.
These cases exhaust the noncentral fiber.  If \(S=tI\), Cayley--Hamilton
gives \(C=tI-A=A^{-1}\), proving the last assertion.
\end{proof}

If \(r_{t,n}(S)\) denotes the representation function in the additive group
of \(M_2(F)\), the theorem gives
\[
 \sum_S r_{t,n}(S)^2
   =|\mathcal F_t(n)|^2\exp(o(n)),
 \qquad
 |\mathcal F_t(n)+\mathcal F_t(n)|
   \ge |\mathcal F_t(n)|^2\exp(-o(n)).
\]
Indeed, the sum \(tI\), arising from the inverse pairs
\((A,A^{-1})\), contributes \(|\mathcal F_t(n)|^2\) to the energy.  Every
other multiplicity is \(\exp(o(n))\), and the total number of
ordered pairs is \(|\mathcal F_t(n)|^2\); Cauchy--Schwarz gives the sumset
bound.  The same statements hold after applying an injective
trace-coordinate map.  Choose
\(B_1,B_2,B_3\in\Delta\) whose traceless parts form a basis and put
\(z(A)=(\Tr(AB_i))_{i=1}^3\).  On \(\Tr A=t\), equality of two sums of
vectors \(z(A)\) is equivalent, by nondegeneracy of the trace pairing, to
equality of the corresponding matrix sums.  The unique coordinate-vector
sum of large
multiplicity is
\(t(\Tr B_1,\Tr B_2,\Tr B_3)\); when the \(B_i\) are conjugate and have
trace \(t\), it is
\(t^2(1,1,1)\), corresponding to the zero matrix
\(A+A^{-1}-tI\).

\subsection{Balanced torus orbits in arithmetic word balls}

We begin with a counting theorem which is independent of hyperbolic surfaces.
It explains the arithmetic content of these fixed trace-coordinate maps in
any rank.

For an algebraic number \(\xi\), write \(H(\xi)\) for its absolute
multiplicative Weil height.  If
\(T=\Gm^r\) acts diagonally on \(K^m\), its weights are the vectors
\(w_i\in\Z^r\) determined by
\[
 (t_1,\ldots,t_r)\cdot(x_1,\ldots,x_m)
 =\bigl(t^{w_1}x_1,\ldots,t^{w_m}x_m\bigr).
\]
The weights are called \emph{balanced} if they span \(\Q^r\) and \(0\) lies
in the interior of their convex hull.  The second condition means that every
direction in the torus expands at least one active coordinate and contracts
at least one other; it is the elementary diagonal-torus form of GIT
stability.

\begin{theorem}[uniform counting in balanced algebraic torus orbits]
\label{thm:balanced-torus}
Let \(K\) be a number field, let \(S\) be a finite set of places containing
the archimedean places, and let \(T=\Gm^r\) act diagonally on \(K^m\) with
balanced weights \(w_1,\ldots,w_m\).  Fix \(C>0\).  For
\(x\in(\cO_{K,S}\setminus\{0\})^m\), put
\[
 \mathcal O_x(K)=\bigl(T(\overline K)\cdot x\bigr)\cap K^m.
\]
Uniformly for \(n\ge2\) and \(H(x_i)\le e^{Cn}\),
\[
 \#\left\{y\in\mathcal O_x(K)\cap\cO_{K,S}^m:
       H(y_i)\le e^{Cn}\ \text{for every }i\right\}
 \le
 \exp\!\left(O_{K,S,C,w_i}\!\left(\frac n{\log n}\right)\right).
\tag{12.8}
\]
The same conclusion holds when some coordinates vanish, provided the weights
on their common nonzero support remain balanced.
\end{theorem}

\begin{proof}
First suppose \(y=t\cdot x\) for \(t=(t_1,\ldots,t_r)\in T(K)\).
At a nonarchimedean place \(\mathfrak p\notin S\), put
\[
 a_{i,\mathfrak p}=v_{\mathfrak p}(x_i),\qquad
 \nu_{\mathfrak p}
 =\bigl(v_{\mathfrak p}(t_1),\ldots,v_{\mathfrak p}(t_r)\bigr).
\]
Both \(x\) and \(y\) are \(S\)-integral.  Therefore
\[
 a_{i,\mathfrak p}\ge0,\qquad
 a_{i,\mathfrak p}
   +\langle w_i,\nu_{\mathfrak p}\rangle\ge0
 \quad(1\le i\le m).
\tag{12.9}
\]

We spell out where balancedness enters.  Because \(0\) is an interior point
of the convex hull, the continuous function
\[
 u\longmapsto\max_i\bigl(-\langle w_i,u\rangle\bigr)
\]
is strictly positive on the unit sphere in \(\R^r\).  Its minimum there is
some \(c_w>0\).  Hence, for every \(u\),
\[
 \max_i\bigl(-\langle w_i,u\rangle\bigr)\ge c_w\lVert u\rVert.
\tag{12.10}
\]
Let \(A_{\mathfrak p}=\sum_i a_{i,\mathfrak p}\).  Inequalities
(12.9)--(12.10) give
\(\lVert\nu_{\mathfrak p}\rVert=O_w(A_{\mathfrak p})\).
Thus \(\nu_{\mathfrak p}\) has at most
\(O_w((1+A_{\mathfrak p})^r)\) possible integral values, and it is zero when
\(A_{\mathfrak p}=0\).

Set
\[
 \mathfrak a=\prod_{\mathfrak p\notin S}
       \mathfrak p^{A_{\mathfrak p}}.
\]
The height bounds imply \(\log N\mathfrak a=O_{K,m,C}(n)\): the left side
is the sum of the outside-\(S\) positive valuations of the coordinates of
\(x\), weighted by \(\log N\mathfrak p\), and this is one of the terms
controlled by their Weil heights.  Consequently the number of possible
collections \((\nu_{\mathfrak p})_{\mathfrak p\notin S}\) is at most
\[
 \prod_{\mathfrak p\mid\mathfrak a}
       O_w((1+A_{\mathfrak p})^r)
 \le \tau_K(\mathfrak a)^{O_w(1)}
 \le\exp\!\left(O\!\left(\frac n{\log n}\right)\right).
\tag{12.11}
\]
Here \(\tau_K(\mathfrak a)\) is the number of ideal divisors of
\(\mathfrak a\); the last inequality is the standard uniform divisor bound
in a fixed number field.

Fix one valuation pattern.  The quotient of any two parameters which realize
it is an \(S\)-unit in each torus coordinate.  Choose \(r\) linearly
independent weights and let \(D\ne0\) be the determinant of their weight
matrix.  The adjugate-matrix identity expresses each \(t_j^D\) as a fixed
Laurent monomial in the ratios \(y_i/x_i\).  Since
\(H(y_i/x_i)\le H(y_i)H(x_i)\le e^{2Cn}\), it follows that
\(H(t_j)\le e^{C'n}\).  Dirichlet's \(S\)-unit theorem identifies the
free part of \(\cO_{K,S}^{\times}\) with a lattice under the logarithmic
embedding.  A box of radius \(O(n)\) contains only \(n^{O_{K,S,r}(1)}\)
lattice points.  Multiplying this polynomial factor by (12.11) proves
(12.8) when \(t\in T(K)\).

Finally suppose only that \(y\in T(\overline K)x\cap K^m\).  The same
adjugate identity shows that every \(t_j^D\) lies in \(K\).  Run the valuation
argument on these \(D\)-th powers, or equivalently on the lattice
\(D^{-1}\Z^r\).  A fixed tuple of powers has at most \(D^r\) preimages acting
on \(x\), so only the implied constant changes.  Zero coordinates are
omitted because a torus orbit preserves its coordinate support.
\end{proof}

Here is the invariant-theoretic form.  If a torus \(T\) acts on an affine
variety \(Y\), the affine quotient
\(Y\sslash T=\operatorname{Spec}K[Y]^T\) records the values of all invariant
regular functions.  Each quotient fiber contains a unique closed orbit.  For
a diagonal action, a point has closed orbit and finite stabilizer exactly
when the weights on its nonzero coordinates are balanced.

\begin{corollary}[stable quotient fibers in word balls]
\label{cor:stable-word-fibers}
Let \(G/K\) be an affine algebraic group, let a \(K\)-torus \(T\) act on
\(G\) by conjugation, and let \(\pi_T:G\to G\sslash T\) be the affine
quotient.  Let \(\Delta<G(K)\) be finitely generated.  On the locus
\(G^{\mathrm{bal}}\) where the active weights in a fixed faithful equivariant
affine embedding are balanced, one has, uniformly in
\(\xi\in(G\sslash T)(K)\),
\[
 \#\{\gamma\in\Delta:|\gamma|\le n,\
       \gamma\in G^{\mathrm{bal}},\ \pi_T(\gamma)=\xi\}
 \le\exp\!\left(O_{G,T,\Delta}\!\left(\frac n{\log n}\right)\right).
\tag{12.12}
\]
\end{corollary}

\begin{proof}
Pass once to a finite extension which splits \(T\), and choose \(S\) so that
the fixed generators of \(\Delta\) and the affine embedding are
\(S\)-integral.  Multiplying at most \(n\) fixed generator matrices makes
every coordinate height at most \(e^{O(n)}\).  Balanced points have closed
torus orbit and finite stabilizer by the diagonal Hilbert--Mumford criterion.
Two such points with the same quotient value therefore lie in the unique
closed orbit in that fiber.  Theorem~\ref{thm:balanced-torus} applies to this
orbit.  There are only finitely many coordinate supports, so summing over
them preserves the bound.
\end{proof}

For \(\SL_d\), the quotient functions can be written explicitly as traces
of alternating products with a fixed regular semisimple matrix.

\begin{corollary}[fibers of centralizer invariants in \(\SL_d\)]
\label{cor:sld-centralizer-fibers}
Let \(\Delta<\SL_d(K)\) be finitely generated, and let
\(B\in\Delta\) be split regular semisimple.  Write
\(P_1(B),\ldots,P_d(B)\) for its spectral projectors.  For every simple
directed cycle \(I=(i_1,\ldots,i_k)\), including loops, put
\[
 c_I(A)=\Tr\bigl(P_{i_1}A P_{i_2}A\cdots P_{i_k}A\bigr).
\tag{12.13}
\]
Let \(\Psi_B\) be the finite vector of all these functions, and let \(U_B\)
be the locus where the directed graph having an edge \(i\to j\) whenever
the \((i,j)\)-entry of \(A\) is nonzero is strongly connected.  Then
\[
 \max_\xi\#\{A\in\Delta:|A|\le n,\ A\in U_B,\ \Psi_B(A)=\xi\}
 \le\exp\!\left(O_{\Delta,B}\!\left(\frac n{\log n}\right)\right).
\tag{12.14}
\]
Every coordinate of \(\Psi_B\) is a fixed linear combination of the
alternating traces
\[
 \Tr(B^{e_1}AB^{e_2}A\cdots B^{e_k}A).
\]
\end{corollary}

\begin{proof}
In a \(B\)-eigenbasis, its connected centralizer modulo the center is the
diagonal torus.  It acts on the entry \(a_{ij}\) with weight \(e_i-e_j\).
An invariant monomial corresponds to a directed multigraph having equal
indegree and outdegree at every vertex.  Such an Eulerian multigraph is a
union of directed cycles, and every directed cycle decomposes into simple
cycles.  Thus the diagonal entries and the simple-cycle monomials generate
the invariant ring.

The expression in (12.13) is exactly
\[
 a_{i_1i_2}a_{i_2i_3}\cdots a_{i_ki_1}.
\]
Strong connectivity says that the active weights span the root lattice.
Moreover every active edge lies on a directed cycle; adding these cycles
gives a strictly positive linear relation among all active weights.  This is
precisely balancedness.  Corollary~\ref{cor:stable-word-fibers} now gives
(12.14).  Finally, Lagrange interpolation writes each \(P_i(B)\) as a
polynomial in \(B\).  Expanding (12.13) therefore gives the asserted finite
linear combination of alternating traces.
\end{proof}

In rank one the entire quotient collapses to two ordinary traces.  If
\(B=\diag(\lambda,\lambda^{-1})\) and
\(A=(\begin{smallmatrix}a&b\\c&d\end{smallmatrix})\), then
\[
 K[\SL_2]^{Z(B)}
 =K[a,d,bc]/(ad-bc-1)=K[a,d]
 =K[\Tr A,\Tr(AB)].
\tag{12.15}
\]
The last equality holds because \(\lambda\ne\lambda^{-1}\), so
\((\Tr A,\Tr(AB))\) is an invertible linear change of \((a,d)\).
The balanced locus is \(bc\ne0\), equivalently
\(\langle A,B\rangle\) is irreducible.

Thus \(\Tr A\) and \(\Tr(AB)\) are literally the affine quotient
coordinates for centralizer conjugation in rank one.  Theorem
\ref{thm:finite-type-fiber} gives the stronger
\(\exp(o(n))\) word-ball bound over every finitely generated
characteristic-zero coefficient field.

The irreducibility hypothesis is essential for a general matrix group.
Indeed, take \(B=\diag(2,1/2)\), \(U(t)=
(\begin{smallmatrix}1&t\\0&1\end{smallmatrix})\), and
\(\Delta=\langle B,U(1)\rangle<\SL_2(\Q)\).  Since
\(BU(t)B^{-1}=U(4t)\), word balls contain exponentially many different
upper-unipotent elements, while all of them have
\(\Tr U(t)=2\) and \(\Tr(U(t)B)=\Tr B\).  Cocompact surface groups avoid
this horospherical phenomenon because a boundary-point stabilizer is cyclic.

\subsection{Trace multiplicity and additive energy}

Low trace growth does force large length multiplicities.  Let \(H\) be a
torsion-free finite-index subgroup of a cocompact Fuchsian lattice.  Write
\(m_t(L)\) for
the number of primitive hyperbolic \(H\)-conjugacy classes of length at most
\(L\) having absolute trace \(t\).

We keep two statistics separate.  The multiplicity second moment
\[
 M_2(L)=\sum_t m_t(L)^2
\]
counts ordered pairs of primitive classes having the same trace.  It is not
additive energy.  For a finite set \(S\subset\R\), put
\[
 r_{S+S}(u)=\#\{(a,b)\in S^2:a+b=u\},
 \qquad
 E_+(S)=\sum_u r_{S+S}(u)^2.
\]
Equivalently, \(E_+(S)\) is the number of ordered quadruples
\((a,b,c,d)\in S^4\) satisfying \(a+b=c+d\).  This second statistic forgets
all geodesic multiplicities and uses each distinct trace value only once.

\begin{proposition}[trace-multiplicity second moment]
\label{prop:collision-energy}
If, for some \(\alpha\ge0\),
\(\#(T(H)\cap[-X,X])=O(X^\alpha)\), then
\[
 \sum_t m_t(L)^2\gg \frac{e^{(2-\alpha/2)L}}{L^2},
 \qquad
 \max_t m_t(L)\gg\frac{e^{(1-\alpha/2)L}}{L}.
\tag{12.16}
\]
In particular, linear trace growth gives the exponents \(3/2\) and \(1/2\).
\end{proposition}

\begin{proof}
The prime geodesic theorem \cite{Roblin,Buser} gives
\(N_H(L)\asymp e^L/L\) primitive classes.  Since
\(|\Tr\gamma|\asymp e^{\ell(\gamma)/2}\), the number \(D_H(L)\) of available
distinct trace values is \(O(e^{\alpha L/2})\).  Cauchy--Schwarz gives
\[
 \sum_t m_t(L)^2\ge \frac{N_H(L)^2}{D_H(L)},
\]
and the bound for the largest fiber follows from
\(N_H(L)\le D_H(L)\max_t m_t(L)\).
\end{proof}

This is a second moment of multiplicities, not yet additive energy of the
distinct trace support.  Fricke does turn an equal-trace pair into an
additive quadruple:
\[
 \Tr A=\Tr A'
 \quad\Longrightarrow\quad
 \Tr(AB)+\Tr(AB^{-1})
 =\Tr(A'B)+\Tr(A'B^{-1}).
\tag{12.17}
\]
To transfer the lower bound in (12.16) to the scalar support, one must prevent
too many group-theoretic pairs from collapsing onto the same scalar
quadruple in (12.17).  The finite-type divisor theorem above gives the
required bound on these fibers.  We first record the geometric normalization
needed to apply it.

Work in the full inverse image \(\widetilde H<\SL_2(\R)\), always choosing
the positive-trace lift of a primitive class.  Fix a hyperbolic
\(B\in\widetilde H\), with \(y=\Tr B\ne0\).  Fix once and for all a compact
fundamental set and a deterministic tie-breaking rule, and conjugate each
axis to meet that set.  Let \(\mathcal P_X\) be the resulting nested family
of representatives of length at most \(2\log X+O(1)\).  Then
\(d(o,Ao)\le\ell(A)+O_H(1)\), so its elements have norm \(O_H(X)\), and
\(|\mathcal P_X|\asymp X^2/\log X\).  Define
\[
 F_B(X)=\max_{x,z}\#\{A\in\mathcal P_X:\Tr A=x,\ \Tr(AB)=z\}.
\tag{12.18}
\]

\begin{theorem}[two trace coordinates on compact Fuchsian lattices]
\label{thm:surface-fiber}
Every hyperbolic \(B\in\widetilde H\) satisfies
\[
 F_B(X)=X^{o(1)}.
\tag{12.19}
\]
No genericity assumption on \(B\) is needed.
\end{theorem}

\begin{proof}
By construction, the axis of every \(A\in\mathcal P_X\) meets one fixed
compact set and \(\ell(A)\le2\log X+O(1)\).  Hence, for a fixed basepoint
\(o\),
\[
 d(o,Ao)\le\ell(A)+O_H(1)=O_H(\log X).
\]
The orbit map of the cocompact action is a quasi-isometry
(the Milnor--Švarc lemma), so the word length of \(A\) in a fixed finite
generating set is \(O_H(\log X)\).

The field generated over \(\Q\) by the entries of a finite generating set
of \(\widetilde H\) is finitely generated.  On the irreducible locus,
Theorem~\ref{thm:finite-type-fiber}, with \(n=O_H(\log X)\), gives
\(\exp(o(\log X))=X^{o(1)}\) possibilities uniformly in the two prescribed
traces.  On the reducible
locus, \(A\) and \(B\) share a boundary fixed point.  A boundary-point
stabilizer in a torsion-free cocompact Fuchsian group is cyclic: there are no
parabolics, and two hyperbolic elements fixing the same endpoint have the
same axis and belong to one maximal cyclic subgroup.  That subgroup has only
its generator and inverse as primitive elements, so it contributes \(O(1)\)
representatives.  Combining the two loci proves (12.19).
\end{proof}

The fiber estimate for
\(A\mapsto(\Tr A,\Tr(AB))\) has the following unconditional consequence,
which does not pass through additive energy.

\begin{corollary}[the sharp universal exponent for distinct lengths]
\label{cor:universal-trace-growth}
For every torsion-free cocompact Fuchsian lattice \(H\),
\[
 D_H(X):=\#\bigl(T(H)\cap[-X,X]\bigr)\ge X^{1-o(1)}.
\]
If \(N_H^*(L)\) denotes the number of distinct primitive closed-geodesic
lengths at most \(L\), then
\[
 N_H^*(L)\ge \exp\bigl((\tfrac12-o(1))L\bigr).
\]
Thus every closed hyperbolic surface has exponentially many distinct
primitive lengths, with the universal exponent \(1/2\).  This exponent is
optimal, because compact arithmetic surfaces have at most
\(O(e^{L/2})\) distinct lengths up to \(L\).
\end{corollary}

\begin{proof}
The map
\[
 \mathcal P_X\longrightarrow T(H)^2,
 \qquad A\longmapsto\bigl(\Tr A,\Tr(AB)\bigr)
\]
has fibers of size at most \(F_B(X)=X^{o(1)}\).  Both coordinates have
absolute value at most \(C_BX\), whereas
\(|\mathcal P_X|\asymp X^2/\log X\).  Therefore
\[
 D_H(C_BX)^2\ge
 \#\{(\Tr A,\Tr(AB)):A\in\mathcal P_X\}
 \ge \frac{|\mathcal P_X|}{F_B(X)}=X^{2-o(1)},
\]
which proves the first assertion after rescaling.

Every nontrivial element of a torsion-free cocompact Fuchsian group is a
power of a unique primitive element, up to inversion.  If such a power has
absolute trace at most \(X\), its translation length is \(O(\log X)\).
The systole of \(H\) is positive, so a fixed primitive element contributes
only \(O_H(\log X)\) power traces in this range.  Hence the number of
distinct primitive absolute traces up to \(X\) is at least
\(D_H(X)/O_H(\log X)=X^{1-o(1)}\).  Finally
\(|\Tr\gamma|=2\cosh(\ell(\gamma)/2)\) converts \(X\) to
\(e^{L/2}\), proving the length statement.  The arithmetic upper bound
follows from linear trace growth.
\end{proof}

Breuillard--de Cornulier--Lubotzky--Meiri prove exponential growth of the
number of characteristic polynomials in word balls of a non-virtually-solvable
linear group \cite[Theorem~1.2]{BCLM}.  In \(\SL_2\) the characteristic
polynomial is determined by the trace; a cocompact Fuchsian lattice is
linear and not virtually solvable, so the theorem applies.  Word length is comparable with
hyperbolic displacement for a cocompact lattice.  Finally, a primitive
element of length at most \(L\) contributes only \(O_H(L)\) powers of length
at most \(L\), because the systole is positive.  Thus their theorem implies
an exponential lower bound for distinct primitive lengths, with a positive
constant that may depend on the surface.  Corollary~\ref{cor:universal-trace-growth} gives the
universal exponent \(1/2\) and, through the arithmetic examples, shows that
this universal value cannot be increased.

\subsection{Trace supports indexed by a finite quotient}

The scalar trace support can be replaced by a finite vector of supports in a
way which retains the group law.  Let \(\Delta<\SL_2(\R)\) have cocompact,
projectively torsion-free image, let
\(\pi:\Delta\twoheadrightarrow Q\) be a finite quotient, and define
\[
 T_r(X)=\{\Tr A:\pi(A)=r,\ |\Tr A|\le X\}\qquad(r\in Q).
\]
These supports may overlap as subsets of \(\R\); their labels remember the
group element, not merely its trace.  Fix \(o\in\mathbb H^2\) and put
\(\|A\|_o=\exp(d(o,\bar A o)/2)\).

\begin{theorem}[growth of \(\Tr(\gamma g\gamma^{-1}B)\) on a conjugacy orbit]
\label{thm:colored-orbit-growth}
Fix noncentral \(g,B\in\Delta\), and write
\[
 a=\pi(g),\quad b=\pi(B),\quad t=\Tr g,\quad y=\Tr B.
\]
For every \(c\in Q\), the set
\[
 \mathcal Z_{g,B,c}(X)=
 \{\Tr(\gamma g\gamma^{-1}B):
       \pi(\gamma)=c,\ d(o,\bar\gamma o)\le\log X\}
\]
satisfies
\[
 \#\mathcal Z_{g,B,c}(X)\ge X^{1-o(1)}.
\]
For a constant \(C=C(g,B,o)\), it also satisfies the two exact component
containments
\[
 \mathcal Z_{g,B,c}(X)\subset T_{cac^{-1}b}(CX),
 \qquad
 ty-\mathcal Z_{g,B,c}(X)\subset T_{cac^{-1}b^{-1}}(CX).
\]
\end{theorem}

\begin{proof}
Let \(\Delta_0=\ker\pi\).  Cocompact orbital counting in each fixed coset
gives
\[
 \#\{\gamma\in\pi^{-1}(c):d(o,\bar\gamma o)\le R\}\asymp e^R.
\]
If two such \(\gamma_i\) give the same conjugate of \(g\), then
\(\gamma_2^{-1}\gamma_1\in Z_\Delta(g)\).  This centralizer is virtually
cyclic and has only \(O(R)\) elements in a displacement ball of radius
\(2R\).  With \(R=\log X\), we therefore obtain \(\gg X/\log X\) distinct
conjugates \(C=\gamma g\gamma^{-1}\).

The triangle inequality gives \(\|C\|_o\ll_gX\), and cocompactness gives
word length \(O(\log X)\).  For prescribed
\(\Tr C=t\) and \(\Tr(CB)=z\), Theorem~\ref{thm:finite-type-fiber} bounds
the irreducible part of the fiber by \(X^{o(1)}\).  In the reducible part,
\(C\) fixes an endpoint of \(B\); the relevant boundary stabilizer is
virtually cyclic and contributes only \(O(\log X)\) elements.  Each
\(z\)-fiber is therefore \(X^{o(1)}\), proving the lower bound.

The image of \(C\) in \(Q\) is \(cac^{-1}\), which proves the first
containment.
The Fricke identity
\[
 \Tr(CB)+\Tr(CB^{-1})=\Tr C\Tr B=ty
\]
proves the reflected containment.
\end{proof}

\begin{corollary}[growth in every quotient component]
\label{cor:colored-growth}
For every \(r\in Q\), one has \(\#T_r(X)\ge X^{1-o(1)}\).
\end{corollary}

\begin{proof}
Fix any noncentral \(g\), with image \(a\in Q\), take \(c=1\), and choose a
noncentral \(B\in\pi^{-1}(a^{-1}r)\).  Such a choice is possible because
each coset of the finite-index kernel contains infinitely many elements.
The first containment in Theorem~\ref{thm:colored-orbit-growth} puts all
\(X^{1-o(1)}\) values of \(\Tr(\gamma g\gamma^{-1}B)\) in the component
indexed by \(r\).
\end{proof}

For finite \(A,B\subset\R\), write
\[
 E_+(A,B)=\#\{(a_1,b_1,a_2,b_2)\in A\times B\times A\times B:
                    a_1+b_1=a_2+b_2\}.
\]

\begin{theorem}[mixed energy between quotient components]
\label{thm:colored-mixed-energy}
Assume that \(\#(T(\Delta)\cap[-X,X])=O(X)\).  For \(q,b\in Q\), choose
a noncentral \(B\in\pi^{-1}(b)\).  Then, for a fixed \(C=C(B,o)\),
\[
 E_+\bigl(T_{qb}(CX),T_{qb^{-1}}(CX)\bigr)\ge X^{3-o(1)}.
\]
Consequently every pair consisting of one quotient component with itself has
mixed energy \(X^{3-o(1)}\).  If \(|Q|\) is odd, the same holds for every
ordered pair of quotient components.
\end{theorem}

\begin{proof}
Let
\[
 \mathcal B_q(X)=\{A\in\Delta:\pi(A)=q,\ \|A\|_o\le X\}.
\]
Finite-index cocompact orbital counting gives
\(|\mathcal B_q(X)|\asymp X^2\): the condition \(\pi(A)=q\) selects one
coset of the lattice \(\ker\pi\), while a displacement ball of radius
\(2\log X\) contains \(\asymp X^2\) orbit points.  For a trace value \(x\),
let \(m_x\)
count the matrices in \(\mathcal B_q(X)\) with trace \(x\), and let \(q_x\)
count the distinct values \(\Tr(AB)\) among them.  Cocompactness identifies
displacement and word length up to multiplicative and additive constants, so
every \(A\in\mathcal B_q(X)\) has word length \(O(\log X)\).  For a fixed
pair of values \(x,z\), Theorem~\ref{thm:finite-type-fiber} bounds the
irreducible matrices with
\(\Tr A=x,\Tr(AB)=z\) by \(X^{o(1)}\).  All reducible matrices lie in the
union of the two boundary-point stabilizers determined by \(B\), and these
contribute \(O(\log X)\) elements in the ball.  If
\(m_x\) is larger than a fixed multiple of \(\log X\), the irreducible
elements therefore give
\(q_x\ge m_xX^{-o(1)}\); if \(m_x=O(\log X)\), the trivial bound
\(q_x\ge1\) gives the same conclusion.  Linear trace growth leaves only
\(O(X)\) possible values of \(x\), because \(|\Tr A|\ll X\), so
\[
 \sum_xm_x^2\ge
 \frac{(\sum_xm_x)^2}{O(X)}\gg X^3.
\]

The element \(B\) is noncentral in a projectively torsion-free cocompact
group, hence hyperbolic, so \(y=\Tr B\ne0\).  Every distinct
value \(z=\Tr(AB)\) gives the ordered representation
\[
 z+\Tr(AB^{-1})=xy.
\]
The two summands belong to the components \(qb\) and \(qb^{-1}\),
respectively.
Distinct \(z\)'s give distinct representations, and distinct \(x\)'s give
distinct sums.  Therefore
\[
 E_+\bigl(T_{qb}(CX),T_{qb^{-1}}(CX)\bigr)
 \ge\sum_xq_x^2\ge X^{-o(1)}\sum_xm_x^2=X^{3-o(1)}.
\]
The same averaging shows that for some exact trace \(x=x(X)\), the set of
values \(\Tr(AB)\) has size \(X^{1-o(1)}\).

For desired components \(r,s\), the displayed pair equals \((r,s)\) exactly
when \(b^2=s^{-1}r\), taking \(q=rb^{-1}\).  The choice \(b=1\) gives
every diagonal pair.  In a finite group of odd order the squaring map is
bijective, so every pair occurs.
\end{proof}

It is therefore useful to retain the entire family of trace sets indexed by
\(Q\), rather than to replace it by their union.  Fricke's identity couples
specified components, and, when \(Q\) is abelian, the finite Fourier
transform diagonalizes translation in the index \(Q\).  This observation
does not invoke modular forms.  The estimates above have not been shown to
be compatible through a tower of quotients associated with a
nondistinguished finite place.  The arithmeticity criterion itself needs
only one component: if the entire set
\[
 \{g+\gamma g\gamma^{-1}-(\Tr g)I:
       \gamma\text{ lies in a fixed coset of }\ker\pi\}
\]
is bounded at all finite places, translation by one coset representative
gives a bounded adjoint orbit for the finite-index subgroup \(\ker\pi\).
Theorem~\ref{thm:one-orbit} makes that subgroup, and hence \(\Delta\),
arithmetic.

\begin{proposition}[a lower bound for additive energy]
\label{prop:energy-extraction}
Let \(S_X=T(H)\cap[-C_BX,C_BX]\), with \(C_B\) large enough.  Then
\[
 E_+(S_X)\gg
 \frac{X^4}{F_B(X)^2|S_X|(\log X)^2}.
\tag{12.20}
\]
In particular, Theorem~\ref{thm:surface-fiber} gives the unconditional
relation \(E_+(S_X)|S_X|\ge X^{4-o(1)}\).
\end{proposition}

\begin{proof}
For a trace value \(x\), let \(m_x\) count its representatives in
\(\mathcal P_X\), and let \(q_x\) count the distinct values
\(z=\Tr(AB)\) among them.  Then \(q_x\ge m_x/F_B(X)\).  Fricke gives, for
each such value,
\[
 z+\Tr(AB^{-1})=xy.
\]
Both summands belong to \(S_X\), after enlarging \(C_B\).  Distinct \(z\)'s
give distinct ordered representations of \(xy\); distinct \(x\)'s give
distinct sums because \(y\ne0\).  Hence
\[
 E_+(S_X)\ge\sum_xq_x^2
 \ge F_B(X)^{-2}\sum_xm_x^2.
\]
There are at most \(|S_X|\) values of \(x\), while
\(\sum_xm_x=|\mathcal P_X|\asymp X^2/\log X\).  Cauchy--Schwarz therefore
gives
\[
 \sum_xm_x^2\ge\frac{|\mathcal P_X|^2}{|S_X|},
\]
which proves (12.20).  Insert (12.19) for the final assertion.
\end{proof}

\begin{corollary}[additive energy under linear trace growth]
\label{cor:unconditional-energy}
For every torsion-free cocompact Fuchsian lattice, linear trace growth
implies
\[
 E_+(S_X)=X^{3-o(1)}.
\tag{12.21}
\]
Consequently \(|S_X|=X^{1-o(1)}\), and \(S_X\) contains subsets of
cardinality \(X^{1-o(1)}\) and doubling \(X^{o(1)}\).
\end{corollary}

\begin{proof}
Linear growth gives \(|S_X|=O(X)\).  Insert (12.19) into (12.20); the
logarithmic factors are \(X^{o(1)}\), so the lower bound is
\(X^{3-o(1)}\).  Since
\(E_+(S_X)\le|S_X|^3\), the first conclusion follows.  Together with the
assumed upper bound \(|S_X|=O(X)\), this says that the energy is within a
factor \(X^{o(1)}\) of its maximum.  The energy form of the
Balog--Szemer\'edi--Gowers theorem then produces the asserted subset.  Here
the doubling constant of a finite set \(A\) is \(|A+A|/|A|\).
\end{proof}

The scalar constraints in (12.12) have a simple geometric description.

\begin{lemma}[centralizer description of fixed-multiplier fibers]
For hyperbolic \(B\), the pair
\((\Tr A,\Tr(AB))\) determines \(A\in\SL_2(\R)\), away from the reducible
locus, up to at most two orbits under conjugation by \(Z_{\SL_2(\R)}(B)\).
\end{lemma}

\begin{proof}
Conjugate \(B\) to \(\diag(\lambda,\lambda^{-1})\), and write
\(A=\left(\begin{smallmatrix}a&b\\c&d\end{smallmatrix}\right)\).  The two
specified traces determine \(a,d\), while determinant one determines
\(bc=ad-1\).  Centralizer conjugation sends
\((b,c)\) to \((s^2b,s^{-2}c)\).  For \(bc\ne0\), only the possible real sign
of \(b\) remains, giving at most two orbits.  The excluded locus \(bc=0\) is
exactly where \(A\) and \(B\) share an eigenline.
\end{proof}

The invariant ring explains why commutator identities do not improve this
lemma.  Over an algebraically closed field of characteristic zero, for a
regular semisimple \(B\),
\[
 k[\SL_2]^{Z(B)}=k[\Tr A,\Tr(AB)].
\]
Indeed, after diagonalizing \(B\), the centralizer acts on the entries of
\(A=(\begin{smallmatrix}a&b\\c&d\end{smallmatrix})\) with weights
\(0,2,-2,0\).  Its invariant ring is
\(k[a,d,bc]/(ad-bc-1)=k[a,d]\), and \(a,d\) are recovered linearly from
\(\Tr A,\Tr(AB)\).  In particular, the trace of every word in \(A,B\),
including every iterated commutator, is constant on the same
centralizer orbit.  Thus no collection of two-generator Fricke identities
can bound the fiber in (12.12); one must add a trace against a third fixed
matrix whose trace functional is linearly independent of the first two.

\begin{proposition}[finite reconstruction from three trace coordinates]
\label{prop:three-trace-reconstruction}
Let \(k\) have characteristic different from two.  Suppose
\(B,C\in\SL_2(k)\) and the annihilator of
\(\Span_k\{I,B,C\}\) under the trace pairing is \(kD\), where
\(\det D\ne0\).  Then
\[
 \Psi_{B,C}:\SL_2\longrightarrow\mathbb A^3,
 \qquad A\longmapsto(\Tr A,\Tr(AB),\Tr(AC))
\]
is finite flat of degree two.  Every geometric fiber has length two and at
most two points.  Its deck involution is the determinant-polar reflection
\[
 \iota_D(A)=A-\frac{\langle A,D\rangle_{\det}}{\det D}D,
 \qquad
 \langle A,D\rangle_{\det}=\Tr(\operatorname{adj}(A)D).
\]
The ramification divisor is
\(\langle A,D\rangle_{\det}=0\); the branch divisor is the discriminant of
the resulting quadratic equation.

If \(\widetilde H<\SL_2(\R)\) is a Fuchsian lattice, one can moreover take
\(C=hBh^{-1}\) for two hyperbolic elements of the same trace in
\(\widetilde H\).
\end{proposition}

\begin{proof}
The three trace functionals have common kernel \(kD\).  Choose a linear
section \(A_0(y)\) of them.  Every point over \(y\in\mathbb A^3\) has the form
\(A_0(y)+sD\), and
\[
 \det(A_0+sD)-1
 =\det D\,s^2+\langle A_0,D\rangle_{\det}s+\det A_0-1.
\]
The leading coefficient is a nonzero constant.  After division by it this is
a monic quadratic in \(s\), so the coordinate ring of \(\SL_2\) is free of
rank two over \(k[y_1,y_2,y_3]\).  The second root above a point \(A\) is
the displayed \(\iota_D(A)\).  The polarization identity shows that
\(\iota_D^2=1\) and preserves the determinant.  Its fixed locus and the
quadratic discriminant give the asserted ramification and branch divisors.

For existence, diagonalize a hyperbolic
\(B=\diag(\lambda,\lambda^{-1})\).  Choose a conjugate
\(C=(\begin{smallmatrix}p&q\\r&s\end{smallmatrix})\) sharing neither
eigenline with \(B\), so \(qr\ne0\).  Such a conjugate exists by
Zariski density (equivalently, choose \(h\) outside the two endpoint
stabilizers).  Direct calculation gives
\[
 \Span\{I,B,C\}^{\perp}
 =k\begin{pmatrix}0&q\\-r&0\end{pmatrix},
 \qquad \det D=qr\ne0.
\]
Conjugation preserves trace, so \(\Tr C=\Tr B\).
\end{proof}

This is a concrete Noether normalization by the three displayed trace
coordinates.
It has a useful higher-rank analogue.  If \(\Gamma<\SL_d(k)\) is Zariski
dense in characteristic zero, there are \(d^2-1\) elements
\(B_j\in\Gamma\) whose trace-pairing annihilator is an invertible line
\(kD\).  Consequently
\[
 A\longmapsto(\Tr(AB_1),\ldots,\Tr(AB_{d^2-1}))
\]
is finite flat of degree \(d\): on the remaining line its equation is the
monic degree-\(d\) polynomial \(\det(A_0+sD)-1\).  To see that suitable
matrices \(B_j\) exist, take \(d^2-1\) conjugates in \(\SL_d\) of one noncentral
trace-zero element which span \(\mathfrak{sl}_d\); the annihilator is then
\(kI\).  The condition that their trace-pairing annihilator be an invertible
line is nonempty and Zariski open, so it meets
\(\Gamma^{d^2-1}\).  For a general embedded reductive group, ordinary
Noether normalization gives the same principle with sufficiently many
matrix coefficients, but the degree and the possibility of choosing all
fixed matrices from the group depend on the representation.

The degree-two result also gives an unconditional energy statement involving
three trace coordinates.  Recall
\(D_H(X)=\#(T(H)\cap[-X,X])\).  For a sufficiently large fixed \(K\), let
\(S_X=T(H)\cap[-KX,KX]\) and
\[
 r_X(s)=\#\{(u,v)\in S_X^2:u+v=s\}.
\]

\begin{corollary}[a fourth moment along multiples of one trace]
\label{cor:trace-coordinate-fourth-moment}
There are a nonzero trace \(y\) and \(K,c>0\), depending only on \(H\), such
that
\[
 \sum_t r_X(ty)^2\ge cX^2,
 \qquad
 \sum_t r_X(ty)^4\ge c\frac{X^4}{D_H(KX)}.
\]
Thus if \(D_H(X)=O(X^\alpha)\), the second lower bound is
\(\gg X^{4-\alpha}\); under linear growth it is \(\gg X^3\).  In addition,
\(D_H(X)\gg X^{2/3}\).  At the extremal exponent, if
\(D_H(X)=X^{2/3+o(1)}\), the sets \(S_X\) contain subsets of size
\(|S_X|^{1-o(1)}\) with doubling \(|S_X|^{o(1)}\).
\end{corollary}

\begin{proof}
Use the conjugate matrices \(B,C\) from
Proposition~\ref{prop:three-trace-reconstruction}, and write
\(y=\Tr B=\Tr C\).  Let
\(\mathcal G_X=\{A\in\widetilde H:\lVert A\rVert\le X\}\).  Cocompact
lattice-point counting gives \(\#\mathcal G_X\asymp X^2\).  If \(n_t\)
counts its elements of trace \(t\), finite reconstruction says that at least
\(n_t/2\) distinct pairs
\[
 (\Tr(AB),\Tr(AC))
\]
occur.  Each coordinate is the first term of a Fricke decomposition of the
same number \(ty\):
\[
 \Tr(AB)+\Tr(AB^{-1})=ty,
 \qquad
 \Tr(AC)+\Tr(AC^{-1})=ty.
\]
All four terms lie in \(S_X\) after fixing \(K\).  Hence
\(r_X(ty)^2\ge n_t/2\).  Summing gives the first claim, while squaring and
using Cauchy--Schwarz gives
\[
 \sum_t r_X(ty)^4\ge\frac14\sum_t n_t^2
 \ge c\frac{X^4}{D_H(KX)}.
\]
Finally, the triple
\((\Tr A,\Tr(AB),\Tr(AC))\in S_X^3\) has fibers of size at most two, so
\(X^2\ll D_H(KX)^3\), proving the last assertion after rescaling \(X\).
If \(D_H(X)=X^{2/3+o(1)}\), the first estimate says
\[
 E_+(S_X)\ge\sum_t r_X(ty)^2\gg X^2=|S_X|^{3-o(1)}.
\]
The energy form of Balog--Szemer\'edi--Gowers gives the final assertion.
\end{proof}

The fourth moment is the strongest conclusion available from finite
reconstruction \emph{alone}.  It does not by itself yield cubic ordinary
additive energy: it only forces
\(E_+(S_X)\ge\sum_t r_X(ty)^2\gg X^2\), which is diagonal order under linear
growth.  Combinatorially, a fiber of \(n_t\) reconstructed matrices could map
to a balanced \(\sqrt{n_t}\)-by-\(\sqrt{n_t}\) grid of the two values
\(\Tr(AB)\) and \(\Tr(AC)\).
Theorem~\ref{thm:finite-type-fiber} rules out this grid after either one of
these coordinates is fixed.  Together with
Theorem~\ref{thm:surface-fiber}, this proves the required incidence
multiplicity estimate and hence Corollary~\ref{cor:unconditional-energy} for
compact Fuchsian lattices.

The next proposition limits what can be deduced from a large exact trace
fiber.  After one independent fixed trace functional is added, the original
fiber splits into almost as many new values as its cardinality permits.

\begin{proposition}[dispersion of a large trace fiber under a fixed multiplier]
\label{prop:fixed-multiplier-dispersion}
Let \(\rho\colon\Pi\to\PSL_2(\R)\) be a faithful cocompact Fuchsian
representation whose absolute trace support has linear growth.  Let
\(\phi\colon\Pi\twoheadrightarrow\Z\), let \(H<\ker\phi\) be finitely
generated and convex cocompact with critical exponent \(\delta_H>1/2\), and
choose \(a\in\phi^{-1}(1)\).  Fix compatible lifts over a finitely generated
characteristic-zero field.  If \(B\in\SL_2(F)\) satisfies
\(\Tr^2B\ne4\), then for every sufficiently large \(R\) there are a
squared-trace value \(t_R\) and a set
\[
 E_R\subset\{h\in H:d(o,\rho(h)o)\le R\}
\]
such that
\[
 |E_R|\ge
 \exp\bigl((\delta_H-\tfrac12-o(1))R\bigr),
 \qquad
 \Tr^2\rho(ah)=t_R\quad(h\in E_R),
\]
whereas, uniformly in \(s\in F\),
\[
 \#\{h\in E_R:\Tr^2(B\rho(a)\rho(h))=s\}=\exp(o(R)).
\]
Consequently the second squared trace assumes at least
\(\exp((\delta_H-\tfrac12-o(1))R)\) distinct values on \(E_R\).
The same lower bound holds for a vector defined by any fixed finite list of
multipliers if at least one multiplier has squared trace different from four.
\end{proposition}

\begin{proof}
Convex-cocompact orbit counting \cite{Roblin} gives
\(\exp((\delta_H+o(1))R)\) elements in the displayed ball.  Every \(ah\)
has translation length in the original Fuchsian representation at most
\(R+O(1)\), and hence absolute trace at most
\(C e^{R/2}\).  Linear trace growth leaves only \(O(e^{R/2})\)
squared-trace values.  Pigeonholing gives \(t_R\) and \(E_R\) with the
asserted cardinality.

Put \(A_0=\rho(a)\) and \(A_1=BA_0\).  Then
\(A_1A_0^{-1}=B\).  The simultaneous-fiber assertion of
Theorem~\ref{thm:two-trace-evaluations}, applied directly to the faithful
convex-cocompact Fuchsian representation \(\rho|_H\), gives
\[
 \sup_{u,v}\#\{h:d(o,\rho(h)o)\le R,
   \ \Tr^2(A_0\rho(h))=u,\ \Tr^2(A_1\rho(h))=v\}
 =\exp(o(R)).
\]
Restricting to \(u=t_R\) proves the uniform fiber bound, and division gives
the asserted number of second-coordinate values.  Adding further coordinates
can only decrease a joint fiber.
\end{proof}

Thus neither the Fricke identities nor conjugation by a fixed independent
multiplier can turn the large first-coordinate fiber into a large
simultaneous fiber.  The exceptional condition \(\Tr^2B=4\) is
exactly where the two-coordinate divisor argument degenerates; central
multipliers merely repeat the first coordinate, while noncentral unipotent
multipliers require a separate argument.  A propagation theorem must
therefore use multipliers that vary with the radius or information from a
nonarchimedean valuation.

Average additive energy does not control all the pairs needed for a trace gap.
Put \(H=\Gamma^{(2)}\), and, for \(\eta>0\), consider the set of all small
differences
\[
 \mathcal C_\eta(H)=
 \{x-y:x,y\in T(H),\ |x-y|<\eta\}.
\]
The gap problem asks whether zero is isolated in this set.  A family
containing a negligible proportion of all pairs can nevertheless contain
every sequence of differences tending to zero.  Cardinality, additive
energy, and Balog--Szemer\'edi--Gowers extraction do not detect such a
family.  Any argument of this kind must therefore control every
sufficiently close pair, rather than almost all pairs.

\begin{problem}[control of every sufficiently small trace difference]
\label{prob:propagation}
Assume that \(\Gamma\) is a cocompact Fuchsian lattice and that
\(T(\Gamma)\) has bounded clustering or linear growth.  Prove, using the
group-level Fricke relations and the fiber theorems for
\((\Tr A,\Tr(AB))\), one of the
following sufficient conclusions for \(H=\Gamma^{(2)}\):
\begin{enumerate}[label=\textup{(\alph*)}]
\item there are \(\eta>0\) and a finite set \(F\subset\R\) such that
\[
 \mathcal C_\eta(H)\subset F+T(H);
\]
\item there are \(\eta>0\), a number field \(k\) with distinguished real
embedding \(v_0\), a fractional ideal \(I\subset k\), and constants
\(C_\sigma\) such that, under \(v_0\),
\[
 \mathcal C_\eta(H)\subset I,
 \qquad |\sigma(d)|\le C_\sigma
 \quad(d\in\mathcal C_\eta(H))
\]
for every nondistinguished archimedean embedding \(\sigma\).
\end{enumerate}
\end{problem}

Conclusion (a) gives a gap by Lemma~\ref{lem:additive-cover-gap}; conclusion
(b) gives one directly from the product formula, by applying the proof of
Lemma~\ref{lem:galois-window} to \(\mathcal C_\eta(H)\).  The finite-type
fiber theorem makes inverse additive theory available, but
Theorem~\ref{thm:same-trace-sum-fiber} shows that, after applying the
injective coordinate map \(A\mapsto(\Tr(AB_i))_{i=1}^3\), an exact trace
fiber in a word ball of radius \(n\) has representation multiplicity
\(\exp(o(n))\) for every sum other than the forced inverse sum.  Thus
ordinary Freiman extraction inside that fiber cannot yield either conclusion, both of which
must hold for every sufficiently small difference.  Problem
\ref{prob:propagation} remains one sufficient scalar approach.

A more direct sufficient condition is furnished by
Theorem~\ref{thm:asymptotic-orbit} and
Lemma~\ref{lem:projective-escape}.  The following theorem quantifies this
expansion in terms of traces.
Theorem~\ref{thm:moving-divisorial-multiplier} attaches to a loxodromic
kernel element \(k\) an exact cancellation exponent
\(\mathfrak c_v(a,k;J)\) and proves
\[
 h_{\Tr}(x)\ge
 \frac{\delta_J}
 {1+\ell_x(k)\mathfrak c_v(a,k;J)/\ell_v(k)}.
\]
Lemma~\ref{lem:loxodromic-in-kernel} shows that an admissible \(k\) exists
at every unbounded divisor.  Thus, when \(\delta_J>1/2\), the remaining
finite-place implication is the concrete inequality
\[
 \frac{\ell_v(k)}{\ell_x(k)}>
 \frac{\mathfrak c_v(a,k;J)}{2\delta_J-1}
\]
for at least one choice of \(J,a,k\), or an argument treating the reverse
inequality.  This statement applies over the
function field of a variety on which all trace collisions at the original
Fuchsian character persist.  It allows the most frequent trace value to
vary with the radius and does not presuppose a number-field form.

\subsection{Volume of a single trace slab}

There is a geometric reason the direct quantitative double-coset approach
stalls in the compact case.

\begin{proposition}[trace-slab volume]\label{prop:slab}
Fix a nonzero \(A\in M_2(\R)\), and let \(\|\cdot\|\) be any matrix norm.  With
respect to Haar measure on \(\SL_2(\R)\), as \(X\to\infty\),
\[
 \operatorname{vol}\{g:\|g\|\le X,\ |\Tr(gA)|\le1\}
 \asymp_A
 \begin{cases}
 X,&\det A\ne0,\\
 X\log X,&\operatorname{rank}A=1.
 \end{cases}
\tag{12.22}
\]
\end{proposition}

\begin{proof}
Use \(KAK\) coordinates
\(g=k_1\diag(r,r^{-1})k_2\), where \(r=e^{t/2}\).  Haar density is comparable
to \(\sinh(t)\,dt\,dk_1\,dk_2=(r-r^{-3})\,dr\,dk_1\,dk_2\), and
\(\|g\|\asymp r\).  Fixing one angular variable and writing
\(C=Ak_1=(c_{ij})\), the other angular variable makes \(\Tr(gA)\) a sinusoid
of amplitude \(D_r\), where
\[
 D_r^2=r^2\|C e_1\|^2+2\det A+r^{-2}\|C e_2\|^2.
\]
Its angular sublevel proportion is exactly
\[
 \frac2\pi\arcsin\!\left(\min\{1,D_r^{-1}\}\right).
\]
If \(A\) is invertible, \(\|Ak_1e_1\|\) is bounded uniformly away from zero,
and averaging gives a proportion comparable to \(r^{-1}\).  If \(A\) has
rank one, the first-column factor has simple angular zeros; splitting into
\(|\cos\theta|\le r^{-2}\) and dyadic complementary regions gives an averaged
proportion comparable to \(r^{-1}\log r\).  Integration against
\((r-r^{-3})dr\), for \(r\ll X\), gives \(X\) and \(X\log X\), respectively.
\end{proof}

For a torsion-free cocompact Fuchsian group, the square-root and secant matrices
attached to positive-trace lifts of hyperbolic elements are invertible.  This
is a continuous Haar-volume calculation, not a lattice-point estimate:
standard mixing errors for a ball of volume \(\asymp X^2\) may dominate an
\(X\)-volume slab.  Thus even an optimal equidistribution estimate for one
such slab would give only the \(O(X)\) bound compatible with bounded
clustering.  A parabolic secant is rank one and gains the logarithm in (12.22),
explaining structurally why unipotent amplification is available in the
nonuniform proofs.

\section{A Noetherian reduction for exceptional compact surfaces}

Hao's 2026 theorem proves superlinear trace growth away from a countable union
of positive-codimension trace-collision loci in Teichm\"uller space
\cite{HaoLengthSets}.  At a point in the exceptional union, Noetherianity
implies that all trace equalities holding there are consequences of finitely
many polynomial equations.

Let \(\mathcal X\) be an affine \(\Q\)-model of the relevant
\(\PSL_2\)-character component, with its open Teichm\"uller component
\(\mathcal T\subset\mathcal X(\R)\).  For a group element \(w\), its
squared-trace function \(\tau_w\in\Q[\mathcal X]\) is regular.  For
\(x\in\mathcal T\), let \(J_x\) be the ideal generated by
\[
 \tau_u-\tau_v
 \quad\text{for all pairs with }\tau_u(x)=\tau_v(x),
\tag{13.1}
\]
and put \(V_x=Z(J_x)\subset\mathcal X\).

\begin{theorem}[persistence of trace collisions]\label{thm:noetherian}
The ideal \(J_x\) is generated by finitely many trace-collision equations.
Every squared-trace equality holding at \(x\) holds identically on \(V_x\).

If \(x\) is a hyperbolic metric with trace-set entropy
\[
 h_{\Tr}(x)=\limsup_{L\to\infty}\frac1L
 \log\#\{\text{distinct primitive absolute traces of length at most }L\}
 \le\frac12,
\tag{13.2}
\]
then, for every \(\epsilon>0\), all hyperbolic metrics \(y\in V_x\) sufficiently
near \(x\) satisfy
\[
 h_{\Tr}(y)\le(1+\epsilon)h_{\Tr}(x)\le\frac{1+\epsilon}{2}.
\tag{13.3}
\]
If the complex algebraic germ of \(V_x\) at \(x\) is zero-dimensional, the
Fricke coordinates of \(x\) are algebraic.  Since \(x\) is a smooth point of the character
variety, it is enough for the gradients of finitely many equal-trace equations
to span the cotangent space at \(x\).  If \(x\) has a transcendental
coordinate, then \(V_x(\R)\cap\mathcal T\) contains a nonconstant
semialgebraic arc through \(x\).
\end{theorem}

\begin{proof}
The coordinate ring \(\Q[\mathcal X]\) is Noetherian, so a finite subfamily
of (13.1) generates \(J_x\).  Every other collision difference belongs to
that ideal and vanishes on \(V_x\).

Hyperbolic lengths for two sufficiently close metrics satisfy, uniformly over
all closed curves,
\[
 \ell_x(\gamma)\le(1+\epsilon)\ell_y(\gamma)
\]
after shrinking the neighborhood.  Every collision at \(x\) persists at
\(y\), so the partition of group elements by trace at \(y\) is coarser than
the partition at \(x\).  Therefore the number of \(y\)-trace values represented
by curves of \(y\)-length at most \(L\) is at most the number of \(x\)-trace
values represented by curves of \(x\)-length at most \((1+\epsilon)L\).
Taking logarithmic upper growth rates proves (13.3).

A zero-dimensional complex germ through \(x\) places \(x\) on a
zero-dimensional irreducible component of the \(\Q\)-scheme \(V_x\); its
residue field is a finite extension of \(\Q\).  The Jacobian criterion at
the smooth Teichm\"uller point gives the stated sufficient condition.
Finally, every isolated real point of a
\(\Q\)-semialgebraic set has algebraic coordinates.  A transcendental \(x\)
is therefore not isolated in \(V_x(\R)\); real curve selection gives a
nonconstant semialgebraic arc, and a short segment remains in the open set
\(\mathcal T\).
\end{proof}

\begin{proposition}[bounded-degree algebraic specializations converging to a transcendental character]
\label{prop:bounded-degree-specialization}
For \(x\in\mathcal T\), let
\[
 I(x)=\{f\in\Q[\mathcal X]:f(x)=0\},\qquad
 Z_x=V(I(x)).
\]
We call \(Z_x\) the \(\Q\)-Zariski closure, or rational closure, of \(x\).
If \(\dim Z_x>0\), there are an integer \(D\) and distinct algebraic
Fuchsian points \(y_j\in Z_x(\R)\), all in the Teichm\"uller component of
\(x\), such that
\[
 y_j\longrightarrow x,\qquad [\Q(y_j):\Q]\le D.
\tag{13.3a}
\]
For every finite \(E\subset\pi_1(S)\),
\[
 \#\{\tau_\gamma(y_j):\gamma\in E\}
 \le
 \#\{\tau_\gamma(x):\gamma\in E\}.
\tag{13.3b}
\]
In particular, if the absolute trace set at \(x\) has linear growth and
\[
 \mathcal B_R^x
 =\{\gamma:d_x(o,\gamma o)\le R\},
\]
then, uniformly in \(j\) and \(R\),
\[
 \#\{\tau_\gamma(y_j):\gamma\in\mathcal B_R^x\}
 =O_x(e^{R/2}).
\tag{13.3c}
\]
After discarding finitely many terms, the \(y_j\) may be taken
nonarithmetic.
\end{proposition}

\begin{proof}
The ideal \(I(x)\) is the kernel of evaluation in \(\R\), hence is prime
and \(Z_x\) is irreducible.  The point \(x\) lies on no proper
\(\Q\)-defined closed subvariety of \(Z_x\).  Since characteristic zero
implies generic smoothness, \(x\) is therefore a smooth point of \(Z_x\).

Put \(d=\dim Z_x\).  Choose a rational linear Noether normalization
\[
 \pi:Z_x\longrightarrow\mathbb A^d
\]
which is generically \'etale.  Its ramification locus is a proper
\(\Q\)-closed subset of \(Z_x\), so the \(\Q\)-generic point \(x\) does
not belong to it.  After shrinking the target away from the discriminant
and nonflat loci, the real inverse-function theorem gives a local
inverse branch through \(x\), and every fiber has total degree
\(\deg\pi\).  Choose distinct rational points tending to \(\pi(x)\) in
this neighborhood and lift them along the branch.  The lifts \(y_j\) are
real algebraic points of degree at most \(\deg\pi\), tend to \(x\), and
remain in the open Teichm\"uller component after the branch is restricted.

If \(\tau_\gamma(x)=\tau_{\gamma'}(x)\), then
\(\tau_\gamma-\tau_{\gamma'}\in I(x)\), so the same equality holds at
every point of \(Z_x\).  Thus the partition at \(y_j\) is coarser than
the partition at \(x\), proving (13.3b).  For
\(\gamma\in\mathcal B_R^x\),
\[
 |\Tr_x\gamma|\le2\cosh(R/2).
\]
Linear growth of the absolute trace set at \(x\), together with (13.3b),
gives (13.3c).

Finally, the Borel--Prasad theorem gives finitely many conjugacy classes of
arithmetic lattices in a fixed semisimple real group whose covolume is at
most a prescribed constant \cite{BorelPrasad}.  Applying it to
\(\PSL_2(\R)\) with that constant equal to \(4\pi(g-1)\) gives, in
particular, only finitely many conjugacy classes having the fixed covolume
\(4\pi(g-1)\).  Each unmarked surface gives one
mapping-class-group orbit in Teichm\"uller space, and proper discontinuity
makes that orbit meet a relatively compact neighborhood of \(x\) in only
finitely many points.  Hence only finitely many members of the convergent
sequence can be arithmetic.
\end{proof}

\begin{theorem}[relative deformations cannot preserve all collisions]
\label{thm:relative-deformation}
Let \(\Pi\) be torsion-free, let \(H<\Pi\) be a subgroup, and let
\[
 \rho_s\colon\Pi\longrightarrow\SL_2(\R),\qquad s\in U,
\]
be a continuous family of lifts of faithful Fuchsian representations,
where \(U\) is connected and contains \(0\).  Suppose that
\(\rho_s|_H=\rho_0|_H\) for every \(s\), and that \(\rho_0(H)\) is convex
cocompact.  For \(a\in\Pi\), put
\[
 W_a=\Span_\R\{\rho_s(a):s\in U\}\subset M_2(\R).
\]
Assume that \(\dim W_a\ge3\) and that every collision in the coset \(aH\)
at \(s=0\) persists throughout the family:
\[
 \Tr^2\rho_0(ah)=\Tr^2\rho_0(ah')
 \quad\Longrightarrow\quad
 \Tr^2\rho_s(ah)=\Tr^2\rho_s(ah')
 \quad(s\in U).
\tag{13.4}
\]
Then every fiber of
\[
 h\longmapsto\Tr^2\rho_0(ah),\qquad h\in H,
\]
has at most three elements.  If \(W_a=M_2(\R)\), every fiber has one
element.

If, in addition, every \(ah\), \(h\in H\), is primitive in \(\Pi\), then
the primitive length-set entropy of \(\rho_0(\Pi)\) is at least the critical
exponent \(\delta_H\) of \(\rho_0(H)\).
\end{theorem}

\begin{proof}
Fix \(h_0\) in one fiber and write \(X_h=\rho_0(h)\).  These matrices do
not depend on \(s\).  The traces in (13.4) never vanish: a nonidentity
element in a torsion-free Fuchsian group is hyperbolic, while the identity
has trace \(2\).  Connectedness of \(U\) therefore gives, for each \(h\) in
the fiber, a sign \(\epsilon_h\in\{\pm1\}\), independent of \(s\), such
that
\[
 \Tr\!\left(\rho_s(a)(X_h-\epsilon_hX_{h_0})\right)=0
 \qquad(s\in U).
\]
For the nondegenerate pairing \((A,B)\mapsto\Tr(AB)\), this says
\[
 X_h-\epsilon_hX_{h_0}\in W_a^\perp.
\tag{13.5}
\]
If \(\dim W_a=4\), the right side is zero.  Projective faithfulness then
gives \(h=h_0\).

Suppose \(\dim W_a=3\), and write \(W_a^\perp=\R D\) with \(D\ne0\).
For a fixed sign, every possible matrix has the form
\[
 X_h=\epsilon X_{h_0}+cD,\qquad
 \det(\epsilon X_{h_0}+cD)=1.
\tag{13.6}
\]
The determinant in (13.6) is a polynomial of degree at most two in \(c\),
and \(c=0\) is a root.  If \(\det D\ne0\), there is at most one further
root for each sign.  If \(\det D=0\), the polynomial is linear unless it is
identically zero.  In the latter case, after writing
\[
 \epsilon X_{h_0}+cD
 =\epsilon X_{h_0}(I+cN),
\]
the identity of determinants implies \(N\ne0\) and \(N^2=0\).  Two distinct values
of \(c\) arising from elements of \(\rho_0(H)\) would then have a quotient
of the form \(I+c'N\), a nontrivial unipotent.  A convex-cocompact
Fuchsian group has no such element.  Thus there is again at most one
nonzero extra root for each sign.  The two roots \(c=0\), one for each
sign, represent the same projective element.  Consequently the projective
fiber has at most three elements.

Finally,
\[
 \#\{h\in H:d(o,\rho_0(h)o)\le R\}
   =e^{(\delta_H+o(1))R}.
\]
Moreover
\(\ell_{\rho_0}(ah)\le d(o,\rho_0(a)o)+d(o,\rho_0(h)o)\).
If all \(ah\) are primitive, the uniform fiber bound just proved therefore
gives \(e^{(\delta_H+o(1))R}\) distinct primitive lengths below
\(R+O(1)\), proving the entropy assertion.
\end{proof}

For a finitely generated group \(G\), write \(\mathfrak X(G)\) for its affine
\(\PSL_2(\C)\)-character variety.  For \(g\in G\), its regular squared-trace
coordinate is
\[
 \tau_g([\sigma])=\Tr^2(\widetilde\sigma(g))
                 =\Tr(\Ad\sigma(g))+1,
\]
where either lift of \(\sigma(g)\) to \(\SL_2(\C)\) gives the same square.

\begin{theorem}[relative fiber bound from two trace evaluations]
\label{thm:two-trace-evaluations}
Let \(H\) be finitely generated, let \(F/\Q\) be a finitely generated field
of characteristic zero, and let
\(\rho\colon H\to\SL_2(F)\).  Assume that there is a faithful torsion-free
convex-cocompact Fuchsian character \(x_H\) such that
\[
 \Tr^2\rho(g)=4\text{ in }F
 \quad\Longrightarrow\quad \tau_g(x_H)=4
 \qquad(g\in H).
\]
Thus every squared-trace identity of the form
\(\Tr^2\rho(g)=4\) remains valid at \(x_H\).  This implication is the only
relation between \(\rho\) and \(x_H\) used below; no homomorphism from the
field \(F\) to \(\R\) is being assumed.
Let \(A_0,A_1\in\SL_2(F)\).  For the conclusion about a fiber of the first
trace coordinate alone, assume in addition that
\[
 \Tr^2(A_0\rho(h))=\Tr^2(A_0\rho(h'))
 \quad\Longrightarrow\quad
 \Tr^2(A_1\rho(h))=\Tr^2(A_1\rho(h')).
\tag{13.7}
\]
Put
\[
 B=A_1A_0^{-1}
 \quad\text{and assume}\quad \Tr^2B\ne4.
\tag{13.8}
\]
Then, for every word metric on \(H\), every simultaneous squared-trace
fiber satisfies
\[
 \sup_{c_0,c_1}\#
 \{h\in H:|h|\le n,\ \Tr^2(A_i\rho(h))=c_i\ (i=0,1)\}
 =\exp(o(n)).
\]
If (13.7) holds, then also
\[
 \sup_t\#\{h\in H:|h|\le n,\ \Tr^2(A_0\rho(h))=t\}
   =\exp(o(n)).
\tag{13.9}
\]
Condition (13.8) says exactly that the relative matrix
\(A_1A_0^{-1}\) has two distinct eigenvalues; no semisimplicity hypothesis
is imposed on either representation as a whole.

In particular, the theorem applies when \(\rho\) itself is a faithful
convex-cocompact Fuchsian representation.  If
\(A_i=\rho_i(a)\) are two evaluations of a deformation fixed on \(H\), all
collisions at \(\rho_0\) persist at \(\rho_1\), and every \(ah\) is
primitive, then the primitive length-set entropy at \(\rho_0\) is at least
\(\delta_H\).
\end{theorem}

\begin{proof}
Fix \(h_*\) in a nonempty simultaneous squared-trace fiber.  After dividing
the fiber into at most four parts according to the two possible signs at
each trace evaluation, the two quantities
\[
 \Tr(A_0\rho(h)),\qquad
 \Tr(A_1\rho(h))
\]
are fixed.  Put \(Y=A_0\rho(h)\).  Since
\(A_1\rho(h)=BY\), both \(\Tr Y\) and \(\Tr(BY)\) are fixed on such
a part.

After a finite extension of \(F\), condition (13.8) diagonalizes \(B\) with
two distinct eigenvalues.  Writing
\[
 Y=\begin{pmatrix}y_{11}&y_{12}\\y_{21}&y_{22}\end{pmatrix},
\]
the two fixed traces determine \(y_{11}\) and \(y_{22}\), while
\(\det Y=1\) becomes
\[
 y_{12}y_{21}=y_{11}y_{22}-1=:q.
\tag{13.10}
\]
The entries of a word of length \(n\), and the fixed linear transforms
used here, belong to one of the finite-type sets \(R_C[n]\) of
Section~12.  More explicitly, in a nonempty radius-\(n\) fiber one may use
\(h_*\) to define the two diagonal entries and \(q\); all have finite-type
complexity \(O(n)\).  The fixed diagonalizing matrix is absorbed into the
same finite extension and finite-type model.  If \(q\ne0\),
Theorem~\ref{thm:finite-type-divisor} gives
\(\exp(o(n))\) possibilities for the ordered factorization in (13.10),
uniformly in \(q\).

If \(q=0\), the solutions form two affine lines, according as
\(y_{12}=0\) or \(y_{21}=0\).  Two distinct matrices \(Y,Y'\) on one of
these lines satisfy \(\det(Y-Y')=0\).  Therefore
\[
 0=\det(Y-Y')
  =2-\Tr(YY'^{-1})
  =2-\Tr\!\left(\rho(h)\rho(h')^{-1}\right).
\]
Thus \(\Tr^2\rho(hh'^{-1})=4\) in \(F\), and the implication in the statement
gives \(\tau_{hh'^{-1}}(x_H)=4\).  Faithfulness and convex cocompactness at
\(x_H\) force \(h=h'\).  Hence the degenerate part contributes at most two
elements for each of the four sign choices.  This proves the simultaneous
fiber assertion.  Under (13.7), every fiber in (13.9) is a simultaneous
fiber, so (13.9) follows.

Word length and orbit displacement on the convex-cocompact group \(H\) are
quasi-isometric.  Dividing its
\(e^{(\delta_H+o(1))R}\) orbit points by the subexponential fiber bound,
and using the primitivity of \(ah\), proves the final assertion.
\end{proof}

The preceding proof only uses finite generation to bound the arithmetic
complexity of the matrices being counted.  We shall need the following
version for a subset of an infinitely generated subgroup.  Word length is
always measured in the fixed ambient group in this statement.

\begin{lemma}[two trace evaluations on subsets of ambient word balls]
\label{lem:ambient-two-trace-evaluations}
Let \(\Pi\) be finitely generated, let \(F/\Q\) be a finitely generated
field of characteristic zero, and let
\(\rho\colon\Pi\to\SL_2(F)\).  Suppose that \(x\) is a faithful,
torsion-free Fuchsian character of \(\Pi\) such that
\[
 \Tr^2\rho(g)=4\text{ in }F\quad\Longrightarrow\quad \tau_g(x)=4.
\]
Fix a word metric on \(\Pi\), a constant \(C\), and subsets
\[
 E_R\subset\{g\in\Pi:|g|_\Pi\le CR+C\}.
\]
If \(A_0,A_1\in\SL_2(F)\) and
\(\Tr^2(A_1A_0^{-1})\ne4\), then, uniformly in \(c_0,c_1\),
\[
 \#\{h\in E_R:\Tr^2(A_i\rho(h))=c_i\ (i=0,1)\}
   =\exp(o(R)).
\]
No finite-generation hypothesis is made on the subgroup containing the
sets \(E_R\).
\end{lemma}

\begin{proof}
Repeat the proof of Theorem~\ref{thm:two-trace-evaluations}, dividing into
the four possible sign classes.  After diagonalizing
\(A_1A_0^{-1}\), the two traces determine the diagonal entries of
\(A_0\rho(h)\), and the determinant leaves one factorization equation
\(y_{12}y_{21}=q\).  The ambient word bound puts every matrix entry in a
finite-type filtration of index \(O(R)\), so the divisor theorem gives
\(\exp(o(R))\) solutions when \(q\ne0\).  When \(q=0\), two points on one
of the two resulting affine lines have relative squared trace \(4\) in
\(F\).  The displayed specialization implication and faithfulness of
\(x\) make them the same projective matrix and then the same element of
\(\Pi\).  The proof therefore does not require a finite generating set for
the union of the \(E_R\).
\end{proof}

The preceding subexponential estimate can be sharpened geometrically in a
fixed Fuchsian representation.  The sharpening also prevents a misleading
commensurator argument: an unbounded two-trace fiber can be produced inside
every lattice by conjugating with powers of the first multiplier.

\begin{theorem}[decomposition and uniform count of internal centralizer orbits]
\label{thm:two-trace-centralizer-orbits}
Let \(\Gamma<\PSL_2(\R)\) be torsion-free and convex cocompact, let
\(q\in\Gamma\) be hyperbolic, and choose compatible lifts to
\(\SL_2(\R)\).  Write
\[
 C_\Gamma(q)=\langle z\rangle,
\]
where \(z\) is primitive.  For \(u,v\in\R\), put
\[
 \mathcal F_{u,v}(q)=
 \{g\in\Gamma:\Tr g=u,\ \Tr(qg)=v\}.
\]
Then every \(\mathcal F_{u,v}(q)\) is a finite union of
\(\langle z\rangle\)-conjugacy orbits.  Moreover, for any fixed word
metric there is \(C=C(\Gamma,q)\), independent of \(u,v\) and of the
orbit, such that
\[
 \#\bigl(\mathcal O\cap B_\Gamma(R)\bigr)\le C(R+1)
\tag{13.10a1}
\]
for every such orbit \(\mathcal O\).  Consequently, if
\(N_{u,v}(R)\) is the number of these orbits which meet \(B_\Gamma(R)\),
then
\[
 \#\bigl(\mathcal F_{u,v}(q)\cap B_\Gamma(R)\bigr)
 \le C(R+1)N_{u,v}(R).
\tag{13.10a2}
\]
The same conclusions, after multiplying \(C\) by at most four, hold when
either trace is prescribed only up to sign and hence for fibers of
\((\Tr^2g,\Tr^2(qg))\).
\end{theorem}

\begin{proof}
Conjugate so that
\[
 z=\begin{pmatrix}\theta&0\\0&\theta^{-1}\end{pmatrix},
 \qquad \theta>1,
\]
and write \(q=z^m\) projectively.  Up to a harmless sign in the chosen
lift, put \(q=\diag(\lambda,\lambda^{-1})\), \(\lambda>1\).  If
\(g=(\begin{smallmatrix}a&b\\c&d\end{smallmatrix})\), the two trace
equations are
\[
 a+d=u,\qquad \lambda a+\lambda^{-1}d=v.
\tag{13.10a3}
\]
They determine \(a,d\), and the determinant equation becomes
\[
 bc=ad-1=:\kappa.
\tag{13.10a4}
\]
If \(\kappa\ne0\), conjugation by \(z^n\) sends
\((b,c)\) to \((\theta^{2n}b,\theta^{-2n}c)\).  Every orbit therefore
has a unique representative, apart from one boundary convention, for which
\[
 1\le |b/c|<\theta^4.
\tag{13.10a5}
\]
For fixed \(u,v\), equations (13.10a3)--(13.10a5) put these representatives
in a compact subset of \(\SL_2(\R)\).  Discreteness of \(\Gamma\) makes
their number finite.  If \(\kappa=0\), a nondiagonal point is triangular
and its commutator with \(z\) is a nonidentity unipotent.  A convex-cocompact
group has no such element.  The degenerate fiber therefore contains at most
one point.

It remains to make the orbit count uniform in \(u,v\).  Positive systole
gives \(\epsilon_\Gamma>0\) such that
\[
 |\Tr\gamma|\ge2+\epsilon_\Gamma
 \qquad(1\ne\gamma\in\Gamma).
\]
Direct multiplication gives
\[
 \Tr[z,g]=2-bc(\theta-\theta^{-1})^2.
\tag{13.10a6}
\]
The commutator is nontrivial in the nondegenerate fiber, so
\[
 |bc|\ge
 \frac{\epsilon_\Gamma}{(\theta-\theta^{-1})^2}=:\kappa_0.
\tag{13.10a7}
\]
For the representative in (13.10a5), both off-diagonal entries consequently
have absolute value at least \(c_0=\theta^{-2}\sqrt{\kappa_0}\).  In a
fixed matrix norm,
\[
 \|z^ngz^{-n}\|\ge c_0\theta^{2|n|}.
\]
On the other hand, \(\|\gamma\|\le AM^{|\gamma|}\) for fixed
\(A,M>1\).  Membership in \(B_\Gamma(R)\) therefore forces
\(|n|=O_{\Gamma,q}(R+1)\), proving (13.10a1) and (13.10a2).
\end{proof}

The linear factor is sharp: for noncommuting hyperbolic \(q,h\in\Gamma\),
all matrices \(q^nhq^{-n}\) have the same two traces
\[
 \Tr(q^nhq^{-n})=\Tr h,
 \qquad
 \Tr(q\,q^nhq^{-n})=\Tr(qh).
\]
Thus even a nonarithmetic cocompact lattice has an infinite two-trace fiber
with \(\asymp R\) points in a radius-\(R\) word ball.  To deduce
arithmeticity, the same transfer parameter must conjugate one fixed
finite-index subgroup back into \(\Gamma\).

\begin{proposition}[transporting one finite-index subgroup forces arithmeticity]
\label{prop:marked-centralizer-commensurator}
Let \(\Gamma<\PSL_2(\R)\) be a lattice, let \(q\in\Gamma\) be hyperbolic,
put \(T=C_{\PSL_2(\R)}(q)^\circ\), and put
\(Z=T\cap\Gamma=\langle z\rangle\).  If \(\Lambda<\Gamma\) has finite
index, define
\[
 \mathcal M_\Lambda(T)=
 \{Zt\in Z\backslash T:t\Lambda t^{-1}\subset\Gamma\}.
\]
If \(\Gamma\) is nonarithmetic, then \(\mathcal M_\Lambda(T)\) is finite.
Equivalently, infinitely many such transfer parameters force
\(\Gamma\) to be arithmetic.  More quantitatively, for every proper
additive coordinate \(s:T\simeq\R\), a nonarithmetic lattice satisfies
\[
 \#\{t\in T:|s(t)|\le R,\ t\Lambda t^{-1}\subset\Gamma\}
 =O_{\Gamma,q,\Lambda}(R).
\tag{13.10a8}
\]
\end{proposition}

\begin{proof}
If \(t\Lambda t^{-1}\subset\Gamma\), then the subgroup on the left is a
lattice and consequently has finite index in the lattice \(\Gamma\).
Since \(\Lambda\) has finite index in \(\Gamma\), this says precisely that
\(t\Gamma t^{-1}\) and \(\Gamma\) are commensurable.  Thus
\(t\in\Comm(\Gamma)\).  For a nonarithmetic Fuchsian lattice, the
commensurator theorem makes \(\Comm(\Gamma)\) discrete.  Hence
\(T\cap\Comm(\Gamma)\) is a discrete subgroup of the one-dimensional
split torus which contains the lattice \(Z\); its quotient by \(Z\) is
finite, and it has only \(O(R)\) elements in a coordinate interval of
radius \(R\).
\end{proof}

The use of one fixed subgroup \(\Lambda\) in
Proposition~\ref{prop:marked-centralizer-commensurator} is essential.  A
point \(tg_0t^{-1}\in\Gamma\) transports only one cyclic
subgroup.  It gives no reason for \(t\) to transport a fixed lattice back
into \(\Gamma\), even if many unrelated values of \(t\) occur.

\begin{theorem}[an entropy bound from a loxodromic element at a divisorial valuation]
\label{thm:moving-divisorial-multiplier}
Let \(\Pi\) be a closed surface group and let
\(x=[\rho_x]\) be a faithful cocompact Fuchsian character.  Let
\(F/\Q\) be a finitely generated field of characteristic zero and let
\(\rho\colon\Pi\to\SL_2(F)\) be a representation with the following two
compatibility properties:
\[
 \begin{split}
 \tau_u(x)=\tau_{u'}(x)
   &\ \Longrightarrow\
   \Tr^2\rho(u)=\Tr^2\rho(u')\quad\text{in }F,\\
 \Tr^2\rho(g)=4\quad\text{in }F
   &\ \Longrightarrow\ \tau_g(x)=4.
 \end{split}
\]
The first property says that all trace collisions at the Fuchsian
character persist in \(\rho\); the second says that trace identities of
\(\rho\) specialize back to \(x\).

Let \(\phi\colon\Pi\twoheadrightarrow\Z\), let
\(J<\ker\phi\) be finitely generated, nonelementary, and convex cocompact
under \(\rho_x\), and write \(\delta_J\) for its critical exponent.  Fix
\(o\in\mathbb H^2\), put
\[
 \mathcal B_J(R)=\{h\in J:d(o,\rho_x(h)o)\le R\},
\]
and choose \(a\in\phi^{-1}(1)\).

Let \(v\) be a divisorial valuation of \(F\).  Suppose that for some
\(k\in\ker\phi\), after a finite field extension and an extension of \(v\),
the matrix \(D=\rho(k)\) has a spectral decomposition
\[
 D=\lambda P_++\lambda^{-1}P_-,
 \qquad P_++P_-=I,\qquad |\lambda|_v>1.
\]
Here and below the same letter \(v\) denotes the extended absolute value.
Set
\[
 A=\rho(a),\qquad
 L_\pm(Y)=\Tr(AP_\pm Y),\qquad
 \ell_v(k)=2\log|\lambda|_v,
\]
and let \(\ell_x(k)\) be the translation length of \(\rho_x(k)\) in
\(\mathbb H^2\).  Define
\[
 \begin{split}
 C_v(R)=\sup\Bigl\{\log^+\left|
 \frac{L_-(\rho(h)-\epsilon\rho(h'))}
      {L_+(\rho(h)-\epsilon\rho(h'))}
 \right|_v:\;&h,h'\in\mathcal B_J(R),\quad
 \epsilon\in\{1,-1\},\\[-2mm]
 &L_+(\rho(h)-\epsilon\rho(h'))\ne0,\\[-2mm]
 &L_-(\rho(h)-\epsilon\rho(h'))\ne0
 \Bigr\},
 \end{split}
\]
where the supremum of the empty set is zero and
\(\log^+t=\max\{0,\log t\}\), and put
\[
 \mathfrak c_v(a,k;J)=
 \limsup_{R\to\infty}\frac{C_v(R)}R.
\]
Then \(\mathfrak c_v(a,k;J)<\infty\), and the primitive trace entropy of
\(x\) satisfies
\[
 h_{\Tr}(x)\ge
 \frac{\delta_J}
 {1+\ell_x(k)\mathfrak c_v(a,k;J)/\ell_v(k)}.
\]
More precisely, if
\(c>\mathfrak c_v(a,k;J)/\ell_v(k)\),
\(n_R=\lceil cR\rceil\), and \(w_R=ak^{n_R}\), then
\[
 \sup_t\#\{h\in\mathcal B_J(R):\tau_{w_Rh}(x)=t\}
 =\exp(o(R)).
\]
\end{theorem}

\begin{proof}
We first justify the finiteness assertion.  Choose a finite symmetric
generating set of \(J\).  Convex cocompactness makes word length on \(J\)
quasi-isometric to the orbit metric, so the word length of every
\(h\in\mathcal B_J(R)\) is \(O(R)\).  Choose a normal finite-type model
of the coefficient field, after the fixed finite extension used to split
\(D\), and the filtration \(R_C[n]\) of (12.1).  Repeated multiplication
by the finitely many generator matrices shows that the entries of
\(\rho(h)\), their differences, and their images under the two fixed
linear forms \(L_\pm\) belong to \(R_{C'}[O(R)]\).

The absolute order at a fixed prime divisor of every nonzero element in
this filtration is \(O(R)\).  This is the valuation estimate used in the
proof of Theorem~\ref{thm:finite-type-divisor}: horizontal orders are
bounded by the degree, vertical orders by the affine coefficient height,
and passage through the fixed finite integral extension changes only the
constant.  Hence the \(v\)-absolute value of every nonzero quotient
occurring in \(C_v(R)\) has logarithm \(O(R)\).  This proves
\(\mathfrak c_v(a,k;J)<\infty\).

Fix \(c>\mathfrak c_v(a,k;J)/\ell_v(k)\).  For every sufficiently large
\(R\),
\[
 C_v(R)<n_R\ell_v(k),\qquad n_R=\lceil cR\rceil.
\]
Suppose that \(h,h'\in\mathcal B_J(R)\) satisfy
\(\tau_{w_Rh}(x)=\tau_{w_Rh'}(x)\).  Persistence turns this into an
equality of squared traces in \(F\), so in characteristic zero there is
an \(\epsilon\in\{1,-1\}\) for which
\[
 \Tr(AD^{n_R}\rho(h))
   =\epsilon\Tr(AD^{n_R}\rho(h')).
\]
With \(\Delta=\rho(h)-\epsilon\rho(h')\), the spectral decomposition of
\(D\) gives
\[
 \lambda^{n_R}L_+(\Delta)
   +\lambda^{-n_R}L_-(\Delta)=0.
\]
If both coefficients were nonzero, then
\[
 \log\left|\frac{L_-(\Delta)}{L_+(\Delta)}\right|_v
 =2n_R\log|\lambda|_v=n_R\ell_v(k),
\]
contrary to the preceding strict inequality.  If one coefficient is zero,
the displayed trace equation makes the other zero as well.  Thus every
collision satisfies
\[
 L_+(\rho(h))=\epsilon L_+(\rho(h')),
 \qquad
 L_-(\rho(h))=\epsilon L_-(\rho(h')).
\]

Choose \(\mu\) in a further fixed finite extension with
\(\mu^2\ne1\), and put
\[
 D_\mu=\mu P_++\mu^{-1}P_-,
 \qquad A_0=A,\qquad A_1=AD_\mu.
\]
The two traces
\[
 \Tr(A_0Y)=L_+(Y)+L_-(Y),\qquad
 \Tr(A_1Y)=\mu L_+(Y)+\mu^{-1}L_-(Y)
\]
are independent linear combinations of \(L_+(Y)\) and \(L_-(Y)\).
Moreover
\[
 A_1A_0^{-1}=AD_\mu A^{-1},\qquad
 \Tr^2(A_1A_0^{-1})=(\mu+\mu^{-1})^2\ne4.
\]
After fixing one element \(h'\) in a trace fiber, the two coefficient
equalities therefore put the whole fiber in one simultaneous
\emph{squared}-trace fiber for the two fixed evaluations
\[
 h\longmapsto
 \bigl(\Tr^2(A_0\rho(h)),\Tr^2(A_1\rho(h))\bigr).
\]
The sign may vary with \(h\), which is harmless because both coordinates
are squared.  The specialization hypothesis in
Theorem~\ref{thm:two-trace-evaluations} is exactly the second compatibility
property in the statement.  That theorem, and the quasi-isometry between
word length and displacement in the Fuchsian representation on \(J\), give the asserted
\(\exp(o(R))\) fiber estimate.

Orbit counting gives
\(\#\mathcal B_J(R)=\exp((\delta_J+o(1))R)\).  Every \(w_Rh\) is primitive:
if \(w_Rh=u^m\) with \(|m|\ge2\), then
\(1=\phi(w_Rh)=m\phi(u)\), which is impossible.  Finally,
\[
 \ell_x(w_Rh)
 \le d(o,\rho_x(a)o)+d(o,\rho_x(k^{n_R})o)
       +d(o,\rho_x(h)o)
 \le (1+c\ell_x(k)+o(1))R,
\]
because \(d(o,\rho_x(k^n)o)=n\ell_x(k)+O(1)\).  Dividing the orbit count
by the subexponential fiber bound, and then letting
\(c\downarrow\mathfrak c_v(a,k;J)/\ell_v(k)\), proves the entropy bound.
\end{proof}

\begin{corollary}[a quantitative criterion excluding an unbounded divisor]
\label{cor:divisorial-half-entropy}
In the setting of
Theorem~\ref{thm:moving-divisorial-multiplier}, suppose that
\(\delta_J>1/2\) and
\[
 \frac{\ell_v(k)}{\ell_x(k)}>
 \frac{\mathfrak c_v(a,k;J)}{2\delta_J-1}.
\]
Then \(h_{\Tr}(x)>1/2\).  Consequently, if \(h_{\Tr}(x)\le1/2\), no
divisor along which all trace collisions at \(x\) persist can satisfy this
inequality.  Equivalently, every unbounded divisor compatible with those
collisions, and every admissible \(v\)-loxodromic \(k\),
must satisfy the exponential-cancellation lower bound
\[
 \mathfrak c_v(a,k;J)\ge
 (2\delta_J-1)\frac{\ell_v(k)}{\ell_x(k)}.
\]
\end{corollary}

\begin{proof}
The displayed inequality says exactly that the lower bound in
Theorem~\ref{thm:moving-divisorial-multiplier} is strictly greater than
\(1/2\).
\end{proof}

The preceding theorem controls every possible cancellation by its largest
valuation.  The collision equation itself is more rigid: it asks whether a
particular quotient is an \emph{exact} even power of \(\lambda\).  The first
exponent for which no such identity occurs gives a sharper estimate.

\begin{theorem}[the first exponent not realized by an exact cancellation identity]
\label{thm:first-missing-resonance}
Retain the hypotheses and notation of
Theorem~\ref{thm:moving-divisorial-multiplier}.  For \(R>0\), let
\(\mathcal E_{v,J}(R)\subset\Z_{\ge0}\) be the set of integers \(n\) for
which there are distinct \(h,h'\in\mathcal B_J(R)\) and
\(\epsilon\in\{1,-1\}\) such that, for
\[
 \Delta=\rho(h)-\epsilon\rho(h'),
\]
both \(L_+(\Delta)\) and \(L_-(\Delta)\) are nonzero and
\[
 -\frac{L_-(\Delta)}{L_+(\Delta)}=\lambda^{2n}.
\tag{13.10a}
\]
Define
\[
 m_{v,J}(R)=
 \min\bigl(\Z_{\ge0}\setminus\mathcal E_{v,J}(R)\bigr),
 \qquad
 \mathfrak m_v(a,k;J)=
 \liminf_{R\to\infty}\frac{m_{v,J}(R)}R.
\tag{13.10b}
\]
Then \(m_{v,J}(R)<\infty\),
\[
 \mathfrak m_v(a,k;J)
 \le \frac{\mathfrak c_v(a,k;J)}{\ell_v(k)},
\tag{13.10c}
\]
and
\[
 h_{\Tr}(x)\ge
 \frac{\delta_J}
 {1+\ell_x(k)\mathfrak m_v(a,k;J)}.
\tag{13.10d}
\]
More precisely, if
\[
 n_R=m_{v,J}(R),\qquad w_R=ak^{n_R},
\]
then
\[
 \sup_t
 \#\{h\in\mathcal B_J(R):\tau_{w_Rh}(x)=t\}
 =\exp(o(R)).
\tag{13.10e}
\]

Consequently, if \(h_{\Tr}(x)\le1/2\) and \(\delta_J>1/2\), then
\[
 \mathfrak m_v(a,k;J)\ge
 \frac{2\delta_J-1}{\ell_x(k)}.
\tag{13.10f}
\]
Equivalently, for every \(0<\eta<2\delta_J-1\) and every sufficiently
large \(R\), each integer
\[
 0\le n<
 \frac{(2\delta_J-1-\eta)R}{\ell_x(k)}
\tag{13.10g}
\]
is realized by an exact identity of the form \((13.10a)\).
\end{theorem}

\begin{proof}
For a fixed triple \((h,h',\epsilon)\), equation (13.10a) can hold for at
most one \(n\), because \(|\lambda|_v>1\).  Thus
\(\mathcal E_{v,J}(R)\) is finite.  More quantitatively, if
\(n\in\mathcal E_{v,J}(R)\), then
\[
 n\ell_v(k)
 =
 \log\left|
 \frac{L_-(\Delta)}{L_+(\Delta)}
 \right|_v
 \le C_v(R).
\]
It follows that
\[
 m_{v,J}(R)
 \le
 1+\left\lfloor\frac{C_v(R)}{\ell_v(k)}\right\rfloor,
\]
which proves finiteness and (13.10c).

Put \(n_R=m_{v,J}(R)\).  Suppose that distinct
\(h,h'\in\mathcal B_J(R)\) satisfy
\(\tau_{w_Rh}(x)=\tau_{w_Rh'}(x)\).  Persistence of trace collisions gives
\[
 \Tr(AD^{n_R}\rho(h))
 =\epsilon\Tr(AD^{n_R}\rho(h'))
\]
for some \(\epsilon\in\{1,-1\}\).  Hence, with
\(\Delta=\rho(h)-\epsilon\rho(h')\),
\[
 \lambda^{n_R}L_+(\Delta)
 +\lambda^{-n_R}L_-(\Delta)=0.
\tag{13.10h}
\]
If both coefficients were nonzero, (13.10h) would say precisely that
\(n_R\in\mathcal E_{v,J}(R)\), contrary to its definition.  If either
coefficient vanishes, (13.10h) forces the other to vanish as well.  Thus
every nontrivial collision satisfies
\[
 L_+(\rho(h))=\epsilon L_+(\rho(h')),
 \qquad
 L_-(\rho(h))=\epsilon L_-(\rho(h')).
\tag{13.10i}
\]

Choose, in a fixed finite extension, an element \(\mu\) with
\(\mu^2\ne1\), and put
\[
 D_\mu=\mu P_++\mu^{-1}P_-,
 \qquad A_0=A,\qquad A_1=AD_\mu.
\]
The two traces against \(A_0,A_1\) are independent linear combinations of
\(L_+,L_-\), and
\[
 \Tr^2(A_1A_0^{-1})
 =\Tr^2(AD_\mu A^{-1})
 =(\mu+\mu^{-1})^2\ne4.
\]
Fixing one member of a trace fiber, (13.10i) places every other member in
the same simultaneous squared-trace fiber for \(A_0,A_1\).
Lemma~\ref{lem:ambient-two-trace-evaluations}, applied to
\(\mathcal B_J(R)\), proves (13.10e).

Every \(w_Rh\) is primitive because \(\phi(w_Rh)=1\).  Moreover
\[
 \ell_x(w_Rh)
 \le R+m_{v,J}(R)\ell_x(k)+O(1).
\tag{13.10j}
\]
Choose a sequence \(R_j\to\infty\) along which
\[
 \frac{m_{v,J}(R_j)}{R_j}
 \longrightarrow \mathfrak m_v(a,k;J).
\]
Convex-cocompact orbit counting gives
\[
 \#\mathcal B_J(R_j)=\exp((\delta_J+o(1))R_j).
\]
Dividing by the fiber bound (13.10e), using (13.10j), and taking the limsup
in the definition of \(h_{\Tr}\) proves (13.10d).  Rearranging (13.10d)
under the assumption \(h_{\Tr}(x)\le1/2\) gives (13.10f).  Since
\(m_{v,J}(R)\) is the first absent integer, (13.10g) follows directly from
the definition of the liminf.
\end{proof}

\begin{theorem}[an internal multiplier produces all initial cancellation exponents]
\label{thm:internal-multiplier-resonance}
Retain the hypotheses and notation of
Theorem~\ref{thm:moving-divisorial-multiplier}, and suppose in addition
that \(k\in J\).  Assume that there are distinct \(u,v\in J\) and a sign
\(\epsilon\in\{1,-1\}\) such that
\[
 \Tr(A\rho(u))=\epsilon\Tr(A\rho(v))
\]
and, for \(\Delta_0=\rho(u)-\epsilon\rho(v)\),
\[
 L_+(\Delta_0)L_-(\Delta_0)\ne0.
\]
Then there is a constant \(C\), independent of \(n\) and \(R\), such that
\[
 \left\{0,1,\ldots,
 \left\lfloor\frac{R-C}{\ell_x(k)}\right\rfloor\right\}
 \subset \mathcal E_{v,J}(R)
\tag{13.10ab}
\]
for all sufficiently large \(R\).  In particular,
\[
 \liminf_{R\to\infty}\frac{m_{v,J}(R)}R
 \ge\frac1{\ell_x(k)}.
\tag{13.10ac}
\]

If \(h_{\Tr}(x)<\delta_J\), then such \(u,v,\epsilon\) necessarily
exist.  In particular they exist when
\(h_{\Tr}(x)\le1/2<\delta_J\).  Consequently, under low trace entropy,
the method of Theorem~\ref{thm:first-missing-resonance} cannot prove an
entropy bound greater than \(1/2\) for a subgroup containing the multiplier
\(k\).
\end{theorem}

\begin{proof}
For \(n\ge0\), put
\[
 h_n=k^{-n}u,\qquad h'_n=k^{-n}v.
\]
Both elements belong to \(J\).  Since
\(D^{-n}=\lambda^{-n}P_++\lambda^nP_-\) and the spectral projectors
commute with \(D\), one has
\[
 \begin{split}
 L_+\bigl(\rho(h_n)-\epsilon\rho(h'_n)\bigr)
   &=\lambda^{-n}L_+(\Delta_0),\\
 L_-\bigl(\rho(h_n)-\epsilon\rho(h'_n)\bigr)
   &=\lambda^nL_-(\Delta_0).
 \end{split}
\]
The assumed trace equality is the identity
\(L_+(\Delta_0)+L_-(\Delta_0)=0\).  The two coefficients are nonzero, so
\[
 -\frac{L_-\bigl(\rho(h_n)-\epsilon\rho(h'_n)\bigr)}
        {L_+\bigl(\rho(h_n)-\epsilon\rho(h'_n)\bigr)}
 =\lambda^{2n}.
\]
Moreover
\[
 d(o,\rho_x(k^{-n}u)o),\quad
 d(o,\rho_x(k^{-n}v)o)
 \le n\ell_x(k)+C,
\]
because \(d(o,\rho_x(k^n)o)=n\ell_x(k)+O(1)\) and \(u,v\) are fixed.
This proves (13.10ab) and (13.10ac).

It remains to obtain the fixed collision from low trace entropy.  Write
\(s=h_{\Tr}(x)<\delta_J\).  Convex-cocompact orbit counting gives
\[
 \#\mathcal B_J(T)=\exp((\delta_J+o(1))T).
\]
Every \(ah\), \(h\in\mathcal B_J(T)\), is primitive because
\(\phi(ah)=1\), and its length in the Fuchsian representation is at most
\(T+O(1)\).  Hence
these elements assume at most \(\exp((s+o(1))T)\) squared-trace values at
\(x\).  A squared-trace fiber \(\mathcal F_T\) consequently has size
\[
 |\mathcal F_T|
 \ge \exp((\delta_J-s-o(1))T).
\tag{13.10ad}
\]

Fix \(v_0\in\mathcal F_T\).  Persistence of collisions implies that, for
each \(u\in\mathcal F_T\), one can choose
\(\epsilon_u\in\{1,-1\}\) with
\[
 \Tr(A\rho(u))=\epsilon_u\Tr(A\rho(v_0)).
\]
Choose in a fixed finite extension an element \(\mu\) with
\(\mu^2\ne1\), and put
\[
 D_\mu=\mu P_++\mu^{-1}P_-,\qquad A_0=A,\qquad A_1=AD_\mu.
\]
Suppose, contrary to the assertion, that every \(u\in\mathcal F_T\)
satisfied
\[
 L_\pm(\rho(u))=\epsilon_uL_\pm(\rho(v_0)).
\]
Then \(\mathcal F_T\) would lie in one simultaneous squared-trace fiber
for \(A_0,A_1\).  Since
\[
 \Tr^2(A_1A_0^{-1})=(\mu+\mu^{-1})^2\ne4,
\]
Lemma~\ref{lem:ambient-two-trace-evaluations} would give
\(|\mathcal F_T|=\exp(o(T))\), contradicting (13.10ad).  Thus for some
\(u\in\mathcal F_T\) at least one of the two coefficient equalities fails.
Take \(v=v_0\) and \(\epsilon=\epsilon_u\).  The trace equality says that
the sum of the two coefficient differences is zero, so if one is nonzero,
then both are nonzero.  This completes the proof.
\end{proof}

\begin{remark}[why the multiplier must lie outside the subgroup]
The obstruction is internal multiplication, not specifically conjugacy.
To exclude it in a way stable under replacing \(k\) by a positive power,
one must at least require
\[
 \langle k\rangle\cap J=\{1\}.
\tag{13.10ae}
\]
Indeed, if \(k^q\in J\), the theorem applies to the equally admissible
multiplier \(k^q\), and both its Fuchsian and its \(v\)-adic translation
length are multiplied by \(q\).  Condition (13.10ae) is therefore a
necessary elementary separation condition for this argument.  By itself,
it says nothing about cancellation identities arising from unrelated pairs.
\end{remark}

This separation is compatible with a critical exponent as close
to one as desired.  We record a form which also separates the chosen
subgroup from Fuchsian conjugacy relations that could imitate multiplication
by \(k\).  For \(w\in\Gamma\), write
\[
 \mathcal L_w(\Gamma)
 =\{w^{-1}\gamma w\gamma^{-1}:\gamma\in\Gamma\}.
\]

\begin{lemma}[separated limit sets and twisted commutators]
\label{lem:finite-twisted-commutators}
Let \(\Gamma<\PSL_2(\R)\) be discrete, let \(w\in\Gamma\) be hyperbolic,
and let \(H<\Gamma\) be nonelementary and convex cocompact.  If
\[
 \Lambda(H)\cap w^{-1}\Lambda(H)=\varnothing,
\tag{13.10af}
\]
then \(H\cap\mathcal L_w(\Gamma)\) is finite.
\end{lemma}

\begin{proof}
Suppose that there are distinct
\(c_i=w^{-1}\gamma_iw\gamma_i^{-1}\in H\), and put
\(b_i=wc_i=\gamma_iw\gamma_i^{-1}\).  The \(b_i\) are distinct
hyperbolic elements, all with translation length \(\ell(w)\).  Their axes
leave every compact subset of \(\mathbb H^2\).  Indeed, if the axes met one
fixed compact set, the displacement of a fixed basepoint under the \(b_i\)
would be uniformly bounded; discreteness would then leave only finitely
many possibilities for \(b_i\).

Write the endpoints of the axis of \(b_i\) as \(x_i,y_i\).  After passage
to a subsequence they converge to \(x,y\in\partial\mathbb H^2\).  If
\(x\ne y\), the geodesics with endpoints \(x_i,y_i\) meet a common compact
neighborhood of the geodesic with endpoints \(x,y\), a contradiction.
Thus \(x=y=\xi\).  A hyperbolic element is determined on the boundary by
its two fixed points and its multiplier.  Since the multiplier here is
fixed, the cross-ratio formula shows that
\[
 b_i z\longrightarrow\xi,
 \qquad b_i^{-1}z\longrightarrow\xi
\]
locally uniformly on
\(\partial\mathbb H^2\setminus\{\xi\}\).  Consequently
\[
 c_i z=w^{-1}b_i z\longrightarrow w^{-1}\xi,
 \qquad
 c_i^{-1}z=b_i^{-1}wz\longrightarrow\xi
\]
away from the corresponding exceptional points.  Applying these limits
to points of the closed, \(H\)-invariant, nonelementary set \(\Lambda(H)\)
shows that both \(\xi\) and \(w^{-1}\xi\) belong to \(\Lambda(H)\).
Therefore
\(w^{-1}\xi\in\Lambda(H)\cap w^{-1}\Lambda(H)\), contradicting
(13.10af).
\end{proof}

\begin{proposition}[high-exponent subgroups transverse to a prescribed multiplier]
\label{prop:high-exponent-transverse-subgroup}
Let \(\Gamma<\PSL_2(\R)\) be a torsion-free cocompact lattice, let
\(\phi:\Gamma\twoheadrightarrow\Z\), and put \(N=\ker\phi\).  Fix
\(a\in\phi^{-1}(1)\) and a nontrivial \(k\in N\), and put
\(w_n=ak^n\).  For every \(\varepsilon>0\) there is a finitely generated
nonelementary convex-cocompact subgroup \(J<N\) such that
\[
 \delta_J>1-\varepsilon,
 \qquad \langle k\rangle\cap J=\{1\},
 \qquad J\cap w_n^{-1}Jw_n=\{1\}
 \quad(n\ge0).
\tag{13.10ag}
\]
It can moreover be chosen so that
\[
 \tau_q(x)\ne\tau_k(x)
 \qquad(q\in J\setminus\{1\}).
\tag{13.10ah}
\]
For each fixed \(n\), the intersection
\(J\cap\mathcal L_{w_n}(\Gamma)\) is finite.  Given any fixed
\(M\), \(J\) may additionally be chosen so that these intersections equal
\(\{1\}\) for every \(0\le n\le M\).
\end{proposition}

\begin{proof}
The quotient \(\Gamma/N\cong\Z\) is amenable, so the
amenable-cover theorem gives \(\delta_N=\delta_\Gamma=1\)
\cite{Brooks,DougallSharp}.  We briefly explain why finitely generated
subgroups of \(N\) approximate this exponent.  On
\(X=N\backslash\mathbb H^2\), amenability gives \(\lambda_0(X)=0\).
Choose compactly supported smooth functions with Rayleigh quotients tending
to zero, and enclose their supports in connected compact subsurfaces
\(K_i\subset X\).  If \(H_i<N\) is the image of \(\pi_1(K_i)\), then
\(K_i\) embeds in the cover
\(H_i\backslash\mathbb H^2\).  Extending the lifted test function by zero
there gives
\(\lambda_0(H_i\backslash\mathbb H^2)\to0\).  The groups \(H_i\) are
eventually nonelementary, are finitely generated, and contain no parabolic
elements; hence they are convex cocompact.  Sullivan's formula
\[
 \lambda_0(H_i\backslash\mathbb H^2)=
 \begin{cases}
  1/4,&\delta_{H_i}\le1/2,\\
  \delta_{H_i}(1-\delta_{H_i}),&\delta_{H_i}>1/2
 \end{cases}
\]
then gives \(\delta_{H_i}\to1\) \cite{SullivanPositivity,BorthwickSpectral}.
Choose one such group \(H<N\) with
\(\delta_H>1-\varepsilon\).

We next choose one boundary interval which is separated from all the
moving translates needed in (13.10ag).  The maps
\(w_n^{-1}=k^{-n}a^{-1}\) converge locally uniformly, away from
\(a(k^+)\), to the constant \(k^-\).  Choose a point
\(\xi\in\partial\mathbb H^2\) different from \(k^-\), from \(a(k^+)\),
and from every fixed point of the countable family \(\{w_n:n\ge0\}\).
For all sufficiently large \(n\), a small closed interval about \(\xi\)
is disjoint from its image under \(w_n^{-1}\); after shrinking it finitely
many more times, we obtain an open interval \(U\) such that
\[
 \overline U\cap\operatorname{Fix}(k)=\varnothing,
 \qquad
 \overline U\cap w_n^{-1}\overline U=\varnothing
 \quad(n\ge0).
\tag{13.10ai}
\]

The normal nonelementary subgroup \(N\) has full limit set, and ordered
pairs of fixed points of its hyperbolic elements are dense off the diagonal.
Choose a hyperbolic \(b\in N\) with \(b^+\in U\) and
\(b^-\notin\Lambda(H)\).  North--south dynamics, uniformly on the compact
set \(\Lambda(H)\), gives
\(b^L\Lambda(H)\subset U\) for all sufficiently large \(L\).  Replace
\(H\) by \(H_1=b^LHb^{-L}\).  Its critical exponent is unchanged and
\(\Lambda(H_1)\subset U\).  The first condition in (13.10ai) implies
\(H_1\cap\langle k\rangle=\{1\}\).  The second implies
\(H_1\cap w_n^{-1}H_1w_n=\{1\}\): a nontrivial element in this
intersection would have both fixed points in the two disjoint limit sets.
Lemma~\ref{lem:finite-twisted-commutators} also makes each
\(H_1\cap\mathcal L_{w_n}(\Gamma)\) finite.

There are only finitely many \(H_1\)-conjugacy classes whose translation
length in the representation \(x\) equals \(\ell_x(k)\), because \(H_1\) is convex
cocompact.  The finitely generated linear group \(H_1\) is residually
finite.  Choose a finite-index subgroup \(J\triangleleft H_1\) which
excludes one representative of every nontrivial such class.  Normality
then excludes the whole class and proves (13.10ah).  If \(M\) was fixed in
advance, also exclude the finitely many nonidentity elements in
\[
 \bigcup_{0\le n\le M}
 H_1\cap\mathcal L_{w_n}(\Gamma).
\]
Passing to finite index preserves the limit set and the critical exponent,
so every assertion follows.
\end{proof}

The next elementary rigidity statement explains the usefulness of
(13.10ah).  It rules out not just propagation by \(k\), but propagation by
any fixed semisimple element of \(J\) whose translation length at \(x\) is different.
The formulation allows the exponents of \(D\) to lie in an arbitrary
four-term arithmetic progression; one multiplication by \(Q^{-1}\) is then
forced to have the same squared trace as multiplication by \(D^r\), where
\(r\) is the common difference.

\begin{lemma}[four cancellation equations determine the multiplier]
\label{lem:coherent-resonance-rigidity}
Let \(K\) be a characteristic-zero valued field.  Write
\[
 D=\lambda P_++\lambda^{-1}P_-,\qquad |\lambda|>1,
 \qquad L_\pm(Y)=\Tr(AP_\pm Y).
\]
Let \(Q\in\SL_2(K)\) have distinct eigenvalues and let
\(\Delta\in M_2(K)\).  Let \(n_0\in\Z\) and \(r\ge1\).  Suppose that,
for \(j=0,1,2,3\), both \(L_+(Q^{-j}\Delta)\) and
\(L_-(Q^{-j}\Delta)\) are nonzero and
\[
 -\frac{L_-(Q^{-j}\Delta)}{L_+(Q^{-j}\Delta)}
   =\lambda^{2(n_0+rj)}.
\tag{13.10aj}
\]
Then
\[
 \Tr^2Q=\Tr^2(D^r).
\tag{13.10ak}
\]
More precisely, if
\(Q=\mu E_++\mu^{-1}E_-\), then either
\(\mu^2=\lambda^{2r}\) or \(\mu^2=\lambda^{-2r}\); the corresponding two
mixed spectral coefficients vanish.
\end{lemma}

\begin{proof}
The assertion may be checked after an algebraic extension over which \(Q\)
is diagonalizable.  Put
\[
 \alpha=L_+(E_+\Delta),\quad \beta=L_+(E_-\Delta),\quad
 \gamma=L_-(E_+\Delta),\quad \zeta=L_-(E_-\Delta).
\]
Set \(\Lambda=\lambda^{2r}\).  Multiplication of (13.10aj) by
\(\mu^j\) gives
\[
 \gamma+\zeta(\mu^2)^j
 +\lambda^{2n_0}\alpha\Lambda^j
 +\lambda^{2n_0}\beta(\Lambda\mu^2)^j=0,
 \qquad j=0,1,2,3.
\tag{13.10al}
\]
If the four bases
\(1,\mu^2,\Lambda,\Lambda\mu^2\) are distinct, their Vandermonde
determinant is nonzero, so all four coefficients vanish, contrary to the
nonvanishing assumption.  Since \(\Lambda\ne1\) and \(Q\) has distinct
eigenvalues, the only possible coincidences are
\(\mu^2=\Lambda\) and \(\mu^2=\Lambda^{-1}\).  In the first case the
three remaining bases in (13.10al) give
\(\gamma=\beta=0\) and
\(\zeta=-\lambda^{2n_0}\alpha\ne0\).  In the second they give
\(\alpha=\zeta=0\) and
\(\gamma=-\lambda^{2n_0}\beta\ne0\).  In either case the
squared traces of \(Q\) and \(D^r\) agree.
\end{proof}

Suppose, under the hypotheses of
Theorem~\ref{thm:first-missing-resonance}, that one fixed
\(q\in J\setminus\{1\}\) and one fixed pair \(u,v\in J\) were to realize
four consecutive exact cancellation equations through
\(h_n=q^{-n}u\) and \(h'_n=q^{-n}v\).  Taking \(r=1\), the lemma would give
\(\Tr^2\rho(q)=\Tr^2\rho(k)\).  Persistence and specialization at the
original Fuchsian character would then give \(\tau_q(x)=\tau_k(x)\),
contradicting (13.10ah).  Thus, if a long interval of cancellation exponents
is realized, the pair realizing the identity must vary with the exponent;
one fixed pair and one fixed element of \(J\) cannot account for the
interval.  This conclusion does not bound identities arising from
unrelated pairs.

The distinction between a fixed pair and varying pairs is essential.  Small sumsets for
many geometrically spaced coefficients do not, by themselves, produce even
two cancellation equations witnessed by the same pair.

\begin{proposition}[small sumsets at geometric coefficients need not have a common collision]
\label{prop:geometric-sumset-no-common-collision}
Fix an integer \(b\ge2\), and for \(q\ge2\) put
\(A_q=\{0,1,\ldots,q-1\}\subset\Q\).  Whenever
\(0\le n\le\lfloor\log_b q\rfloor\),
\[
 \#(A_q+b^nA_q)=(b^n+1)(q-1)+1\le(b^n+1)\#A_q.
\tag{13.10ala}
\]
Nevertheless, for two distinct ordered pairs
\((x,y),(x',y')\in A_q^2\), the equality
\[
 x+b^ny=x'+b^ny'
\tag{13.10alb}
\]
can hold for at most one integer \(n\ge0\).
This example occurs inside determinant-one matrices: the map
\[
 (x,y)\longmapsto
 M(x,y)=\begin{pmatrix}x&1\\xy-1&y\end{pmatrix}
 \quad\text{has}\quad \det M(x,y)=1,
\]
and, over \(\Q(\sqrt b)\), if
\(D=\operatorname{diag}(b^{-1/2},b^{1/2})\), then
\[
 \Tr(D^nM(x,y))=b^{-n/2}(x+b^ny).
\]
Thus the obstruction persists for the same trace observables that arise from
a diagonal one-parameter subgroup of \(\SL_2\).
\end{proposition}

\begin{proof}
For fixed \(y\), the sums \(x+b^ny\) fill the integer interval
\([b^ny,b^ny+q-1]\).  Since \(b^n\le q\), the intervals for consecutive
values of \(y\) overlap or touch, and their union is
\([0,(b^n+1)(q-1)]\cap\Z\).  This proves (13.10ala).
If (13.10alb) held for two different integers \(m,n\), subtraction would
give \((b^m-b^n)(y-y')=0\).  Hence \(y=y'\), and then \(x=x'\), contrary
to the assumed distinction.  The determinant and trace identities follow
by direct calculation.
\end{proof}

It is tempting to replace \(J\) by the full kernel \(\ker\phi\), whose
critical exponent is \(1\) because the quotient is amenable
\cite{DougallSharp}.  The following proposition shows that this
cannot improve the argument: conjugation alone fills exactly the initial
interval that would be needed to cross the exponent \(1/2\).

\begin{proposition}[conjugation realizes the initial cancellation exponents in the full homology kernel]
\label{prop:kernel-resonance-saturation}
Retain the preceding hypotheses and put \(N=\ker\phi\).  Define
\(\mathcal E_{v,N}(R)\) exactly as in (13.10a), but with
\[
 \mathcal B_N(R)=
 \{h\in N:d(o,\rho_x(h)o)\le R\}
\]
in place of \(\mathcal B_J(R)\).  Then, for every \(\eta\in(0,1)\) and
every sufficiently large \(R\),
\[
 \left\{0,1,\ldots,
 \left\lfloor\frac{(1-\eta)R}{\ell_x(k)}\right\rfloor
 \right\}
 \subset\mathcal E_{v,N}(R).
\tag{13.10k}
\]
Consequently, if
\[
 m_{v,N}(R)=
 \min\bigl(\Z_{\ge0}\setminus\mathcal E_{v,N}(R)\bigr),
\]
then
\[
 \liminf_{R\to\infty}\frac{m_{v,N}(R)}R
 \ge\frac1{\ell_x(k)}.
\tag{13.10l}
\]
Thus replacing the finitely generated subgroup \(J\) by the full homology
kernel cannot make the argument of
Theorem~\ref{thm:first-missing-resonance} prove trace entropy greater than
\(1/2\).
\end{proposition}

\begin{proof}
Choose once and for all a finitely generated nonelementary
convex-cocompact subgroup \(J_0<N\), for example the subgroup \(J\) in
Theorem~\ref{thm:first-missing-resonance}, and let \(\delta_0>0\) be its
critical exponent.  For \(n\ge0\), put
\[
 w_n=ak^n,\qquad h_0=w_n^{-1}a=k^{-n}.
\]
For \(\gamma\in J_0\), define
\[
 h_\gamma=w_n^{-1}\gamma a\gamma^{-1}.
\tag{13.10m}
\]
Both \(h_0\) and \(h_\gamma\) belong to \(N\), and
\[
 w_nh_0=a,\qquad
 w_nh_\gamma=\gamma a\gamma^{-1}.
\tag{13.10n}
\]
Their traces are therefore equal identically in every representation of
\(\Pi\).

There is a constant \(C\), independent of \(n,R,\gamma\), such that
\[
 d(o,\rho_x(h_\gamma)o)
 \le
 n\ell_x(k)+2d(o,\rho_x(\gamma)o)+C.
\tag{13.10o}
\]
Fix \(\eta\in(0,1)\), and suppose that
\[
 0\le n\le\frac{(1-\eta)R}{\ell_x(k)}.
\]
For large \(R\), put
\[
 S_{R,n}=\frac{R-n\ell_x(k)-C}{2}.
\]
Then \(S_{R,n}\ge\eta R/3\).  Every
\(\gamma\in\mathcal B_{J_0}(S_{R,n})\) gives
\(h_\gamma\in\mathcal B_N(R)\), and convex-cocompact orbit counting gives
\[
 \#\mathcal B_{J_0}(S_{R,n})
 =\exp((\delta_0+o(1))S_{R,n})
 \ge \exp(c_\eta R)
\tag{13.10p}
\]
for some \(c_\eta>0\), uniformly over the indicated integers \(n\).

The map \(\gamma\mapsto h_\gamma\) is injective.  Indeed, equality of two
values in (13.10m) says that \(\gamma_2^{-1}\gamma_1\) centralizes \(a\).
Centralizers of nontrivial elements in a closed surface group are cyclic.
Moreover \(\phi(a)=1\) implies that \(a\) is not a proper power and that
\[
 C_\Pi(a)\cap N=\{1\}.
\]
Since \(\gamma_2^{-1}\gamma_1\in J_0<N\), one obtains
\(\gamma_1=\gamma_2\).

Apply (13.10n) in the representation \(\rho\).  With
\[
 \Delta_\gamma=\rho(h_\gamma)-\rho(h_0),
\]
the spectral expansion gives
\[
 \lambda^nL_+(\Delta_\gamma)
 +\lambda^{-n}L_-(\Delta_\gamma)=0.
\tag{13.10q}
\]
The elements \(h_\gamma\) all have ambient \(\Pi\)-word length \(O(R)\).
Those for which
\[
 L_+(\Delta_\gamma)=L_-(\Delta_\gamma)=0
\tag{13.10r}
\]
lie in one simultaneous squared-trace fiber for
\[
 A_0=A,\qquad A_1=AD_\mu,
\]
where \(D_\mu=\mu P_++\mu^{-1}P_-\) and \(\mu^2\ne1\).
Lemma~\ref{lem:ambient-two-trace-evaluations} bounds the number satisfying
(13.10r) by \(\exp(o(R))\), uniformly in \(n\) and in the resulting trace
values.  This is negligible compared with (13.10p).  Hence, for every
integer \(n\) in the indicated interval, some \(\gamma\) has both
coefficients in (13.10q) nonzero.  Equation (13.10q) then becomes
\[
 -\frac{L_-(\Delta_\gamma)}{L_+(\Delta_\gamma)}
 =\lambda^{2n},
\]
so \(n\in\mathcal E_{v,N}(R)\).

There are only \(O(R)\) integers in (13.10k), and all estimates above are
uniform in \(n\); consequently the conclusion holds simultaneously for
all of them once \(R\) is sufficiently large.  This proves (13.10k), and
letting \(\eta\downarrow0\) gives (13.10l).
\end{proof}

\begin{remark}
Proposition~\ref{prop:kernel-resonance-saturation} explains why the
subgroup in Theorem~\ref{thm:first-missing-resonance} cannot be replaced
by the full normal kernel.  The conjugacy families (13.10m) in the full
kernel realize every exponent up to the limit imposed by their Fuchsian
translation length.  Thus the limitation comes from an explicit
group-theoretic construction, not from the valuation estimate.
\end{remark}

\begin{lemma}[an unbounded representation has a kernel element of positive local translation length]
\label{lem:loxodromic-in-kernel}
Let \(K\) be a complete valued field of characteristic zero of
archimedean or discrete nonarchimedean type, and let
\(r\colon\Pi\to\PSL_2(K)\) be Zariski dense.  Suppose that
\(J<\ker\phi\) has Zariski-dense image, where
\(\phi\colon\Pi\twoheadrightarrow\Z\).  If \(r(\Pi)\) is unbounded, then
\(r(\ker\phi)\) contains an element which, after a finite extension of
\(K\), has eigenvalues \(\lambda,\lambda^{-1}\) with
\(|\lambda|>1\).
\end{lemma}

\begin{proof}
Put \(N=\ker\phi\).  Its image is Zariski dense because it contains
\(r(J)\).  The elliptic-subgroup dichotomy for an action on a rank-one
symmetric space or a tree says that a subgroup containing no element of
positive translation length either has a global fixed point or fixes a
point at infinity; finite generation is not required.  The second
alternative puts its Zariski closure in a proper parabolic and is
impossible here.  Thus, if \(r(N)\) had no such element, its fixed-point
set would be nonempty.

This fixed-point set is bounded.  Indeed, \(J\) is finitely generated and
Zariski dense, so it contains two elements with no common eigenline.  In
rank one, the intersection of the fixed subspaces of two elliptic elements
is unbounded only if they have a common endpoint at infinity, equivalently
a common eigenline.  Hence \(\operatorname{Fix}(r(J))\) is bounded, and
\(\operatorname{Fix}(r(N))\subset\operatorname{Fix}(r(J))\) is bounded as
well.  Normality of \(N\) makes \(\operatorname{Fix}(r(N))\) invariant
under \(r(\Pi)\).  Its circumcenter in the symmetric space, or its center
in the tree after subdividing edges if necessary, is consequently fixed
by \(r(\Pi)\), contrary to the assumed unboundedness.  Hence \(r(N)\)
contains an element of positive translation length.  Its two fixed points
and eigenvalues are defined after a finite extension, and positive
translation length is equivalent to \(|\lambda|>1\).
\end{proof}

\begin{remark}
Both compatibility properties in
Theorem~\ref{thm:moving-divisorial-multiplier} occur naturally.  If \(x\)
is algebraic and \(\rho\) is a lift of its algebraic representation, they
follow from injectivity of the distinguished field embedding.  If
\(\rho\) is the generic representation of an irreducible component of
the variety \(V_x\) on which the trace collisions at \(x\) persist, the first property is the
definition of \(V_x\).  For the second, an identity
\(\Tr^2\rho(g)=4\) says that the regular function
\(\tau_g|_{V_x}\) is identically \(4\); evaluating that regular identity
at \(x\) gives \(\tau_g(x)=4\).  No homomorphism from the abstract
function field to \(\R\) is being asserted.
Lemma~\ref{lem:loxodromic-in-kernel} shows that every unbounded
divisor in either setting admits such an element \(k\); the remaining
condition is now the explicit, scale-invariant inequality in
Corollary~\ref{cor:divisorial-half-entropy}.

Replacing \(k\) by \(k^m\) does not improve that inequality:
\(\ell_x(k^m)\) and \(\ell_v(k^m)\) both multiply by \(m\), while the
spectral projectors and \(\mathfrak c_v(a,k;J)\) do not change.  Thus the
criterion compares the translation length in the original Fuchsian
representation, the translation length at \(v\), and the rate of
cancellation.  Taking a large power changes neither side of this comparison.

For an algebraic character, the same proof works at a nondistinguished
archimedean embedding.  Use the ordinary absolute value, choose a proximal
kernel element with \(|\lambda|>1\), and define
\(\mathfrak c_\sigma\) by the same formula.  Linear word-height bounds give
\(\mathfrak c_\sigma<\infty\), and every subsequent cancellation and
two-evaluation step is unchanged.  Thus the theorem also gives a
quantitative criterion for compactness of a nondistinguished archimedean realization.
The divisorial formulation was stated separately because it also applies
over function fields of varieties on which a prescribed collection of
trace collisions persists.
\end{remark}

We now determine when a single scalar trace fiber already admits a second,
independent trace functional that is constant on the fiber.  The result also
identifies the exceptional case in which the group structure or a valuation
must provide additional information.

\begin{proposition}[linear observables constant on one squared-trace fiber]
\label{prop:signed-collision-span}
Let \(k\) be a field of characteristic zero, let \(A\in\SL_2(k)\), and
let \(\mathcal E\subset\SL_2(k)\) be a nonempty set on which
\(X\mapsto\Tr^2(AX)\) has one nonzero value.  Fix
\(X_0\in\mathcal E\).  For each \(X\in\mathcal E\), choose the unique
sign \(\epsilon_X\in\{1,-1\}\) for which
\[
 \Tr(A\widehat X)=\Tr(AX_0),
 \qquad \widehat X=\epsilon_XX,
\]
and take \(\epsilon_{X_0}=1\).  Put
\[
 W_{\mathcal E}=
 \Span_k\{\widehat X-X_0:X\in\mathcal E\}\subset M_2(k).
\]
For the nondegenerate trace pairing
\(\langle B,C\rangle=\Tr(BC)\), one has
\[
 W_{\mathcal E}\subset A^\perp,
 \qquad
 \{B:\Tr(B\widehat X)\text{ is independent of }X\in\mathcal E\}
 =W_{\mathcal E}^\perp.
\tag{13.10s}
\]
Consequently, a linear trace observable independent of the original one
exists exactly when \(\dim W_{\mathcal E}\le2\).  If
\(\dim W_{\mathcal E}=3\), then
\[
 W_{\mathcal E}=A^\perp,
 \qquad W_{\mathcal E}^\perp=kA,
\]
so every constant linear trace observable is a scalar multiple of the
original one.

Assume \(\dim W_{\mathcal E}\le2\).  Exactly one of the following occurs.
\begin{enumerate}[label=\textup{(\roman*)}]
\item There is \(B\in W_{\mathcal E}^\perp\) with \(\det B\ne0\) and
\[
 \Delta(BA^{-1})
 :=\Tr^2(BA^{-1})-4\det(BA^{-1})\ne0.
\tag{13.10t}
\]
After adjoining \(\sqrt{\det B}\) and rescaling \(B\), the set
\(\mathcal E\) is contained in a simultaneous squared-trace fiber for
two matrices in \(\SL_2\) whose relative matrix has two distinct
eigenvalues.
\item One has \(\dim W_{\mathcal E}=2\) and
\[
 \{BA^{-1}:B\in W_{\mathcal E}^\perp\}
   =kI\oplus kN
\tag{13.10u}
\]
for a nonzero nilpotent matrix \(N\) with \(N^2=0\).  This is the only
exception.
\end{enumerate}
In case \textup{(ii)}, the matrices \(Y_X=A\widehat X\) have both
\(\Tr Y_X\) and \(\Tr(NY_X)\) fixed.  After extending \(k\), if
necessary, and conjugating so that \(N=E_{12}\), write
\(c=\Tr(NY_X)=(Y_X)_{21}\).  If \(c\ne0\), all the \(Y_X\) lie on the
single unipotent conjugacy orbit
\[
 \{(I+sN)Y_{X_0}(I-sN):s\in k\}.
\tag{13.10v}
\]
If \(c=0\), the corresponding curve is the union of at most two affine
lines; two distinct points on either line have relative trace \(2\).
\end{proposition}

\begin{proof}
The signs give \(\Tr(A(\widehat X-X_0))=0\), hence
\(W_{\mathcal E}\subset A^\perp\).  A matrix \(B\) has constant trace on
the signed set precisely when it annihilates every difference
\(\widehat X-X_0\), which proves (13.10s).  Nondegeneracy of the trace
pairing gives
\(\dim W_{\mathcal E}^\perp=4-\dim W_{\mathcal E}\), proving the first
claims.

Put
\[
 U=\{BA^{-1}:B\in W_{\mathcal E}^\perp\}\subset M_2(k).
\]
This subspace contains \(I\).  The discriminant
\(\Delta(M)=\Tr^2M-4\det M\) is unchanged upon adding a scalar matrix, so
it descends to the nondegenerate quadratic form on
\(M_2(k)/kI\simeq\mathfrak{sl}_2(k)\).  A totally isotropic subspace for
a nondegenerate quadratic form in three variables has dimension at most
one.  Thus, if \(\dim U\ge3\), some \(M\in U\) has
\(\Delta(M)\ne0\).  Replacing it by \(M+\alpha I\), for a suitable
\(\alpha\in k\), makes it invertible without changing its discriminant,
and \(B=MA\) proves \textup{(i)}.  If \(\dim U=2\), the same conclusion
holds unless \(U/kI\) is an isotropic line.  In the latter case,
subtracting the repeated eigenvalue of a nonscalar generator gives a
nonzero matrix \(N\) with \(N^2=0\), proving (13.10u).

In the exceptional case, \(NA\in W_{\mathcal E}^\perp\), so
\(\Tr(NY_X)\) is fixed.  With \(N=E_{12}\), every such matrix has the form
\[
 Y(u)=
 \begin{pmatrix}
  u&\bigl(u(t-u)-1\bigr)/c\\ c&t-u
 \end{pmatrix},
 \qquad t=\Tr Y_X,
\]
when \(c\ne0\).  Conjugation by \(I+sN\) replaces \(u\) by \(u+sc\),
which proves (13.10v).  If \(c=0\), the determinant equation is
\(u(t-u)=1\); each of its at most two roots leaves one off-diagonal entry
free.  Two points on the same resulting line have singular difference,
and
\(\det(P-Q)=2-\Tr(PQ^{-1})\) gives relative trace \(2\).
\end{proof}

\begin{theorem}[subexponential fibers unless the signed differences fill the trace hyperplane]
\label{thm:collision-span-fiber-bound}
Let \(H\) be finitely generated, let \(F/\Q\) be finitely generated of
characteristic zero, and let \(\rho\colon H\to\SL_2(F)\) satisfy the
Fuchsian specialization hypothesis of
Theorem~\ref{thm:two-trace-evaluations}.  Use the finite-type filtration of
Section~12.  Let \(A_n\in\SL_2(F)\) have complexity \(O(n)\), and let
\(\mathcal E_n\) be contained in the word ball of radius \(n\) in \(H\)
and in one nonzero squared-trace fiber of
\[
 h\longmapsto\Tr^2(A_n\rho(h)).
\]
Form the signed span \(W_{\mathcal E_n}\) from the matrices
\(\rho(h)\) as in Proposition~\ref{prop:signed-collision-span}.  If
\(\dim W_{\mathcal E_n}\le2\) and alternative \textup{(ii)} of that
proposition does not occur, then
\[
 |\mathcal E_n|=\exp(o(n)),
\]
uniformly over the fiber.  Hence an exponentially large squared-trace
fiber must either satisfy
\[
 \Span_F\{\epsilon_h\rho(h)-\epsilon_{h_0}\rho(h_0):
                 h\in\mathcal E_n\}=A_n^\perp
\tag{13.10w}
\]
or lie in the explicitly described unipotent-conjugation exception.
\end{theorem}

\begin{proof}
Choose at most two signed differences which form a basis of
\(W_{\mathcal E_n}\).  A basis of their common annihilator can be written
using minors of this fixed-size linear system; no division is needed, and
all entries have finite-type complexity \(O(n)\).  Unless the exceptional
alternative holds, the discriminant is not identically zero on this
annihilator after right multiplication by \(A_n^{-1}\).  One basis vector
or the sum of two basis vectors therefore has nonzero discriminant: if the
quadratic form vanished on every basis vector and every such sum, its
polarization would vanish on the entire space.  We obtain
\(B_n\in W_{\mathcal E_n}^\perp\) such that
\[
 \Delta(B_nA_n^{-1})\ne0
\]
and all its entries have complexity \(O(n)\).

Put
\[
 Y=A_n\widehat X,\qquad M=B_nA_n^{-1},\qquad
 t=\Tr Y,\quad u=\Tr(MY),
\]
and take traceless parts
\[
 Z=Y-\frac t2I,
 \qquad S=M-\frac{\Tr M}{2}I.
\]
Set
\[
 q=\frac{t^2}{4}-1,\qquad
 d=-\det S=\frac{\Delta(M)}4,\qquad
 r=u-\frac{t\Tr M}{2}.
\]
Then \(d\ne0\), \(Z^2=qI\), \(S^2=dI\), and \(\Tr(SZ)=r\).  Therefore
\[
 V=2dZ-rS
\]
satisfies
\[
 \Tr(SV)=0,
 \qquad \frac12\Tr(V^2)=d(4dq-r^2)=:\mathcal N.
\tag{13.10x}
\]
All entries and scalars in these formulas have complexity \(O(n)\).

Write
\[
 S=\begin{pmatrix}a&b\\c&-a\end{pmatrix},
 \qquad
 V=\begin{pmatrix}x&y\\z&-x\end{pmatrix},
 \qquad d=a^2+bc.
\]
If \(b\ne0\), the equation \(\Tr(SV)=0\) and (13.10x) give the injective
assignment
\[
 V\longmapsto(bx-ay,y),
 \qquad (bx-ay)^2-dy^2=b^2\mathcal N.
\]
If \(b=0\) and \(c\ne0\), use instead
\[
 V\longmapsto(cx-az,z),
 \qquad (cx-az)^2-dz^2=c^2\mathcal N.
\]
If \(b=c=0\), then \(x=0\) and (13.10x) is the factorization equation
\(yz=\mathcal N\).  When \(\mathcal N\ne0\),
Lemma~\ref{lem:quadratic-norm-divisor}, or
Theorem~\ref{thm:finite-type-divisor} in the last case, gives
\(\exp(o(n))\) possibilities.  When \(\mathcal N=0\), after a quadratic
extension the solution is the union of at most two affine lines.  Two
distinct group matrices on one line have singular difference and hence
relative squared trace \(4\).  The Fuchsian specialization hypothesis
makes them the same projective group element.  Thus the degenerate case
contributes only \(O(1)\) points.  The two possible signs change the bound
by only a constant factor.
\end{proof}

The full-rank alternative in (13.10w) occurs generically in the scheme
defined by the scalar collision equations; it is not an artifact of the
proof.

\begin{proposition}[geometry of the universal scalar-collision scheme]
\label{prop:universal-collision-scheme}
Let \(k\) be a field of characteristic zero.  For \(m\ge1\), consider
\[
 \mathcal C_m=
 \{(A,X_0,\ldots,X_m)\in\SL_2^{m+2}:
   \Tr(AX_i)=\Tr(AX_0)\ (1\le i\le m)\}.
\]
Then \(\mathcal C_m\) is an integral complete intersection of codimension
\(m\).  Its projection to \((A,X_0)\) is flat, and its fiber is the
\(m\)-fold product of the geometrically integral affine quadric
\[
 Q_{A,t}=\{X\in\SL_2:\Tr(AX)=t\},
 \qquad t=\Tr(AX_0).
\]
The affine span of \(Q_{A,t}\) is exactly the hyperplane
\(\{X\in M_2:\Tr(AX)=t\}\).

For \(m\ge3\), at the generic point of \(\mathcal C_m\),
\[
 \Span\{X_i-X_0:1\le i\le m\}=A^\perp.
\tag{13.10y}
\]
Consequently, even if \(B\) is allowed to have coefficients in the full
function field \(k(\mathcal C_m)\),
\[
 \Tr(BX_i)=\Tr(BX_0)\quad(1\le i\le m)
 \quad\Longrightarrow\quad B\in k(\mathcal C_m)A.
\tag{13.10z}
\]
The same conclusion holds componentwise for fixed-sign components of the
equations \(\Tr^2(AX_i)=\Tr^2(AX_0)\); on the open set where the common
trace is nonzero, the signs are unique.
\end{proposition}

\begin{proof}
Left multiplication by \(A\) identifies \(Q_{A,t}\) with
\(Q_t=\{Y\in\SL_2:\Tr Y=t\}\).  Writing
\(Y=(\begin{smallmatrix}u&v\\w&t-u\end{smallmatrix})\) gives
\[
 k[Q_t]=k[u,v,w]/\bigl(u(t-u)-vw-1\bigr).
\]
After any field extension, the displayed polynomial is primitive in
\(k[u,w][v]\), linear and irreducible over \(k(u,w)\), so Gauss's lemma
proves geometric integrality, including when \(t^2=4\).  Its principal
ideal contains no nonzero affine-linear polynomial.  Hence \(Q_t\)
affinely spans the entire trace hyperplane.

The trace morphism \(\SL_2\to\mathbb A^1\) is flat because its coordinate
ring is torsion-free over the principal ideal domain \(k[t]\).  The
projection in the statement is its \(m\)-fold base change and has
geometrically integral fibers over the integral base \(\SL_2^2\); its
total space is therefore integral.  Every fiber has dimension \(2m\), so
\(\mathcal C_m\) has codimension \(m\) in the smooth variety
\(\SL_2^{m+2}\).  The displayed \(m\) equations consequently form a
regular sequence.

Because the quadric affinely spans its trace hyperplane, three differences
from four generic points span its three-dimensional translation space
\(A^\perp\).  Linear independence is a nonempty open condition, which
proves (13.10y).  The trace pairing is nondegenerate, so the annihilator of
this space is the line spanned by \(A\), proving (13.10z).
\end{proof}

The proposition shows that the scalar collision equations alone cannot
imply the desired second trace relation.  It does not assert that a
full-rank fiber of exponential size occurs inside a word ball of a Fuchsian
group.  To exclude that possibility one must use the group law or arithmetic
information not present in the universal collision scheme.  In
Theorem~\ref{thm:persistent-local-expansion-bound}, that information comes
from reduced divisors and congruence expansion applied to a varying family
of group elements.  Differentiating a collision gives
\[
 \Tr\bigl((dA)(X_i-X_0)\bigr)
 +\Tr\bigl(A(dX_i-dX_0)\bigr)=0.
\]
For a deformation of a representation, the second term depends on the
group element and disappears only when the restricted cocycle is a
coboundary removed by infinitesimal conjugation, exactly the vertical
situation treated in Theorem~\ref{thm:restriction-unramified}.  Hasse
derivatives leave two free matrix directions at every order on the generic
smooth trace quadric.  Likewise, an identity in an associated graded or Rees
algebra gives a second identity in the original fiber only if it can be
lifted or separated by a valuation.  Theorem
\ref{thm:moving-divisorial-multiplier} provides such a valuation argument.

Finally, the square-root construction itself cannot evade this linear-span
alternative by producing a new conjugation-covariant matrix on the same
fiber.

\begin{proposition}[rational conjugation covariants in rank one]
\label{prop:rank-one-rational-covariants}
Let \(k\) be algebraically closed of characteristic zero.  Every rational
conjugation-equivariant map
\[
 \Phi\colon\SL_2\dashrightarrow M_2
\]
has the form
\[
 \Phi(g)=\alpha(\Tr g)I+\beta(\Tr g)g,
 \qquad \alpha,\beta\in k(t).
\tag{13.10aa}
\]
The same statement holds after a Kummer extension of the
adjoint-quotient field \(k(t)\), with \(\alpha\) and \(\beta\) in that
extension.  Consequently, on a fixed trace quadric every such construction
is affine-linear in \(g\) and preserves the full-affine-span obstruction of
Proposition~\ref{prop:signed-collision-span}.  In particular,
\[
 R(g)=\frac{I+g}{\sqrt{2+\Tr g}}
\]
does not by itself produce a second linear trace observable on one of its
scalar trace fibers.
\end{proposition}

\begin{proof}
For regular semisimple \(g\), equivariance makes \(\Phi(g)\) fixed by the
centralizer of \(g\).  The commutant of that torus in \(M_2\) is the
quadratic algebra \(k[g]=kI\oplus kg\).  Thus
\(\Phi(g)=\alpha(g)I+\beta(g)g\) on a dense open set.  Uniqueness of the
two coefficients and equivariance make them conjugation-invariant, while
the rational invariant field is
\(k(\SL_2)^{\PGL_2}=k(\Tr g)\).  This proves (13.10aa).  The same argument
after base change proves the Kummer assertion.
\end{proof}

Passing to a finite quotient and a vector-valued representation has a
similar limitation.  If \(q\colon\Pi\to Q\) is finite and
\(\sigma\colon Q\to\GL(V)\), then
\[
 \Tr((\rho\otimes\sigma)(g))
   =\Tr\rho(g)\,\chi_\sigma(q(g)).
\]
On a fixed \(q\)-coset this is only a fixed scalar multiple of the original
trace.  The partial trace over the two-dimensional factor is
\(\Tr\rho(g)\sigma(q(g))\), again a fixed matrix multiple on that coset.
Thus deck-character twists, Kummer sign characters, and the regular
representation of a finite quotient do not by themselves produce a second
left multiplier.  To obtain one, a vector-valued argument must use either a
different representation of the subgroup or the valuation separation used
above.

\subsection{The trace recurrence and its height at all places}

There is another way to retain more than one scalar trace without adding a
new representation.  Fix one hyperbolic element \(k\), and record the two
coordinates
\[
 \bigl(\Tr g,\Tr(kg)\bigr).
\]
All later coordinates \(\Tr(k^ng)\) are obtained from these two by one
linear recurrence.  The expansion rate of that recurrence at all places is
a global height, and not merely the translation length at the distinguished
real place.

\begin{lemma}[the two-trace recurrence]
\label{lem:two-trace-recurrence}
Let \(k,g\in\SL_2(\C)\) and put \(s=\Tr k\).  Define
\[
 q_0=0,\qquad q_1=1,\qquad q_{n+1}=s q_n-q_{n-1}.
\]
Then, for every \(n\ge1\),
\[
 k^n=q_n(s)k-q_{n-1}(s)I,
 \qquad
 \Tr(k^ng)=q_n(s)\Tr(kg)-q_{n-1}(s)\Tr(g).
\tag{13.10am}
\]
If
\[
 \alpha_n=-\frac{q_{n-1}(s)}{q_n(s)},\qquad
 \beta_n=-\alpha_n^2,
\tag{13.10an}
\]
then \(\beta_n\in\Q(s^2)\).  Moreover
\[
 q_n(s)^2-q_{n-1}(s)q_{n+1}(s)=1.
\tag{13.10ao}
\]
\end{lemma}

\begin{proof}
Cayley--Hamilton gives \(k^2-sk+I=0\).  Induction proves the matrix
identity in (13.10am), and taking the trace after multiplication by \(g\)
proves the scalar identity.  The polynomial \(q_n\) has the parity of
\(n-1\), so the square of \(q_{n-1}/q_n\) is rational in \(s^2\).
The left side of (13.10ao) is constant in \(n\) by the recurrence and is
equal to one at \(n=1\).
\end{proof}

\begin{proposition}[all cyclic trace vectors have one determinantal modulus]
\label{prop:cyclic-trace-determinantal-ideal}
Put \(q_0(T)=0\), \(q_1(T)=1\), and
\(q_{j+1}(T)=Tq_j(T)-q_{j-1}(T)\) for every \(j\in\mathbb Z\);
thus \(q_{-j}=-q_j\).  Let \(J\subset\mathbb Z\) be finite with at least
two elements, and put
\[
 C_J(T)=\bigl(-q_{j-1}(T),q_j(T)\bigr)_{j\in J}
 \in M_{J\times2}(\mathbb Z[T]).
\]
If
\[
 g_J=\gcd\{|i-j|:i,j\in J,\ i\ne j\},
\]
then the ideal generated by the maximal minors of \(C_J(T)\) is exactly
\[
 I_2(C_J)=(q_{g_J}(T))\subset\mathbb Z[T].
\tag{13.10ao1}
\]
This equality survives every specialization \(T\mapsto s\) to a
commutative ring.  In particular, if \(K\) is a commutative ring,
\(A\in\SL_2(K)\),
\(s=\Tr A\), and \(X\in M_2(K)\), then
\[
 \Tr(A^jX)
 =-q_{j-1}(s)\Tr X+q_j(s)\Tr(AX),
\tag{13.10ao2}
\]
and the common determinantal divisor of any finite vector of such trace
coordinates is the single Chebyshev term \(q_{g_J}(s)\).
\end{proposition}

\begin{proof}
Cayley--Hamilton gives
\(A^j=q_j(s)A-q_{j-1}(s)I\) for \(j\ge0\).
The identity \(A^{-1}=sI-A\), or the same recurrence run backwards,
proves this for every \(j\in\mathbb Z\) and gives \(q_{-j}=-q_j\).
For any \(i,j\in\mathbb Z\), the corresponding minor is
\[
 q_iq_{j-1}-q_{i-1}q_j=q_{j-i}.
\tag{13.10ao3}
\]
The addition and subtraction identities
\[
 q_{u+v}=q_uq_{v+1}-q_{u-1}q_v,\qquad
 q_{u-v}=q_uq_{v+1}-q_{u+1}q_v
\tag{13.10ao4}
\]
hold in \(\mathbb Z[T]\).  Consequently
\((q_u,q_v)=(q_v,q_{u-v})\) for \(u\ge v>0\).
The Euclidean algorithm on the indices gives
\[
 (q_u,q_v)=(q_{\gcd(u,v)})
\]
without localization or exceptional primes.  Apply this repeatedly to the
differences \(i-j\) in (13.10ao3).  Since the resulting ideal equality is
already over \(\mathbb Z[T]\), evaluation in any commutative ring preserves
it.
\end{proof}

Thus adding finitely many positions in one orbit of the cyclic group
\(\langle A\rangle\) never creates independent recurrence moduli.  If the
index differences generate \(\mathbb Z\), the determinantal ideal is the
unit ideal.  If the indices form a primitive congruence pattern modulo
\(n\), meaning that the greatest common divisor of their differences is
exactly \(n\), the common modulus is exactly the same \(q_n(s)\).
Finite-character twists multiply these rows by scalars from a fixed finite
set.  After zero coordinates are discarded and the fixed coefficient field
and bad-prime set are enlarged, they introduce only factors independent of
the recurrence indices; they do not produce a new growing Chebyshev modulus.
Hence a finite vector of cyclic trace coordinates cannot replace the radical
of one Lucas--Chebyshev term by a product of independent factors.  To obtain
independent growing moduli, one must use either a noncommuting group element
or congruence information modulo the full prime powers.

\begin{proposition}[the first noncommuting recurrence determinant]
\label{prop:noncommuting-recurrence-determinant}
Let \(K\) have characteristic zero, let \(a,b\in\SL_2(K)\), and assume
that \(a\) is regular semisimple and that \(a,b\) have no common invariant
line over \(\overline K\).  Then
\[
 \{I,a,b,ab\}
\]
is a basis of \(M_2(K)\).  Put
\[
 Q_m=q_m(\Tr a),\qquad R_n=q_n(\Tr b).
\]
Relative to any fixed coordinates on \(M_2(K)\),
\[
 \det(I,a^m,b^n,a^mb^n)
 =Q_m^2R_n^2\det(I,a,b,ab).
\tag{13.10ao5}
\]
Equivalently, the four trace functionals
\[
 X\longmapsto
 \bigl(\Tr X,\Tr(a^mX),\Tr(b^nX),\Tr(a^mb^nX)\bigr)
\]
have determinant equal to a fixed nonzero factor times \(Q_m^2R_n^2\).
\end{proposition}

\begin{proof}
Write \(P_m=q_{m-1}(\Tr a)\) and
\(S_n=q_{n-1}(\Tr b)\).  Cayley--Hamilton gives
\[
 a^m=Q_ma-P_mI,\qquad b^n=R_nb-S_nI.
\]
In the ordered basis \((I,a,b,ab)\), the change-of-basis matrix to
\((I,a^m,b^n,a^mb^n)\) is triangular, with diagonal
\[
 1,\quad Q_m,\quad R_n,\quad Q_mR_n.
\]
This proves (13.10ao5).  To see that the fixed determinant is nonzero,
diagonalize \(a\) over \(\overline K\), say
\(a=\operatorname{diag}(\lambda,\lambda^{-1})\), and write the
off-diagonal entries of \(b\) as \(u,v\).  Direct expansion gives the
fixed determinant, up to the sign imposed by the coordinate order, as
\[
 uv(\lambda-\lambda^{-1})^2
 =2-\Tr[a,b].
\]
Regularity makes the last eigenvalue factor nonzero, and absence of a
common invariant line makes \(uv\ne0\).  Finally, the trace pairing on
\(M_2(K)\) is nondegenerate, so the determinant of the four trace
functionals differs from the matrix-vector determinant by one fixed
nonzero factor.
\end{proof}

Proposition~\ref{prop:noncommuting-recurrence-determinant} identifies the
first point at which two independent recurrence moduli occur.  The
determinant identity alone does not give a useful count.  One group ball has
\(e^{R+o(R)}\) elements, whereas four scalar trace supports of size
\(e^{R/2+o(R)}\) form an ambient grid of size \(e^{2R+o(R)}\).  A direct
Cartesian-support comparison therefore loses a factor \(e^{R-o(R)}\).
Thus a direct comparison with the Cartesian product is too weak by a factor
of \(e^{R-o(R)}\).  Any application of the determinant identity must use
additional incidence relations among the four trace coordinates.

\begin{lemma}[the recurrence dilates eventually generate the squared-trace field]
\label{lem:recurrence-field-stabilization}
Let \(s\) be algebraic, suppose that one real embedding of \(s^2\) is
greater than \(4\), and put
\[
 K=\Q(s^2),\qquad
 \beta_n=-\frac{q_{n-1}(s)^2}{q_n(s)^2}\quad(n\ge2).
\]
Then
\[
 \Q(\beta_n)=K
\]
for all but finitely many \(n\).
\end{lemma}

\begin{proof}
Put \(u=s^2\).  If \(z+z^{-1}=u-2\), the closed formula for the recurrence
gives
\[
 \beta_n=f_n(u)
 =-z\left(\frac{z^{n-1}-1}{z^n-1}\right)^2.
\tag{13.10ap}
\]
The expression is invariant under \(z\mapsto z^{-1}\), and is therefore
rational in \(u\).

A number field has only finitely many subfields.  If the assertion failed,
one proper subfield \(F\subsetneq K\) would contain \(f_n(u)\) for
infinitely many \(n\).  Fix the distinguished embedding
\(\sigma_0:K\hookrightarrow\C\), for which \(\sigma_0(u)>4\).
The restriction \(\sigma_0|_F\) has two distinct extensions
\(\sigma_0,\tau:K\hookrightarrow\C\).  Hence
\[
 f_n(\sigma_0u)=f_n(\tau u)
\tag{13.10aq}
\]
for infinitely many \(n\).

Choose \(z_0>1\) with
\(z_0+z_0^{-1}=\sigma_0u-2\), and choose \(w\) with
\(|w|\ge1\) and \(w+w^{-1}=\tau u-2\).  No denominator in
(13.10ap) vanishes.  Indeed, \(q_n(s)\) is a nonzero algebraic number
because its distinguished conjugate is positive.  Substitution in
(13.10aq) and clearing denominators gives the exponential polynomial
\[
 z_0(z_0^{n-1}-1)^2(w^n-1)^2
 -w(w^{n-1}-1)^2(z_0^n-1)^2.
\tag{13.10ar}
\]
By the Skolem--Mahler--Lech theorem \cite{Lech}, infinitely many zeros
force (13.10ar) to vanish for every \(n\) in an infinite arithmetic
progression.

If \(|w|>1\), passage to the limit along that progression gives
\(-z_0^{-1}=-w^{-1}\).  Thus \(w=z_0\) and
\(\tau(u)=\sigma_0(u)\).  Since \(u\) generates \(K\), this contradicts
\(\tau\ne\sigma_0\).

It remains to consider \(|w|=1\).  If \(w\ne1\) is a root of unity, then
\(w^m=1\) for some \(m\), and the conjugate \(q_m(\tau s)\) vanishes,
again contradicting \(q_m(s)\ne0\).  If \(w=1\), then
\(\tau(u)=4\); irreducibility of the minimal polynomial of \(u\) would
force \(u=4\), contrary to \(\sigma_0(u)>4\).  Thus \(w\) is not a root
of unity.  On a progression \(n=a+mj\), the powers \(w^{a+mj}\) are dense
in the unit circle, whereas \(f_{a+mj}(\sigma_0u)\to-z_0^{-1}\).
But
\[
 f_n(\tau u)
 =-\frac1w\left(\frac{w^n-w}{w^n-1}\right)^2.
\]
Subsequences converging away from the pole at \(1\) would make the
nonconstant rational function
\[
 t\longmapsto-\frac1w\left(\frac{t-w}{t-1}\right)^2
\]
constant on the unit circle, a contradiction.
\end{proof}

For an algebraic number \(\xi\), put \(F=\Q(\xi)\) and define the
Conlon--Lim height
\[
 \mathsf H(\xi)=
 N_{F/\Q}(\mathfrak D_{\xi;F})
 \prod_{\sigma:F\hookrightarrow\C}(1+|\sigma\xi|),
 \qquad
 \mathfrak D_{\xi;F}
 =\{a\in\mathcal O_F:a\xi\in\mathcal O_F\}.
\tag{13.10as}
\]
Every complex embedding occurs separately in the product.  Thus, for
\(\xi=p/q\in\Q\) in lowest terms,
\(\mathsf H(\xi)=|p|+|q|\).  Conlon and Lim prove that every finite
\(A\subset\C\) satisfies
\[
 |A+\xi A|\ge
 \bigl(\mathsf H(\xi)-o_\xi(1)\bigr)|A|
 \qquad(|A|\to\infty).
\tag{13.10at}
\]
Their theorem is sharp \cite{ConlonLimAlgebraicDilates}.

We now compute this height for the recurrence dilates.  For a place \(v\)
of \(K=\Q(s^2)\), choose a root \(z_v\) of
\[
 z_v+z_v^{-1}=s^2-2
\]
in an algebraic closure of \(K_v\), and set
\[
 \chi_v(s)=\frac12
 \log\max\{|z_v|_v,|z_v|_v^{-1}\},\qquad
 \mathcal J_K(s)=\sum_{v\in M_K}[K_v:\Q_v]\chi_v(s).
\tag{13.10au}
\]
These quantities are independent of the choice \(z_v\leftrightarrow
z_v^{-1}\).  At the distinguished real place,
\(\chi_{v_0}(s)=\ell(k)/2\).  The sum \(\mathcal J_K(s)\) measures the
growth, over all places, of the recurrence determined by \(s^2-2\);
equivalently, it is its Mahler measure in logarithmic form.  The
distinguished translation length accounts for only one of its nonnegative
summands.

\begin{lemma}[global height of the recurrence dilates]
\label{lem:recurrence-height}
With the preceding notation,
\[
 h(\alpha_n)
 =\frac{n}{[K:\Q]}\mathcal J_K(s)+O_s(\log n).
\tag{13.10av}
\]
Whenever \(\Q(\beta_n)=K\),
\[
 \log\mathsf H(\beta_n)
 =2n\mathcal J_K(s)+O_s(\log n).
\tag{13.10aw}
\]
\end{lemma}

\begin{proof}
Let \(L=\Q(s)\).  For each place \(w\) of \(L\), choose
\(\lambda_w+\lambda_w^{-1}=s\).  The closed formula
\[
 q_n(s)=\frac{\lambda_w^n-\lambda_w^{-n}}
              {\lambda_w-\lambda_w^{-1}}
\]
and (13.10ao) give
\[
 \log\max\{|q_{n-1}(s)|_w,|q_n(s)|_w\}
 =n\chi_w(s)+O_{s,w}(\log n).
\tag{13.10ax}
\]
At a nonarchimedean place, if \(|s|_w>1\), then
\(|q_n(s)|_w=|s|_w^{n-1}\).  If \(|s|_w\le1\), the recurrence bounds both
adjacent terms by one and (13.10ao) prevents both from having norm less
than one.  At an archimedean place the closed formula proves
(13.10ax); a logarithmic error is needed only at a repeated root.

The projective-height formula, followed by summation of local degrees over
the places of \(L\) above a place of \(K\), proves (13.10av).  Also
\(h(\beta_n)=2h(\alpha_n)\).  If \(K=\Q(\beta)\), then
\[
 \log\mathsf H(\beta)
 =\sum_{v\nmid\infty}[K_v:\Q_v]\log^+|\beta|_v
  +\sum_{v\mid\infty}[K_v:\Q_v]\log(1+|\beta|_v),
\]
and hence
\[
 0\le
 \log\mathsf H(\beta)-[K:\Q]h(\beta)
 \le [K:\Q]\log2.
\]
Together with (13.10av), this proves (13.10aw).
\end{proof}

The recurrence becomes useful only after one deals with a combinatorial
point which is absent from the usual sum-of-dilates problem.  The pairs
\((\Tr g,\Tr(kg))\) need not contain a large Cartesian product.  The next
two lemmas show that a dense bipartite graph of such pairs is already
enough.  We include the details because this is where an unjustified
``dense graph contains a large rectangle'' assertion would otherwise enter.

\begin{lemma}[a difference-set bound on a bipartite graph]
\label{lem:graph-ruzsa}
Let \(U,V\) be finite subsets of a field, let
\(G\subset U\times V\), and let \(\alpha\ne0\).  For
\(u\in U\), write
\[
 N_G(u)=\{v\in V:(u,v)\in G\}.
\]
For every \(d\in U-\alpha^2U\), choose one representation
\(d=u_d-\alpha^2u'_d\), and put
\(b(d)=\#(N_G(u_d)\cap N_G(u'_d))\).  Then
\[
 \sum_{d\in U-\alpha^2U}b(d)
 \le
 \#\{u+\alpha v:(u,v)\in G\}
 \#\{v+\alpha u:(u,v)\in G\}.
\tag{13.10ay}
\]
In particular, if every two vertices of \(U\), including two equal
vertices, have at least \(r\) common neighbours, then
\[
 \#(U-\alpha^2U)
 \le \frac{1}{r}
 \#\{u+\alpha v:(u,v)\in G\}
 \#\{v+\alpha u:(u,v)\in G\}.
\tag{13.10az}
\]
\end{lemma}

\begin{proof}
For each \(d\) and each
\(v\in N_G(u_d)\cap N_G(u'_d)\), consider
\[
 (d,v)\longmapsto
 (u_d+\alpha v,\ v+\alpha u'_d).
\]
This map is injective.  Indeed, its image \((x,y)\) first recovers
\(d=x-\alpha y\); the chosen representative \((u_d,u'_d)\) is then
known, and \(v=(x-u_d)/\alpha\).  Counting its domain and codomain proves
(13.10ay), and (13.10az) follows immediately.
\end{proof}

\begin{lemma}[dependent random choice with an explicit loss]
\label{lem:drc-explicit}
Let \(G\subset U\times V\) be a bipartite graph, where
\(\#U,\#V\le N\), and write
\[
 \#G=N^2e^{-a},\qquad a\ge0.
\]
For every positive integer \(t\), there is \(W\subset U\) such that
\[
 \#W\ge\frac12Ne^{-at}
\tag{13.10ba}
\]
and every two vertices of \(W\), including two equal vertices, have at
least
\[
 r=\left\lfloor e^{-a}N(2N)^{-1/t}\right\rfloor
\tag{13.10bb}
\]
common neighbours.
\end{lemma}

\begin{proof}
Pad the two vertex classes by isolated vertices so that both have size
\(N\).  Choose \(t\) vertices of \(V\), independently and with
replacement, and let \(X\) be their common neighbourhood.  Convexity gives
\[
 \mathbb E\#X
 =\sum_{u\in U}\left(\frac{\deg u}{N}\right)^t
 \ge N\left(\frac{\#G}{N^2}\right)^t
 =Ne^{-at}.
\]
Call an unordered pair, allowing a repeated vertex, bad if its common
neighbourhood has size less than \(r\).  The expected number \(B\) of bad
pairs contained in \(X\) is at most
\[
 N^2(r/N)^t\le\frac12Ne^{-at}.
\]
Some choice therefore has \(\#X-B\ge Ne^{-at}/2\).  Repeatedly remove one
vertex from each remaining bad pair.  At most \(B\) vertices are removed,
and the surviving set \(W\) has (13.10ba) and (13.10bb).
\end{proof}

\begin{theorem}[large trace sets with simultaneous recurrence control]
\label{thm:drc-trace-recurrence}
Let \(\Gamma<\PSL_2(\R)\) be a torsion-free cocompact lattice and suppose
that its signed trace set has linear growth.  Choose a lift
\(\widehat\Gamma<\SL_2(\R)\), a hyperbolic
\(k\in\widehat\Gamma\) with \(s=\Tr k>2\), and a basepoint on the axis of
\(k\).  Let \(\ell=\ell(k)\), and let \(\beta_n\) be as in
(13.10an).  There are finite sets \(W_R\subset\R\) and numbers
\(\Delta_R=o(R)\), depending on \(k\) but not on \(n\), such that
\[
 \#W_R=\exp\bigl((\tfrac12-o(1))R\bigr)
\tag{13.10bc}
\]
and, simultaneously for every integer \(n\ge2\),
\[
 \#(W_R+\beta_nW_R)
 \le \exp(n\ell+\Delta_R)\#W_R.
\tag{13.10bd}
\]

More explicitly, define the simple bipartite graph
\[
 G_R=\{(\Tr g,\Tr(kg)):
          g\in\widehat\Gamma,\ d(o,go)\le R\}.
\tag{13.10be}
\]
For a sufficiently large fixed \(C\), put
\[
 N_R=\lceil Ce^{R/2}\rceil,
 \qquad a_R=\log\frac{N_R^2}{\#G_R},
 \qquad L_R=\log(2N_R).
\]
Then \(a_R=o(R)\).  If
\[
 \bar a_R=\max\{a_R,L_R^{-1}\},
 \qquad
 t_R=\left\lceil\sqrt{L_R/\bar a_R}\right\rceil,
\tag{13.10bf}
\]
one may take
\[
 \Delta_R
 \le a_R+a_Rt_R+L_R/t_R+O_{\Gamma,k}(1)
 =O\bigl(1+a_R+\sqrt{a_RR}\bigr)=o(R).
\tag{13.10bg}
\]
\end{theorem}

\begin{proof}
The two coordinate projections of \(G_R\) have at most \(N_R\) elements:
matrix norm estimates put both traces in an interval of size \(O(e^{R/2})\),
and the assumed linear trace growth bounds the number of values in that
interval.  Cocompact orbit counting gives \(e^{R+O(1)}\) group elements in
the ball.  Lemma~\ref{lem:ambient-two-trace-evaluations}, with
\(A_0=I\) and \(A_1=k\), bounds every fiber of the two trace coordinates
by \(e^{o(R)}\).  Hence
\[
 \#G_R=e^{R-o(R)},
 \qquad a_R=o(R).
\tag{13.10bh}
\]

Apply Lemma~\ref{lem:drc-explicit} with \(t=t_R\), and call the resulting
left vertex set \(W_R\).  Formula (13.10ba), together with
\(a_Rt_R=o(R)\), proves (13.10bc).  The quantity inside the floor in
(13.10bb) tends to infinity; consequently, after changing an absolute
constant,
\[
 r_R\ge
 \tfrac12 e^{-a_R}N_R(2N_R)^{-1/t_R}.
\tag{13.10bi}
\]

Fix \(n\ge2\), and put \(\alpha_n=-q_{n-1}(s)/q_n(s)\).  The recurrence
and its inverse give
\[
 \begin{aligned}
 q_n(s)(\Tr g+\alpha_n\Tr(kg))
   &=\Tr(k^{-(n-1)}g),\\
 q_n(s)(\Tr(kg)+\alpha_n\Tr g)
   &=\Tr(k^ng).
 \end{aligned}
\tag{13.10bj}
\]
The basepoint lies on the axis of \(k\), so the elements on the right have
displacement at most \(R+n\ell\).  Linear trace growth therefore bounds
each of the two graph projections in Lemma~\ref{lem:graph-ruzsa} by
\[
 O_{\Gamma,k}\bigl(e^{(R+n\ell)/2}\bigr).
\tag{13.10bk}
\]
Since \(\beta_n=-\alpha_n^2\), Lemma~\ref{lem:graph-ruzsa}, (13.10ba),
and (13.10bi) now give
\[
 \frac{\#(W_R+\beta_nW_R)}{\#W_R}
 \le
 \exp\bigl(n\ell+a_R+a_Rt_R+L_R/t_R+O_{\Gamma,k}(1)\bigr).
\]
This is (13.10bd)--(13.10bg).  Notice that \(W_R\) and \(\Delta_R\) were
chosen before \(n\); the estimate holds simultaneously for all
recurrence indices.
\end{proof}

Theorem~\ref{thm:drc-trace-recurrence} requires no large Cartesian product
inside the trace-pair graph.  It gives one large set on which every
recurrence dilate has a controlled sumset, with a loss of only
\(\exp(o(R))\).

\begin{theorem}[Chebyshev recurrence and Chinese-remainder uncertainty on a curve]
\label{thm:chebyshev-crt-curve}
Let \(C/k\) be a smooth projective geometrically integral curve over a
field of characteristic zero, let \(s\in k(C)\) be nonconstant, and let
\(D_0\) be an effective divisor containing \((s)_\infty\).  Put
\[
 q_0(T)=0,\quad q_1(T)=1,\quad
 q_{n+1}(T)=Tq_n(T)-q_{n-1}(T),\qquad
 \beta_n=-\frac{q_{n-1}(s)^2}{q_n(s)^2}.
\]
For \(n\ge2\), let \(A_n\) be the part of the zero divisor
\((q_n(s)^2)_0\) supported outside \(|D_0|\).  Let \(I\) be a finite set
of pairwise coprime integers at least \(2\).  Suppose that
\[
 U\subset\Gamma(C\setminus|D_0|,\mathcal O_C)
\]
is finite and nonempty and that an effective divisor \(D\), supported on
\(D_0\),
satisfies
\[
 (u-u')_\infty\le D\qquad(u,u'\in U).
\tag{13.10bh1}
\]
If
\[
 \sum_{n\in I}\deg A_n>\deg D,
\tag{13.10bh2}
\]
then
\[
 \#U\le
 \prod_{n\in I}\frac{\#(U+\beta_nU)}{\#U}.
\tag{13.10bh3}
\]
Writing \(d=\deg(s)_\infty\), the divisors \(A_n\), \(n\in I\), have
pairwise disjoint supports and
\[
 2(n-1)d-C_{D_0,s}\le\deg A_n\le2(n-1)d,
\tag{13.10bh4}
\]
where the constant is independent of \(n\).  One may take
\[
 C_{D_0,s}
 =2\sum_{\substack{P\in|D_0|\\s(P)\ne\infty}}
 e_P(s)\deg P,
\]
after extending the constant field if necessary; degrees are unchanged by
that extension.
\end{theorem}

\begin{proof}
As polynomials,
\[
 q_n(T)=U_{n-1}(T/2),\qquad
 \gcd(q_m,q_n)=q_{\gcd(m,n)}.
\tag{13.10bh5}
\]
Thus pairwise coprime \(m,n\) give coprime polynomials
\(q_m,q_n\), and their zero divisors after substitution \(T=s\) are
disjoint.  The polynomial \(q_n\) is monic of degree \(n-1\), so the full
zero divisor of \(q_n(s)^2\) has degree \(2(n-1)d\).
All roots of \(q_n\) are simple in characteristic zero.  At a fixed
point \(P\in|D_0|\) where \(s(P)\) is finite, a zero therefore has order
\(2\operatorname{ord}_P(s-s(P))\), independently of \(n\); at a pole of
\(s\) there is no zero.  Removing the fixed set \(|D_0|\) consequently
costs at most a constant, proving (13.10bh4).

Let \(\pi_n\) denote reduction modulo \(A_n\), and choose
\(T_n\subset U\) containing one representative of each class in
\(\pi_n(U)\).  Cassini's identity makes \(q_{n-1}(s)\) a unit modulo
\(A_n\).  The map
\[
 T_n\times U\longrightarrow q_n(s)^2(U+\beta_nU),\qquad
 (t,u)\longmapsto q_n(s)^2u-q_{n-1}(s)^2t
\tag{13.10bh6}
\]
is injective.  Indeed, reduction modulo \(A_n\) first recovers the class
of \(t\), hence \(t\) itself, and then equality in the function field
recovers \(u\).  Therefore
\[
 \frac{\#(U+\beta_nU)}{\#U}\ge\#\pi_n(U).
\tag{13.10bh7}
\]

The Chinese remainder theorem applies to the pairwise disjoint divisors
\(A_n\).  The joint reduction map
\[
 U\longrightarrow\prod_{n\in I}
 \Gamma(A_n,\mathcal O_{A_n})
\]
is injective: otherwise a nonzero difference \(u-u'\) would have at
least \(\sum_{n\in I}\deg A_n\) zeros, counted with multiplicity, but by
(13.10bh1) at most \(\deg D\) poles.  Hence
\(\#U\le\prod_{n\in I}\#\pi_n(U)\), and multiplication of (13.10bh7)
proves (13.10bh3).
\end{proof}

\begin{corollary}[entropy bound on an algebraic curve]
\label{cor:chebyshev-crt-entropy}
In the setting of Theorem~\ref{thm:chebyshev-crt-curve}, suppose that
finite sets \(U_R\) satisfy, for constants \(\sigma,c,H>0\),
\[
 \log\#U_R\ge(\sigma-o(1))R,
\tag{13.10bh8}
\]
that every difference of two elements of \(U_R\) has a common pole bound
of degree at most \((H+o(1))R\), and, uniformly for \(n=o(R)\),
\[
 \log\frac{\#(U_R+\beta_nU_R)}{\#U_R}
 \le cn+\Delta_R,\qquad \Delta_R=o(R).
\tag{13.10bh9}
\]
Then
\[
 \sigma\le\frac{cH}{2\deg(s)_\infty}.
\tag{13.10bh10}
\]
\end{corollary}

\begin{proof}
Choose \(N_R=o(R)\) so that
\[
 N_R\gg\Delta_R,\qquad
 N_R^2\gg R\log N_R.
\]
Write the degree of the common pole bound as at most \(HR+e_R\), where
\(e_R=o(R)\).  Choose \(\eta_R\downarrow0\) so slowly that
\[
 \eta_RR\gg |e_R|+R/N_R+N_R.
\]
Greedily select primes in \([N_R,2N_R]\) until
\[
 \sum_{n\in I_R}n
 \ge\frac{HR+\eta_RR}{2\deg(s)_\infty}.
\tag{13.10bh11}
\]
The overshoot is at most \(2N_R=o(R)\).  The prime number theorem gives
enough primes: the required number is \(O(R/N_R)\), while the interval
contains \(\gg N_R/\log N_R\) primes and its total prime sum is
\(\gg N_R^2/\log N_R\gg R\).  Moreover
\(\#I_R=O(R/N_R)\), so (13.10bh4) gives
\[
 \sum_{n\in I_R}\deg A_n
 \ge2\deg(s)_\infty\sum_{n\in I_R}n
     -O(R/N_R)>HR+e_R.
\]
Thus (13.10bh2) holds.  Apply
(13.10bh3) and (13.10bh9).  Since
\(\#I_R=O(R/N_R)\),
\[
 (\sigma-o(1))R
 \le c\sum_{n\in I_R}n+\#I_R\Delta_R
 =\frac{cHR}{2\deg(s)_\infty}+o(R).
\]
Division by \(R\) proves (13.10bh10).
\end{proof}

\begin{lemma}[sharpness of the function-field radical bound]
\label{lem:chebyshev-radical-curve}
In the notation of Theorem~\ref{thm:chebyshev-crt-curve}, let \(R_s\) be
the ramification divisor of the morphism
\(s:C\to\mathbb P^1\), and let \(B_n\) be the reduced part, outside
\(|D_0|\), of \((q_n(s))_0\).  Then
\[
 (n-1)d-\deg R_s-\deg(D_0)_{\mathrm{red}}
 \le\deg B_n\le(n-1)d,\qquad d=\deg(s)_\infty.
\tag{13.10bha}
\]
The leading constant is best possible: for
\(C=\mathbb P^1\), \(s=T\), and \(D_0=[\infty]\), one has
\(\deg B_n=n-1\) for every \(n\ge2\).
\end{lemma}

\begin{proof}
After extending the constant field, write
\[
 (q_n(s))_0=\sum_{q_n(a)=0}s^*[a].
\]
The \(n-1\) roots of \(q_n\) are distinct, so this divisor has degree
\((n-1)d\).  Passing to its reduced divisor loses
\(\sum(e_P(s)-1)\deg P\) over only some ramification points and hence at
most \(\deg R_s\).  Deleting the fixed support of \(D_0\) loses at most
\(\deg(D_0)_{\mathrm{red}}\).  Degrees are unchanged by the constant-field
extension.  In the stated example \(s\) is unramified and \(q_n(T)\) is
square-free.
\end{proof}

This is the leading term given by Mason--Stothers.  It cannot sharpen
(13.10bh10): Theorem~\ref{thm:chebyshev-crt-curve} already uses the
nonreduced divisor of \(q_n(s)^2\), whose degree is the exact maximum
\(2(n-1)d\).  Replacing it by the radical gives only
\((n-1)d+O(1)\), and the example in the lemma rules out a universal strict
improvement.  Thus the possible loss on a curve comes from the coefficients
and from the alignment of their poles, not from a weaker-than-necessary
function-field radical estimate.

For the sets given by Theorem~\ref{thm:drc-trace-recurrence}, the
natural values are \(\sigma=1/2\) and \(c=\ell_x(k)\).  The corollary then
gives the sharp necessary inequality
\[
 \deg(s)_\infty\le\ell_x(k)H.
\tag{13.10bh12}
\]
This is a product-formula estimate on a curve along which the prescribed
trace collisions persist, but it is not by itself a contradiction: the
common pole rate \(H\) can be governed by
different boundary valuations for different words.  A strict improvement
of the coefficient \(c\), together with one word nearly maximizing
\(\deg(s)_\infty/\ell_x(k)\) relative to that same pole bound, would exclude
the curve.  The near-maximizing condition is additional and is not implied
merely by taking a supremum over unrelated boundary valuations.

\begin{theorem}[a lower sumset bound would eliminate every nondistinguished expanding place]
\label{thm:moving-height-criterion}
Retain the hypotheses and sets of
Theorem~\ref{thm:drc-trace-recurrence}.  Suppose that \(s\) is algebraic,
put \(K=\Q(s^2)\), and choose integers \(n_R\to\infty\) such that
\[
 n_R=o(R),\qquad \Delta_R=o(n_R).
\tag{13.10bl}
\]
Such a choice always exists; for example one may take the integer part of
\(\sqrt{R(1+\Delta_R)}\), after an immaterial monotone adjustment.  Suppose
that the following lower bound holds along this sequence:
\[
 \log\frac{\#(W_R+\beta_{n_R}W_R)}{\#W_R}
 \ge \log\mathsf H(\beta_{n_R})-o(n_R).
\tag{13.10bm}
\]
Then
\[
 \mathcal J_K(s)=\frac{\ell(k)}2.
\tag{13.10bn}
\]
Equivalently, \(\chi_v(s)=0\) at every place of \(K\) other than the
distinguished real place.  In particular, \(s^2\) is an algebraic integer
and every nonidentity conjugate of \(s^2\) lies in \([0,4]\).
\end{theorem}

\begin{proof}
Lemma~\ref{lem:recurrence-field-stabilization} shows that
\(\Q(\beta_{n_R})=K\) for all sufficiently large \(R\).  Combining
(13.10bd), (13.10bm), and Lemma~\ref{lem:recurrence-height} gives
\[
 2n_R\mathcal J_K(s)+o(n_R)
 \le n_R\ell(k)+\Delta_R.
\]
Thus \(\mathcal J_K(s)\le\ell(k)/2\).  The distinguished real place alone
contributes \(\ell(k)/2\), and every local term is nonnegative, proving
(13.10bn).  At a finite place, \(\chi_v(s)=0\) says that the two roots of
\(z+z^{-1}=s^2-2\) are units; hence \(s^2\) is integral there.  At an
archimedean place it says that those roots have modulus one, so the
corresponding conjugate of \(s^2\) is real and belongs to \([0,4]\).
\end{proof}

One can instead count the distribution of the two trace coordinates modulo
the recurrence ideal.  This avoids both the loss from dependent random
choice and the need for a general sum-of-dilates theorem.

\begin{lemma}[counting trace pairs by congruence classes]
\label{lem:congruence-wedge}
Let \(R\) be a Dedekind domain, let \(P,Q\in R\setminus\{0\}\) generate
coprime ideals, and let \(G\subset U\times V\), where \(U,V\subset R\)
are finite.  Put
\[
 \begin{split}
 X_-&=\{Qu-Pv:(u,v)\in G\},\\
 X_+&=\{Qv-Pu:(u,v)\in G\},
 \end{split}
\]
and
\[
 M_{(PQ)}(V)=
 \max_{c\in R/(PQ)}\#\bigl(V\cap(c+(PQ))\bigr).
\]
Then
\[
 \#X_-\,\#X_+
 \ge \frac{1}{M_{(PQ)}(V)}
       \sum_{v\in V}\deg_G(v)^2
 \ge \frac{(\#G)^2}{\#V\,M_{(PQ)}(V)}.
\tag{13.10bp0}
\]
\end{lemma}

\begin{proof}
A wedge is a triple \((u,v,u')\) with
\((u,v),(u',v)\in G\).  Map it to
\[
 (x,y)=(Qu-Pv,\ Qv-Pu')\in X_-\times X_+.
\]
For fixed \((x,y)\), reduction of these equations gives
\[
 v\equiv-P^{-1}x\pmod{(Q)},
 \qquad
 v\equiv Q^{-1}y\pmod{(P)}.
\]
The displayed inverses exist because \((P)+(Q)=R\).  The Chinese remainder
theorem confines \(v\) to one class modulo \((PQ)\).  Once \(v,x,y\) are
fixed, the two equations determine \(u,u'\) uniquely in the fraction field.
Thus every fiber of the wedge map has size at most \(M_{(PQ)}(V)\).
There are \(\sum_v\deg_G(v)^2\) wedges, and Cauchy--Schwarz proves the
second inequality.
\end{proof}

For \(P=q_{n-1}(s)\), \(Q=q_n(s)\), the two sets in this lemma are not
formal sumsets.  Equation (13.10bj) identifies them, up to signs, with
actual trace supports
\[
 \{\Tr(k^{-(n-1)}g):(\Tr g,\Tr(kg))\in G\},
 \quad
 \{\Tr(k^ng):(\Tr g,\Tr(kg))\in G\}.
\tag{13.10bp1}
\]
Thus linear trace growth bounds their product by
\(\exp(R+n\ell(k)+O(1))\) when \(G\subset G_R\).

There is an essential normalization when the coefficients are not integral.
Let \(F\) be a number field containing \(s\) and the trace vertices, put
\(K=\Q(s^2)\), and set
\(\mathcal J_F(s)=[F:K]\mathcal J_K(s)\).  For
\[
 P=q_{n-1}(s),\qquad Q=q_n(s),
\]
define the integral ideal
\[
 \mathfrak m_n=
 \prod_{\mathfrak p}
 \mathfrak p^{|v_{\mathfrak p}(P)-v_{\mathfrak p}(Q)|}.
\tag{13.10bp2}
\]
Thus \(\mathfrak m_n\) is the product of the coprime numerator and
denominator ideals of \(P/Q\).  If \(\Lambda\) is a fractional
\(\mathcal O_F\)-ideal containing a common rescaling of the vertex sets,
the two congruences in the proof of Lemma~\ref{lem:congruence-wedge}
confine the middle vertex to one coset of \(\Lambda\mathfrak m_n\).
Indeed, locally they say
\[
 v_{\mathfrak p}(v-v')-v_{\mathfrak p}(\Lambda)
 \ge
 \max\{v_{\mathfrak p}(P)-v_{\mathfrak p}(Q),
       v_{\mathfrak p}(Q)-v_{\mathfrak p}(P)\}.
\]
This formulation includes primes where the original representation is not
integral and is independent of the common rescaling.

At an archimedean embedding, choose an algebraic characteristic root
\(\lambda\) with \(\lambda+\lambda^{-1}=s\), and put \(z=\lambda^2\).
If \(|\lambda|\ne1\), then \(P/Q\) and its inverse tend to nonzero limits.
If \(z=1\), the repeated-root formula
\(q_n(\mathord\pm2)=(\mathord\pm1)^{n-1}n\) gives polynomial growth.
Otherwise \(z\) is not a root of unity: if \(z^m=1\), then the conjugate of
\(q_m(s)\) at this embedding would vanish, although its distinguished real
conjugate is positive.
Baker's lower bound for linear forms in logarithms gives a constant
\(C_z>0\) such that
\[
 |z^j-1|\ge j^{-C_z}\qquad(j\ge2)
\]
\cite{BakerLinearForms}.  Applying this to the closed formulas for
\(q_{n-1}(s)\) and \(q_n(s)\) shows that \(P/Q\) and its inverse are
\(n^{O_s(1)}\).  Formula (13.10av) and the product formula therefore give
\[
 \log N_{F/\Q}\mathfrak m_n
 =2n\mathcal J_F(s)+O_s(\log n).
\tag{13.10bp2a}
\]
This is the correct exponent.  Taking an unnormalized full ideal norm in
the quadratic sign extension can incorrectly double it.

\begin{theorem}[a uniform congruence-count criterion]
\label{thm:congruence-occupancy-bridge}
Let \(G_R\) be the trace-pair graph (13.10be), with vertex classes
\(U_R,V_R\).  Write
\[
 \#G_R=e^{R-\eta_R},\qquad \eta_R=o(R).
\]
Choose integers \(n_R\to\infty\) such that
\[
 n_R=O(R),\qquad \eta_R=o(n_R).
\tag{13.10bp3}
\]
Suppose that, after a common rescaling into a fractional ideal \(\Lambda_R\),
\[
 \max_c\#\bigl(V_R\cap(c+\Lambda_R\mathfrak m_{n_R})\bigr)
 \le
 \frac{\#V_R}{N\mathfrak m_{n_R}}e^{o(n_R)}.
\tag{13.10bp4}
\]
Then
\[
 \mathcal J_F(s)=\frac{\ell(k)}2.
\]
Consequently \(F=K\), \(\mathcal J_K(s)=\ell(k)/2\), and all
nondistinguished local expansion exponents vanish.
\end{theorem}

\begin{proof}
The two vertex classes have size at most \(e^{R/2+O(1)}\).  Apply
Lemma~\ref{lem:congruence-wedge}, in the fractional-lattice form just
proved, with \(n=n_R\).  Equations (13.10bp2a)--(13.10bp4) give
\[
 \begin{split}
 \log(\#X_-\,\#X_+)
 &\ge 2\log\#G_R-2\log\#V_R
       +\log N\mathfrak m_{n_R}-o(n_R)\\
 &\ge R+2n_R\mathcal J_F(s)-o(n_R).
 \end{split}
\]
On the other hand, (13.10bp1) and linear trace growth give
\[
 \log(\#X_-\,\#X_+)\le R+n_R\ell(k)+O(1).
\]
Divide by \(n_R\).  This yields
\(2\mathcal J_F(s)\le\ell(k)\); the distinguished real place already
contributes \(\ell(k)/2\), so equality holds and every other nonnegative
local term is zero.  Finally,
\(\mathcal J_F(s)=[F:K]\mathcal J_K(s)\) and
\(\mathcal J_K(s)\ge\ell(k)/2\), so equality forces \([F:K]=1\).
\end{proof}

Condition (13.10bp4) asks for a uniform upper bound on the number of trace
values in each residue class modulo the recurrence ideal.  It can be
deduced from a corresponding count of group elements.  Remove from \(V_R\)
vertices of degree
less than \(\#G_R/(4\#V_R)\), and make the analogous removal on the other
side.  A positive proportion of the edges remains, while every surviving
trace value is represented by at least \(\#G_R/(4\#V_R)\) group elements.
Consequently a uniform estimate
\[
 \max_c\#\{g:d(o,go)\le R+O(1),\ \Tr g\equiv c
                    \pmod{\Lambda_R\mathfrak m_{n_R}}\}
 \le \frac{e^{R+o(n_R)}}{N\mathfrak m_{n_R}}
\tag{13.10bp5}
\]
implies (13.10bp4) for the remaining graph and hence gives the theorem.
Property \(\tau\), together with uniform lattice-point counting, is the
natural source of (13.10bp5)~\cite{GorodnikNevoCounting}.

\subsection{The congruence criterion over the full trace field}

We now put the preceding implication in the form needed for the compact
conjectures.  This also removes two possible ambiguities.  The coefficient
field must be the full trace field of the subgroup generated by squares, and
the finitely many places at which a matrix model has denominators must be
deleted before ideal norms are taken.

Let \(H=\Gamma^{(2)}\), let \(\widetilde H\) be its inverse image in
\(\SL_2(\R)\), and suppose for the moment that
\[
 k=\Q(\Tr\widetilde H)
\tag{13.10bpa}
\]
is a number field.  The \(k\)-algebra generated by \(\widetilde H\) is its
invariant quaternion algebra \(B/k\).  Choose a finite set \(S\) of finite
places which contains the ramified finite places of \(B\) and clears the
denominators of a finite generating set.  Thus, for an
\(\cO_{k,S}\)-order \(\mathcal B\subset B\),
\[
 \widetilde H\subset \SL_1(\mathcal B)(\cO_{k,S}).
\tag{13.10bpb}
\]
Enlarge \(S\), once and for all, so that it is saturated over a finite set
of rational primes containing the denominators and the bad-reduction and
model discriminants of the restriction-of-scalars realization used below:
if one place above such a rational prime is removed, all places above it are
removed.  This harmless enlargement is needed when a number-field ideal is
compared with a rational congruence modulus.
Norms of ideals of \(\cO_{k,S}\) will mean norms of their prime-to-\(S\)
parts.

\begin{lemma}[the full trace field gives the full restriction-of-scalars closure]
\label{lem:full-trace-field-res-closure}
The connected \(\Q\)-Zariski closure of the restriction-of-scalars image of
\(\widetilde H\) is
\[
 \operatorname{Res}_{k/\Q}\SL_1(B).
\tag{13.10bpba}
\]
\end{lemma}

\begin{proof}
The group \(\widetilde H\) is non-elementary and is Zariski dense in
\(\SL_1(B)\) over \(k\).  After extension to \(\overline\Q\), the ambient
restriction of scalars is a product of \(\SL_2\)-factors indexed by the
embeddings \(\sigma:k\hookrightarrow\overline\Q\), and the connected
closure projects onto every factor.

A connected subdirect algebraic subgroup of a product of simple
\(A_1\)-groups is a product of diagonal blocks.  This follows, for example,
by applying the Lie-algebra version of Goursat's lemma successively: a
surjective subalgebra of a product of simple Lie algebras either contains a
factor or identifies it with another factor by an isomorphism.  Suppose a
block joined two different embeddings \(\sigma\ne\tau\).  After passing to
the finite-index subgroup \(\widetilde H_0\) lying in the connected
closure, the two projections would be identified by an automorphism of
\(\SL_2\).  Every such automorphism preserves trace, so
\[
 \sigma(\Tr h)=\tau(\Tr h)\qquad(h\in\widetilde H_0).
\tag{13.10bpbb}
\]
The invariant trace field of a non-elementary Fuchsian group is unchanged
on passage to finite index \cite[Chapter~3]{MaclachlanReid}.  Since \(H=\Gamma^{(2)}\), this invariant trace
field is precisely the full field \(k\) in (13.10bpa).  Equation
(13.10bpbb), also applied to squares, would therefore make \(\sigma\) and
\(\tau\) agree on \(k\), a contradiction.  Every diagonal block is a
singleton, proving (13.10bpba).
\end{proof}

Fix \(a\in\widetilde H\) with \(s=\Tr a>2\), and put
\(\ell=\ell(a)\).  For the recurrence in
Lemma~\ref{lem:two-trace-recurrence}, set
\[
 P_n=q_{n-1}(s),\qquad Q_n=q_n(s),\qquad
 \mathfrak n_n=(P_nQ_n)\cO_{k,S}\quad(n\ge2).
\tag{13.10bpc}
\]
The two factors in (13.10bpc) are coprime in \(\cO_{k,S}\), by Cassini's
identity (13.10ao).

\begin{lemma}[norm of the prime-to-\(S\) recurrence modulus]
\label{lem:prime-to-S-recurrence-norm}
With the preceding notation,
\[
 \log N\mathfrak n_n
 =2n\mathcal J_k(s)+O_{s,S}(\log n).
\tag{13.10bpd}
\]
Before \(S\) is chosen, the intrinsic version is obtained by dividing the
two principal fractional ideals by their common fractional-ideal factor.
\end{lemma}

\begin{proof}
Put \(\alpha_n=P_n/Q_n\).  The local calculation in the proof of
Lemma~\ref{lem:recurrence-height}, now summed over the places of \(k\), gives
\[
 [k:\Q]h(\alpha_n)=n\mathcal J_k(s)+O_s(\log n).
\tag{13.10bpe}
\]
We explain why deleting \(S\) loses only the displayed logarithmic error.
At a fixed archimedean place, choose an algebraic characteristic root
\(\lambda\) and set \(z=\lambda^2\).  If the two characteristic roots have
different moduli, the quotient of consecutive recurrence terms converges to
a nonzero limit.  If \(z=1\), use
\(q_n(\mathord\pm2)=(\mathord\pm1)^{n-1}n\).  In the remaining unit-circle
case, \(z\) is not a root of unity: otherwise \(q_m(s)\) would vanish under
this embedding for some \(m\), and injectivity of the embedding would give
\(q_m(s)=0\), contrary to its positive distinguished conjugate.  Baker's
linear-forms theorem gives \(|z^n-1|\ge n^{-C_z}\)
\cite{BakerLinearForms}.  The closed formula for \(q_n\) now shows that both
\(\log^+|\alpha_n|_v\) and
\(\log^+|\alpha_n^{-1}|_v\) are \(O_s(\log n)\).

The same assertion holds at every fixed nonarchimedean place in \(S\).
If \(|s|_v>1\), then
\(|q_n(s)|_v=|s|_v^{n-1}\), so the valuation of the quotient of two
consecutive terms is constant.  Suppose that \(|s|_v\le1\), and pass to a
finite extension containing a characteristic root \(\lambda\).  Then
\(\lambda\) and \(z=\lambda^2\) are units, and, unless \(z=1\),
\[
 v(q_n(s))=v(z^n-1)-v(z-1).
\tag{13.10bpe-local}
\]
We use the following elementary local valuation lemma.  If \(E/\Q_p\) is
finite and \(z\in\mathcal O_E^\times\) is not a root of unity, then
\[
 v_E(z^n-1)\le C_z+v_E(n).
\tag{13.10bpe-lte}
\]
Indeed, let \(r\) be the order of the residue class of \(z\).  Choose \(j\)
so that \(w=z^{rp^j}\) lies in a principal-unit neighborhood on which the
\(p\)-adic logarithm preserves valuations.  If \(r\nmid n\), the left side
is zero.  If \(n=rp^au\), with \(p\nmid u\), then for \(a<j\) it is bounded
by the maximum over the finitely many values \(a=0,\ldots,j-1\), because a
binomial expansion gives
\(v_E((z^{rp^a})^u-1)=v_E(z^{rp^a}-1)\).  For \(a\ge j\), the logarithm gives
\[
 v_E(z^n-1)=v_E(w-1)+v_E(p^{a-j}u)
           =v_E(w-1)+(a-j)v_E(p).
\]
This proves (13.10bpe-lte).  Equation (13.10bpe-local) therefore gives
\(v(q_n(s))=O_{s,v}(\log n)\).  When \(z=1\), the repeated-root formula
gives the same bound directly.  No other root of unity can occur, since
\(z^m=1\), \(z\ne1\), would force \(q_m(s)=0\).
Thus the archimedean and \(S\)-adic parts of the heights of both
\(\alpha_n\) and \(\alpha_n^{-1}\) are \(O_{s,S}(\log n)\).

Outside \(S\), coprimality says that the denominator ideals of
\(\alpha_n\) and \(\alpha_n^{-1}\) are respectively \((Q_n)\) and
\((P_n)\).  The ideal-height formula applied to (13.10bpe) in the two
directions gives
\[
 \log N(Q_n)=n\mathcal J_k(s)+O_{s,S}(\log n),\qquad
 \log N(P_n)=n\mathcal J_k(s)+O_{s,S}(\log n).
\]
Adding proves (13.10bpd).  In particular, even a place at which the original
representation is unbounded costs only \(O(\log n)\) after the fixed bad set
has been deleted.
\end{proof}

Choose a basepoint \(o\) on the axis of \(a\), put
\[
 \mathcal B_R=\{g\in\widetilde H:d(o,\bar g o)\le R\},
\]
and form the simple graph
\[
 \mathcal G_R=
 \{(\Tr g,\Tr(ag)):g\in\mathcal B_R\}.
\tag{13.10bpf}
\]
Write \(\#\mathcal G_R=e^{R-\eta_R}\).  Under linear trace growth,
orbit counting and Lemma~\ref{lem:ambient-two-trace-evaluations} give
\(\eta_R=o(R)\), while both vertex classes have size at most
\(e^{R/2+O_a(1)}\).

\begin{theorem}[full-trace-field congruence criterion]
\label{thm:full-trace-field-congruence}
Suppose that the signed trace set of \(\Gamma\) has linear growth.  Let
\(n_R\to\infty\) satisfy
\[
 n_R=O(R),\qquad \eta_R=o(n_R).
\tag{13.10bpg}
\]
Suppose, uniformly in
\(c\in\cO_{k,S}/\mathfrak n_{n_R}\), that
\[
 \#\{g\in\mathcal B_R:
       \Tr(ag)\equiv c\pmod{\mathfrak n_{n_R}}\}
 \le
 \frac{\exp(R+\epsilon_R)}{N\mathfrak n_{n_R}},
 \qquad \eta_R+\epsilon_R=o(n_R).
\tag{13.10bph}
\]
Then
\[
 \mathcal J_k(s)=\frac{\ell(a)}2.
\tag{13.10bpi}
\]
Consequently \(s\) is an algebraic integer and every nondistinguished
embedding \(\sigma:k\hookrightarrow\C\) satisfies
\(\sigma(s)\in[-2,2]\).
\end{theorem}

\begin{proof}
Let \(U_R,V_R\) be the two vertex classes of (13.10bpf).  Delete only those
vertices \(v\in V_R\) for which
\[
 \deg(v)<d_R:=\frac{\#\mathcal G_R}{4\#V_R}.
\]
Fewer than one quarter of the edges are deleted.  If
\(\mathcal G'_R\subset U_R\times V'_R\) is the remaining graph, then
\[
 \#\mathcal G'_R\ge\tfrac34\#\mathcal G_R,
 \qquad
 \#V'_R\ge
 \frac{\#\mathcal G'_R}{\#U_R}
 =\exp(R/2-\eta_R-O_a(1)).
\tag{13.10bpj}
\]
Choose one group element representing each edge.  Every retained right
vertex occurs for at least \(d_R\) of these distinct representatives, so
(13.10bph) implies
\[
 \max_c\#\bigl(V'_R\cap(c+\mathfrak n_{n_R})\bigr)
 \le
 \frac{\#V'_R}{N\mathfrak n_{n_R}}
 \exp(2\eta_R+\epsilon_R+O_a(1)).
\tag{13.10bpk}
\]
This pruning is necessary: equidistribution of group elements is a weighted
statement and does not directly imply flatness of the unweighted trace
support.

Apply Lemma~\ref{lem:congruence-wedge} to \(\mathcal G'_R\), with
\(P=P_{n_R}\) and \(Q=Q_{n_R}\).  The two images are actual trace
supports, because
\[
 \begin{aligned}
 Q\Tr g-P\Tr(ag)&=\Tr(a^{-(n_R-1)}g),\\
 Q\Tr(ag)-P\Tr g&=\Tr(a^{n_R}g).
 \end{aligned}
\tag{13.10bpl}
\]
Linear trace growth therefore gives
\[
 \#X_-\,\#X_+\le\exp(R+n_R\ell+O_a(1)).
\tag{13.10bpm}
\]
On the other hand, (13.10bpj), (13.10bpk), and the wedge injection give
\[
 \#X_-\,\#X_+
 \ge \exp(R-o(n_R))N\mathfrak n_{n_R}.
\tag{13.10bpn}
\]
Lemma~\ref{lem:prime-to-S-recurrence-norm}, followed by division by \(n_R\),
now gives \(2\mathcal J_k(s)\le\ell\).  The distinguished real place alone
contributes \(\ell/2\), and all local terms are nonnegative.  Equality
follows and every other local term vanishes.  At a finite place this says
\(|s|_v\le1\); at an archimedean place it says that the roots of
\(X^2-\sigma(s)X+1\) have modulus one, whence
\(\sigma(s)\in[-2,2]\).
\end{proof}

\begin{corollary}[the algebraic compact conjectures from recurrence counting]
\label{cor:algebraic-compact-from-congruence}
Let \(\Gamma<\PSL_2(\R)\) be torsion-free and cocompact.  Assume that its
invariant traces are algebraic and that its trace set has linear growth.
If (13.10bpg)--(13.10bph) hold for every hyperbolic element of
\(\widetilde{\Gamma^{(2)}}\), then \(\Gamma\) is arithmetic.  Thus this one
congruence-counting statement implies both the algebraic compact Schmutz
conjecture and, since bounded clustering implies linear growth, the
algebraic compact Sarnak conjecture.
\end{corollary}

\begin{proof}
A finite set of character coordinates generates all traces integrally.
Their algebraicity makes \(k\) in (13.10bpa) a number field.  Apply
Theorem~\ref{thm:full-trace-field-congruence} to every nonidentity element
of \(\widetilde{\Gamma^{(2)}}\).  All its traces are algebraic integers, and
every nondistinguished embedding sends every trace into \([-2,2]\); in
particular \(k\) is totally real.  Takeuchi's criterion \cite{Takeuchi} makes
\(\Gamma^{(2)}\), and therefore its finite-index overgroup \(\Gamma\),
arithmetic.
\end{proof}

Every count in this proof is initially a count of group elements: this is
true for the graph (13.10bpf), for cocompact orbit counting, and for the
congruence estimate.  Sparsity concerns only the set of trace values.  In a
torsion-free Fuchsian group, every
nonidentity element is a power of a unique primitive element up to inversion,
and \(|\Tr g|=2\cosh(\ell(g)/2)\).  This is exactly the trace/primitive-length
conversion in the standard formulations.  The two choices of lift change
cardinalities by at most two.

\begin{lemma}[quotient mixing implies the required trace count]
\label{lem:quotient-mixing-trace-count}
Let \(\mathfrak n\) be prime to \(S\).  Suppose strong approximation in the
form of Weisfeiler's theorem \cite{WeisfeilerStrong} gives
a bounded-index image
\(Q_{\mathfrak n}\subset
\SL_1(\mathcal B)(\cO_{k,S}/\mathfrak n)\), uniformly in
\(\mathfrak n\), and suppose
\[
 \max_{x\in Q_{\mathfrak n}}
 \#\{g\in\mathcal B_R:\pi_{\mathfrak n}(g)=x\}
 \le \frac{\exp(R+o(R))}{\#Q_{\mathfrak n}}.
\tag{13.10bpo}
\]
Then, uniformly in \(c\),
\[
 \#\{g\in\mathcal B_R:
       \Tr(ag)\equiv c\pmod{\mathfrak n}\}
 \le \frac{\exp(R+o(R))}{N\mathfrak n}.
\tag{13.10bpp}
\]
The errors are uniform in any family for which
\(\log N\mathfrak n=O(R)\) and (13.10bpo) is uniform.
\end{lemma}

\begin{proof}
Outside \(S\), the group scheme is \(\SL_2\).  If a prime
\(\mathfrak p\) has residue cardinality \(q\), then, for
\(A=\cO_k/\mathfrak p^e\),
\[
 \max_t\#\{x\in\SL_2(A):\Tr x=t\}\le(e+1)q^{2e}.
\tag{13.10bpq}
\]
Indeed, after choosing the upper-left entry \(u\), the determinant equation
is \(bc=u(t-u)-1\).  Splitting \(b\) according to its valuation gives at
most \((e+1)q^e\) pairs \((b,c)\).  Since
\(\#\SL_2(A)=q^{3e}(1-q^{-2})\), the Chinese remainder theorem gives
\[
 \frac{\max_t\#\{x:\Tr x=t\}}
      {\#\SL_1(\mathcal B)(\cO_{k,S}/\mathfrak n)}
 \le \frac{(N\mathfrak n)^{o(1)}}{N\mathfrak n}.
\tag{13.10bpr}
\]
A bounded-index image changes only the implied constant.  Multiplication by
\(a\) permutes that image, so (13.10bpo)--(13.10bpr) prove (13.10bpp).
\end{proof}

A uniform spectral gap for all the congruence quotients, combined with
uniform lattice-point counting on the corresponding compact covers, gives
(13.10bpo) in a range
\(\log N\mathfrak n\le cR\): more precisely the count in a quotient coset
has a main term \(e^R/\#Q_{\mathfrak n}\) and an error
\(O(e^{(1-\delta)R})\), with \(\delta>0\) independent of the modulus
\cite{BrooksTower,Burger,GorodnikNevoCounting}.  Equations
(13.10bpq)--(13.10bpr) absorb
the error after \(c>0\) is chosen sufficiently small.  Therefore full
arbitrary-ideal superapproximation for the group in (13.10bpb) proves
(13.10bph), with \(n_R=\lfloor\vartheta R\rfloor\) and
\(\vartheta>0\) sufficiently small, and hence proves the algebraic compact
Sarnak and Schmutz conjectures.

This last spectral estimate is not currently available in the required
generality.  Salehi Golsefidy proves expansion for moduli \(m^E\), where
\(m\) is square-free and the exponent \(E\) is fixed in advance
\cite{SalehiSuperII}.  The claimed
arbitrary-prime-power extension in \cite{SalehiSuperI,SalehiSuperII} has an
unresolved issue, recorded explicitly in
\cite[Remark~2.8]{TangZhang}, so no such extension is used here.  Varj\'u
treats square-free ideals in the split \(\SL_2\) case
\cite{VarjuSquarefree}.  He and de Saxc\'e prove all
integer moduli for integral groups with absolutely simple Zariski closure
\cite{HeDeSaxce}; their published theorem therefore does not include
\(\operatorname{Res}_{k/\Q}\SL_1(B)\) when \([k:\Q]>1\), because that
group becomes a product over \(\overline\Q\), and it also does not state the
needed \(S\)-integer
version.  Tang and Zhang prove all integer moduli, including fixed
denominators, for the split non-simple cases
\(\SL_2\times\SL_2\) and \(\operatorname{ASL}_2\)
\cite{TangZhang}; their theorem does not cover a general twisted
restriction-of-scalars form.  The general \(S\)-integer arbitrary-modulus
assertion remains open.  The induction arguments of
\cite{SalehiZhangInducing} transfer expansion in certain extensions, but do
not prove the twisted restriction-of-scalars assertion required here.  Thus Corollary
\ref{cor:algebraic-compact-from-congruence} is a complete implication
theorem, not a claim that this external spectral problem has already been
settled.

\begin{proposition}[the congruence criterion on balls from another Fuchsian metric]
\label{prop:fixed-ball-specialization}
Let \(x\) be a faithful cocompact Fuchsian character of a closed surface
group \(\Pi\), and assume that its signed trace set has linear growth.  Let
\(Z\subset V_x\) be a \(\Q\)-subvariety through \(x\), and let
\(y\in Z(\overline\Q)\) also be a faithful cocompact Fuchsian character.
Fix a lift \(a\in\widetilde{\Pi^{(2)}}\) of a nonidentity element of
\(\Pi^{(2)}\), put
\[
 s_y=\Tr_y(a),\qquad
 \mathcal B_R^x=
 \{g\in\widetilde{\Pi^{(2)}}:d_x(o,\bar g o)\le R\},
\]
and choose \(o\) on the \(x\)-axis of \(a\).  Form
\[
 \mathcal G_{R,y}
 =\{(\Tr_y g,\Tr_y(ag)):g\in\mathcal B_R^x\},
 \qquad
 \#\mathcal G_{R,y}=e^{R-\eta_{R,y}}.
\tag{13.10bpr1}
\]
Let \(k_y\) be the full trace field of the specialized subgroup generated
by squares, choose its finite bad set \(S_y\), and let
\(\mathfrak n_{n,y}\) be the recurrence ideal (13.10bpc) formed from
\(s_y\).  Suppose that, for
\(n_R=\lfloor\vartheta R\rfloor\), where \(\vartheta>0\) is fixed, one
has uniformly in \(c\)
\[
 \#\{g\in\mathcal B_R^x:
   \Tr_y(ag)\equiv c\pmod{\mathfrak n_{n_R,y}}\}
 \le
 \frac{\exp(R+\epsilon_R)}{N\mathfrak n_{n_R,y}},
 \qquad
 \eta_{R,y}+\epsilon_R=o(R).
\tag{13.10bpr2}
\]
Then
\[
 2\mathcal J_{k_y}(s_y)\le\ell_x(a).
\tag{13.10bpr3}
\]
\end{proposition}

\begin{proof}
Every equality of squared traces at \(x\) holds on \(V_x\), hence at
\(y\).  An equality of signed traces at \(x\) can split into at most the
two choices of sign at \(y\).  Linear trace growth at \(x\) therefore
bounds both coordinate projections of (13.10bpr1) by
\(e^{R/2+O_a(1)}\).

The representation at \(y\) is faithful Fuchsian.  The two-trace
finite-type estimate, Lemma~\ref{lem:ambient-two-trace-evaluations},
applied to the fixed \(x\)-ball, bounds every fiber of the two coordinates
by \(e^{o(R)}\); word length and \(x\)-displacement are comparable.
Cocompact orbit counting in the metric \(x\) now gives
\[
 \#\mathcal G_{R,y}=e^{R-o(R)},\qquad \eta_{R,y}=o(R).
\tag{13.10bpr4}
\]

The recurrence identities hold in the representation \(y\):
\[
 \begin{aligned}
 q_n(s_y)\Tr_y g-q_{n-1}(s_y)\Tr_y(ag)
   &=\Tr_y(a^{-(n-1)}g),\\
 q_n(s_y)\Tr_y(ag)-q_{n-1}(s_y)\Tr_y g
   &=\Tr_y(a^ng).
 \end{aligned}
\tag{13.10bpr5}
\]
The elements on the right lie in the \(x\)-ball of radius
\(R+n\ell_x(a)+O_a(1)\).  Persistence again bounds their \(y\)-trace
supports, up to the two signs, by the corresponding \(x\)-trace supports.
Consequently
\[
 \#X_-\,\#X_+
 \le\exp\bigl(R+n\ell_x(a)+O_a(1)\bigr).
\tag{13.10bpr6}
\]
The pruning and congruence-class counting proof of
Theorem~\ref{thm:full-trace-field-congruence}, using (13.10bpr2), is
otherwise unchanged and gives
\[
 \#X_-\,\#X_+
 \ge
 \exp(R-o(n_R))N\mathfrak n_{n_R,y}
 =
 \exp\bigl(R+2n_R\mathcal J_{k_y}(s_y)-o(n_R)\bigr).
\tag{13.10bpr7}
\]
Comparison of (13.10bpr6) and (13.10bpr7), followed by division by
\(n_R\), proves (13.10bpr3).
\end{proof}

There is no change in the spectral interpretation of the hypothesis in
Proposition~\ref{prop:fixed-ball-specialization}.  A congruence kernel
defined by the algebraic representation \(y\) is a finite-index abstract
subgroup of \(\Pi\), hence also defines a cover of the fixed surface
\((S,x)\).  Expansion of the \(y\)-congruence quotients gives a uniform
spectral gap on these \(x\)-covers by the Brooks--Burger transfer
\cite{BrooksTower,Burger}, and
uniform lattice counting then proves (13.10bpr2) for sufficiently small
\(\vartheta\).  No displacement ball for the metric \(y\) is used.

\pagebreak[3]
\begin{lemma}[local expansion and marked Fuchsian rigidity]
\label{lem:local-expansion-fuchsian-rigidity}
Let \(x\) and \(y\) be faithful cocompact Fuchsian characters of the same
closed surface group \(\Pi\), lying in the same Fuchsian component, and
assume that \(y\) is algebraic.  Let \(k_y\) be the full trace field of the
subgroup generated by squares.  Suppose that every nonidentity
\(b\in\widetilde{\Pi^{(2)}}\) satisfies
\[
 \mathcal J_{k_y}(\Tr_y b)\le \frac{\ell_x(b)}2.
\tag{13.10bpr9}
\]
Then \(y=x\).
\end{lemma}

\begin{proof}
For \(\gamma\in\Pi\setminus\{1\}\), choose a lift \(b\) of
\(\gamma^2\).  Since
\(\mathcal J_{k_y}(-t)=\mathcal J_{k_y}(t)\), the choice of lift does not
matter.  At the distinguished real embedding,
\[
 |\Tr_y b|=2\cosh\!\left(\frac{\ell_y(\gamma^2)}2\right),
\]
so this place contributes \(\ell_y(\gamma^2)/2\) to the nonnegative sum
\(\mathcal J_{k_y}(\Tr_y b)\).  Hence
\[
 \ell_y(\gamma)
 =\frac{\ell_y(\gamma^2)}2
 \le \mathcal J_{k_y}(\Tr_y b)
 \le \frac{\ell_x(\gamma^2)}2
 =\ell_x(\gamma).
\]
Thurston's minimal-stretch theorem identifies the logarithm of the least
Lipschitz constant in the marked homotopy class with
\[
 d_{\rm Th}(x,y)
 =\log\sup_{\gamma\ne1}\frac{\ell_y(\gamma)}{\ell_x(\gamma)}.
\]
The common Fuchsian component identifies \(x\) and \(y\) with marked points
of the same Teichm\"uller space.  The displayed formula defines a
nonnegative separating asymmetric metric on that space
\cite{ThurstonStretch}.  The preceding length domination gives
\(d_{\rm Th}(x,y)\le0\).  Therefore \(d_{\rm Th}(x,y)=0\), and separation
gives \(y=x\).
\end{proof}

For comparison, we formulate a general superapproximation hypothesis that
would also imply both compact conjectures, including the case of a
transcendental character.  Let
\(\mathrm{SA}_{\mathrm{all}}\) denote the following assertion:
for every number field \(k\), quaternion algebra \(B/k\), finite set
\(S\), and finitely generated Zariski-dense subgroup
\[
 \Delta<\SL_1(B)(\cO_{k,S})
\]
whose full trace field is \(k\), the Cayley graphs of its reduction images
modulo all ideals prime to \(S\) have a uniform spectral gap.

\begin{theorem}[arbitrary-ideal superapproximation implies both compact conjectures]
\label{thm:superapproximation-implies-compact}
Assume \(\mathrm{SA}_{\mathrm{all}}\).  Then every torsion-free cocompact
Fuchsian lattice whose signed trace set has linear growth is arithmetic.
Consequently both the compact Schmutz linear-growth conjecture and the
compact Sarnak bounded-clustering conjecture hold.
\end{theorem}

\begin{proof}
Let \(x\) be the Fuchsian character.  If \(x\) is algebraic, the expansion
hypothesis, Lemma~\ref{lem:quotient-mixing-trace-count}, and
Corollary~\ref{cor:algebraic-compact-from-congruence} prove arithmeticity.

Suppose that \(x\) is transcendental.  Its rational Zariski closure has
positive dimension.  Proposition~\ref{prop:bounded-degree-specialization}
gives an algebraic Fuchsian point \(y\ne x\), arbitrarily close to \(x\)
in the same Fuchsian component, at which all trace collisions at \(x\)
persist.  The hypothesis \(\mathrm{SA}_{\mathrm{all}}\)
and the preceding spectral interpretation verify (13.10bpr2).  Thus, for
every nonidentity element of \(\Pi^{(2)}\) and either lift
\(a\in\widetilde{\Pi^{(2)}}\),
\[
 \mathcal J_{k_y}(\Tr_ya)
 \le\tfrac12\ell_x(a).
\tag{13.10bpr8}
\]
Lemma~\ref{lem:local-expansion-fuchsian-rigidity} now gives \(y=x\),
contrary to the choice of \(y\).

Thus \(x\) is algebraic and the first paragraph applies.  Finally, bounded
clustering implies linear trace growth, so the Sarnak implication follows
from the Schmutz implication.
\end{proof}

Nor can estimates for the prime-power factors simply be multiplied.  The
following finite model records the exact obstruction.

\begin{proposition}[flat marginals need not be flat modulo a product]
\label{prop:crt-diagonal-obstruction}
Let \(p\) be prime and let
\[
 V=\{(t,t):t\in\F_p\}\subset\F_p\times\F_p.
\]
Each coordinate projection is perfectly flat: every residue class has
exactly
\[
 \frac{\#V}{p}=1
\]
preimage.  But the joint reduction occupies only \(p\) of the \(p^2\)
classes, and its largest fiber is \(1\), rather than the formal product
prediction \(\#V/p^2=1/p\).  Equivalently, if
\[
 E_i(V)=\sum_{a\in\F_p}\#\{v\in V:v_i=a\}^2,
\]
then \(E_i(V)=\#V^2/p=p\) for each factor, whereas the joint collision
energy is \(\#V=p\), not \(\#V^2/p^2=1\).
\end{proposition}

\begin{proof}
Both coordinate maps restrict to bijections on \(V\), while the joint map is
the inclusion of the diagonal.  The formulas follow immediately.
\end{proof}

If \(\log p\) is proportional to the recurrence index, the discrepancy in
Proposition~\ref{prop:crt-diagonal-obstruction} is exponential.  Therefore
one-prime or one-prime-power trace estimates do not imply (13.10bph) for a
mixed recurrence ideal.  One would need either a spectral theorem directly
for the product quotient or an estimate that rules out correlations between
the different prime-power coordinates.

The available expansion theorems nevertheless give an unconditional
restriction on the argument with a fixed group element whenever a
nondistinguished place contributes positive recurrence height.  For a fixed
integer \(D\ge1\),
write
\[
 \mathfrak c_{n,D}
 =\prod_{\mathfrak p\notin S}
   \mathfrak p^{\min\{v_{\mathfrak p}(\mathfrak n_n),D\}},
 \qquad
 \mathfrak h_{n,D}=\mathfrak n_n\mathfrak c_{n,D}^{-1}.
\tag{13.10bps}
\]
Thus \(\mathfrak c_{n,D}\) is obtained by truncating every local exponent at
\(D\), while \(\mathfrak h_{n,D}\) contains the remaining powers.

\begin{theorem}[large prime-power contributions from a nondistinguished expanding place]
\label{thm:high-power-concentration}
Under the algebraic and linear-growth hypotheses above, for every fixed
hyperbolic \(a\in\widetilde H\) and every fixed \(D\),
\[
 \limsup_{n\to\infty}\frac1n\log N\mathfrak c_{n,D}
 \le \ell(a).
\tag{13.10bpt}
\]
Consequently
\[
 \liminf_{n\to\infty}\frac1n\log N\mathfrak h_{n,D}
 \ge 2\mathcal J_k(s)-\ell(a).
\tag{13.10bpu}
\]
If \(a\) has any nondistinguished local expansion, the right side of
(13.10bpu) is positive.  Hence, for every fixed cutoff \(D\), a linear
amount of the recurrence height is forced into prime powers of exponent
strictly greater than \(D\).
\end{theorem}

\begin{proof}
After restriction of scalars, (13.10bpb) embeds the finitely generated group
in \(\GL_N(\Z[1/q_0])\) for one fixed \(q_0\).  If
\(m\) is the square-free product of the rational primes below
\(\mathfrak c_{n,D}\), then
\[
 m^D\cO_{k,S}\subset\mathfrak c_{n,D}.
\tag{13.10bpv}
\]
Indeed, at a prime \(\mathfrak p\mid p\), the exponent on the left is
\(D e(\mathfrak p/p)\), which is at least
\(\min\{v_{\mathfrak p}(\mathfrak n_n),D\}\); recall that larger ideal
exponents mean smaller ideals.  The rational prime \(p\) is prime to
\(q_0\) by the saturation convention on \(S\).  Reduction modulo
\(\mathfrak c_{n,D}\) is therefore a quotient of reduction modulo the
bounded power \(m^D\).

By Lemma~\ref{lem:full-trace-field-res-closure}, the connected
\(\Q\)-Zariski closure is the full group
\(\operatorname{Res}_{k/\Q}\SL_1(B)\).  It is semisimple and therefore
perfect; moreover, it is connected and simply connected, since both
properties are preserved by restriction of scalars.  Salehi Golsefidy's
theorem for square-free moduli raised to a fixed power
\cite{SalehiSuperII} applies; absolute simplicity is not required in this
theorem.  Weisfeiler's strong-approximation theorem
\cite{WeisfeilerStrong}, applied after restriction of scalars and after enlarging
the fixed bad set, gives a uniformly bounded-index congruence image and hence
\(\#Q_c\gg(Nc)^3\) for \(c=\mathfrak c_{n,D}\).  The spectral gap
descends from modulus \(m^D\) to its quotient at modulus \(c\).

For completeness, uniform lattice-point counting on the congruence covers
gives, for some \(\delta_D>0\), a count per quotient element bounded by
\[
 \frac{e^{R+o(R)}}{\#Q_c}+O(e^{(1-\delta_D)R}).
\]
The local trace-fiber estimate (13.10bpq), the Chinese remainder theorem,
and the ideal-divisor bound
\(\prod_{\mathfrak p^e\Vert c}(e+1)=(Nc)^{o(1)}\) therefore give
\[
 \#\{g\in\mathcal B_R:\Tr(ag)\equiv t\pmod c\}
 \ll (Nc)^{o(1)}\left\{\frac{e^R}{Nc}
             +(Nc)^2e^{(1-\delta_D)R}\right\}.
\tag{13.10bpw}
\]
Lemma~\ref{lem:prime-to-S-recurrence-norm} gives
\(\log Nc\le2\mathcal J_k(s)n+O(\log n)\).  Take
\(R=C_{D,a,\Gamma}n\), with the fixed constant so large that
\(\delta_DC_{D,a,\Gamma}>6\mathcal J_k(s)\).  Then the second term in
(13.10bpw) is at most \(e^{R+o(n)}/Nc\), and \(\eta_R=o(n)\).
Repeat the proof of
Theorem~\ref{thm:full-trace-field-congruence}, with
\(\mathfrak c_{n,D}\) in place of \(\mathfrak n_n\).  Comparison of
(13.10bpm) and (13.10bpn) gives
\[
 \log N\mathfrak c_{n,D}\le n\ell(a)+o(n),
\]
which proves (13.10bpt).  Subtract this from
Lemma~\ref{lem:prime-to-S-recurrence-norm} to obtain (13.10bpu).
\end{proof}

\begin{theorem}[a one-term Chebyshev radical bound]
\label{thm:single-term-radical-upper}
Retain the algebraic and linear-trace-growth hypotheses of
Theorem~\ref{thm:high-power-concentration}.  Fix a hyperbolic
\(a\in\widetilde H\), choose its lift so that
\(s=\Tr a>2\), and put
\[
 \ell=\ell(a),\qquad P_n=q_{n-1}(s),\qquad Q_n=q_n(s).
\]
For \(n\ge2\), define the prime-to-\(S\) ideals
\[
 \mathfrak r_n
 =\prod_{\substack{\mathfrak p\notin S\\
                    v_{\mathfrak p}(Q_n)>0}}\mathfrak p,
 \qquad
 \mathfrak d_n=\mathfrak r_n^2.
\tag{13.10bpw1}
\]
Thus \(\mathfrak d_n\) is the ideal square of the radical of the
\emph{single} recurrence term \(q_n(s)\); it is not the radical of
\(q_n(s)^2\).  Then
\[
 \limsup_{n\to\infty}\frac1n\log N\mathfrak r_n
 \le\frac{\ell(a)}2.
\tag{13.10bpw2}
\]
Consequently, if
\(\mathfrak q_n=(Q_n)\mathfrak r_n^{-1}\) is the part of the principal
prime-to-\(S\) ideal left after one copy of each prime divisor is removed,
then
\[
 \liminf_{n\to\infty}\frac1n\log N\mathfrak q_n
 \ge \mathcal J_k(s)-\frac{\ell(a)}2.
\tag{13.10bpw3}
\]
In particular, any nondistinguished expanding place forces a positive
linear amount of norm into repeated prime factors of every individual term
\(q_n(s)\).
\end{theorem}

\begin{proof}
Let \(m_n\) be the square-free product of the rational primes lying below
the prime ideals occurring in \(\mathfrak r_n\).  The rational saturation
of \(S\) gives
\[
 m_n^2\cO_{k,S}\subseteq\mathfrak d_n.
\tag{13.10bpw4}
\]
To check the direction of this inclusion, fix
\(\mathfrak p\mid p\).  The exponent of \(\mathfrak p\) on the left is
\(2e(\mathfrak p/p)\).  It is at least \(2\) if
\(\mathfrak p\mid\mathfrak d_n\), while the required exponent is zero at
the other primes.  Since larger ideal exponents give smaller ideals,
(13.10bpw4) follows.  Hence
\[
 \cO_{k,S}/m_n^2\cO_{k,S}\twoheadrightarrow
 \cO_{k,S}/\mathfrak d_n,
\]
so a spectral gap modulo the bounded power \(m_n^2\) descends to the
quotient modulo \(\mathfrak d_n\).

After restriction of scalars, the connected Zariski closure is the perfect
semisimple group \(\operatorname{Res}_{k/\Q}\SL_1(B)\), by
Lemma~\ref{lem:full-trace-field-res-closure}.  Salehi Golsefidy's theorem
for fixed powers of square-free rational moduli, used here with the square,
therefore applies \cite{SalehiSuperII}.  Combining it with uniform lattice
counting and the local trace-fiber estimate (13.10bpq) gives, for some
fixed \(\delta_2>0\),
\[
 \max_c\#\{g\in\mathcal B_R:
             \Tr g\equiv c\pmod{\mathfrak d_n}\}
 \ll
 (N\mathfrak d_n)^{o(1)}
 \left\{
   \frac{e^R}{N\mathfrak d_n}
   +(N\mathfrak d_n)^2e^{(1-\delta_2)R}
 \right\}.
\tag{13.10bpw5}
\]
The factor \((N\mathfrak d_n)^{o(1)}\) comes from the ideal-divisor
factor in the local trace count.  Since
\(\log N\mathfrak d_n=O(n)\), it is \(\exp(o(n))\); the
lattice-counting error denoted by \(o(R)\) below is also \(o(n)\) after
\(R=Cn\).  Since
\((Q_n^2)\subseteq\mathfrak d_n\), the one-factor norm estimate in the
proof of Lemma~\ref{lem:prime-to-S-recurrence-norm} gives
\[
 \log N\mathfrak d_n
 \le 2n\mathcal J_k(s)+O_{s,S}(\log n).
\tag{13.10bpw6}
\]
Choose a fixed \(C>6\mathcal J_k(s)/\delta_2\) and put \(R=Cn\).
The ratio of the error term in braces in (13.10bpw5) to its main term is
\[
 (N\mathfrak d_n)^{3+o(1)}e^{-\delta_2R}
 \le
 \exp\bigl((6\mathcal J_k(s)-\delta_2C)n+o(n)\bigr)=o(1).
\]
Thus
\[
 \max_c\#\{g\in\mathcal B_R:
             \Tr g\equiv c\pmod{\mathfrak d_n}\}
 \le\frac{e^{R+o(n)}}{N\mathfrak d_n}.
\tag{13.10bpw7}
\]

Use the set \(W_R\) produced in the proof of
Theorem~\ref{thm:drc-trace-recurrence}, as a subset of the left vertex
class \(\{\Tr g:g\in\mathcal B_R\}\).  Lemma~\ref{lem:drc-explicit}
allows the two selected vertices to be equal.  Applied to the repeated
pair \((u,u)\), its common-neighbour conclusion says that every
\(u\in W_R\) has degree at least \(r_R\).  Equations (13.10ba) and
(13.10bi) consequently give
\[
 \#W_R\,r_R
 \ge \tfrac14N_R^2
       \exp\bigl(-a_R(t_R+1)-L_R/t_R\bigr)
 =\exp(R-o(R)).
\tag{13.10bpw8}
\]
Choose one group element representing each edge incident to \(W_R\);
different edges have different representatives.  If \(W_R\) meets
\(M_{R,n}\) residue classes modulo \(\mathfrak d_n\), all these
representatives lie in the union of those \(M_{R,n}\) residue classes.
In particular,
\[
 M_{R,n}\!
 \max_c\#\{g\in\mathcal B_R:
             \Tr g\equiv c\pmod{\mathfrak d_n}\}
 \ge \#W_R\,r_R.
\]
Equations (13.10bpw7)--(13.10bpw8), and \(R=Cn\), now give the explicit
support estimate
\[
 M_{R,n}
 \ge
 \frac{\#W_R\,r_R}{e^{R+o(n)}/N\mathfrak d_n}
 \ge N\mathfrak d_n\,e^{-o(n)}.
\tag{13.10bpw9}
\]
No second application of dependent random choice and no further degree
pruning are needed.

Let \(T_R\subset W_R\) contain exactly one representative from each
occupied residue class modulo \(\mathfrak d_n\), and recall that
\(\beta_n=-P_n^2/Q_n^2\).  The map
\[
 T_R\times W_R\longrightarrow
 Q_n^2(W_R+\beta_nW_R),\qquad
 (t,u)\longmapsto Q_n^2u-P_n^2t
\tag{13.10bpw10}
\]
is injective.  Indeed, reduction of an equality of two images modulo
\(\mathfrak d_n\) kills the \(Q_n^2\)-terms.  Because
\(\mathfrak d_n\) is supported on primes dividing \(Q_n\), Cassini's
identity makes \(P_n\) a unit modulo \(\mathfrak d_n\), so it recovers the residue class
of \(t\).  The definition of \(T_R\) then gives \(t=t'\), and equality in
the number field gives \(u=u'\).  Therefore (13.10bpw9) implies
\[
 \frac{\#(W_R+\beta_nW_R)}{\#W_R}
 \ge N\mathfrak d_n\,e^{-o(n)}.
\tag{13.10bpw11}
\]
The simultaneous upper bound (13.10bd), at \(R=Cn\), is
\[
 \frac{\#(W_R+\beta_nW_R)}{\#W_R}
 \le\exp(n\ell+o(n)).
\]
Comparison gives
\(\log N\mathfrak d_n\le n\ell+o(n)\).  Since
\(N\mathfrak d_n=(N\mathfrak r_n)^2\), this proves (13.10bpw2).
Finally,
\[
 \log N\bigl((Q_n)\bigr)
 =n\mathcal J_k(s)+O_{s,S}(\log n)
 =\log N\mathfrak r_n+\log N\mathfrak q_n,
\]
which proves (13.10bpw3).
\end{proof}

\begin{proposition}[residue-class deficit modulo a complete recurrence term]
\label{prop:full-depth-support-saturation}
Retain the hypotheses and notation of
Theorem~\ref{thm:single-term-radical-upper}.  Let \(W_R\) and
\(\Delta_R\) be given by
Theorem~\ref{thm:drc-trace-recurrence}, and put
\[
 \mathfrak m_n=(q_n(s)^2)\cO_{k,S},\qquad
 M_{R,n}=\#\operatorname{im}
 \bigl(W_R\longrightarrow\cO_{k,S}/\mathfrak m_n\bigr).
\tag{13.10bpw12}
\]
Then, for every \(R\) and \(n\ge2\),
\[
 M_{R,n}\le
 \frac{\#(W_R+\beta_nW_R)}{\#W_R}
 \le \exp(n\ell(a)+\Delta_R).
\tag{13.10bpw13}
\]
Consequently, for every fixed \(C>0\),
\[
 \log\frac{N\mathfrak m_n}{M_{Cn,n}}
 \ge (2\mathcal J_k(s)-\ell(a))n-o(n).
\tag{13.10bpw14}
\]
In particular, if along an unbounded sequence
\[
 M_{Cn,n}\ge N\mathfrak m_n\,e^{-o(n)},
\tag{13.10bpw15}
\]
then \(2\mathcal J_k(s)\le\ell(a)\).  The contribution of the
distinguished real place gives the reverse inequality, so equality holds and every
nondistinguished local expansion exponent vanishes.
\end{proposition}

\begin{proof}
Write \(P_n=q_{n-1}(s)\), \(Q_n=q_n(s)\), and choose
\(T_{R,n}\subset W_R\) with one representative of every class modulo
\(\mathfrak m_n\).  The map
\[
 T_{R,n}\times W_R\longrightarrow Q_n^2(W_R+\beta_nW_R),
 \qquad (t,u)\longmapsto Q_n^2u-P_n^2t,
\]
is injective.  Reducing an equality modulo \((Q_n^2)\) recovers the
residue class of \(t\), because Cassini's identity makes \(P_n\) a unit
modulo \((Q_n^2)\).  The choice of representatives then gives \(t=t'\),
and equality in the number field gives \(u=u'\).  This proves the first
inequality in (13.10bpw13); the second is (13.10bd).  Finally,
\[
 \log N\mathfrak m_n
 =2n\mathcal J_k(s)+O_{s,S}(\log n),
\]
and \(\Delta_{Cn}=o(n)\), which proves (13.10bpw14) and the final
assertion.
\end{proof}

This criterion asks less than superapproximation for arbitrary ideals.  It
requires only that the particular set \(W_{Cn}\), obtained by dependent
random choice, occupy almost all residue classes modulo the particular
principal ideal \((q_n(s)^2)\).  Conversely, (13.10bpw14) says that a
nondistinguished place with positive recurrence height forces this set to
miss an exponentially large proportion of that quotient.  Bounded-power
expansion proves the required estimate only after \((q_n(s)^2)\) is replaced
by the square of its radical.  Theorem
\ref{thm:single-term-radical-upper} measures exactly the resulting loss.

The same bounded-power estimate gives a square-divisor sieve when the group
element, rather than the recurrence index, is allowed to vary.  We first
record the exact local calculation.  It also explains why divisibility by
\(\mathfrak p^2\) is not a condition of codimension two modulo
\(\mathfrak p\): it is a second-order congruence along a smooth divisor.

\begin{lemma}[local density of a Chebyshev power divisor]
\label{lem:chebyshev-local-power-density}
Fix \(m\ge2\).  Let \(\mathfrak p\) be a prime ideal at which the group
scheme is split, put \(q=N\mathfrak p\), and assume that the residue
characteristic does not divide \(2m\).  Define
\[
 r_+(m,q)=\frac{\gcd(2m,q-1)-2}{2},\qquad
 r_-(m,q)=\frac{\gcd(2m,q+1)-2}{2}.
\tag{13.10bpw16}
\]
For every \(e\ge1\),
\[
 \frac{\#\{x\in\SL_2(\cO_k/\mathfrak p^e):
                  q_m(\Tr x)=0\}}
      {\#\SL_2(\cO_k/\mathfrak p^e)}
 =\frac{r_+(m,q)(q+1)+r_-(m,q)(q-1)}
        {q^{e-1}(q^2-1)}.
\tag{13.10bpw17}
\]
In particular, this density is
\[
 \frac{r_+(m,q)+r_-(m,q)}{q^e}+O_m(q^{-e-1}),
\tag{13.10bpw18}
\]
and the set counted in the numerator on the left of (13.10bpw17) has
\(O_m(q^{2e})\) elements in
\(\SL_2(\cO_k/\mathfrak p^e)\).

Moreover, for \(x\in\SL_2(\cO_k/\mathfrak p^e)\),
\[
 q_m(\Tr x)=0\quad\Longrightarrow\quad x^m\in\{I,-I\}.
\tag{13.10bpw19}
\]
Conversely, if the reduction of \(x\) modulo \(\mathfrak p\) is not
central, then \(x^m\) scalar implies \(q_m(\Tr x)=0\).  On the two central
residue discs, a scalar \(m\)-th power occurs only for \(x=I\) and
\(x=-I\).  Thus the scalar-power locus consists of the divisor counted in
(13.10bpw17), together with these two central points.
\end{lemma}

\begin{proof}
Write \(t=\lambda+\lambda^{-1}\).  Away from \(t=\pm2\),
\[
 q_m(t)=\frac{\lambda^m-\lambda^{-m}}
              {\lambda-\lambda^{-1}}.
\tag{13.10bpw20}
\]
Hence a split root over \(\F_q\) is obtained from an element
\(\lambda\in\F_q^\times\) satisfying \(\lambda^{2m}=1\), with
\(\lambda=\pm1\) deleted and \(\lambda\) identified with
\(\lambda^{-1}\).  This gives \(r_+(m,q)\) roots.  The nonsplit roots are
obtained in the same way from the norm-one subgroup of
\(\F_{q^2}^\times\), whose order is \(q+1\), and their number is
\(r_-(m,q)\).  The assumption on the residue characteristic makes all
these roots simple.  It also excludes \(\pm2\), since
\[
 q_m(2)=m,\qquad q_m(-2)=(-1)^{m-1}m.
\]

For \(t\in\F_q\), direct counting in the determinant equation gives
\[
 \#\{x\in\SL_2(\F_q):\Tr x=t\}
 =q\bigl(q+\chi(t^2-4)\bigr),
\tag{13.10bpw21}
\]
where \(\chi\) is the quadratic character.  Indeed, after choosing the
upper-left entry \(a\), one counts pairs \((b,c)\) with
\(bc=a(t-a)-1\).  A split root in (13.10bpw20) has
\(\chi(t^2-4)=1\), and a nonsplit root has character \(-1\).
Every root of \(q_m\) has one Hensel lift modulo \(\mathfrak p^e\).
The trace morphism on \(\SL_2\) is smooth above a regular trace, so every
point of its fibre modulo \(\mathfrak p\) has \(q^{2(e-1)}\) lifts in the
fibre of a prescribed lifted trace.  Since
\[
 \#\SL_2(\cO_k/\mathfrak p^e)
 =q^{3(e-1)}q(q^2-1),
\]
equation (13.10bpw17) follows.

Cayley--Hamilton gives
\[
 x^m=q_m(\Tr x)x-q_{m-1}(\Tr x)I.
\tag{13.10bpw22}
\]
If the first coefficient vanishes, Cassini's identity gives
\(q_{m-1}(\Tr x)^2=1\), which proves (13.10bpw19).  Conversely, if
\(x^m\) is scalar and \(x\) is noncentral modulo \(\mathfrak p\), an
off-diagonal entry or the difference of the two diagonal entries in
(13.10bpw22) shows that \(q_m(\Tr x)=0\) modulo \(\mathfrak p^e\).
Finally, write an element in a central residue disc as \(x=\epsilon u\),
where \(\epsilon\in\{I,-I\}\) and \(u\equiv I\pmod{\mathfrak p}\).
If \(x^m\) is scalar, then \(u^m\) is a scalar congruent to \(I\) modulo
\(\mathfrak p\).  Its determinant is one, so it equals \(I\), because the
residue characteristic is odd.  The kernel of reduction is a finite group
of residue-characteristic power order, whereas that characteristic does
not divide \(m\); hence \(u=I\).
\end{proof}

We now state the precise quotient-counting hypothesis used in the sieve.
For each fixed integer \(E\ge2\), suppose that there are
\(\delta_E>0\) and \(C_E>0\) such that, for every good prime ideal
\(\mathfrak p\), every \(1\le e\le E\), and every residue class \(x\) in
the congruence image,
\[
 \frac{\#\{g\in\mathcal B_R:g\equiv x\pmod{\mathfrak p^e}\}}
      {\#\mathcal B_R}
 \le \frac{C_0}{\#\SL_2(\cO_k/\mathfrak p^e)}
       +C_Ee^{-\delta_E R}.
\tag{13.10bpw23}
\]
Here \(C_0\) is independent of \(E\), \(e\), and \(\mathfrak p\), and
\(\#\mathcal B_R=e^{R+O(1)}\).  A uniformly bounded congruence-image
index gives the constant \(C_0\).  In the present setting,
(13.10bpw23) follows, after enlarging the fixed set of bad primes, from
superapproximation for these fixed powers and uniform lattice-point counting
\cite{SalehiSuperII,GorodnikNevoCounting}.  Only a fixed value of \(E\)
is used at any one time.  In particular, no assertion about an unbounded
prime-power exponent is being invoked.

\begin{theorem}[localization of repeated prime factors from bounded powers]
\label{thm:bounded-power-squareful-localization}
Fix \(m\ge2\), enlarge \(S\) by the primes below \(2m\), and put
\[
 A_m(g)=q_m(\Tr g).
\]
Assume (13.10bpw23).  For a fixed \(D\ge2\), define
\[
 \mathcal E^{(D)}_{m,z}(g)=
 \sum_{\substack{\mathfrak p\notin S\\N\mathfrak p\le z}}
 \min\{(v_{\mathfrak p}(A_m(g))-1)_+,D-1\}
       \log N\mathfrak p.
\tag{13.10bpw24}
\]
If \(0<\alpha<\delta_D/(2D+1)\), then
\[
 \frac1{\#\mathcal B_R}\sum_{g\in\mathcal B_R}
 \mathcal E^{(D)}_{m,e^{\alpha R}}(g)
 =O_{m,D,k}(1).
\tag{13.10bpw25}
\]
The average is over group elements, counted with multiplicity; it is not
an average over distinct trace values.  It follows that a proportion
\(1-o(1)\) of the elements satisfy
\[
 \mathcal E^{(D)}_{m,e^{\alpha R}}(g)=o(R).
\tag{13.10bpw26}
\]

Suppose in addition that the standard height estimate
\[
 \sum_{\mathfrak p\notin S}
 v_{\mathfrak p}(A_m(g))\log N\mathfrak p
 \le H_mR+O_m(1)\qquad(g\in\mathcal B_R)
\tag{13.10bpw27}
\]
holds.  This estimate is automatic for a fixed polynomial in the matrix
entries of a finitely generated \(S\)-integral group.  Then, for every
\(\epsilon>0\), there is a fixed \(D=D(\epsilon,m,k)\), followed by a
fixed \(\alpha=\alpha(\epsilon,m,k,\Gamma)>0\), such that
\[
 \limsup_{R\to\infty}\frac1{R\#\mathcal B_R}
 \sum_{g\in\mathcal B_R}
 \sum_{\substack{\mathfrak p\notin S\\
                  N\mathfrak p\le e^{\alpha R}}}
 (v_{\mathfrak p}(A_m(g))-1)_+\log N\mathfrak p
 <\epsilon.
\tag{13.10bpw28}
\]
Consequently, after first making the right side smaller than
\(\epsilon^2\), Markov's inequality shows that a proportion at least
\(1-\epsilon-o(1)\) of the group elements have repeated-prime-factor
contribution at most \(\epsilon R\) from primes of norm at most
\(e^{\alpha R}\).
\end{theorem}

\begin{proof}
Lemma~\ref{lem:chebyshev-local-power-density} and (13.10bpw23) give,
uniformly for \(2\le j\le D\),
\[
 \frac{\#\{g\in\mathcal B_R:
          v_{\mathfrak p}(A_m(g))\ge j\}}
      {\#\mathcal B_R}
 \le C_0C_m(N\mathfrak p)^{-j}
       +C_{m,D}(N\mathfrak p)^{2j}e^{-\delta_D R},
\tag{13.10bpw29}
\]
where \(C_m\) is independent of \(D\) and \(j\).
The elementary identity
\[
 \min\{(v-1)_+,D-1\}=\sum_{j=2}^D\boldsymbol 1_{v\ge j}
\]
therefore reduces the average in (13.10bpw25) to a sum of (13.10bpw29).
The main term is bounded because
\[
 \sum_{\mathfrak p}\frac{\log N\mathfrak p}
                          {(N\mathfrak p)^j}<\infty
 \qquad(j\ge2),
\]
and prime-ideal counting bounds the total error by
\[
 O_{m,D,k}\bigl(e^{-\delta_D R}z^{2D+1}\bigr).
\]
This proves (13.10bpw25).  Taking the exceptional threshold to be
\(\sqrt R\), for example, proves (13.10bpw26) by Markov's inequality.

For the second assertion use (13.10bpw23) only through the fixed depth
\(D+1\).  The part above the cutoff in (13.10bpw24) is bounded pointwise,
using (13.10bpw27), by
\[
 (H_mR+O_m(1))
 \boldsymbol 1_{\{\exists\mathfrak p:\,
        N\mathfrak p\le z,\ v_{\mathfrak p}(A_m(g))\ge D+1\}}.
\tag{13.10bpw30}
\]
The union bound and (13.10bpw29), now at depth \(D+1\), show that if
\(0<\alpha<\delta_{D+1}/(2D+3)\), then
\[
 \begin{split}
 &\frac{\#\{g\in\mathcal B_R:\exists\mathfrak p,\,
       N\mathfrak p\le e^{\alpha R},\,
       v_{\mathfrak p}(A_m(g))\ge D+1\}}
       {\#\mathcal B_R}\\
 &\hspace{25mm}\le C_0C_m
 \sum_{\mathfrak p\notin S}(N\mathfrak p)^{-(D+1)}+o(1).
\end{split}
\tag{13.10bpw31}
\]
The prime-ideal sum tends to zero as the fixed integer \(D\) tends to
infinity.  Choose \(D\) first so that (13.10bpw30)--(13.10bpw31) contribute
less than \(\epsilon R\) on average, then choose \(\alpha\), and only then
let \(R\) tend to infinity.  The truncated sum is \(o(R)\) on average by
(13.10bpw25).  This proves (13.10bpw28).  Notice that no uniform lower
bound for \(\delta_D\), and no choice \(D=D(R)\), is required.
\end{proof}

The order of quantifiers in Theorem
\ref{thm:bounded-power-squareful-localization} is essential.  The exponent
\(\alpha\) may decrease when \(D\) is increased.  Moreover, (13.10bpw27)
only confines a prime which occurs to exponent at least two to
\[
 N\mathfrak p\le \exp(H_mR/2+O_m(1)).
\]
Thus the theorem leaves the interval
\[
 e^{\alpha R}<N\mathfrak p
 \le \exp(H_mR/2+O_m(1))
\tag{13.10bpw32}
\]
untouched.  A single prime in this interval can account for a positive
multiple of \(R\) in the difference between the logarithmic norm and the
logarithm of the radical.

The following elementary family shows that the uncontrolled interval cannot
be removed by a formal refinement of the preceding estimate.  Fix
\[
 0<\alpha<\beta<H/2,\qquad 0<h<\beta.
\]
For all sufficiently large \(R\), the prime number theorem gives
\(\lfloor e^{hR}\rfloor\) distinct primes
\(p_i\in[e^{\beta R},2e^{\beta R}]\).  The integers
\(a_i=p_i^2\) satisfy
\[
 \log a_i\le HR+O(1),\qquad
 \log\frac{a_i}{\operatorname{rad}(a_i)}
   =\beta R+O(1),
\]
but no prime at most \(e^{\alpha R}\) square-divides any \(a_i\), and
the repeated-prime supports of distinct \(a_i\) are disjoint.  Hence even
perfect information below \(e^{\alpha R}\), together with all estimates
for pairs or larger tuples sharing the same repeated prime, does not force
one member of such a family to have a nearly full radical.  If one also
requires compatibility with the Chebyshev congruence, choose
\(p_i\equiv1\pmod{2m}\); equation (13.10bpw20) and Hensel lifting then
give regular classes \(t_i\pmod{p_i^2}\) for which
\(p_i^2\mid q_m(t_i)\).  Removing the interval (13.10bpw32) therefore
requires an additional large-prime argument, not a further consequence of
superapproximation for fixed powers of square-free moduli.

We avoid this large-prime problem by fixing the recurrence index and varying
the group element whose trace enters the recurrence.  Roth's theorem then
gives a lower bound for the reduced divisor.  Expansion modulo squares is
enough to choose the varying element so that no prime occurs to a large
power.  Thus the available expansion theorem for fixed powers suffices.

We begin with the precise Diophantine statement.  Absolute values on a number
field \(k\) are normalized to extend the standard absolute values on
\(\Q\), and \(n_v=[k_v:\Q_v]\).  Thus
\(\sum_v n_v\log|a|_v=0\), and
\[
 [k:\Q]h(a)=\sum_v n_v\log^+|a|_v
 =\frac12\sum_v n_v|\log|a|_v|.
\]

\begin{lemma}[Roth's theorem for a reduced numerator--denominator divisor]
\label{lem:roth-reduced-divisor}
Let \(k\) be a number field, let \(S\) be a finite set of finite places,
and let \(f\in k(T)\) have degree \(d\ge2\).  Suppose that every zero and
every pole of \(f\) on \(\mathbb P^1_{\overline k}\) is simple.  For
\(t\in k\) which is neither a zero nor a pole, put
\[
 \mathfrak m_{f,S}(t)=
 \prod_{\mathfrak p\notin S}
 \mathfrak p^{|\operatorname{ord}_{\mathfrak p}f(t)|}.
\]
For every \(\epsilon>0\), outside a finite set of \(t\in k\),
\[
 \log N\mathfrak m_{f,S}(t)
 \ge (2d-2-\epsilon)[k:\Q]h(t)-O_{f,S,\epsilon}(1).
\tag{13.10mov1}
\]
The exceptional set and the implied constant may depend on
\(f,S,\epsilon\), which remain fixed while \(t\) varies.
\end{lemma}

\begin{proof}
Let \(T\) consist of \(S\) and all archimedean places.  The product formula
and height functoriality give
\[
 \sum_{v\in M_k}n_v|\log|f(t)|_v|
 =2[k:\Q]h(f(t))
 =2d[k:\Q]h(t)+O_f(1).
\tag{13.10mov2}
\]
At a place \(v\in T\), the quantity
\(|\log|f(t)|_v|\), up to a bounded local term, is a Weil proximity
function to one member of the finite set consisting of the zeros and poles
of \(f\).  Simplicity is important here: every target occurs with
multiplicity one.  Pass once to a fixed finite extension \(L/k\) which
contains every zero and pole, and sum over all extensions \(w\mid v\) of
each \(v\in T\).  Roth's theorem over \(L\), divided by \([L:k]\),
retains the coefficient \(2\) with the base-field normalization:
\([L:\Q]h(t)/[L:k]=[k:\Q]h(t)\) for \(t\in k\), and the local degrees
above each \(v\) sum to \([L:k]n_v\).  There are only finitely many ways to choose one such
target at each place of \(L\) above \(T\).  The number-field Roth theorem, applied to
each choice, gives
\[
 \sum_{v\in T}n_v|\log|f(t)|_v|
 \le (2+\epsilon)[k:\Q]h(t)+O_{f,S,\epsilon}(1)
\tag{13.10mov3}
\]
outside a finite set.  We use here the simultaneous-place form of Roth's
theorem, for example \cite[Theorem~6.4.1]{BombieriGubler}.  The finite-place
part outside \(S\) of the left side of (13.10mov2) is exactly
\(\log N\mathfrak m_{f,S}(t)\).  Subtracting (13.10mov3) from
(13.10mov2), and replacing \(\epsilon\) by a smaller positive number if
necessary, proves (13.10mov1).
\end{proof}

For \(N\ge3\), put
\[
 P_N(T)=q_{N-1}(T),\qquad Q_N(T)=q_N(T),\qquad
 F_N(T)=P_N(T)Q_N(T).
\]
The rational function \(P_N/Q_N\) has degree \(N-1\).  Its finite zeros
and poles are simple and disjoint, and infinity is one further simple zero.
Indeed, the Chebyshev polynomials \(q_j\) are square-free in characteristic
zero, while Cassini's identity makes two consecutive terms coprime.

\begin{corollary}[the reduced divisor of two consecutive recurrence terms]
\label{cor:roth-chebyshev-pair}
Let \(k,S,N\) be fixed, with \(N\ge3\), and enlarge \(S\) so that the
arguments under consideration are \(S\)-integral.  For every
\(\epsilon>0\), outside a finite set of \(t\in k\),
\[
 \log N\bigl(F_N(t)\cO_{k,S}\bigr)
 \ge(2N-4-\epsilon)\mathcal J_k(t)-O_{N,S,\epsilon}(1).
\tag{13.10mov4}
\]
Here and below the norm of a principal \(S\)-ideal means the norm of its
prime-to-\(S\) part.
\end{corollary}

\begin{proof}
Cassini's identity says that \(P_N(t)\) and \(Q_N(t)\) generate the unit
ideal of \(\cO_{k,S}\).  Consequently
\(\mathfrak m_{P_N/Q_N,S}(t)=F_N(t)\cO_{k,S}\).  Apply
Lemma~\ref{lem:roth-reduced-divisor} with \(d=N-1\).  The local definition
(13.10au) gives, uniformly in \(t\in k\),
\[
 \bigl|\mathcal J_k(t)-[k:\Q]h(t)\bigr|
 \le[k:\Q]\log2:
\]
at a nonarchimedean place the two sides have the same local contribution,
and at each archimedean place their difference is bounded.  This converts
(13.10mov1) into (13.10mov4).
\end{proof}

We next use expansion modulo squares to sieve the group elements.  The order
of quantifiers is important.  Let \(x\) be a faithful cocompact Fuchsian character of
a closed surface group \(\Pi\), and let \(y\) be an algebraic faithful
cocompact Fuchsian character in the rational closure of \(x\).  Thus every
squared-trace equality at \(x\) also holds at \(y\).  Write
\(\Delta=\widetilde{\Pi^{(2)}}\), let \(k_y\) be the full trace field of
the representation \(y\), and use a fixed finite generating set of
\(\Delta\) for word length.

\begin{lemma}[selection of elements with bounded recurrence exponents]
\label{lem:moving-bounded-recurrence-exponents}
Fix a noncentral \(b\in\Delta\), an integer \(N\ge3\), and a sufficiently
small fixed \(\theta>0\).  There are a rationally saturated finite set
\(S=S(N,\theta,b,y)\) of finite places of \(k_y\), an integer
\(D=D(N,\theta,b,y)\), and, for every sufficiently large \(m\), an element
\[
 a_m=b^m h_m,\qquad |h_m|\le\theta m,
\tag{13.10mov5}
\]
with the following properties.  If \(s_m=\Tr_y(a_m)\), then
\[
 \mathcal J_{k_y}(s_m)
 \ge m\mathcal J_{k_y}(\Tr_y b)-O_{b,y}(\theta m)-O_{b,y}(1),
\tag{13.10mov6}
\]
and, for a fixed basepoint \(o\) on the \(x\)-axis of \(b\),
\[
 D_x(a_m):=d_x(o,a_mo)
 \le m\ell_x(b)+O_{b,x}(\theta m)+O_{b,x}(1).
\tag{13.10mov7}
\]
Moreover, \(F_N(s_m)\ne0\), every exponent in the prime-to-\(S\) ideal
\[
 \mathfrak c_m=F_N(s_m)\cO_{k_y,S}
\tag{13.10mov8}
\]
is at most \(D\), and
\[
 \log N\mathfrak c_m=O_{N,b,y}(m).
\tag{13.10mov9}
\]
Only expansion modulo the square of a rational prime is used in choosing
\(a_m\).  In particular, neither \(D\)-th-power expansion nor a lower
bound for a spectral gap depending on \(D\) is used.
\end{lemma}

\begin{proof}
First enlarge a rationally saturated set \(S_0\) so that the representation
is integral, the restriction-of-scalars model has good reduction outside
\(S_0\), and the primes below \(2N(N-1)\) are removed.  Let \(h\) be
sampled from the distribution of the lazy random walk on \(\Delta\) after
\(L=\lfloor\theta m\rfloor\) steps, and consider \(a=b^mh\).

There are only finitely many places \(v\) of \(k_y\) for which
\(\chi_v(\Tr_yb)>0\).  Over one fixed finite extension which splits the
quaternion algebra and the characteristic polynomial of \(b\), choose at
each such place the eigenvalue \(\lambda_v\) with \(|\lambda_v|_v>1\).
Then
\[
 \Tr_y(b^mh)=\lambda_v^mL_{v,+}(h)
              +\lambda_v^{-m}L_{v,-}(h),
\tag{13.10mov10}
\]
where \(L_{v,+}\) and \(L_{v,-}\) are fixed regular linear functions of
the matrix entries of \(h\), and \(L_{v,+}(I)=1\).  Take the product of
the Galois norms of all the \(L_{v,+}\).  After clearing a fixed
denominator this is a nonzero regular function \(\Phi\), defined over
\(\Q\), on the restriction-of-scalars group.

Choose one sufficiently large good rational prime \(r\).  The Lang--Weil
estimate \cite{LangWeil} and the bounded index of the congruence image make the proportion of its
reduction on which \(\Phi=0\) smaller than \(1/8\).  Enlarge \(S_0\) by
all places above \(r\).  Mixing on this one fixed quotient then gives
\[
 \Pr(\Phi(h)=0)<\frac14
\tag{13.10mov11}
\]
for all sufficiently large \(m\).  If \(\Phi(h)\ne0\), the height of each
nonzero \(L_{v,+}(h)\), and the height of \(L_{v,-}(h)\), is \(O(L)\).
The product formula therefore bounds the negative logarithm of the former
and the positive logarithm of the latter at \(v\) by \(O(L)\).  For
\(\theta\) sufficiently small, the first term of (13.10mov10) dominates
the second and gives
\[
 \chi_v(\Tr_y(b^mh))
 \ge m\chi_v(\Tr_yb)-O_{b,y}(L)-O_{b,y}(1).
\]
Summing over the finitely many expanding places proves (13.10mov6).
The triangle inequality and the bounded displacement of the fixed
generators give (13.10mov7).

It remains to impose bounded exponents without destroying (13.10mov11).
For a value of \(h\) with \(\Phi(h)\ne0\), formula (13.10mov6) already makes
\(\mathcal J_{k_y}(\Tr_y(b^mh))\) grow linearly.  Thus, for large \(m\),
its trace avoids the fixed zero set of \(F_N\).
The word complexity of \(b^mh\) is \(O(m)\), so the usual height estimate
for the fixed polynomial \(F_N(\Tr_y(\,\cdot\,))\) gives a constant \(A_N\)
such that
\[
 \log N\bigl(F_N(\Tr_y(b^mh))\cO_{k_y,S_0}\bigr)
 \le A_Nm+O_N(1)
\tag{13.10mov12}
\]
whenever the value is nonzero.  Choose a large constant \(Z\), and enlarge
\(S_0\) by every place above a rational prime at most \(Z\); call the
result \(S\).

Outside \(S\), the polynomial \(F_N=P_NQ_N\) is square-free.  The local
calculation of Lemma~\ref{lem:chebyshev-local-power-density}, applied to
its two disjoint factors and followed by Hensel lifting, shows that for a
good rational prime \(p\)
\[
 \Pr\bigl(\mathfrak p^2\mid
 F_N(\Tr_y(b^mh))\text{ for some }\mathfrak p\mid p\bigr)
 \le C_Np^{-2}+p^B e^{-\delta L}.
\tag{13.10mov13}
\]
Here the union over all prime ideals above \(p\) is taken inside the single
restriction-of-scalars quotient modulo \(p^2\).  More explicitly,
\(p^2\cO_{k_y}\subseteq\mathfrak p^2\), so reduction modulo rational
\(p^2\) surjects onto reduction modulo every \(\mathfrak p^2\) above it.
The constants
\(B,C_N,\delta>0\) are fixed.  Salehi Golsefidy's expansion theorem for
moduli \(p^2\) applies to that quotient, starting the walk at the arbitrary point
\(b^m\); summing the pointwise mixing error over the local divisor costs
only the fixed power \(p^B\), and the bounded index of the congruence image
is absorbed into \(C_N\).  Lemma
\ref{lem:full-trace-field-res-closure} establishes the required connected
perfect Zariski closure.  No modulus \(p^e\) with \(e>2\) occurs here.

If an exponent in (13.10mov8) is greater than \(D\), (13.10mov12) shows
that the rational prime below it is at most
\[
 Y_m=\exp\left(\frac{A_Nm+O_N(1)}{D+1}\right).
\]
Indeed,
\((D+1)\log N\mathfrak p\le\log N\mathfrak c_m
\le A_Nm+O_N(1)\), and the rational prime \(p\) below
\(\mathfrak p\) satisfies \(p\le N\mathfrak p\le Y_m\).
The union bound in (13.10mov13) gives
\[
 \begin{split}
 \Pr(&\text{some prime-to-}S\text{ exponent exceeds }D)\\
 &\le C_N\sum_{p>Z}p^{-2}
 +\exp\left(
   \left(\frac{(B+1)A_N}{D+1}-\delta\theta\right)m+o(m)
          \right).
 \end{split}
\tag{13.10mov14}
\]
Choose \(Z\) first so that the first term is less than \(1/8\), and then
choose the fixed integer \(D\) so that the exponential term tends to zero.
For large \(m\), (13.10mov14) is less than \(1/4\).  Its union with the
event in (13.10mov11) still has probability less than one, so there is a
value \(h_m\) satisfying all the assertions.  Formula (13.10mov6) tends
to infinity linearly; hence \(s_m\) eventually avoids the roots of \(F_N\)
and any prescribed finite set.  Finally (13.10mov9) follows from
(13.10mov12).
\end{proof}

The next lemma gives a fiber bound uniform in the varying multiplier
\(a_m\).  Its proof avoids diagonalizing \(a_m\) over a field that depends
on \(m\).

\begin{lemma}[uniform fibers for a varying pair of trace coordinates]
\label{lem:moving-two-trace-fibers}
Retain \(x,y,\Delta\) as above.  Suppose that \(A_m\in\Delta\) has word
complexity \(O(m)\) and is noncentral.  For each fixed \(C>0\), uniformly
in \(\xi,\zeta\in k_y\),
\[
 \#\{H\in\Delta:d_x(o,Ho)\le Cm,
       \ \Tr_yH=\xi,\ \Tr_y(A_mH)=\zeta\}
 =\exp(o_C(m)).
\tag{13.10mov15}
\]
The same conclusion holds for squared traces, with only the four possible
choices of signs added.
\end{lemma}

\begin{proof}
Pass once to a fixed finite splitting field for the invariant quaternion
algebra and use the fixed arithmetic filtration of Section~12.  Write
\[
 A=\begin{pmatrix}a&b\\c&d\end{pmatrix},\qquad
 H=\begin{pmatrix}r&s\\u&\xi-r\end{pmatrix},\qquad
 w=\zeta-d\xi,
\]
and put
\[
 \Delta_A=(a+d)^2-4,\qquad
 y_0=2br-(a-d)s,\qquad
 K=b\{\xi(a-d)-2w\}.
\]
Since \(A\) is a noncentral element of a cocompact Fuchsian group,
\(\Delta_A\ne0\).  If \(b\ne0\), direct elimination using
\(\det H=1\) and \(\Tr(AH)=\zeta\) gives
\[
 (\Delta_As+K)^2-\Delta_A(y_0-b\xi)^2
 =K^2-\Delta_Ab^2(\xi^2-4).
\tag{13.10mov16}
\]
For fixed \(A,\xi,\zeta\), the map
\[
 H\longmapsto(Y,V)=(y_0-b\xi,\Delta_As+K)
\]
is injective: \(V\) recovers \(s\), \(Y\) recovers \(r\), and the trace
equation recovers \(u\).  All coefficients and all possible solutions in
(13.10mov16) have complexity \(O_C(m)\).

The right side of (13.10mov16) is
\[
 4b^2\bigl(\Tr[A,H]-2\bigr),
\tag{13.10mov17}
\]
by the Fricke identity.  If it is nonzero,
Lemma~\ref{lem:quadratic-norm-divisor} gives \(\exp(o_C(m))\) possible
pairs \((Y,V)\), and hence as many \(H\).  If it is zero, then
\(\Tr[A,H]=2\).  The faithful cocompact Fuchsian image has no nontrivial
parabolic or unipotent element, so \([A,H]=I\).  The centralizer of a
nonidentity element of a surface group is cyclic; fixing \(\Tr H=\xi\)
leaves \(O(1)\) choices.

If \(b=0\) and \(c\ne0\), transpose the calculation.  If \(b=c=0\), the
two trace equations determine both diagonal entries of \(H\), and
\(\det H=1\) fixes the product \(su\).  A nonzero product has
\(\exp(o_C(m))\) factorizations by
Theorem~\ref{thm:finite-type-divisor}.  A zero product would make the
commutator unipotent; cocompactness again forces commutation and leaves
only \(O(1)\) choices.  This proves (13.10mov15).  Replacing two squared
trace conditions by their four signed versions proves the last assertion.
\end{proof}

We now combine the bounded exponent condition with the congruence count in
Lemma~\ref{lem:congruence-wedge}.  The basepoint remains fixed throughout;
neither the axis nor the conjugation varies.

\begin{theorem}[upper bound for bounded-exponent recurrence ideals]
\label{thm:moving-congruence-wedge}
Retain \(x,y,\Delta\), and suppose that the signed trace set at \(x\) has
linear growth.  Fix \(N\ge3\).  Let \(a_m\in\Delta\) have word complexity
\(O(m)\), assume that \(a_m\) is noncentral for all sufficiently large
\(m\), put \(s_m=\Tr_y(a_m)\), and suppose that
\[
 D_x(a_m)=d_x(o,a_mo)=O(m).
\]
Let \(S\) and \(D\) be fixed, with \(S\) rationally saturated and containing
all bad and ramified primes, and assume that the nonzero ideals
\[
 \mathfrak c_m=
 \bigl(q_{N-1}(s_m)q_N(s_m)\bigr)\cO_{k_y,S}
\]
have every prime exponent at most \(D\) and satisfy
\(\log N\mathfrak c_m=O(m)\).  Then
\[
 \log N\mathfrak c_m
 \le\frac{2N-1}{2}D_x(a_m)+o(m).
\tag{13.10mov18}
\]
The error may depend on \(N,S,D,x,y\) and on fixed implied constants in
the word-complexity, displacement, and ideal-norm bounds for \(a_m\).  All
of these are fixed before \(m\to\infty\).
\end{theorem}

\begin{proof}
Let \(M_m\) be the square-free product of the rational primes below
\(\mathfrak c_m\).  Rational saturation and the removal of ramified primes
give
\[
 M_m^D\cO_{k_y,S}\subseteq\mathfrak c_m.
\tag{13.10mov19}
\]
Thus the quotient modulo \(\mathfrak c_m\) is a quotient of the reduction
modulo the fixed power \(M_m^D\) of a square-free rational modulus.
Lemma~\ref{lem:full-trace-field-res-closure} and bounded-power
superapproximation \cite{SalehiSuperII}, together with the
Brooks--Burger transfer and uniform lattice counting used after
Proposition~\ref{prop:fixed-ball-specialization}, imply the following.
There is a sufficiently large fixed \(C\), depending also on the fixed
implied complexity and norm bounds for the sequence \(a_m\), such that, for
\(R=Cm\),
\[
 \max_c\#\{g\in\mathcal B_R^x:
       \Tr_y g\equiv c\pmod{\mathfrak c_m}\}
 \le\frac{e^{R+o(m)}}{N\mathfrak c_m}.
\tag{13.10mov20}
\]
Here is the quantitative point in this invocation.  The local trace fiber
has relative size
\((N\mathfrak c_m)^{-1+o(1)}\), because all local exponents are at most
\(D\).  The expander error is at most a fixed power of
\(N(M_m^D\cO_{k_y,S})\) times \(e^{-\delta_DR}\).  Both this norm and
\(N\mathfrak c_m\) have logarithm \(O_D(m)\), so choosing \(C\) after
\(D\) absorbs the error.  No uniformity of \(\delta_D\) as \(D\to\infty\)
is asserted or needed.

Form the simple bipartite graph, with the trace involving \(a_m\) in the
first coordinate,
\[
 \mathcal G_m=
 \{(u,v)=(\Tr_y(a_mg),\Tr_yg):g\in\mathcal B_R^x\}.
\tag{13.10mov21}
\]
Cocompact orbit counting gives \(e^{R+O(1)}\) elements in the ball, and
Lemma~\ref{lem:moving-two-trace-fibers} bounds every fiber of the two
coordinates by \(e^{o(m)}\).  Hence
\[
 \#\mathcal G_m=e^{R-o(m)}.
\tag{13.10mov22}
\]
Let \(V_m\) be the second vertex class.  Persistence of squared-trace
collisions from \(x\) to \(y\), with the two possible signs restored, and
linear trace growth at \(x\) give
\[
 \#V_m\le e^{R/2+O(1)}.
\tag{13.10mov23}
\]

Delete the vertices of \(V_m\) whose degree is less than
\(\#\mathcal G_m/(4\#V_m)\), and let \(\mathcal G'_m\) be the remaining
graph.  At most one quarter of the edges are deleted.  Choose one group
element representing each remaining edge.  Equation (13.10mov20) then
implies
\[
 M_{\mathfrak c_m}(V'_m):=
 \max_c\#\{v\in V'_m:v\equiv c\pmod{\mathfrak c_m}\}
 \le\frac{e^{R/2+o(m)}}{N\mathfrak c_m}.
\tag{13.10mov24}
\]
Indeed, every retained \(v\) occurs in at least
\(\#\mathcal G_m/(4\#V_m)=e^{R/2-o(m)}\) distinct chosen representatives
to the same congruence class.

Apply Lemma~\ref{lem:congruence-wedge} over \(\cO_{k_y,S}\) to
\(\mathcal G'_m\), with
\[
 P=q_{N-1}(s_m),\qquad Q=q_N(s_m).
\]
Cassini's identity makes \(P,Q\) coprime.  Equations
(13.10mov22)--(13.10mov24) give
\[
 \#X_-\,\#X_+
 \ge\frac{(\#\mathcal G'_m)^2}
          {\#V'_mM_{\mathfrak c_m}(V'_m)}
 \ge e^{R-o(m)}N\mathfrak c_m.
\tag{13.10mov25}
\]
The recurrence identities identify the two images with actual traces:
\[
 \begin{aligned}
 Q\Tr_y(a_mg)-P\Tr_yg&=\Tr_y(a_m^Ng),\\
 Q\Tr_yg-P\Tr_y(a_mg)&=\Tr_y(a_m^{-(N-1)}g).
 \end{aligned}
\tag{13.10mov26}
\]
Since the basepoint is fixed,
\[
 d_x(o,a_m^jgo)\le R+jD_x(a_m)\qquad(j\ge0),
\]
and the analogous inequality holds for negative powers.  Persistence and
linear trace growth therefore give
\[
 \#X_-\,\#X_+
 \le
 \exp\left(R+\frac{2N-1}{2}D_x(a_m)+O(1)\right).
\tag{13.10mov27}
\]
Comparison of (13.10mov25) and (13.10mov27) proves
(13.10mov18).
\end{proof}

As \(N\to\infty\), the coefficient in Roth's lower bound and the
coefficient in the geometric upper bound have the same leading term.  Their
comparison gives the following unconditional estimate.

\begin{theorem}[recurrence height at algebraic specializations preserving trace equalities]
\label{thm:persistent-local-expansion-bound}
Let \(x\) be a faithful cocompact Fuchsian character of a closed surface
group \(\Pi\), and suppose that its signed trace set has linear growth.
Let \(y\) be any algebraic faithful cocompact Fuchsian character in the
rational closure of \(x\), and let \(k_y\) be the full trace field of the
subgroup generated by squares.  Then every nonidentity
\(b\in\widetilde{\Pi^{(2)}}\) satisfies
\[
 \mathcal J_{k_y}(\Tr_yb)\le\frac{\ell_x(b)}2.
\tag{13.10mov28}
\]
\end{theorem}

\begin{proof}
If \(b\) is central, then \(b=\pm I\),
\(\mathcal J_{k_y}(\Tr_yb)=0=\ell_x(b)\), and there is nothing to prove.
Assume henceforth that \(b\) is noncentral.
Fix in this order an integer \(N\ge3\), positive numbers
\(\epsilon,\theta\), with \(\theta\) sufficiently small, and then the set
\(S\), exponent \(D\), and sequence \(a_m\) given by
Lemma~\ref{lem:moving-bounded-recurrence-exponents}.  Finally choose the
constant \(C\) in Theorem~\ref{thm:moving-congruence-wedge}.  All of these
data remain fixed while \(m\to\infty\).  Formula (13.10mov6) makes
\(s_m\) leave every finite set, so Corollary~\ref{cor:roth-chebyshev-pair}
applies for all sufficiently large \(m\).  Together with (13.10mov18) it
gives
\[
 (2N-4-\epsilon)
 \bigl(m\mathcal J_{k_y}(\Tr_yb)-O(\theta m)-O(1)\bigr)
 \le
 \frac{2N-1}{2}
 \bigl(m\ell_x(b)+O(\theta m)+O(1)\bigr)+o(m).
\tag{13.10mov29}
\]
Let \(m\to\infty\), then \(\theta\to0\), then
\(\epsilon\to0\), and finally \(N\to\infty\).  This proves
(13.10mov28).  Equivalently, the complete quantifier order is: fix
\(N,\epsilon,\theta\); choose the fixed sieve cutoff, \(S,D,C\); let
\(m\to\infty\); then let \(\theta,\epsilon\to0\); and only then let
\(N\to\infty\).
\end{proof}

\begin{theorem}[trace sparsity forces arithmeticity]
\label{thm:unconditional-compact-trace-sparsity}
Let \(\Gamma<\PSL_2(\R)\) be torsion-free and cocompact.  If its signed
trace set has linear growth, then \(\Gamma\) is arithmetic.  Consequently
the compact Schmutz linear-growth conjecture and the compact Sarnak
bounded-clustering conjecture hold.
\end{theorem}

\begin{proof}
Suppose first that the Fuchsian character \(x\) is algebraic and take
\(y=x\).  Theorem~\ref{thm:persistent-local-expansion-bound} gives the
upper bound (13.10mov28), while the distinguished real place contributes
\(\ell_x(b)/2\) and every other local contribution is nonnegative.  Thus
equality holds, all finite-place
contributions vanish, and all nondistinguished archimedean contributions
vanish.  Apply this to every nonidentity element of
\(\widetilde{\Gamma^{(2)}}\).  All its traces are algebraic integers, its
full trace field is totally real, and every nondistinguished embedding
sends every trace into \([-2,2]\).  Takeuchi's criterion \cite{Takeuchi} makes
\(\Gamma^{(2)}\), and hence \(\Gamma\), arithmetic.

Suppose now that \(x\) is transcendental.  Proposition
\ref{prop:bounded-degree-specialization} gives an algebraic faithful
Fuchsian character \(y\ne x\), arbitrarily close to \(x\) in the same
Fuchsian component, on which all squared-trace collisions at \(x\)
persist.  Theorem~\ref{thm:persistent-local-expansion-bound} gives, for
every nonidentity \(b\in\widetilde{\Pi^{(2)}}\),
\[
 \mathcal J_{k_y}(\Tr_yb)\le\frac{\ell_x(b)}2.
\tag{13.10mov30}
\]
Lemma~\ref{lem:local-expansion-fuchsian-rigidity} therefore gives \(y=x\),
contradicting the choice of \(y\).  Thus the transcendental case cannot
occur, and the algebraic argument proves arithmeticity.  Bounded clustering
implies linear growth, which gives the Sarnak conclusion from the Schmutz
conclusion.
\end{proof}

\begin{remark}[the independent pointwise radical problem]
\label{rem:pointwise-radical-still-open}
Theorem~\ref{thm:unconditional-compact-trace-sparsity} does not prove the
universal pointwise assertion that every fixed Chebyshev sequence has the
full radical rate.  That assertion still contains the Mersenne and
non-Wieferich problem described below.  The proof instead uses the group
family in an essential way: for each fixed recurrence index, it chooses a
nearby multiplier \(b^mh_m\) whose reduced recurrence divisor has bounded
exponents.  The pointwise radical results and examples below therefore
remain separate results, but the proof of the compact trace-sparsity
conjectures does not depend on them.
\end{remark}

Primitive-divisor theorems for Lucas sequences guarantee new prime factors,
but they do not show that the bounded-exponent part in (13.10bps) captures
\(e^{(2\mathcal J_k(s)-o(1))n}\) of the norm.  Exponential control of the
powerful part is closely related to unresolved radical and Wieferich-type
questions.  Theorem~\ref{thm:high-power-concentration} gives additional
information about this problem for the recurrence of a fixed element without
assuming a stronger Lucas-sequence estimate.

The recurrence permits a further sharpening.  Except at finitely many
places, divisibility acquired after a prime first appears is accounted for
by the ordinary prime divisors of the index.  Those later contributions have
only logarithmic total norm.  Any exponential powerful part must therefore
come from primes which already divide the recurrence term to unusually high
order when they first appear.

\begin{lemma}[local valuation formula and the Lucas--Wieferich excess]
\label{lem:lucas-wieferich-excess}
Let \(s\in\cO_{k,S}\), let \(q_0=0,q_1=1\), and let
\(q_{m+1}=sq_m-q_{m-1}\).  Enlarge \(S\) by the prime ideals dividing
\(2(s^2-4)\).  For a prime ideal \(\mathfrak p\notin S\) which divides
some \(q_m(s)\), let \(r_{\mathfrak p}\) be the least positive integer
\(r\) for which \(\mathfrak p\mid q_r(s)\), and put
\[
 w_{\mathfrak p}=v_{\mathfrak p}(q_{r_{\mathfrak p}}(s)).
\]
Then
\[
 \mathfrak p\mid q_m(s)\ \Longleftrightarrow\
 r_{\mathfrak p}\mid m,\qquad
 v_{\mathfrak p}(q_m(s))
 =w_{\mathfrak p}+v_{\mathfrak p}(m/r_{\mathfrak p}).
\tag{13.10bpx}
\]
For \(D\ge1\), define
\[
 \mathfrak W_{m,D}:=
 \prod_{\substack{\mathfrak p\notin S\\r_{\mathfrak p}\mid m}}
 \mathfrak p^{\max\{w_{\mathfrak p}-D,0\}},
 \qquad
 \mathfrak h^{(q)}_{m,D}:=
 \prod_{\mathfrak p\notin S}
 \mathfrak p^{\max\{v_{\mathfrak p}(q_m(s))-D,0\}}.
\]
Then
\[
 \log N\mathfrak h^{(q)}_{m,D}
 \le \log N\mathfrak W_{m,D}
      +[k:\Q]\log m+O_{s,S}(1).
\tag{13.10bpy}
\]
If the finitely many newly excluded places are retained instead, their
total contribution is \(O_{s,S}(\log m)\).
\end{lemma}

\begin{proof}
Let \(L=k(\lambda)\), where
\(\lambda+\lambda^{-1}=s\).  At a place outside \(S\), the extension
\(L/k\) is unramified, and \(\lambda\) and \(\lambda^2-1\) are units.
The closed formula is
\[
 q_m(s)=\lambda^{-(m-1)}
         \frac{\lambda^{2m}-1}{\lambda^2-1}.
\tag{13.10bpz}
\]
Thus \(r_{\mathfrak p}\) is the multiplicative order of the reduction of
\(\lambda^2\) at any place of \(L\) above \(\mathfrak p\).  In the split
case the two reductions are inverse and have the same order.  The local
lifting-the-exponent formula, in odd residue characteristic, applied after
writing \(m=r_{\mathfrak p}u\), proves (13.10bpx).

For every prime in the product,
\[
 \max\{w_{\mathfrak p}+v_{\mathfrak p}(u)-D,0\}
 \le \max\{w_{\mathfrak p}-D,0\}+v_{\mathfrak p}(m).
\]
Summing against \(\log N\mathfrak p\) proves (13.10bpy), since
\(\log N(m\cO_k)=[k:\Q]\log m\).  At each of the finitely many excluded
places, (13.10bpz) and the corresponding local
lifting-the-exponent estimate give
\(v_{\mathfrak p}(q_m(s))=O_{s,\mathfrak p}(\log m)\).
\end{proof}

\begin{theorem}[a nondistinguished expanding place forces large prime-power multiplicities]
\label{thm:wieferich-mass}
Use the notation of Theorem~\ref{thm:high-power-concentration}, and fix
\(D\ge1\).  If the signed trace set has linear growth, then every
hyperbolic \(a\), with \(s=\Tr a\), satisfies
\[
 \liminf_{n\to\infty}\frac1n
 \left(\log N\mathfrak W_{n,D}
       +\log N\mathfrak W_{n-1,D}\right)
 \ge 2\mathcal J_k(s)-\ell(a).
\tag{13.10bpza}
\]
In particular, if \(a\) has a nondistinguished expanding place, then, for
every fixed \(D\), the primes whose initial multiplicities
\(w_{\mathfrak p}\) exceed \(D\) carry positive exponential norm.
Conversely, if for some fixed \(D\)
\[
 \log N\mathfrak W_{n,D}
 +\log N\mathfrak W_{n-1,D}=o(n)
\tag{13.10bpzb}
\]
along an unbounded sequence, then every nondistinguished local expansion
exponent of \(s\) vanishes.
\end{theorem}

\begin{proof}
Cassini's identity makes \((q_n(s))\) and \((q_{n-1}(s))\) coprime.
Apply Lemma~\ref{lem:lucas-wieferich-excess} to both factors.  It gives
\[
 \log N\mathfrak h_{n,D}
 \le \log N\mathfrak W_{n,D}
    +\log N\mathfrak W_{n-1,D}+O_{s,S}(\log n).
\]
The lower bound (13.10bpu) now proves (13.10bpza).  Conversely,
(13.10bpzb) and the same estimate show that
\(\log N\mathfrak h_{n,D}=o(n)\).  Hence
\[
 \log N\mathfrak c_{n,D}
 =2n\mathcal J_k(s)-o(n)
\]
along the chosen sequence.  Comparing this with (13.10bpt) yields
\(2\mathcal J_k(s)\le\ell(a)\).  The contribution of the distinguished
real place gives the reverse inequality, so equality holds and every other nonnegative
local term is zero.
\end{proof}

\begin{proposition}[the one-term radical bound at an algebraic specialization]
\label{prop:fixed-ball-single-term-radical}
Let \(x\) be a faithful cocompact Fuchsian character whose signed trace
set has linear growth, and let \(y\) be an algebraic Fuchsian point in
\(V_x\).  For a nonidentity element \(\bar a\in\Pi^{(2)}\), choose the
lift \(a\in\widetilde{\Pi^{(2)}}\) for which
\[
 s_y=\Tr_y(a)>2,\qquad Q_{n,y}=q_n(s_y).
\]
Let \(k_y\) be the full trace field at \(y\), choose a rationally saturated
bad set \(S_y\), and set
\[
 \mathfrak r_{n,y}
 =\prod_{\substack{\mathfrak p\notin S_y\\
                    v_{\mathfrak p}(Q_{n,y})>0}}\mathfrak p .
\tag{13.10bpzc}
\]
Then
\[
 \limsup_{n\to\infty}\frac1n\log N\mathfrak r_{n,y}
 \le\frac{\ell_x(\bar a)}2.
\tag{13.10bpzcb}
\]
All constants may depend on the fixed specialization \(y\).
\end{proposition}

\begin{proof}
Use the graph \(\mathcal G_{R,y}\) of
Proposition~\ref{prop:fixed-ball-specialization}, with \(a\) as chosen
above.  Persistence of squared-trace collisions at \(x\) bounds each
coordinate set at \(y\) by \(e^{R/2+O_a(1)}\): a signed trace value at
\(x\) can split into at most its two signs at \(y\).  The two-trace
finite-type estimate at \(y\), together with cocompact orbit counting in
the metric \(x\), gives
\[
 \#\mathcal G_{R,y}=e^{R-o(R)}.
\tag{13.10bpzcc1}
\]
Apply Lemma~\ref{lem:drc-explicit} to this graph with the parameters used
in the proof of Theorem~\ref{thm:drc-trace-recurrence}.  It produces a set
\[
 W_{R,y}\subset\{\Tr_y g:g\in\mathcal B_R^x\}
\]
of \(e^{R/2-o(R)}\) left vertices, each incident to
\(e^{R/2-o(R)}\) edges.  Indeed, repeated pairs are allowed in
Lemma~\ref{lem:drc-explicit}, and for the pair \((u,u)\) the common
neighbours are exactly the neighbours of \(u\).  If this minimum degree is
denoted by \(r_{R,y}\), then
\[
 \#W_{R,y}\,r_{R,y}=e^{R-o(R)}.
\tag{13.10bpzcc1a}
\]
Repeating the graph--Ruzsa calculation with the
recurrence identities at \(y\) gives, simultaneously for every \(n\ge2\),
\[
 \frac{\#(W_{R,y}+\beta_{n,y}W_{R,y})}{\#W_{R,y}}
 \le\exp\bigl(n\ell_x(\bar a)+o(R)\bigr),
 \qquad
 \beta_{n,y}=-\frac{q_{n-1}(s_y)^2}{q_n(s_y)^2}.
\tag{13.10bpzcc2}
\]
Indeed, the recurrence images are \(y\)-traces of elements in \(x\)-balls
of radius \(R+n\ell_x(\bar a)+O_a(1)\), and persistence transfers the
linear trace-support bound from \(x\) to \(y\).  This reconstructs the
dependent-random-choice argument in the fixed metric \(x\); it does not
silently change the metric in Theorem~\ref{thm:drc-trace-recurrence}.

Put \(\mathfrak d_{n,y}=\mathfrak r_{n,y}^2\), and let \(m_{n,y}\) be the
square-free product of the rational primes below \(\mathfrak r_{n,y}\).
Rational saturation gives
\[
 m_{n,y}^2\cO_{k_y,S_y}\subseteq\mathfrak d_{n,y}.
\tag{13.10bpzcc3}
\]
Thus reduction modulo \(\mathfrak d_{n,y}\) is a quotient of reduction
modulo the fixed square of the square-free rational modulus \(m_{n,y}\).
Salehi Golsefidy's theorem for fixed powers of square-free rational moduli
applies to the algebraic representation at \(y\)
\cite{SalehiSuperII}.  Its spectral gap transfers to the corresponding covers of the
fixed surface \((S,x)\) by the Brooks--Burger comparison
\cite{BrooksTower,Burger}.  Uniform lattice counting in the metric \(x\)
therefore gives, for some \(\delta_y>0\),
\[
 \max_c\#\{g\in\mathcal B_R^x:
   \Tr_y g\equiv c\pmod{\mathfrak d_{n,y}}\}
 \ll
 (N\mathfrak d_{n,y})^{o(1)}
 \left\{
 \frac{e^R}{N\mathfrak d_{n,y}}
 +(N\mathfrak d_{n,y})^2e^{(1-\delta_y)R}
 \right\}.
\tag{13.10bpzcc4}
\]
The one-factor recurrence norm estimate gives
\[
 \log N\mathfrak d_{n,y}
 \le2n\mathcal J_{k_y}(s_y)+O_y(\log n).
\]
Choose a fixed \(C_y\) with
\(\delta_yC_y>6\mathcal J_{k_y}(s_y)\), and put \(R=C_yn\).  The second
term in (13.10bpzcc4) is then absorbed by the first.

Choose one group representative for every edge incident to \(W_{R,y}\).
Representatives of distinct graph edges are distinct.  If this set of left
vertices occupies \(M_{R,n,y}\) residue classes modulo
\(\mathfrak d_{n,y}\), then (13.10bpzcc1a) and (13.10bpzcc4) give
\[
 M_{R,n,y}
 \max_c\#\{g\in\mathcal B_R^x:
   \Tr_y g\equiv c\pmod{\mathfrak d_{n,y}}\}
 \ge \#W_{R,y}\,r_{R,y}.
\]
Since \(R=C_yn\), it follows that
\[
 M_{R,n,y}\ge N\mathfrak d_{n,y}\,e^{-o(n)}.
\tag{13.10bpzcc5}
\]
Let \(T_{R,n,y}\subset W_{R,y}\) contain one representative from each
occupied class.  Exactly as in (13.10bpw10), the map
\[
 T_{R,n,y}\times W_{R,y}\longrightarrow
 Q_{n,y}^2(W_{R,y}+\beta_{n,y}W_{R,y}),\qquad
 (t,u)\longmapsto
 Q_{n,y}^2u-q_{n-1}(s_y)^2t
\]
is injective.  The ideal \(\mathfrak d_{n,y}\) is supported on primes
dividing \(Q_{n,y}\), so \(Q_{n,y}^2\) vanishes modulo it while Cassini's
identity makes \(q_{n-1}(s_y)\) a unit there.  Consequently
\[
 \frac{\#(W_{R,y}+\beta_{n,y}W_{R,y})}{\#W_{R,y}}
 \ge N\mathfrak d_{n,y}\,e^{-o(n)}.
\]
Comparison with (13.10bpzcc2) gives
\[
 \log N\mathfrak d_{n,y}
 \le n\ell_x(\bar a)+o(n).
\]
Since \(\mathfrak d_{n,y}=\mathfrak r_{n,y}^2\), this is
(13.10bpzcb).
\end{proof}

We next record the exact consequence of the number-field \(abc\)
conjecture.  For a number field \(L\), a finite set \(S_L\) of finite
places, and an \(S_L\)-integer \(\alpha\), write
\[
 \operatorname{rad}_{S_L}(\alpha)
 =\prod_{\substack{\mathfrak P\notin S_L\\
                    v_{\mathfrak P}(\alpha)>0}}\mathfrak P.
\]
We use the standard number-field \(abc\) assertion in the following
fixed-field normalization: for every \(\epsilon>0\) and every nonzero
solution \(A+B=C\),
\[
 [L:\Q]h(A:B:C)
 \le(1+\epsilon)
 \sum_{\substack{\mathfrak P\\
  v_{\mathfrak P}(A),v_{\mathfrak P}(B),v_{\mathfrak P}(C)
  \text{ are not all equal}}}
 \log N\mathfrak P+C_{L,\epsilon}.
\tag{13.10bpzcc}
\]
Deleting a fixed finite set \(S_L\) changes only the constant.  This is the
usual fixed-number-field form of the Masser--Oesterl\'e--Vojta
conjecture~\cite[p.~84]{VojtaDiophantine}; a discriminant term in a
field-uniform formulation is constant here.

\begin{lemma}[\(abc\) gives the full Chebyshev radical]
\label{lem:abc-chebyshev-radical}
Let \(k\) be a number field, let \(s\in\cO_{k,S}\) have a distinguished
real embedding with \(s>2\), and let
\(q_{n+1}=sq_n-q_{n-1}\), \(q_0=0,q_1=1\).  Assume the number-field
\(abc\) conjecture for \(L=k(\lambda)\), where
\(\lambda+\lambda^{-1}=s\).  After enlarging \(S\) by a fixed finite set,
\[
 \log N_k\operatorname{rad}_S(q_n(s))
 =n\mathcal J_k(s)+o_s(n).
\tag{13.10bpzd}
\]
Equivalently, the quotient of the principal ideal \((q_n(s))\) by its
radical has logarithmic norm \(o_s(n)\).
\end{lemma}

\begin{proof}
Put \(z=\lambda^2\).  The distinguished embedding shows that \(z\) is not
a root of unity.  The product formula and the definition (13.10au) give
the global identity
\[
 [L:\Q]h(z)=[L:k]\mathcal J_k(s).
\tag{13.10bpze}
\]
When \(z\notin k\), the two roots \(z,z^{-1}\) give this identity
place by place after extension to \(L\); when \(z\in k\), it follows after
summing because the positive and negative logarithms have equal total mass.

Enlarge \(S\) so that \(\lambda\), \(z-1\), and their inverses are units
above its complement, and so that every prime ramified in \(L/k\) is
removed.  In the group application this enlargement is made after adjoining
the model and denominator primes and is then saturated over rational
primes.  Removing this fixed set changes every logarithmic radical by only
\(O_{s,S}(1)\).  Apply \(abc\) over \(L\) to
\[
 z^n+(-1)=z^n-1.
\]
Outside the enlarged fixed set, only \(z^n-1\) contributes to the radical,
whereas
\[
 h(z^n:-1:z^n-1)=nh(z)+O_z(1).
\]
Thus, for every \(\epsilon>0\),
\[
 \log N_L\operatorname{rad}_{S_L}(z^n-1)
 \ge\frac{n[L:k]\mathcal J_k(s)}{1+\epsilon}
      -O_{s,S,\epsilon}(1).
\tag{13.10bpzf}
\]
The closed formula
\[
 q_n(s)=\lambda^{-(n-1)}\frac{z^n-1}{z-1}
\]
identifies the two radicals outside \(S_L\).  Since \(q_n(s)\in k\) and
\(L/k\) is unramified there, every prime above a prime divisor of \(q_n(s)\)
occurs, and
\[
 \log N_L\operatorname{rad}_{S_L}(q_n(s)\cO_L)
 =[L:k]\log N_k\operatorname{rad}_S(q_n(s)).
\tag{13.10bpzg}
\]
Combine (13.10bpzf)--(13.10bpzg), then let
\(\epsilon\downarrow0\), to obtain the lower bound in (13.10bpzd).  The
local norm calculation in
Lemma~\ref{lem:prime-to-S-recurrence-norm}, applied to one recurrence
factor, gives
\[
 \log N_k(q_n(s))=n\mathcal J_k(s)+O_{s,S}(\log n).
\]
The radical divides this principal ideal, which gives the matching upper
bound and the assertion about the quotient.
\end{proof}

\begin{conjecture}[the one-term Chebyshev radical hypothesis]
\label{conj:chebyshev-radical}
Let \(k\) be a number field, let \(S\) be a finite set of finite places,
and let \(s\in\cO_{k,S}\) have a real embedding at which \(s>2\).  Then
\[
 \limsup_{n\to\infty}\frac1n
 \log N_k\operatorname{rad}_S(q_n(s))
 \ge\mathcal J_k(s).
\tag{\(\mathrm{CR}_{\mathrm{Ch}}\)}
\]
The reverse inequality follows from the one-factor recurrence norm formula
\[
 \log N_k\bigl((q_n(s))\bigr)
 =n\mathcal J_k(s)+O_{s,S}(\log n),
\]
so the assertion is equivalently equality of the two exponential rates.
It is unchanged by enlarging \(S\) by a fixed finite set.
\end{conjecture}

\begin{theorem}[the one-term radical hypothesis implies the compact conjectures]
\label{thm:chebyshev-radical-implies-compact}
Assume Conjecture~\ref{conj:chebyshev-radical}.  Then every torsion-free
cocompact Fuchsian lattice whose signed trace set has linear growth is
arithmetic.  Consequently the compact Schmutz linear-growth conjecture and
the compact Sarnak bounded-clustering conjecture hold.
\end{theorem}

\begin{proof}
First suppose that the Fuchsian character \(x\) is algebraic.  For every
nontrivial projective element of \(\Gamma^{(2)}\), choose its
positive-trace lift \(a\), put \(s=\Tr a>2\), and use the full trace field
\(k\).  Conjecture~\ref{conj:chebyshev-radical} and
Theorem~\ref{thm:single-term-radical-upper} give
\[
 \mathcal J_k(s)
 \le
 \limsup_{n\to\infty}\frac1n
 \log N_k\operatorname{rad}_S(q_n(s))
 \le\frac{\ell(a)}2.
\tag{13.10rad1}
\]
The distinguished real place already contributes \(\ell(a)/2\).
Therefore equality holds and every other local expansion exponent
vanishes.  This holds for every projectively nontrivial element of
\(\Gamma^{(2)}\).  Takeuchi's criterion \cite{Takeuchi} makes
\(\Gamma^{(2)}\) arithmetic, and arithmeticity passes to its finite-index
overgroup \(\Gamma\).

Suppose next that \(x\) is transcendental.  Let \(y_j\to x\) be the
bounded-degree algebraic Fuchsian specializations given by
Proposition~\ref{prop:bounded-degree-specialization}.  Since
\(Z_x\subset V_x\), the squared-trace equalities needed below persist at
every \(y_j\); the proposition also places all \(y_j\) in the same
Teichm\"uller component as \(x\).  For each fixed
\(j\), Proposition~\ref{prop:fixed-ball-single-term-radical} and
Conjecture~\ref{conj:chebyshev-radical} give, for every nonidentity
\(\bar a\in\Pi^{(2)}\),
\[
 \mathcal J_{k_{y_j}}(\Tr_{y_j}a)
 \le\frac{\ell_x(\bar a)}2,
\tag{13.10rad2}
\]
where the lift is chosen to have positive trace at \(y_j\).  No constants
need be uniform in \(j\), because the limit in \(n\) is taken with \(j\)
fixed.
Because \(\mathcal J_k(-s)=\mathcal J_k(s)\), the choice of lift in
(13.10rad2) is immaterial.  Taking any \(y_j\ne x\), the inequality holds
for every nonidentity element of the subgroup generated by squares.
Lemma~\ref{lem:local-expansion-fuchsian-rigidity} gives \(y_j=x\), a
contradiction.  Hence only the algebraic case occurs.  Finally, bounded
clustering implies linear growth.
\end{proof}

\begin{proposition}[specialization to Mersenne numbers]
\label{prop:mersenne-radical-boundary}
Let
\[
 k=\Q(\sqrt2),\qquad s=\frac3{\sqrt2},
\]
and let \(S\) contain the unique prime of \(k\) above \(2\).  Then
\[
 q_n(s)=2^{-(n-1)/2}(2^n-1)
\tag{13.10mer1}
\]
and
\[
 N_k\operatorname{rad}_S(q_n(s))
 =\operatorname{rad}(2^n-1)^2,
 \qquad
 \mathcal J_k(s)=2\log2.
\tag{13.10mer2}
\]
Consequently Conjecture~\ref{conj:chebyshev-radical} at this single value
of \(s\) is exactly
\[
 \limsup_{n\to\infty}\frac1n
 \log\operatorname{rad}(2^n-1)=\log2.
\tag{13.10mer3}
\]
In particular, (13.10mer3) implies that there are infinitely many
non-Wieferich primes to base \(2\).
\end{proposition}

\begin{proof}
Take \(\lambda=\sqrt2\), so that
\(\lambda+\lambda^{-1}=3/\sqrt2\).  The closed recurrence formula gives
(13.10mer1).  The power of \(2\) is an \(S\)-unit.  Every odd rational
prime is unramified in \(k\), and the product of the norms of all primes
above it is its square.  This proves the first identity in (13.10mer2).
Moreover, for the prime-to-\(S\) principal ideal,
\[
 \log N_k\bigl((q_n(s))\bigr)
 =2\log(2^n-1)=2n\log2+O(1).
\]
The one-factor recurrence norm formula therefore gives
\(\mathcal J_k(s)=2\log2\), proving the equivalence with (13.10mer3).

Suppose that only finitely many odd primes \(E\) are non-Wieferich to base
\(2\).  If \(p\notin E\) divides \(2^n-1\), put
\(r_p=\operatorname{ord}_p(2)\).  Then \(r_p\mid n\) and
\(r_p\mid p-1\).  Since \(p\nmid(p-1)/r_p\), the
lifting-the-exponent formula gives
\[
 v_p(2^{r_p}-1)=v_p(2^{p-1}-1)\ge2.
\]
A second application, using \(r_p\mid n\), gives
\(v_p(2^n-1)\ge2\).  Hence
\[
 \log\operatorname{rad}(2^n-1)
 \le\frac12\log(2^n-1)+O_E(1),
\]
contradicting (13.10mer3).  Thus this special case already implies the
classical infinitude of non-Wieferich primes; number-field \(abc\)
implies it by Lemma~\ref{lem:abc-chebyshev-radical}, in agreement with
Silverman's theorem \cite{SilvermanWieferich}.
\end{proof}

Known unconditional recurrence results are far below the exponential rate
in Conjecture~\ref{conj:chebyshev-radical}.  In the classical integral
setting, Stewart's lower bound for a simple nondegenerate recurrence with a
dominant root has the shape
\[
 \operatorname{rad}(u_n)>
 n^{C\log_2 n/\log_3 n},
\]
where \(\log_2 n=\log\log n\) and
\(\log_3 n=\log\log\log n\) \cite{StewartSquarefree}.  Its logarithm is
still \(o(n)\).  His later lower bound for the greatest prime divisor of a
Lucas--Lehmer cyclotomic factor is likewise subexponential on the scale
needed here \cite{StewartLucasDivisors}.  Primitive-divisor theorems guarantee
a new prime at large indices, but one new prime per index does not give a
linear lower bound for the logarithm of the radical.  Thus the Mersenne
specialization is not hiding a known Lucas-sequence theorem.

\begin{proposition}[a nondistinguished expanding place contradicts fixed-field \(abc\)]
\label{prop:hidden-place-abc-quality}
Assume the algebraic linear-trace-growth hypotheses of
Theorem~\ref{thm:high-power-concentration}.  Fix a hyperbolic element
with positive-trace lift \(a\), put \(s=\Tr a>2\), choose
\(\lambda+\lambda^{-1}=s\), and set
\[
 L=k(\lambda),\qquad d=[L:k],\qquad z=\lambda^2.
\]
For \(m\ge1\), define the height and radical of the \(abc\) triple
\(z^m+(-1)=z^m-1\) by
\[
 \begin{aligned}
 H_m&=[L:\Q]h(z^m:-1:z^m-1),\\
 R_m&=\sum_{\substack{\mathfrak P\subset\cO_L\\
 v_{\mathfrak P}(z^m),v_{\mathfrak P}(-1),
 v_{\mathfrak P}(z^m-1)\text{ are not all equal}}}
 \log N\mathfrak P.
 \end{aligned}
\tag{13.10bpzj}
\]
With the convention \(H_m/R_m=+\infty\) when \(R_m=0\), one has
\[
 \liminf_{m\to\infty}\frac{H_m}{R_m}
 \ge \frac{2\mathcal J_k(s)}{\ell(a)}.
\tag{13.10bpzk}
\]
Consequently, if \(a\) has a nondistinguished expanding place, then the
fixed number field \(L\) and the explicit sequence
\((z^m,-1,z^m-1)\) contradict the fixed-field \(abc\) conjecture
(13.10bpzcc).
\end{proposition}

\begin{proof}
Use the same enlarged set \(S\) as in
Lemma~\ref{lem:abc-chebyshev-radical}: above its complement, \(z\) and
\(z-1\) are units and \(L/k\) is unramified.  At every place outside
\(S_L\),
\[
 z^m-1=\lambda^{m-1}(z-1)q_m(s)
\]
shows that the prime support of the triple is exactly that of \(q_m(s)\).
Because \(q_m(s)\in k\) and the extension is unramified there,
\[
 \operatorname{rad}_{S_L}(q_m(s)\cO_L)
 =\operatorname{rad}_{S}(q_m(s))\cO_L.
\]
The finitely many places in \(S_L\) contribute only \(O_{s,S}(1)\) to a
radical, regardless of the sizes of their valuations.  Therefore
\[
 R_m=d\log N_k\operatorname{rad}_{S}(q_m(s))+O_{s,S}(1).
\tag{13.10bpzl}
\]
Likewise, the ultrametric triangle inequality at the finite places and
the ordinary triangle inequality at the archimedean places give
\[
 H_m=[L:\Q]h(z^m)+O_L(1)
     =dm\mathcal J_k(s)+O_{s,L}(1),
\tag{13.10bpzm}
\]
where the last equality is (13.10bpze).

Write
\(r_m=\log N_k\operatorname{rad}_{S}(q_m(s))\).
Theorem~\ref{thm:single-term-radical-upper} gives directly
\[
 r_m\le\frac{m\ell(a)}2+o(m).
\tag{13.10bpzn}
\]
Together with (13.10bpzl)--(13.10bpzm), this gives
\[
 \frac{H_m}{R_m}
 \ge
 \frac{dm\mathcal J_k(s)+O(1)}
      {dm\ell(a)/2+o(m)}
\]
whenever \(R_m>0\); if \(R_m=0\), the asserted inequality holds by
convention.  Taking the limit inferior proves (13.10bpzk) for every
individual recurrence term, without averaging consecutive indices.

Finally suppose \(2\mathcal J_k(s)/\ell(a)>1\), and choose
\[
 0<\epsilon<\frac{2\mathcal J_k(s)}{\ell(a)}-1.
\]
There are an unbounded sequence of indices \(m\) and a fixed \(\delta>0\)
for which
\(H_m\ge(1+\epsilon+\delta)R_m\).  If \(R_m\) is unbounded on a
subsequence, then \(H_m-(1+\epsilon)R_m\to\infty\) there; if it is
bounded on a subsequence, the same follows from (13.10bpzm).  Thus no
constant \(C_{L,\epsilon}\) can make (13.10bpzcc) hold for these triples.
This is a counterexample to the fixed-field \(abc\) conjecture over \(L\).
\end{proof}

\begin{theorem}[number-field \(abc\) implies compact Schmutz and Sarnak]
\label{thm:abc-implies-compact}
Assume the number-field \(abc\) conjecture for every number field.  Then
every torsion-free cocompact Fuchsian lattice whose signed trace set has
linear growth is arithmetic.  Consequently the compact Schmutz
linear-growth conjecture and the compact Sarnak bounded-clustering
conjecture both hold.
\end{theorem}

\begin{proof}
For every fixed \(k,S,s\), Lemma~\ref{lem:abc-chebyshev-radical} shows
that number-field \(abc\) over \(k(\lambda)\), where
\(\lambda+\lambda^{-1}=s\), implies
Conjecture~\ref{conj:chebyshev-radical}.  The assumed family of
fixed-field \(abc\) conjectures therefore establishes that hypothesis for
every algebraic specialization needed here.  Apply
Theorem~\ref{thm:chebyshev-radical-implies-compact}.
\end{proof}

There is a coefficient-uniform way to prove the lower bound (13.10bm) if
the trace set is sufficiently spread out in residue classes.  Unlike the
general sum-of-dilates theorem, it has no error depending on the size of the
coefficients.

\begin{lemma}[coprime dilates and residue classes]
\label{lem:coprime-dilate-residues}
Let \(L\) be a number field, let \(A\subset\mathcal O_L\) be finite, and
let \(a,b\in\mathcal O_L\setminus\{0\}\) generate coprime ideals.  Denote
by \(\nu_{\mathfrak m}(A)\) the number of residue classes modulo
\(\mathfrak m\) met by \(A\).  Then
\[
 \#(aA+bA)\ge
 \max\{\nu_{(a)}(A),\nu_{(b)}(A)\}\#A.
\tag{13.10bp}
\]
More generally, put
\[
 E_{\mathfrak m}(A)=
 \sum_{\xi\in\mathcal O_L/\mathfrak m}
       \#(A\cap(\xi+\mathfrak m))^2.
\]
Then
\[
 \#(aA+bA)
 \ge \frac{(\#A)^3}
 {\min\{E_{(a)}(A),E_{(b)}(A)\}}.
\tag{13.10bq}
\]
The same statements hold for arbitrary \(a,b\in L^\times\): divide their
principal fractional ideals by their greatest common ideal and multiply the
sumset by a common scalar.
\end{lemma}

\begin{proof}
Choose one element \(u\in A\) from each residue class modulo \((a)\).
The sets \(aA+bu\) are pairwise disjoint.  Indeed, a collision would give
\(b(u-u')\in(a)\), and coprimality would give \(u-u'\in(a)\).  This proves
the first bound using \((a)\); interchange \(a,b\) for the other one.

For the second bound, let \(\mathcal E\) count solutions of
\(ax+by=ax'+by'\) in \(A^4\).  Coprimality implies
\(x\equiv x'\pmod{(b)}\) and
\(y\equiv y'\pmod{(a)}\).  Once \((x,x',y)\) is fixed, there is at most
one \(y'\), and symmetrically.  Hence
\[
 \mathcal E\le
 \#A\min\{E_{(a)}(A),E_{(b)}(A)\}.
\]
Cauchy--Schwarz gives \(\#(aA+bA)\ge(\#A)^4/\mathcal E\).
\end{proof}

For the recurrence, this residue modulus has exactly the required
exponential size.  Work in \(K=\Q(s^2)\), and divide the two principal
fractional ideals generated by \(q_{n-1}(s)^2\) and \(q_n(s)^2\) by their
greatest common ideal.  Denote the resulting coprime integral ideals by
\(\mathfrak a_n\) and \(\mathfrak b_n\), respectively.  Cassini's identity
(13.10ao) is the algebraic reason for this coprimality.  If
\(\Q(\beta_n)=K\), then
\[
 \log N\mathfrak a_n
 =\log N\mathfrak b_n
 =2n\mathcal J_K(s)+O_s(\log n).
\tag{13.10br}
\]
To see this, use (13.10ap) at every archimedean embedding.  When the
corresponding algebraic root \(z\) has modulus different from one,
\(\beta_n\) tends to a nonzero limit.  If \(|z|=1\), then \(z\) is not a
root of unity by Lemma~\ref{lem:recurrence-field-stabilization}, and Baker's
linear-forms theorem supplies \(C_z>0\) such that
\(|z^n-1|\ge n^{-C_z}\) for \(n\ge2\)
\cite{BakerLinearForms}.  Thus the archimedean
factors in the heights of both \(\beta_n\) and \(\beta_n^{-1}\) are
\(e^{O_s(\log n)}\).  Lemma~\ref{lem:recurrence-height} then puts the full
exponential height into the denominator ideals of \(\beta_n\) and its
inverse, which are \(\mathfrak b_n\) and \(\mathfrak a_n\).

We state a sufficient group-theoretic hypothesis for applying
Lemma~\ref{lem:coprime-dilate-residues}.  This
formulation permits a finite extension \(L/K\), since the trace values need
not themselves lie in \(K\).

\begin{proposition}[congruence spreading is sufficient for recurrence expansion]
\label{prop:congruence-spreading-sufficient}
Retain the algebraic hypotheses of
Theorem~\ref{thm:moving-height-criterion}.  Suppose all trace values under
consideration lie in a fixed number field \(L\supset K\).  Fix a sufficiently
small \(\vartheta>0\), put \(n_R=\lfloor\vartheta R\rfloor\), and let
\(\mathfrak m_R\) be either
\(\mathfrak a_{n_R}\mathcal O_L\) or
\(\mathfrak b_{n_R}\mathcal O_L\).  After multiplying \(W_R\) by a common
scalar to make it integral, assume that
\[
 \nu_{\mathfrak m_R}(W_R)
 \ge
 \exp\bigl(2n_R\mathcal J_K(s)-o(R)\bigr).
\tag{13.10bs}
\]
Then \(\mathcal J_K(s)=\ell(k)/2\), so every nondistinguished local
expansion of the recurrence vanishes.

It is enough, more concretely, to have the following uniform counting
estimate for one of the two ideals.  Let \(G_R(W_R)\) be the edges of
(13.10be) whose first coordinate belongs to \(W_R\), and choose one group
element representing each such edge.  If, uniformly in residue classes
\(c\),
\[
 \#\{g\text{ among these representatives}:\Tr g\equiv c
          \pmod{\mathfrak m_R}\}
 \le
 \frac{\exp(R+o(R))}
      {\exp(2n_R\mathcal J_K(s))},
\tag{13.10bt}
\]
then (13.10bs) follows.
\end{proposition}

\begin{proof}
Since \(\beta_n=-q_{n-1}(s)^2/q_n(s)^2\), multiplication by a common
scalar identifies \(W_R+\beta_nW_R\) with a sum of the two coprime
dilates defining \(\mathfrak a_n,\mathfrak b_n\).  Equations
(13.10bp), (13.10br), and (13.10bs) give
\[
 \#(W_R+\beta_{n_R}W_R)
 \ge
 \exp\bigl(2n_R\mathcal J_K(s)-o(R)\bigr)\#W_R.
\]
The upper bound (13.10bd) is
\(\exp(n_R\ell(k)+o(R))\#W_R\).  Divide logarithms by
\(n_R\sim\vartheta R\) to obtain
\(2\mathcal J_K(s)\le\ell(k)\); the distinguished-place contribution
gives the reverse inequality.

For the final assertion, the dependent-random-choice construction gives
minimum degree at least \(r_R\), so (13.10ba) and (13.10bi) imply
\[
 \#G_R(W_R)\ge\#W_R\,r_R=\exp(R-o(R)).
\]
If \(W_R\) met \(M\) residue classes, all chosen edge representatives
would lie in the union of the corresponding \(M\) trace congruence
classes.  Estimate (13.10bt) therefore forces
\(M\ge\exp(2n_R\mathcal J_K(s)-o(R))\), which is (13.10bs).
\end{proof}

This proposition makes the remaining arithmetic issue concrete.  Expansion
theorems give estimates of the shape (13.10bt) for several important
families of congruence moduli: square-free ideals, bounded powers of
square-free rational moduli, and special groups for which arbitrary-modulus
theorems are known, subject to the usual good-reduction hypotheses.  The
ideals \(\mathfrak a_n,\mathfrak b_n\), however, can have
simultaneously many prime factors and unbounded prime-power exponents.
Moreover, at a finite place where the original representation is unbounded,
there is no fixed integral reduction map.  Neither issue may be suppressed:
one needs either an arbitrary-ideal trace equidistribution theorem in the
range \(\log N\mathfrak m=O(R)\), or a recurrence-specific substitute which
retains the full norm in (13.10br).

The pointwise theorem of Conlon and Lim, (13.10at), does not prove
(13.10bm): its error is uniform in the size of the finite set only after
the algebraic dilate has been fixed, whereas here \(n_R\to\infty\).
This order of quantifiers is substantive.  The desired estimate would have
to be uniform for algebraic dilates of bounded degree whose heights grow
with \(R\), for the particular sets \(W_R\) furnished by the trace graph.
It may be written in the more
informative equivalent form
\[
 \#(W_R+\beta_{n_R}W_R)
 \ge
 \exp\bigl(2n_R\mathcal J_K(s)-o(n_R)\bigr)\#W_R.
\tag{13.10bo}
\]

\begin{corollary}[the arithmetic consequence of uniform recurrence expansion]
\label{cor:moving-height-arithmeticity}
Let \(\Gamma<\PSL_2(\R)\) be cocompact with linearly growing trace set.
Assume that all squared traces of \(\Gamma\) are algebraic and that
(13.10bm) holds for every hyperbolic element of the subgroup generated by
squares, with the corresponding sets furnished by
Theorem~\ref{thm:drc-trace-recurrence}.  Then \(\Gamma\) is arithmetic.
Consequently, the stated uniform recurrence-expansion estimate, together
with the exclusion of transcendental exceptional characters, is another
sufficient condition for both compact arithmeticity conclusions.
\end{corollary}

\begin{proof}
The invariant trace field is a number field because a finite collection of
character coordinates generates it and those coordinates are algebraic.
Theorem~\ref{thm:moving-height-criterion}, applied to every hyperbolic
element of \(\Gamma^{(2)}\), says that every squared trace is an algebraic
integer and that all of its nondistinguished conjugates lie in \([0,4]\).
Takeuchi's arithmeticity criterion \cite{Takeuchi} now applies to \(\Gamma^{(2)}\), and
arithmeticity passes to its finite-index overgroup \(\Gamma\).  Bounded
clustering implies linear growth, so the same implication handles both
hypotheses.
\end{proof}

\begin{proposition}[coefficient extraction for a pure cyclic twist]
\label{prop:pure-twist-coefficients}
Let \(k\) be a field of characteristic zero, let
\[
 D(z)=\begin{pmatrix}z&0\\0&z^{-1}\end{pmatrix},\qquad
 B=\begin{pmatrix}a&b\\c&d\end{pmatrix}\in\SL_2(k),
\]
and assume \(bc\ne0\) and \((a,d)\ne(0,0)\).  Let
\(\Delta<\SL_2(k)\) contain no nonidentity unipotent.  Then every fiber on
\(\Delta\) of
\[
 X\longmapsto
 \Tr^2\!\bigl(D(z)BD(z)^{-1}X\bigr)\in k[z,z^{-1}]
\]
has at most two projective elements, and hence at most four elements of
\(\Delta\).

Consequently, suppose that \(\rho_z\) is an algebraic pure-twist family and
that, for some \(w\in\Pi\) and \(H<\Pi\), lifts can be chosen so that
\[
 \widetilde\rho_z(wh)=
 D(z)BD(z)^{-1}\widetilde\rho_1(h)\qquad(h\in H).
\]
Assume that \(\rho_1|_H\) is faithful and convex cocompact and that every
squared-trace collision among the elements \(wh\), \(h\in H\), at \(z=1\)
persists identically in \(z\).  Then every fiber of
\(h\mapsto\Tr^2\widetilde\rho_1(wh)\) has at most two projective elements.
If, moreover, \(H\subset\ker\phi\) and \(\phi(w)=1\) for a homomorphism
\(\phi\colon\Pi\to\Z\), then all \(wh\) are primitive and
\[
 h_{\Tr}(\rho_1)\ge\delta_H.
\]
\end{proposition}

\begin{proof}
For \(X=(\begin{smallmatrix}p&q\\r&s\end{smallmatrix})\), direct
multiplication gives
\[
 \Tr\!\bigl(D(z)BD(z)^{-1}X\bigr)
 =ap+ds+z^2br+z^{-2}cq.
\]
If the squares of this Laurent polynomial for \(X\) and \(X_0\) agree,
the two polynomials differ by one sign
\(\epsilon\in\{1,-1\}\).  Comparing coefficients gives
\[
 q=\epsilon q_0,\qquad r=\epsilon r_0,\qquad
 ap+ds=\epsilon(ap_0+ds_0).
\]
After replacing \(X_0\) by \(\epsilon X_0\), fix \(q,r\) and put
\(P=1+qr\).  The remaining equations are
\[
 ps=P,\qquad ap+ds=L.
\]
If \(ad\ne0\), elimination gives a quadratic equation and hence at most
two solutions.  If exactly one of \(a,d\) is nonzero, the equations give a
unique solution except when their solution set is an affine line.  Two
distinct matrices on that exceptional line have a nonidentity upper or
lower unipotent quotient.  The line therefore meets \(\Delta\) in at most
one point.  Negating \(X\) exchanges the two signs without changing its
projective class, proving the first assertion.

A faithful convex-cocompact Fuchsian group has no nonidentity unipotent.
Persistence turns every trace fiber at \(z=1\) into a fiber of the Laurent
polynomial map just considered.  Finally, \(\phi(wh)=1\), whereas the
\(\phi\)-value of a proper \(m\)-th power is divisible by
\(|m|\ge2\).  Thus every \(wh\) is primitive.  Orbit counting in \(H\),
together with \(\ell(wh)\le d(o,\rho_1(w)o)+d(o,\rho_1(h)o)\), gives the
entropy bound.
\end{proof}

\begin{proposition}[coefficients of one trace functional do not determine a matrix jet]
\label{prop:one-trace-jet-obstruction}
Let \(k\) be algebraically closed of characteristic zero, let \(H\) be
finitely generated, and let
\[
 \rho\colon H\longrightarrow\SL_2(k[[t]])
\]
have Zariski-dense reduction
\(\bar\rho\colon H\to\SL_2(k)\).  Write
\[
 \rho(h)=(I+t z(h)+O(t^2))\bar\rho(h),
\]
and assume that the class of \(z\) in
\(H^1(H,\mathfrak{sl}_2{}_{\Ad\bar\rho})\) is nonzero.  For every
\(r\ge1\), the Zariski closure of the reduction of \(\rho(H)\) in
\[
 J_r\SL_2=\SL_2(k[t]/(t^{r+1}))
\]
is the whole jet group.

For fixed \(A(t)\in M_2(k[[t]])\), consider
\[
 \Phi_r\colon J_r\SL_2\longrightarrow\mathbb A^{r+1},\qquad
 X(t)\longmapsto
 \bigl([t^j]\Tr(A(t)X(t))\bigr)_{0\le j\le r}.
\]
Every irreducible component of every nonempty fiber of \(\Phi_r\) has
dimension at least
\[
 3(r+1)-(r+1)=2(r+1).
\]
The same conclusion holds for a Laurent matrix after multiplication by a
fixed power of \(t\) and reindexing its coefficients.

This loss of two thirds of the dimensions is already exact at the matrix
level.  Put \(R=k[[t]]\) and
\(A_m=\diag(t^{-m},t^m)\), where \(m\ge1\).  Fix
\(p_0,s_0\in R\).  For arbitrary \(s\in R\) and \(q\in R^\times\), set
\[
 p=p_0+t^{2m}(s_0-s),\qquad
 r=\frac{ps-1}{q},\qquad
 X_{q,s}=\begin{pmatrix}p&q\\r&s\end{pmatrix}.
\]
Then \(X_{q,s}\in\SL_2(R)\) and
\[
 \Tr(A_mX_{q,s})=t^{-m}p_0+t^m s_0
\]
for every \((q,s)\), while their reductions fill a two-dimensional open
subset of the surface \(\bar p=\bar p_0\) in \(\SL_2(k)\).
\end{proposition}

\begin{proof}
Let \(L_1\) be the Zariski closure of the first-jet image.  Its projection
to \(\SL_2\) is onto.  The kernel is an
\(\Ad\SL_2\)-invariant subspace of \(\mathfrak{sl}_2\), hence is either
zero or all of \(\mathfrak{sl}_2\).  In the first case \(L_1\) is the
graph of an algebraic \(1\)-cocycle on \(\SL_2\).  The vanishing of
algebraic \(H^1(\SL_2,\mathfrak{sl}_2)\) makes that cocycle a coboundary,
contrary to the hypothesis on \(z\).  Hence \(L_1=J_1\SL_2\).

Let \(L_r\) be the closure at order \(r\).  Truncation sends \(L_r\) onto
\(L_1=J_1\SL_2\), because the image of a Zariski closure under the
surjective homomorphism of jet groups contains the dense first-jet image.
Its differential is therefore onto.  In the Lie algebra of \(L_r\), choose
lifts of the degree-one vectors \(tX\), \(X\in\mathfrak{sl}_2\).
Brackets of two such lifts have initial term \(t^2[X,Y]\).  Since
\([\mathfrak{sl}_2,\mathfrak{sl}_2]=\mathfrak{sl}_2\), all degree-two
congruence directions occur.  Bracketing a degree-one lift with one whose
initial term has degree \(j\) produces every direction of degree \(j+1\).
Induction gives all congruence layers through degree \(r\), so
\(\operatorname{Lie}L_r=\operatorname{Lie}J_r\SL_2\).  The jet group is
connected, and therefore \(L_r=J_r\SL_2\).

The jet group has dimension \(3(r+1)\), while a fiber of \(\Phi_r\) is cut
out by \(r+1\) scalar equations.  Krull's principal ideal theorem gives
the stated lower bound.  Finally, the displayed \(X_{q,s}\) has determinant
one, direct substitution gives the fixed Laurent trace, and reduction
modulo \(t\) leaves the two free parameters \((\bar q,\bar s)\), with
\(\bar q\ne0\).
\end{proof}

The last proposition concerns dimension counting; it does not construct a
large fiber inside a discrete group.  A Zariski-dense subset need not be
dense in any specified fiber.  Deforming the
representation introduces new jet variables at the same rate as new trace
coefficients.  A discrete estimate must therefore use the distribution of
words, as in Theorem~\ref{thm:moving-divisorial-multiplier}, or a second
independent trace observable.

\begin{proposition}[the limit of pigeonholing with two trace observables]
\label{prop:two-observable-budget}
Let \(\rho_0\colon\Pi\to\PSL_2(\R)\) be Fuchsian, let
\(\phi\colon\Pi\twoheadrightarrow\Z\), and let
\(J<\ker\phi\) be finitely generated and convex cocompact, with critical
exponent \(\delta_J\).  Fix \(o\in\mathbb H^2\) and put
\[
 \mathcal B_J(R)=\{h\in J:d(o,\rho_0(h)o)\le R\}.
\]
Suppose that \(h_{\Tr}(\rho_0)\le s\).  For \(j=1,2\), let
\(w_{j,R}\in\phi^{-1}(1)\) and assume
\[
 d(o,\rho_0(w_{j,R})o)\le(a_j+o(1))R,
 \qquad a_j\ge0.
\]
Then the map
\[
 h\longmapsto
 \bigl(\Tr^2\rho_0(w_{1,R}h),
       \Tr^2\rho_0(w_{2,R}h)\bigr),
 \qquad h\in\mathcal B_J(R),
\]
has at most
\[
 \exp\!\left(\bigl(s(2+a_1+a_2)+o(1)\bigr)R\right)
\]
values.  Consequently, pigeonholing alone guarantees a joint fiber only
of size
\[
 \exp\!\left(
   \bigl(\delta_J-s(2+a_1+a_2)-o(1)\bigr)R\right).
\]
In particular, when \(s=1/2\), this exponent is nonpositive because
\(\delta_J\le1\).  Scalar trace sparsity, even in finite covers for which
\(\delta_J\) tends to \(1\), therefore cannot by itself force a collision
of two independent trace observables.
\end{proposition}

\begin{proof}
The equality \(\phi(w_{j,R}h)=1\) shows that every \(w_{j,R}h\) is
primitive.  Moreover,
\[
 \ell_{\rho_0}(w_{j,R}h)
 \le d(o,\rho_0(w_{j,R})o)+d(o,\rho_0(h)o)
 \le(1+a_j+o(1))R.
\]
The definition of \(h_{\Tr}\), with an arbitrarily small fixed enlargement
of \(s\), bounds the number of values of the \(j\)-th coordinate by
\(\exp((s(1+a_j)+o(1))R)\).  Multiplying the two bounds proves the first
assertion.  Convex-cocompact orbit counting gives
\(\#\mathcal B_J(R)=\exp((\delta_J+o(1))R)\), and division gives the joint
fiber bound.  Finally, a convex-cocompact subgroup of
\(\PSL_2(\R)\) has critical exponent at most the ambient exponent \(1\).
\end{proof}

\begin{remark}[a sharp finite model]
For a prime power \(q\), put \(U_q=\mathbb F_q^2\), and define
\[
 f_t(x,y)=x+ty\quad(t\in\mathbb F_q),\qquad
 f_\infty(x,y)=y.
\]
Every \(f_t\) has exactly \(q\) values and every fiber has \(q\) elements,
whereas
\((f_s,f_t)\colon U_q\to\mathbb F_q^2\) is a bijection whenever
\(s\ne t\).  With \(R=2\log q\), this gives
\(\#U_q=e^R\) and scalar supports of size \(e^{R/2}\), but no nontrivial
simultaneous fiber for any two independent observables.  Thus even an
arbitrarily large collection of individually large fibers need not contain
two aligned partitions.
\end{remark}

\begin{proposition}[the valuation cost of extracting several Cartan coefficients]
\label{prop:multi-cartan-budget}
Let \(K\) be a nonarchimedean field of characteristic zero, let
\(\rho_v\colon\Pi\to\SL_2(K)\), fix \(a\in\Pi\), and let
\(k_1,\ldots,k_m\in\ker\phi\) be split regular semisimple elements.  Write
\[
 A=\rho_v(a),\qquad
 \rho_v(k_i)=\lambda_iP_i^++\lambda_i^{-1}P_i^-,
 \qquad |\lambda_i|_v>1,
\]
and put \(\ell_v(k_i)=2\log|\lambda_i|_v\).  For
\(\epsilon=(\epsilon_1,\ldots,\epsilon_m)\in\{+,-\}^m\), set
\[
 L_\epsilon(Y)=
 \Tr\bigl(AP_1^{\epsilon_1}\cdots P_m^{\epsilon_m}Y\bigr),
 \qquad
 \alpha_\epsilon(\mathbf n)=
 \prod_{i=1}^m\lambda_i^{\epsilon_i n_i},
\]
where a sign in the exponent means \(+1\) or \(-1\).  Then
\[
 \Tr\!\left(
 A\rho_v(k_1)^{n_1}\cdots\rho_v(k_m)^{n_m}Y\right)
 =
 \sum_{\epsilon\in\{+,-\}^m}
 \alpha_\epsilon(\mathbf n)L_\epsilon(Y).
\]
In particular, regardless of \(m\), the functionals in this expansion span
a space of dimension at most four.

Let \(\mathcal L=(L_1,\ldots,L_q)\) be \(q\) nonzero linearly independent
functionals remaining after only coefficient-independent identifications
among these terms, with coefficients
\(\alpha_1,\ldots,\alpha_q\).  Thus each
\(|\alpha_r|_v\) is one of the displayed monomial absolute values; no
additional valuation separation caused by cancellation of equal-size
monomials is used.  Necessarily \(q\le4\).  For
\[
 \mathcal D_R=\{\rho_v(h)-\eta\rho_v(h'):
   h,h'\in\mathcal B_J(R),\ \eta\in\{1,-1\}\},
\]
define
\[
 C_{\mathcal L}(R)=
 \sup_{\substack{\Delta\in\mathcal D_R,\ i,j\\
                   L_i(\Delta)L_j(\Delta)\ne0}}
 \left|\log\left|\frac{L_i(\Delta)}{L_j(\Delta)}\right|_v\right|.
\]
Order the coefficients so that
\(b_1>\cdots>b_q\), where \(b_r=\log|\alpha_r|_v\).  If
\[
 b_r-b_{r+1}>C_{\mathcal L}(R)\qquad(1\le r<q),
\]
then, uniformly for \(\Delta\in\mathcal D_R\),
\[
 \sum_{r=1}^q\alpha_rL_r(\Delta)=0
 \quad\Longrightarrow\quad
 L_1(\Delta)=\cdots=L_q(\Delta)=0.
\]
Every choice satisfying these coefficient-separation inequalities obeys
the necessary valuation budget
\[
 \sum_{i=1}^m n_i\ell_v(k_i)>
 (q-1)C_{\mathcal L}(R).
\]
This is a necessary cost for this uniform leading-term argument, not a
claim about every possible arithmetic proof of the same conclusion.
\end{proposition}

\begin{proof}
Multiplying the spectral decompositions proves the expansion.  All the
\(L_\epsilon\) lie in the four-dimensional space \(M_2(K)^*\).

Suppose that the displayed sum vanishes and not every
\(L_r(\Delta)\) is zero.  Let \(r\) be the least index with
\(L_r(\Delta)\ne0\).  For every \(j>r\) with
\(L_j(\Delta)\ne0\), the definitions give
\[
 \log|\alpha_rL_r(\Delta)|_v-
 \log|\alpha_jL_j(\Delta)|_v
 \ge b_r-b_j-C_{\mathcal L}(R)>0.
\]
Thus \(\alpha_rL_r(\Delta)\) is the unique term of largest absolute
value, contradicting the ultrametric triangle inequality.  This proves the
coefficient extraction.

Adding the \(q-1\) separation inequalities gives
\(b_1-b_q>(q-1)C_{\mathcal L}(R)\).  Since the coefficients are pure
Cartan weights,
\[
 b_1-b_q\le
 2\sum_i n_i\log|\lambda_i|_v
 =\sum_i n_i\ell_v(k_i),
\]
which proves the cost bound.
\end{proof}

\begin{lemma}[Fuchsian displacement of a product of large powers]
\label{lem:visible-power-cost}
Retain the preceding notation and suppose that
\(\rho_0\colon\Pi\to\PSL_2(\R)\) is faithful, with torsion-free discrete
image.  Let
\(Q_i^\pm\) be the spectral projectors of \(\rho_0(k_i)\).  If
\[
 Q_i^+Q_{i+1}^+\ne0\qquad(1\le i<m),
\]
then, as \(\min_i n_i\to\infty\),
\[
 d\!\left(o,\rho_0(k_1^{n_1}\cdots k_m^{n_m})o\right)
 =\sum_{i=1}^m n_i\ell_0(k_i)+O(1).
\]
If the projector condition fails, the two adjacent elements share a
boundary fixed point.  They are then commensurable, and combining their
powers shows that they do not provide two independent spectral directions
under any representation of \(\Pi\).
\end{lemma}

\begin{proof}
Write
\[
 \rho_0(k_i)^{n_i}
 =\mu_i^{n_i}Q_i^++\mu_i^{-n_i}Q_i^-,
 \qquad |\mu_i|=e^{\ell_0(k_i)/2}.
\]
The projector assumption makes the leading product
\(Q_1^+\cdots Q_m^+\) nonzero.  Every other term is smaller by an
exponential factor as \(\min_i n_i\to\infty\), and hence
\[
 \bigl\|\rho_0(k_1^{n_1}\cdots k_m^{n_m})\bigr\|
 \asymp\prod_i|\mu_i|^{n_i}.
\]
The standard comparison
\(d(o,go)=2\log\|g\|+O(1)\) proves the formula.

The equality \(Q_i^+Q_{i+1}^+=0\) says that the attracting fixed point of
\(k_{i+1}\) is the repelling fixed point of \(k_i\).  In a torsion-free
discrete Fuchsian group, two hyperbolic elements sharing one endpoint share
both and belong to the same cyclic subgroup.  Their abstract
commensurability persists under every representation, proving the last
assertion.
\end{proof}

\begin{corollary}[the cost of extracting four coefficients from two Cartan factors]
\label{cor:four-coordinate-cost}
For two factors, suppose that all four functionals
\(L_{\epsilon_1,\epsilon_2}\) are linearly independent.  Put
\(x_i=n_i\ell_v(k_i)\), and assume \(x_1\ge x_2\).  The four logarithmic
weights, in decreasing order, are
\[
 \frac{x_1+x_2}{2},\quad
 \frac{x_1-x_2}{2},\quad
 \frac{-x_1+x_2}{2},\quad
 \frac{-x_1-x_2}{2}.
\]
Their consecutive gaps are \(x_2,x_1-x_2,x_2\).  Uniform extraction with
the bound \(C_{\mathcal L}(R)\) therefore requires
\[
 x_2>C_{\mathcal L}(R),\qquad
 x_1-x_2>C_{\mathcal L}(R),
\]
and in particular
\[
 x_1+x_2>3C_{\mathcal L}(R).
\]
If the corresponding word in the Fuchsian representation satisfies
Lemma~\ref{lem:visible-power-cost} and
\[
 r_*=\min_i\frac{\ell_0(k_i)}{\ell_v(k_i)},
\]
the Fuchsian displacement contributed by the prefix is greater than
\(3r_*C_{\mathcal L}(R)+O(1)\).
\end{corollary}

More generally, if \(q\) extracted coordinates have subexponential common
fibers and
\(c_{\mathcal L}=\limsup C_{\mathcal L}(R)/R\), this particular
leading-term scheme can certify no more than the entropy exponent
\[
 \frac{\delta_J}{1+(q-1)r_*c_{\mathcal L}}.
\]
This is a bound on the method, not an upper bound for the actual trace
entropy.  Extracting three or four coordinates improves on the
two-coordinate argument only if the corresponding cancellation exponent
drops by more than a factor of two or three.  A product of many Cartan
powers therefore does not generically amortize cancellation.

\begin{remark}[the exceptional two-dimensional family]
\label{rem:unipotent-pencil}
Suppose \(W=\Span_\R\{A_s:s\in U\}\) has dimension two.  If the determinant
quadratic form restricted to \(W\) is nondegenerate, any nonconstant
connected subfamily contains a pair satisfying (13.8), so
Theorem~\ref{thm:two-trace-evaluations} applies.  If it is degenerate, write
\(W=\Span\{A_0,D\}\).  Infinitely many points of \(W\) have determinant
one, so the polynomial \(\det(A_0+cD)\) is identically one along the
relevant affine line.  Consequently
\[
 A_s=A_0(I+c(s)N),\qquad N\ne0,\quad N^2=0.
\tag{13.11}
\]
Thus (13.11) is the only two-dimensional family of matrices \(A_s\) not
covered by the two theorems.  Relative Fenchel--Nielsen twists are
motions in the semisimple centralizer of a hyperbolic boundary element and
are not of the form (13.11).
\end{remark}

\begin{theorem}[constant restriction on a locus preserving trace collisions forces trace entropy at least \(\delta_H\)]
\label{thm:relative-fiber-rigidity}
Let \(\Pi\) be a closed surface group, let
\(\rho_0\colon\Pi\to\PSL_2(\R)\) be a faithful cocompact Fuchsian
representation, and let \(H<\Pi\) be nonelementary and convex cocompact
under \(\rho_0\).  Suppose that \(H\subset\ker\phi\) for an epimorphism
\(\phi\colon\Pi\to\Z\), and write \(\delta_H\) for the critical exponent of
\(\rho_0(H)\).

Let \(\mathcal C\) be an irreducible algebraic subvariety, through
\([\rho_0]\), of the \(\PSL_2(\C)\)-character component containing
the Fuchsian characters.  Assume that:
\begin{enumerate}[label=\textup{(\alph*)}]
\item the restriction map \(\mathcal C\to\mathfrak X(H)\) is constant;
\item every equality between two squared traces at \(\rho_0\) holds
identically on \(\mathcal C\).
\end{enumerate}
If \(\dim\mathcal C>0\), then the primitive trace entropy at \(\rho_0\)
is at least \(\delta_H\).  In particular, if
\(\delta_H>1/2\) and that entropy is at most \(1/2\), then
\(\mathcal C\) is zero-dimensional.
\end{theorem}

\begin{proof}
Suppose that \(\dim\mathcal C>0\), and choose
\([\rho_1]\in\mathcal C\) distinct from \([\rho_0]\).  Because
\(\rho_0|_H\) is absolutely irreducible, equality of the restricted
characters means conjugacy of the restricted representations.  Conjugate
\(\rho_1\) so that its projective restriction to \(H\) equals that of
\(\rho_0\):
\[
 \rho_1|_H=\rho_0|_H.
\tag{13.12}
\]
Here equality is in \(\PSL_2(\C)\).  Since \(H\) is finitely generated and
lies in \(\ker\phi\), it has infinite index in \(\Pi\).  The covering
surface corresponding to \(H\) is noncompact, and finite generation gives
a compact core with nonempty boundary.  Such a surface retracts onto a
finite graph, so \(H\) is free (of rank at least two because it is
nonelementary).  We may therefore choose one lift
\(\widetilde\rho_H\colon H\to\SL_2(\R)\) of this common projective
restriction.

Consider
\[
 (\rho_1,\rho_0)\colon
 \Pi\longrightarrow\PSL_2(\C)\times\PSL_2(\C).
\]
Both projections are Zariski dense because their restrictions to \(H\)
are conjugate to the nonelementary Fuchsian group \(\rho_0(H)\).
Algebraic Goursat's lemma says that the Zariski closure of the paired image
is either the full product or the graph of an algebraic automorphism of
\(\PSL_2\).  In the graph case, (13.12) forces that automorphism to fix the
Zariski-dense subgroup \(\rho_0(H)\) pointwise; it is therefore the
identity.  This would make the two characters equal, contrary to the choice
of \(\rho_1\).  Hence the paired image is Zariski dense in the product.

Put \(K=\ker\phi\).  It contains \(H\), so the projections of the Zariski
closure of \((\rho_1,\rho_0)(K)\) are both \(\PSL_2\).  That closure is
normal in the full product.  Since \(\PSL_2\) is adjoint and simple, the
only such normal algebraic subgroup is the full product.  Every coset of
\(K\), in particular \(\phi^{-1}(1)\), consequently has Zariski-dense
image under \((\rho_1,\rho_0)\).

The condition
\[
 \Tr^2\!\left(\rho_1(w)\rho_0(w)^{-1}\right)\ne4
\tag{13.13}
\]
is a nonempty Zariski-open condition on the product.  Choose
\(w\in\phi^{-1}(1)\) satisfying (13.13).  Choose arbitrary lifts
\(A_i\in\SL_2(\C)\) of \(\rho_i(w)\), with \(A_0\in\SL_2(\R)\).
The squared trace in (13.13) is independent of these two sign choices.
The entries of \(A_0,A_1\) and of the images of a finite free basis of
\(H\) generate a finitely generated field \(F/\Q\).
Hypothesis (b), restricted to the
elements \(wh\), \(h\in H\), gives implication (13.7) in
Theorem~\ref{thm:two-trace-evaluations}, with
\(\rho=\widetilde\rho_H\); (13.13) gives its condition (13.8).  That theorem
gives subexponential fibers for
\[
 h\longmapsto\Tr^2\rho_0(wh)
\]
in word balls of \(H\).

Every \(wh\) is primitive in \(\Pi\).  Indeed, if \(wh=u^m\) for an integer
\(|m|\ge2\), then \(1=\phi(wh)=m\phi(u)\), which is impossible in \(\Z\).
Convex-cocompact orbit counting and the subexponential fiber bound therefore
give primitive trace entropy at least \(\delta_H\).  The final assertion
follows.
\end{proof}

\begin{corollary}[fibers of restriction to Hao's complementary subgroup]
\label{cor:hao-restriction-fibers}
Let \(x\) be a closed hyperbolic surface of genus \(g\ge6\) with
\(h_{\Tr}(x)\le1/2\), and let \(H\) be the complementary subgroup in the
surface decomposition used in Hao's proof of
\cite[Theorem~A, Proposition~3.2, and (5.1)]{HaoLengthSets}.  Thus
\(\delta_H>1/2\), and \(H\) lies in the
kernel of a homomorphism from the surface group onto \(\Z\).
For the variety \(V_x\), on which all trace collisions at \(x\) persist,
the fiber through \(x\) of the
restriction map
\[
 r_H\colon V_x\longrightarrow\mathfrak X(H),\qquad
 [\rho]\longmapsto[\rho|_H],
\]
is zero-dimensional at \(x\).  In particular, every positive-dimensional
algebraic germ in \(V_x\) through \(x\) must vary the character of this
complementary subgroup.
\end{corollary}

\begin{proof}
The required homomorphism is explicit in the standard generators of the
splitting.  In the one-holed-torus case set
\(\phi(a_1)=1\) and set \(\phi\) equal to zero on the remaining standard
generators, which generate \(H\).  In the pair-of-pants case set
\(\phi(a_2)=1\) and set it equal to zero on the displayed generators of
\(H\) (and on \(a_1\)).  Thus \(H<\ker\phi\) in both cases.

If the fiber had positive dimension at \(x\), take a reduced irreducible
component through \(x\) of
\[
 \bigl(V_x\cap r_H^{-1}(r_H(x))\bigr)_{\mathrm{red}}.
\]
Every squared-trace equality at \(x\) would persist on that component, while
its restricted \(H\)-character would be constant.
Theorem~\ref{thm:relative-fiber-rigidity} rules this out.
\end{proof}

The preceding argument excludes an algebraic family contained in a
restriction fiber.  Persistence in fact excludes even an infinitesimal
vertical direction.  This distinction matters at a ramification point: a
morphism can have a zero-dimensional fiber and nevertheless have a
noninjective differential there.

\begin{theorem}[restriction has injective differential on a trace-collision locus]
\label{thm:restriction-unramified}
Let \(\Pi\) be a closed surface group, let
\(x=[\rho]\) be a faithful cocompact Fuchsian character in the liftable
\(\PSL_2(\C)\)-character component, and suppose that
\(h_{\Tr}(x)\le 1/2\).  Let
\(\phi\colon\Pi\twoheadrightarrow\Z\), and let
\(H<\ker\phi\) be finitely generated, nonelementary, and convex cocompact
under \(\rho\), with critical exponent \(\delta_H>1/2\).

If \(Z\) is any reduced algebraic subvariety of \(V_x\), where all trace
collisions at \(x\) persist, and \(Z\) contains \(x\), then restriction
has injective differential at \(x\):
\[
 d(r_H)_x\colon T_xZ\longrightarrow
 T_{[\rho|_H]}\mathfrak X(H)
 \quad\hbox{is injective}.
\tag{13.14}
\]
In particular, this conclusion also rules out vertical tangent vectors at a
singular point of \(Z\); no integrability assumption on the tangent vector
is needed.  If \(H\) is free of rank \(r\ge2\), then the restricted
Fuchsian character is a smooth point of \(\mathfrak X(H)\), of dimension
\(3r-3\).  Consequently
\[
 \dim_{\C}T_xZ\le 3r-3.
\]
\end{theorem}

\begin{proof}
We first translate a putative tangent vector into a cocycle.  The Fuchsian
representation \(\rho\) is irreducible and its centralizer in
\(\PSL_2(\C)\) is trivial.  The usual tangent-space description of the
character variety at such a smooth surface-group character is
\[
 T_x\mathfrak X(\Pi)=
 H^1\!\left(\Pi,\mathfrak{sl}_2(\C)_{\Ad\rho}\right).
\tag{13.15}
\]
For clarity, a representative \(z\) of a class on the right satisfies
\[
 z(uv)=z(u)+\Ad\rho(u)z(v).
\tag{13.16}
\]
It gives the dual-number deformation
\[
 \rho_\epsilon(u)=
 \bigl(I+\epsilon z(u)\bigr)\rho(u),
 \qquad \epsilon^2=0.
\tag{13.17}
\]
Coboundaries are precisely infinitesimal conjugations.  Formula (13.15)
can also be obtained directly by differentiating the surface relation; the
vanishing of the infinitesimal centralizer identifies the quotient by
infinitesimal conjugation with \(H^1\).

Suppose that \(0\ne v\in T_xZ\) belongs to the kernel in (13.14), and let
\(z\) represent its nonzero cohomology class.  The vanishing of the
restricted tangent class says that \(z|_H\) is a coboundary.  Subtracting
the corresponding coboundary on all of \(\Pi\), we may and shall assume
\[
 z(h)=0\qquad(h\in H).
\tag{13.18}
\]
The resulting cocycle is not a coboundary on \(\Pi\), because \(v\ne0\).

We next find a group element whose infinitesimal displacement is regular
semisimple.  Put
\[
 G=\PSL_2(\C),\qquad \mathfrak g=\mathfrak{sl}_2(\C),
\]
and give \(G\ltimes\mathfrak g\) the multiplication
\[
 (g,U)(g',U')=(gg',U+\Ad(g)U').
\]
Equation (13.16) says that
\[
 \widehat\rho\colon\Pi\longrightarrow G\ltimes\mathfrak g,
 \qquad
 \widehat\rho(u)=(\rho(u),z(u)),
\tag{13.19}
\]
is a homomorphism.  Let \(L\) be the Zariski closure of its image.  Its
projection to \(G\) is onto, since the Fuchsian group \(\rho(\Pi)\) is
Zariski dense.  The subgroup
\(L\cap\mathfrak g\) is a vector subspace invariant under the adjoint
action of \(G\).  The adjoint module of \(\PSL_2\) is irreducible, so this
intersection is either \(0\) or all of \(\mathfrak g\).

If \(L\cap\mathfrak g=0\), projection identifies \(L\) with \(G\), and
\(L\) is the graph of an algebraic cocycle
\(c\colon G\to\mathfrak g\).  Such a cocycle is a coboundary.  Indeed, its
differential at the identity is a Lie-algebra cocycle
\(\mathfrak g\to\mathfrak g\); Whitehead's lemma makes that differential a
coboundary.  After subtracting the corresponding algebraic coboundary, the
cocycle has zero differential.  Since \(G\) is connected in characteristic
zero, it is then constant, and its value at the identity is zero.  It
follows that \(z\) is a coboundary on the Zariski-dense subgroup
\(\rho(\Pi)\), contrary to the choice of \(v\).  Consequently
\[
 L=G\ltimes\mathfrak g.
\tag{13.20}
\]
Here we used that a subgroup which contains the full translation group
\(\mathfrak g\) and projects onto \(G\) is the whole semidirect product:
multiply any lift of \(g\in G\) by a suitable translation to obtain
\((g,0)\).

Let \(K=\ker\phi\), and let \(M\) be the Zariski closure of
\(\widehat\rho(K)\).  Since \(K\) is normal in \(\Pi\), (13.20) makes
\(M\) a normal algebraic subgroup of \(G\ltimes\mathfrak g\).  Its
projection to \(G\) is onto: it contains
\((\rho(h),0)\) for \(h\in H\), and a nonelementary subgroup of
\(\PSL_2\) is Zariski dense.  As above,
\(M\cap\mathfrak g\) is either \(0\) or \(\mathfrak g\).  In the first
case \(M\) is the graph of an algebraic cocycle on \(G\).  That cocycle
vanishes on the Zariski-dense set \(\rho(H)\), by (13.18), and hence
vanishes identically.  Thus the first case would give
\[
 M=G\times\{0\},\qquad z|_K=0.
\tag{13.21}
\]
Choose \(w_0\in\Pi\) with \(\phi(w_0)=1\).  For every \(k\in K\), the
cocycle identity and (13.21) give
\[
 0=z(w_0kw_0^{-1})
   =z(w_0)-\Ad\rho(w_0kw_0^{-1})z(w_0).
\]
The group \(\rho(K)\) is Zariski dense because it contains \(\rho(H)\),
and the adjoint representation has no nonzero invariant vector.  Hence
\(z(w_0)=0\).  Since \(K\) together with \(w_0\) generates \(\Pi\), this
would make \(z=0\), again a contradiction.  We have therefore proved
\[
 M=G\ltimes\mathfrak g.
\tag{13.22}
\]

Every coset of \(K\), and in particular \(\phi^{-1}(1)\), now has
Zariski-dense image under \(\widehat\rho\).  The complement of the
nilpotent cone
\[
 \{U\in\mathfrak{sl}_2:\det U=0\}
\]
is nonempty and Zariski open.  We can consequently choose
\[
 w\in\phi^{-1}(1)
 \quad\hbox{such that}\quad
 N:=z(w)\ \hbox{satisfies}\ \det N\ne0.
\tag{13.23}
\]
Since \(\Tr N=0\), the matrix \(N\) has two distinct eigenvalues.

We use the two first-order trace coordinates in (13.23).  Fix a
lift of \(\rho\) to \(\SL_2(\R)\), use the same notation for this lift,
and put
\[
 Y_h=\rho(w)\rho(h),\qquad h\in H.
\]
If \(h,h'\in H\) satisfy
\(\Tr^2Y_h=\Tr^2Y_{h'}\), then
\(\tau_{wh}-\tau_{wh'}\) is one of the defining collision functions of
\(V_x\).  It vanishes on \(Z\), so every tangent vector to \(Z\) at \(x\),
including \(v\), annihilates its differential.  This is true even when
\(v\) is only a dual-number tangent and is not tangent to an algebraic or
analytic arc.  By (13.18),
\[
 z(wh)=z(w),
 \qquad
 \left.\frac{d}{d\epsilon}\right|_{\epsilon=0}
 \Tr\rho_\epsilon(wh)=\Tr(NY_h).
\tag{13.24}
\]

Fix \(h_*\) in one squared-trace fiber.  Every \(wh\) is nontrivial---in
fact it is primitive by the argument below---so its trace at the Fuchsian
point is nonzero.  There is therefore a sign
\(\varepsilon_h\in\{\pm1\}\) such that
\[
 \Tr Y_h=\varepsilon_h\Tr Y_{h_*}.
\]
Differentiating
\(\Tr^2Y_h=\Tr^2Y_{h_*}\) along \(v\), and using (13.24), gives
\[
 \Tr(NY_h)=\varepsilon_h\Tr(NY_{h_*}).
\tag{13.25}
\]
Thus, after dividing the fiber into its two sign classes, both
\(\Tr Y_h\) and \(\Tr(NY_h)\) are fixed.

Diagonalize \(N\) over a finite extension of the field generated by its
entries:
\[
 N=\begin{pmatrix}\lambda&0\\0&-\lambda\end{pmatrix},
 \qquad \lambda\ne0.
\]
Writing \(Y_h=(y_{ij})\) in this fixed basis, the two quantities in
(13.25) determine \(y_{11}\) and \(y_{22}\), while determinant one gives
\[
 y_{12}y_{21}=y_{11}y_{22}-1=:q.
\tag{13.26}
\]
We spell out why Theorem~\ref{thm:finite-type-divisor} applies here despite
the use of a tangent vector.  Let \(F/\Q\) be the field generated by the
entries of \(\rho\) on a finite generating set of \(\Pi\), the finitely
many entries of \(N=z(w)\), and the entries of the fixed diagonalizing
matrix.  This is a finitely generated field of characteristic zero.  Choose
for \(F\) a normal finite-type model and a filtration \(R_C[n]\) as in
(12.1).  If \(h\) has word length at most \(n\), repeated multiplication
by the fixed generator matrices shows that every entry of \(Y_h\), after
the fixed change of basis, belongs to \(R_C[n]\) for one constant \(C\)
independent of \(h\) and \(n\).  The entries of the dual-number cocycle do
not grow here: only the single fixed matrix \(N=z(w)\) occurs in (13.24).
After increasing \(C\), the fixed diagonal entries and the right side
\(q\) of (13.26) belong to the same filtration.  If \(q\ne0\), the
finite-type divisor bound therefore leaves only
\(\exp(o(n))\) possibilities for the ordered pair
\((y_{12},y_{21})\), uniformly in the fixed trace fiber.

If \(q=0\), the solutions in (13.26) are the two triangular affine lines
\(y_{12}=0\) and \(y_{21}=0\).  Each line contains at most one matrix
\(Y_h\) arising from \(H\).  Indeed, two distinct such matrices on the
same line would satisfy \(\det(Y_h-Y_{h'})=0\).  For determinant-one
matrices,
\[
 \det(U-V)=2-\Tr(UV^{-1}),
\]
and hence
\[
 2=\Tr(Y_hY_{h'}^{-1})
   =\Tr\rho(hh'^{-1}).
\]
Faithfulness and torsion-free convex cocompactness make this impossible
when \(h\ne h'\).  Combining the two signs and the two triangular lines,
we conclude that every fiber of
\[
 h\longmapsto\Tr^2\rho(wh)
\tag{13.27}
\]
in an \(H\)-word ball of radius \(n\) has size \(\exp(o(n))\).  Word
length and orbit displacement are quasi-isometric on the convex-cocompact
group \(H\), so the same statement is \(\exp(o(R))\) in a displacement
ball of radius \(R\).

Finally, orbit counting gives
\(\exp((\delta_H+o(1))R)\) elements \(h\in H\) with
\(d(o,\rho(h)o)\le R\).  Every \(wh\) is primitive: if
\(wh=u^m\) with \(|m|\ge2\), then
\(1=\phi(wh)=m\phi(u)\), which is impossible.  Their lengths are at most
\(R+O(1)\).  The hypothesis \(h_{\Tr}(x)\le1/2\) says that they assume at
most \(\exp((1/2+o(1))R)\) distinct absolute-trace values.  Pigeonholing
therefore produces a fiber of (13.27) of size
\[
 \exp\bigl((\delta_H-\tfrac12-o(1))R\bigr),
\]
contradicting the subexponential bound because \(\delta_H>1/2\).  The
kernel in (13.14) is consequently zero.
\end{proof}

\begin{theorem}[restriction to the subgroup has degree one]
\label{thm:restriction-birational}
Let \(\Pi\) be a closed surface group and let
\(x=[\rho_x]\) be a faithful cocompact Fuchsian character.  Let
\(\phi\colon\Pi\twoheadrightarrow\Z\), and let
\(H<\ker\phi\) be a finitely generated subgroup for which
\(\rho_x(H)\) is nonelementary and convex cocompact with critical exponent
\(\delta_H>1/2\).  Assume \(h_{\Tr}(x)\le1/2\).  Let \(Z\) be an irreducible
positive-dimensional \(\Q\)-subvariety of \(V_x\), on which all trace
collisions at \(x\) persist, and let
\[
 r_H\colon Z\longrightarrow\mathfrak X(H)
\]
be restriction of characters.  If \(r_H\) is generically finite onto its
image \(Y=\overline{r_H(Z)}\), then it has degree one:
\[
 \Q(Z)=\Q(Y).
\]
\end{theorem}

\begin{proof}
Put \(K=\Q(Z)\) and \(L=\Q(Y)\).  Suppose \([K:L]>1\).  Characteristic
zero makes this extension separable, so there is a nonidentity
\(L\)-embedding \(\sigma\colon K\to\overline K\).

The generic absolutely irreducible \(\PSL_2\)-character of \(Z\) can be
represented after a finite extension of \(K\): equivalently, split its
degree-two central simple algebra and choose matrices for a finite
generating set.  Extending \(\sigma\) to this finite coefficient field
gives two generic representations
\[
 \rho_0,\rho_1\colon\Pi\longrightarrow\PSL_2(\Omega),
\]
where \(\Omega\) is an algebraically closed extension of a finitely
generated characteristic-zero field.  They correspond to the identity and
\(\sigma\) embeddings of \(K\).  Since \(\sigma\) fixes \(L\), their
restrictions to \(H\) have the same character.  Their generic restricted
character is Zariski dense.  Indeed, it is faithful.  For each fixed
nonidentity \(h\in H\), the regular function \(\tau_h-4\) on \(Z\) is
nonzero at \(x\), because \(\rho_x\) is faithful and torsion-free Fuchsian.
It is therefore nonzero in the function field \(K\), so \(h\) is not the
identity under the generic projective representation.  This holds
simultaneously for all \(h\ne1\): the function-field point evaluates every
nonzero element of the coordinate ring as a nonzero element, and no
intersection of countably many Zariski-open sets is being asserted.  Now
\(H\) is a nonabelian free group, whereas every proper algebraic subgroup
of \(\PSL_2(\C)\) is virtually solvable.  A faithful image of \(H\) cannot
lie in such a subgroup.  Hence the generic restricted image is Zariski
dense.  In particular the restriction is absolutely
irreducible, so we may conjugate \(\rho_1\) and arrange the projective
equality
\[
 \rho_1|_H=\rho_0|_H
\tag{13.28}
\]
in \(\PSL_2(\Omega)\).  The group \(H\) is a finitely generated
infinite-index subgroup of a closed surface group and is therefore free by
the same compact-core argument used above.  Choose one common lift
\[
 r\colon H\longrightarrow\SL_2(\Omega)
\]
of the projective representation in (13.28).

If \(\rho_0\) and \(\rho_1\) were projectively equal on \(\Pi\), then
\(\sigma\) would fix every squared-trace coordinate on \(Z\).  Those
coordinates generate \(K\): in the adjoint representation,
\(\Tr(\Ad \rho(g))=\Tr^2(\widetilde\rho(g))-1\), and traces of words in a
faithful matrix representation generate the character function field.
Thus \(\sigma\) would be the identity.  Hence the representations are not
projectively equal.  Algebraic Goursat's lemma,
followed by the normal-subgroup argument in the proof of
Theorem~\ref{thm:relative-fiber-rigidity}, shows that
\[
 (\rho_1,\rho_0)(\ker\phi)
\]
is Zariski dense in \(\PSL_2\times\PSL_2\).  We may consequently choose
\(w\in\phi^{-1}(1)\) and arbitrary lifts \(A_i\in\SL_2(\Omega)\) of
\(\rho_i(w)\) such that
\[
 B=A_1A_0^{-1}
 \quad\text{has}\quad \Tr^2B\ne4.
\tag{13.29}
\]

There are \(e^{(\delta_H+o(1))R}\) elements
\(h\in H\) with \(d(o,\rho_x(h)o)\le R\), and every \(wh\) is primitive
because \(\phi(wh)=1\).  Their \(x\)-lengths are at most \(R+O(1)\), so
for every \(\epsilon>0\), the hypothesis \(h_{\Tr}(x)\le1/2\) leaves at
most \(e^{(1/2+\epsilon)R}\) squared-trace values for all sufficiently
large \(R\).  Therefore one value is assumed by at least
\[
 e^{(\delta_H-1/2-\epsilon-o(1))R}
\tag{13.30}
\]
of these elements.

All collisions at \(x\) persist on \(Z\).  Thus the squared-trace functions
of the elements in (13.30) are equal as elements of \(K\).  Evaluating
these identities first at the identity generic point and then at its
\(\sigma\)-conjugate shows that the elements \(h\) form one simultaneous
squared-trace fiber for
\[
 h\longmapsto
 \bigl(\Tr^2(A_0r(h)),\Tr^2(A_1r(h))\bigr).
\]
The coefficient field generated by the finitely many matrices involved is
finitely generated over \(\Q\).  For each \(g\in H\), the element
\(\Tr^2r(g)\) is the regular squared-trace coordinate \(\tau_g\) at the
generic point of \(Y\).  If \(\Tr^2r(g)=4\), then \(\tau_g-4\) is identically
zero on \(Y\), and hence it also vanishes at \(x|_H\).  Thus the common
\(H\)-representation and the Fuchsian character \(x|_H\) satisfy the exact
implication in Theorem~\ref{thm:two-trace-evaluations}.  That theorem and
(13.29) bound
this simultaneous fiber by \(e^{o(R)}\).  Here convex cocompactness gives
\(|h|_H=O(R)\), and the fixed matrices and diagonalizer are absorbed into
one finite extension and finite-type model.  Choosing
\(0<\epsilon<\delta_H-1/2\) contradicts (13.30).  Hence \([K:L]=1\).
\end{proof}

\begin{theorem}[the adjoint trace field is already generated on the
complementary subgroup]
\label{thm:algebraic-restriction-field}
Let \(\Pi\) be a closed surface group and let \(x=[\rho_x]\) be a faithful
cocompact Fuchsian character which is an algebraic point of the liftable
\(\PSL_2(\C)\)-character component.  Let
\(\phi\colon\Pi\twoheadrightarrow\Z\), and let \(H<\ker\phi\) be finitely
generated, nonelementary, and convex cocompact under \(\rho_x\), with
critical exponent \(\delta_H\).  Suppose that
\(h_{\Tr}(x)<\delta_H\).  Define the adjoint trace field and the
corresponding field for the restriction to \(H\) by
\[
 K_x=\Q\bigl(\tau_g(x):g\in\Pi\bigr),
 \qquad
 L_H=\Q\bigl(\tau_h(x):h\in H\bigr).
\]
Then
\[
 K_x=L_H.
\]
Equivalently, no nonidentity embedding of the full adjoint trace field can
fix all squared traces on \(H\).
\end{theorem}

\begin{proof}
Both fields are number fields and \(L_H\subset K_x\).  Suppose that the
inclusion is proper.  Separability in characteristic zero gives a
nonidentity \(L_H\)-embedding
\(\sigma\colon K_x\hookrightarrow\overline{\Q}\).  Applying \(\sigma\) to
the character coordinates of \(x\) gives an algebraic character
\(x^\sigma\).  It differs from \(x\), since the squared-trace coordinates
generate \(K_x\), but its restriction to \(H\) has exactly the same
character as \(x|_H\).  The latter is absolutely irreducible.  After
realizing the two characters over \(\overline{\Q}\) and conjugating the
second realization, we may therefore write
\[
 \rho^\sigma|_H=\rho_x|_H
 \quad\hbox{in }\PSL_2(\overline{\Q}).
\]
The cover corresponding to the finitely generated infinite-index subgroup
\(H\) is noncompact and has a compact core with nonempty boundary.  That
core retracts onto a finite graph, so \(H\) is free.  Choose a common lift
\(r\colon H\to\SL_2(\overline{\Q})\).

The paired representation \((\rho^\sigma,\rho_x)\) is Zariski dense in
\(\PSL_2\times\PSL_2\).  Indeed, both factors are Zariski dense.  If its
closure were the graph of an automorphism, equality on the Zariski-dense
subgroup \(r(H)\) would make that automorphism the identity, and the two
global characters would coincide, contrary to \(x^\sigma\ne x\).  The
normal-subgroup argument from
Theorem~\ref{thm:relative-fiber-rigidity} then shows that the paired image of
\(\ker\phi\) is also Zariski dense in the product.  Hence one can choose
\(w\in\phi^{-1}(1)\) and lifts \(A_0,A_1\in\SL_2(\overline{\Q})\) of
\(\rho_x(w),\rho^\sigma(w)\), respectively, such that
\[
 \Tr^2(A_1A_0^{-1})\ne4.
\]

There are \(\exp((\delta_H+o(1))R)\) elements \(h\in H\) whose displacement
in the Fuchsian representation is at most \(R\), and every \(wh\) is primitive because
\(\phi(wh)=1\).  Choose
\(\epsilon>0\) with
\(h_{\Tr}(x)+\epsilon<\delta_H\).  By the definition of trace entropy,
some squared trace is shared by
\[
 \exp\bigl((\delta_H-h_{\Tr}(x)-\epsilon-o(1))R\bigr)
\]
of these elements.  Equality of their squared traces at \(x\) remains an
equality after applying \(\sigma\).  They consequently lie in one
simultaneous fiber of
\[
 h\longmapsto
 \bigl(\Tr^2(A_0r(h)),\Tr^2(A_1r(h))\bigr).
\]
Theorem~\ref{thm:two-trace-evaluations} bounds this simultaneous fiber by
\(\exp(o(R))\).  (Its specialization hypothesis is immediate here because
\(r\) is the faithful convex-cocompact Fuchsian restriction itself.)  This
contradicts
\(\delta_H-h_{\Tr}(x)-\epsilon>0\), and proves \(K_x=L_H\).
\end{proof}

\begin{corollary}[the invariant quaternion algebra is generated on the
complementary subgroup]
\label{cor:algebra-generated-on-H}
Under the hypotheses of
Theorem~\ref{thm:algebraic-restriction-field}, let
\(\mathcal A_x\) be the invariant quaternion algebra of
\(\rho_x(\Pi)\) over \(K_x\): after choosing lifts, this is the
\(K_x\)-subalgebra generated by \(\rho_x(\Pi^{(2)})\).  Then the
\(K_x\)-subalgebra generated by \(\rho_x(H^{(2)})\) is all of
\(\mathcal A_x\).
\end{corollary}

\begin{proof}
The subgroup \(H^{(2)}\) has finite index in the finitely generated group
\(H\), because the quotient is the finite elementary abelian group
\(H_1(H,\mathbb F_2)\).  It is therefore nonelementary, and its
two-dimensional representation is absolutely irreducible.  Its image lies
in \(\mathcal A_x\), since \(H^{(2)}<\Pi^{(2)}\).  The
\(K_x\)-subalgebra it generates becomes all of \(M_2\) after extending
scalars to an algebraic closure, by Burnside's theorem.  It consequently
has dimension four over \(K_x\).  The ambient quaternion algebra
\(\mathcal A_x\) also has dimension four over \(K_x\), so the inclusion
is an equality.
\end{proof}

\begin{corollary}[birational determination by the complementary subgroup]
\label{cor:hao-birational}
Under the hypotheses of
Corollary~\ref{cor:hao-restriction-fibers}, restriction to \(H\) is
birational onto its image on every irreducible positive-dimensional
\(\Q\)-subvariety of \(V_x\) through \(x\).  Its differential at \(x\)
is injective on the tangent space of every reduced such subvariety.
\end{corollary}

\begin{proof}
Corollary~\ref{cor:hao-restriction-fibers} says that the fiber dimension at
\(x\) is zero.  Upper semicontinuity of fiber dimension therefore makes the
restriction morphism generically finite on \(Z\).  Equivalently, the induced
extension of function fields is finite.  Theorem
\ref{thm:restriction-birational} makes its degree one.  The final assertion
is Theorem~\ref{thm:restriction-unramified}.
\end{proof}

For completeness, the following direct coordinate calculation verifies the
same conclusion on the specific twist subspaces used in Hao's proof.

\begin{corollary}[explicit calculation in Hao's Fenchel--Nielsen coordinates]
\label{cor:hao-transversality}
Let \(x\) be a closed hyperbolic surface of genus \(g\ge6\) with
\(h_{\Tr}(x)\le1/2\), and let \(V_x\) be the variety on which all trace
collisions at \(x\) persist.
For the one-holed-torus or pair-of-pants surface decomposition used in Hao's proof of
\cite[Theorem~A, Proposition~3.2, and (5.1)]{HaoLengthSets}, the complementary
subgroup \(H\) has
\(\delta_H>1/2\).  The real germ of \(V_x\) at \(x\) has
zero-dimensional intersection with the associated Fenchel--Nielsen twist
line in the one-holed-torus case, and with the full two-twist plane in the
pair-of-pants case.
\end{corollary}

\begin{proof}
We give the matrix calculation because it is where the algebraic locus
preserving the trace collisions meets the geometry of the splitting.  In the
one-holed-torus case, the deformation fixes \(H\) and has
\[
 \rho_s(a_1)=D(s)\rho_0(a_1)D(-s),\qquad
 D(s)=\diag(\kappa^s,\kappa^{-s}).
\]
Hao's normal form has
\(a_{22}=-\kappa a_{11}\) and
\(a_{11}a_{12}a_{21}\ne0\).  The diagonal part and the two off-diagonal
matrix units therefore show that
\(\Span\{\rho_s(a_1)\}\) has dimension three.  The homology class of
\(a_1h\) has \(a_1\)-coordinate one, so \(a_1h\) cannot be a proper power.
Theorem~\ref{thm:relative-deformation} would force primitive length entropy
at least \(\delta_H>1/2\) on any nonconstant piece of
\(V_x\) in this twist line, contrary to \(h_{\Tr}(x)\le1/2\).

In the pair-of-pants case, use twist parameters \(s,t\), put
\(u=e^s\), \(v=e^t\), and choose the standard matrices so that
\[
 \rho_{s,t}(a_1^{-1}a_2)=E(-s)NE(t)P^{-1},
 \qquad E(r)=\diag(e^r,e^{-r}),
\]
where every entry \(n_{ij}\) of \(N\) is nonzero.  Before the fixed right
factor \(P^{-1}\), the four entries are
\[
 n_{11}\frac vu,\qquad
 \frac{n_{12}}{uv},\qquad
 n_{21}uv,\qquad
 n_{22}\frac uv.
\tag{13.31}
\]
After projectivizing and multiplying all coordinates by \(uv\), (13.31)
becomes the diagonally rescaled Segre map
\[
 (u^2,v^2)\longmapsto
 [\,n_{11}v^2:n_{12}:n_{21}u^2v^2:n_{22}u^2\,].
\]
Writing these four homogeneous coordinates as
\([X_{11}:X_{12}:X_{21}:X_{22}]\), direct multiplication gives the exact
equation
\[
 n_{12}n_{21}X_{11}X_{22}
   =n_{11}n_{22}X_{12}X_{21}.
\]
All \(n_{ij}\) are nonzero, so this is a smooth quadric in
\(\mathbb P^3\).  Its two families of projective lines are obtained by
holding \(u^2\) constant or holding \(v^2\) constant.  Hence, on any
irreducible curve on which both \(s\) and \(t\) vary, the matrices in
(13.31) span at least three dimensions: otherwise the projective image
would lie on a line, forcing \(s\) or \(t\) to be constant.  Every
\(a_1^{-1}a_2h\) is primitive because its two indicated homology
coordinates are \((-1,1)\).  Theorem~\ref{thm:relative-deformation} excludes
such a curve from \(V_x\).

It remains to treat the two rulings.  If \(t=t_0\), choose \(Z\in H\) so
that
\[
 M=\rho_0(a_1^{-1}Za_1)
\]
has two nonzero off-diagonal entries and a nonzero diagonal part.  Such
\(Z\) exists by Zariski density of the non-elementary group \(\rho_0(H)\).
For \(a_Z=a_1^{-1}Za_1a_2^{-1}\), the varying matrix is
\[
 \rho_{s,t_0}(a_Z)
 =E(-s)ME(s)B_2(-t_0)\rho_0(a_2)^{-1}.
\]
Its constant diagonal part and its two independently scaled off-diagonal
parts span three dimensions.  Also \(a_Zh\) has \(a_2\)-homology
coordinate \(-1\), so it is primitive.  The other ruling is identical
after interchanging the two boundary generators and using the
\(B_2\)-eigenbasis.  Thus no nonconstant real-analytic curve in the
two-twist plane can lie in \(V_x\).  Real curve selection now gives the
claimed zero-dimensional real germs.
\end{proof}

\begin{theorem}[trace functions forced to coincide in every finite cover]
\label{thm:cover-persistent-collapse}
Let \(x\) be the point of Teichm\"uller space represented by a closed
hyperbolic surface \(S\), let
\(\Gamma=\pi_1(S)\), and let \(\rho\) be its holonomy representation.
Count distinct
absolute traces by
\[
 D_x(T)=
 \#\{|\Tr\widetilde\rho(\gamma)|:
          \gamma\in\Gamma,\ |\Tr\widetilde\rho(\gamma)|\le T\},
\]
where the choice of lift to \(\SL_2(\R)\) is immaterial.  Assume
\(D_x(T)\le C(1+T)\).

Let \(S_G\to S\) be any connected finite cover of genus \(G\ge6\), and
write \(x_G\) for the restricted Fuchsian character.  Let
\(\mathfrak X_G\) be the \(\Q\)-defined \(\PSL_2\)-character component
containing its Teichm\"uller space.  Define
\[
 I(x_G)=\{f\in\Q[\mathfrak X_G]:f(x_G)=0\},
 \qquad
 Z_G=V(I(x_G))=\overline{\{x_G\}}^{\,\Q}.
\]
Then the surface decomposition in Hao's proof of
\cite[Theorem~A, Proposition~3.2, and (5.1)]{HaoLengthSets} gives a
convex-cocompact subgroup
\(H=\pi_1(\Sigma')\), an element \(A\in\pi_1(S_G)\), and a number
\(\delta_H\ge\delta(G)>1/2\) with the following property.  There are
constants \(c>0\) and \(R_0\), depending on this cover and metric, such
that, for every \(R\ge R_0\), there are \(h_1,\ldots,h_M\in H\) for which
\[
 M\ge c\exp\bigl((\delta_H-\tfrac12)R\bigr),
\tag{13.32}
\]
and:
\begin{enumerate}[label=\textup{(\roman*)}]
\item \(d(o,\rho(h_i)o)\le R\), and every \(Ah_i\) is primitive;
\item the squared-trace functions
\[
 \tau_{Ah_i}\colon\mathfrak X_G\longrightarrow\mathbb A^1,
 \qquad \tau_w([\sigma])=\Tr^2\sigma(w),
\]
are pairwise distinct regular functions;
\item all restrictions \(\tau_{Ah_i}|_{Z_G}\) are the same regular
function on \(Z_G\).
\end{enumerate}
\end{theorem}

\begin{proof}
Convex-cocompact orbit counting gives
\[
 \#\{h\in H:d(o,\rho(h)o)\le R\}
   \ge c_0e^{\delta_HR}
\]
for all large \(R\).  Hao's primitivity lemma for this decomposition says
that every \(Ah\), \(h\in H\), is primitive
\cite[Lemma~4.7]{HaoLengthSets}.  If
\(D_A=d(o,\rho(A)o)\), then
\[
 \ell_x(Ah)\le D_A+d(o,\rho(h)o)\le D_A+R.
\]
Since \(|\Tr\widetilde\rho(w)|=2\cosh(\ell_x(w)/2)\), all these elements
have absolute trace at most \(C_Ae^{R/2}\).  The assumed linear bound on
the base trace set applies because the cover group is a subgroup of the
base surface group.  It therefore leaves at most \(C_1e^{R/2}\) possible
absolute traces.  Pigeonholing produces a subset \(E_R\subset H\) of
cardinality
\[
 |E_R|\ge(c_0/C_1)e^{(\delta_H-1/2)R}
\]
on which all squared traces agree at \(x_G\).

For \(u,v\in E_R\),
\(\tau_{Au}(x_G)-\tau_{Av}(x_G)=0\); hence
\(\tau_{Au}-\tau_{Av}\in I(x_G)\).  By the definition of \(Z_G\), their
restrictions to \(Z_G\) are equal.  It remains to ensure that the functions
on the full component are distinct.  Partition \(E_R\) by equality of the
regular functions \(\tau_{Ah}\) on \(\mathfrak X_G\).  Teichm\"uller space
is Zariski dense in this component, so these are precisely the universal
trace-equivalence classes considered in Hao's proof.  His matrix
calculation bounds every such class by three in the one-holed-torus
decomposition and by one in the pair-of-pants decomposition
\cite[Lemmas~4.2--4.3 and Corollary~4.6]{HaoLengthSets}.  Choosing one
representative from each class loses a factor of at most three and proves
(13.32) and (i)--(iii).
\end{proof}

The lower exponent in (13.32) can approach \(1/2\).  More explicitly, Hao
obtains
\[
 \delta(G)=\frac12\left(
 1+\sqrt{1-\frac{2.728\,H(G)}{(4G-6)\pi}}\right)>\frac12,
\]
where
\[
 H(G)=4\log G+5.362-4\log2+
 4\log\left(1+\frac{4(G-1)\pi}{2\log G+2.681}\right).
\]
For all sufficiently large \(G\),
\[
 \delta(G)>1-\frac{5.456\log G}{(2G-3)\pi},
\]
so \(\delta(G)\to1\).

\begin{lemma}[transcendence degree is unchanged by a finite cover]
\label{lem:finite-cover-trdeg}
Let \(\Gamma'<\Gamma\) have finite index, and for a character \(x\) put
\[
 K_\Gamma=\Q(\tau_\gamma(x):\gamma\in\Gamma),\qquad
 K_{\Gamma'}=\Q(\tau_h(x):h\in\Gamma').
\]
Then \(K_\Gamma/K_{\Gamma'}\) is algebraic.  Consequently the rational
Zariski closures of \(x\) in the full character components for \(\Gamma\)
and \(\Gamma'\) have the same dimension.
\end{lemma}

\begin{proof}
For every \(\gamma\in\Gamma\), left multiplication by
\(\langle\gamma\rangle\) on the finite set of left cosets
\(\Gamma/\Gamma'\) gives an integer
\(1\le n\le[\Gamma:\Gamma']\) with \(\gamma^n\in\Gamma'\).  The
Chebyshev trace identity has the form
\[
 \tau_{\gamma^n}=Q_n(\tau_\gamma)
\]
for a nonconstant polynomial \(Q_n\in\Z[T]\).  Thus every
\(\tau_\gamma(x)\) is algebraic over \(K_{\Gamma'}\), proving the first
assertion.  A character-coordinate field is algebraic over its
squared-trace field: choose finitely many trace generators and adjoin their
square roots.  Its transcendence degree, which is the dimension of the
rational Zariski closure of the point, is therefore unchanged.
\end{proof}

\begin{corollary}[at least \(e^{(1/2-\eta)R}\) distinct trace functions agree after restriction]
\label{cor:cover-collapse-half}
If the field generated by the squared-trace coordinates of \(x\) in
Theorem~\ref{thm:cover-persistent-collapse} has positive transcendence
degree over \(\Q\), then \(Z_G\) is positive-dimensional in every finite
cover.
For every \(\eta>0\), one can choose one cyclic cover and a constant
\(c_\eta>0\) such that, for every sufficiently large \(R\), at least
\[
 c_\eta e^{(1/2-\eta)R}
\]
pairwise distinct squared-trace functions on the full character component
restrict to one and the same function on \(Z_G\).
\end{corollary}

\begin{proof}
A cyclic cover of degree \(m\) has genus
\(G=1+m(g-1)\), which tends to infinity with \(m\).  Choose \(m\) so that
\(\delta(G)>1-\eta\), and apply (13.32).  Positivity of
\(\dim Z_G\) follows from Lemma~\ref{lem:finite-cover-trdeg}.
\end{proof}

The cover is chosen before \(R\) tends to infinity, and the constants
\(c,R_0,D_A\) may depend on that cover.  The pairwise distinct functions
live on the full character variety of the cover; its twist deformations need
not descend to the base surface.  Consequently the theorem does not by
itself rule out a transcendental character.  It shows that any
transcendental character with linear trace growth would produce, in every
sufficiently large cover, almost \(e^{R/2}\) distinct trace polynomials that
agree identically on one positive-dimensional algebraic locus.

\begin{proposition}[arbitrary nonconstant loci can have trace-function fibers of size \(e^{(1/2-o(1))R}\)]
\label{prop:no-uniform-locus-gap}
For every \(\eta>0\), there are a closed arithmetic hyperbolic surface
\((S,x)\), a positive-dimensional nonconstant algebraic locus
\(W\) through \(x\) in its character variety, and constants \(c>0\) and
\(R_0\) such that, for every \(R\ge R_0\), at least
\[
 c e^{(1/2-\eta)R}
\]
pairwise distinct squared-trace functions of primitive elements of length at
most \(R+O(1)\) restrict to one and the same constant function on \(W\).
Consequently there is no cover- and locus-independent estimate
\(e^{(1/2-\epsilon)R+o(R)}\), with fixed \(\epsilon>0\), for fibers of trace
functions on arbitrary nonconstant loci.
\end{proposition}

\begin{proof}
Start with a torsion-free compact arithmetic surface \(S_0\) of sufficiently
large genus.  Choose disjoint embedded one-holed tori \(T_1,T_2\), with
geodesic boundary, such that
\[
 C=\overline{S_0\setminus(T_1\cup T_2)}
\]
is connected and has positive genus; one may first pass to a finite cover of
an arithmetic surface to arrange this topology.  Choose an epimorphism
\(\psi\colon\pi_1(S_0)\to\Z\) which vanishes on both torus subgroups and whose
restriction to \(\pi_1(C)\) is still onto.  For example, use a primitive
cohomology class dual to a nonseparating curve contained in the interior of
\(C\).  Let \(S_m\to S_0\) be the cyclic cover defined by the reduction of
\(\psi\) modulo \(m\), equip it with the lifted hyperbolic metric, and denote
its Fuchsian character by \(x_m\).  The inverse image of \(C\) is connected because
\(\psi(\pi_1(C))=\Z\), whereas each \(T_i\) has \(m\) isometric lifts.  Choose
one lift \(\widetilde T_{i,m}\) of each torus and set
\[
 Y_m=S_m\setminus\bigl(\operatorname{int}\widetilde T_{1,m}
                   \cup\operatorname{int}\widetilde T_{2,m}\bigr),
 \qquad K_m=\pi_1(Y_m).
\]
The surface \(Y_m\) is connected.  It is the convex core of the complete
convex-cocompact surface
\[
 X_m=K_m\backslash\mathbb H^2.
\]
Moreover
\[
 \operatorname{area}(Y_m)
 =m\operatorname{area}(S_0)-\operatorname{area}(T_1)-\operatorname{area}(T_2),
\]
whereas \(\operatorname{length}(\partial Y_m)\) is independent of \(m\).
Extend the function equal to one on \(Y_m\) by a fixed compactly supported
cutoff across the two funnels of \(X_m\).  Its Dirichlet energy is bounded by
a constant times \(\operatorname{length}(\partial Y_m)\), while its squared
\(L^2\)-norm is at least \(\operatorname{area}(Y_m)\).  The Rayleigh quotient
therefore gives \(\lambda_0(X_m)\to0\).  For a convex-cocompact hyperbolic
surface, once \(\lambda_0(X_m)<1/4\),
\[
 \lambda_0(X_m)=\delta_m(1-\delta_m),
\]
where \(\delta_m\) is the critical exponent of \(K_m\)
\cite{BorthwickSpectral}.  It follows that
\[
 \delta_m\longrightarrow1.
\]

Put
\(H_m=\pi_1(S_m\setminus\operatorname{int}\widetilde T_{1,m})\).  Keep the
\(H_m\)-character fixed and vary the hyperbolic structure on
\(\widetilde T_{1,m}\) with
its boundary holonomy and gluing fixed.  A one-holed torus with fixed
boundary has positive-dimensional Teichm\"uller space.  The Zariski closure
of this deformation lies in a positive-dimensional component \(W_m\) of
the algebraic restriction fiber, and every squared-trace function
\(\tau_w\), \(w\in H_m\), is constant on \(W_m\).

Choose a standard generator \(A_m\) of the other chosen one-holed torus
\(\widetilde T_{2,m}\), and choose a symplectic homology basis of \(S_m\)
adapted to that torus.  Every loop in \(Y_m\) has zero \(A_m\)-coordinate in
this basis.  Hence, for \(h\in K_m\), the element \(A_mh\) lies in \(H_m\)
and has \(A_m\)-coordinate one.  If \(A_mh=u^q\) with \(|q|\ge2\), applying
this integral homology coordinate would give \(1=q\,[u]_{A_m}\), which is
impossible.  Thus every \(A_mh\) is primitive.
Convex-cocompact orbit counting gives \(c'_m e^{\delta_mR}\) such
elements of length at most \(R+O_m(1)\).  The surface \(S_m\) is arithmetic,
so its trace set has linear growth; these elements therefore assume at most
\(C_me^{R/2}\) squared-trace values.  One value is assumed at least
\[
 c''_m e^{(\delta_m-1/2)R}
\]
times.

It remains to discard possible identities between trace functions on the
full character variety.  Twist the chosen torus \(\widetilde T_{2,m}\) along
its separating boundary while fixing the representation of its complement,
and hence fixing \(K_m\).  In an eigenbasis for the boundary holonomy this
deformation has
\[
 \rho_s(A_m)=D(s)\rho_0(A_m)D(-s),
 \qquad D(s)=\operatorname{diag}(e^s,e^{-s}).
\]
Write \(\rho_0(A_m)=(a_{ij})\).  Both off-diagonal entries are nonzero: if
one vanished, \(A_m\) and the boundary holonomy would share a fixed point,
which in a discrete torsion-free Fuchsian group would force them to have the
same two fixed points, contrary to the topology of the one-holed torus.  Its
diagonal part is also nonzero, since a determinant-one matrix with zero
diagonal has projective order two.  Consequently the constant diagonal part
and the two independently scaled off-diagonal matrix units span a
three-dimensional subspace of \(M_2(\R)\).

If two squared-trace functions \(\tau_{A_mh}\) agree on the full character
component, their equality persists along this twist.  Applying
Theorem~\ref{thm:relative-deformation} with \(H=K_m\) shows that each such
class contains at most three elements.  Choosing one representative from
each class loses only this factor.  All remaining functions are globally
distinct, while their restrictions to \(W_m\) are the same constant because
the elements \(A_mh\) lie in \(H_m\).  Finally choose \(m\) with
\(\delta_m>1-\eta\), take \(c=c''_m/3\), and relabel
\((S_m,x_m,W_m)\) as \((S,x,W)\).  The constants are then fixed before
\(R\) tends to infinity, as required.
\end{proof}

The loci \(W_m\) in Proposition~\ref{prop:no-uniform-locus-gap} need not
preserve collisions involving the torus that moves.  The proposition
therefore does not weaken the results above concerning trace equalities that
hold identically on a fixed locus.  It shows
instead that any argument through these loci must exploit the full collision ideal,
or another condition that forces the varying part of the representation to
be seen by the chosen trace functions.

Although these results are not used in the recurrence proof of arithmeticity,
they give a concrete dichotomy for a locus \(V_x\) with low trace entropy:
either the Fricke coordinates are algebraic, or a real semialgebraic
deformation preserves all trace equalities that hold at \(x\).  Classifying the latter loci
means determining when exponentially many equalities between individual
closed-curve length functions can hold identically on one algebraic locus;
it is a structural question about the exceptional locus, not a remaining
step in the compact conjectures.

\section{Quantitative separation from the product formula}

The product formula converts upper bounds at all nondistinguished places into
a lower bound at the distinguished place.

\begin{theorem}[adelic separation]\label{thm:adelic}
Let \(k\) be a number field, \(v_0\) a distinguished archimedean place of local
degree \(n_0\in\{1,2\}\).  Write \(n_v=[k_v:\Q_v]\) for the usual local
degree (so a real place has degree one and a complex place degree two), and
normalize absolute values by
\(\prod_v|z|_v^{n_v}=1\).  Let \(S_T\subset k\) be a family such that for every
\(T\ge2\) and every distinct \(x,y\in S_T\), with \(z=x-y\),
\[
 |z|_v\le C_vT^{a_v}\qquad(v\ne v_0),
\tag{14.1}
\]
where \(C_v\ge1\), \(a_v\ge0\), and \(C_v=1,a_v=0\) outside a fixed finite
set.  Put
\[
 A=\sum_{v\ne v_0}n_va_v,
 \qquad C=\prod_{v\ne v_0}C_v^{n_v}.
\]
Then
\[
 |x-y|_{v_0}\ge C^{-1/n_0}T^{-A/n_0}.
\tag{14.2}
\]
Consequently a real unit interval, or a complex unit disk, contains
\(O(T^A)\) points of \(S_T\).  Explicitly, the total counts are
\[
 \#(S_T\cap[-T,T])=O(T^{1+A}),\qquad
 \#\{s\in S_T:|s|\le T\}=O(T^{2+A})
\tag{14.3}
\]
at a real and a complex distinguished place, respectively.
\end{theorem}

\begin{proof}
The product formula and (14.1) give
\[
 1\le |z|_{v_0}^{n_0}CT^A,
\]
which is (14.2).  One-dimensional interval packing and two-dimensional disk
packing give the local bounds; covering the radius-\(T\) region by
\(O(T^{n_0})\) unit regions gives (14.3).
\end{proof}

The theorem gives an upper bound controlled by the \emph{sum} of the growth
exponents at the nondistinguished places; it does not say that this bound is
sharp.  In a derived arithmetic group, all nondistinguished archimedean
coordinates are bounded and all finite-place denominators are bounded, so
\(A=0\).  For semi-arithmetic groups,
modular embeddings control archimedean escape and lead to the recent weak
clustering results of Belolipetsky--Cosac--D\'oria--Teixeira Paula
\cite{BCDPmultiplicities}.  For quasi-arithmetic or \(S\)-arithmetic behavior, a separate
estimate relating Fuchsian translation length to Bruhat--Tits displacement
is indispensable; the size of the real trace alone does not control
denominators.

\section{Why iterated square roots do not prove the compact conjectures}

Apply Theorem~\ref{thm:root-commensurator} to the finite-index lattice
\(H=\Gamma^{(2)}\).  Since \(H\) is normal in \(\Gamma\), every
\(\gamma\in\Gamma\) commensurates \(H\), and \(\gamma^2\in H\).  Among
primitive \(H\)-classes represented by such literal squares and having length
at most \(L\), there are at most, up to bounded finite-index class splitting,
as many as \(\Gamma\)-classes of length at most \(L/2\), namely only
\[
 e^{hL/2+o(L)}
\]
classes.  This is exactly the critical exponent allowed for a nonarithmetic
\(H\) when \(m=2\).

There is no further iteration inside this fixed subgroup: for \(r\ge2\), the
element \(\gamma^{2^r}\) is already a proper power in \(H\), so it is not
primitive there.  One may instead define the characteristic tower
\[
 H_0=\Gamma,\qquad H_{r+1}=H_r^{(2)}.
\]
Then \(\gamma^{2^r}\in H_r\), while \(\gamma\in\Comm_G(H_r)\), and the number
of possibly primitive \(H_r\)-classes obtained this way below length \(L\) is
at most \(e^{hL/2^r+o(L)}\).  Theorem~\ref{thm:root-commensurator} with
\(m=2^r\) permits exactly that exponent for a nonarithmetic lattice.  Thus the
iterated tower stays at, rather than crosses, the critical threshold.

Nor can an arbitrary element of \(\Gamma^{(2)}\) be replaced by a bounded
product of literal squares.  For a closed surface group the square word is
proper, and its verbal subgroup has infinite width by
Myasnikov--Nikolaev~\cite{MyasnikovNikolaev}.  At the algebro-geometric level,
the Kummer obstruction
\(\delta_2(h)\in H^1(k,\cJ_h[2])\) lives in a centralizer that varies with \(h\),
and has no canonical multiplicative law across products of noncommuting
elements.  Therefore one cannot prove (9.1) simply by taking square roots
after passing to the subgroup generated by squares.

\section{The main argument and the auxiliary results}

We summarize the logical role of the principal results.

\begin{enumerate}[label=\textbf{\arabic*.},leftmargin=*]
\item \emph{Incidence estimates and universal growth.}
Theorems~\ref{thm:finite-type-fiber} and~\ref{thm:surface-fiber} give the
sharp universal bounds of Corollary~\ref{cor:universal-trace-growth}.
Theorem~\ref{thm:colored-orbit-growth} strengthens this to every component of
every finite quotient, and Theorem~\ref{thm:colored-mixed-energy} gives the
corresponding mixed correlations under linear growth.  No algebraicity
assumption on the lattice is used.

\item \emph{Why ordinary additive closure is insufficient.}
Theorem~\ref{thm:same-trace-sum-fiber} proves that, after recording
\((\Tr(AB_i))_{i=1}^3\) with fixed \(B_i\) whose trace functionals recover
\(A\), every matrix sum except the forced inverse sum has representation
multiplicity \(\exp(o(n))\) on a word ball of radius \(n\).  Thus the energy
estimates cannot by themselves make the full trace fiber bounded at every finite place.  This is
a proved obstruction to that method, not an omitted estimate.

\item \emph{Simultaneous control of the recurrence dilates.}
The two trace coordinates \((\Tr g,\Tr(ag))\) form a graph with
\(e^{R-o(R)}\) edges and two vertex classes of size \(O(e^{R/2})\).
Theorem~\ref{thm:drc-trace-recurrence} extracts one set of
\(e^{(1/2-o(1))R}\) vertices on which all Chebyshev recurrence
coefficients are controlled simultaneously.  On a nonconstant curve,
Theorem~\ref{thm:chebyshev-crt-curve} combines the coprime recurrence zero
divisors by the Chinese remainder theorem and gives the exact bound
(13.10bh10).

\item \emph{The congruence criterion for a fixed element.}
Lemma~\ref{lem:congruence-wedge} does not require a Cartesian product in the
trace-pair graph.  It requires only a uniform bound for the number of trace
values in each residue class modulo
\[
 \mathfrak n_n=(q_{n-1}(\Tr a)q_n(\Tr a)).
\]
Theorem~\ref{thm:full-trace-field-congruence} turns that estimate into
\(\mathcal J_k(\Tr a)=\ell(a)/2\).  If it holds for every
\(a\in\widetilde{\Gamma^{(2)}}\), Takeuchi's theorem \cite{Takeuchi} proves arithmeticity.
Lemma~\ref{lem:quotient-mixing-trace-count} derives the estimate from
uniform expansion and lattice counting in all these mixed congruence
quotients.

\item \emph{The proof using varying group elements.}
Fix the recurrence index \(N\), instead of sending it to infinity along one
fixed element.  Corollary~\ref{cor:roth-chebyshev-pair} gives a
\((2N-4-\epsilon)\)-multiple of the global trace height in the reduced
divisor of two consecutive recurrence terms.  Lemma
\ref{lem:moving-bounded-recurrence-exponents} uses only expansion modulo
rational \(p^2\) to choose \(a_m=b^mh_m\) for which that divisor has bounded
prime exponents.  Lemma~\ref{lem:moving-two-trace-fibers} is uniform in
this moving multiplier, and Theorem~\ref{thm:moving-congruence-wedge}
uses bounded-power superapproximation to give the geometric upper
coefficient \((2N-1)/2\).  Letting \(N\to\infty\) proves
Theorem~\ref{thm:persistent-local-expansion-bound}:
\[
 \mathcal J_{k_y}(\Tr_yb)\le\frac{\ell_x(b)}2
\]
at every algebraic faithful cocompact Fuchsian character \(y\) in the
rational closure of \(x\).

\item \emph{Arithmeticity and transcendental characters.}
At an algebraic original character, the distinguished real place gives
the reverse inequality, and Takeuchi's theorem \cite{Takeuchi} gives arithmeticity.
If the original character were transcendental,
Proposition~\ref{prop:bounded-degree-specialization} would give a distinct
nearby algebraic Fuchsian character in the same component and in its
rational closure.  The preceding inequality makes every marked length at
the specialization at most the corresponding original marked length.
Lemma~\ref{lem:local-expansion-fuchsian-rigidity}, which applies Thurston's
asymmetric metric, then forces the two characters to coincide.  This
contradiction proves
Theorem~\ref{thm:unconditional-compact-trace-sparsity}, and hence both
compact conjectures, without arbitrary-ideal expansion or \(abc\).

\item \emph{Separate pointwise results for a fixed recurrence.}
Theorem~\ref{thm:single-term-radical-upper} gives
\[
 \limsup_n\frac1n
 \log N\operatorname{rad}_S(q_n(\Tr a))\le\frac{\ell(a)}2
\]
for every hyperbolic \(a\).  Thus a nondistinguished place with positive
local expansion forces the repeated prime factors in every individual
recurrence term to have positive linear logarithmic norm, not merely on
average over two consecutive terms.  Theorem
\ref{thm:high-power-concentration} also forces every hypothetical
nondistinguished expanding place in the fixed-element argument to put positive linear
logarithmic norm above each fixed prime-power cutoff, and
Theorem~\ref{thm:wieferich-mass} locates
the whole possible loss in exceptional initial Lucas--Wieferich
multiplicity.  The matching lower radical rate is isolated as
Conjecture~\ref{conj:chebyshev-radical};
Theorem~\ref{thm:chebyshev-radical-implies-compact} shows that it implies
both compact conjectures.  Number-field \(abc\) gives the lower rate,
whereas Proposition~\ref{prop:mersenne-radical-boundary} shows why it is
not currently an unconditional Lucas-sequence theorem.  Proposition
\ref{prop:crt-diagonal-obstruction} proves that flatness at the separate
prime-power factors does not formally glue to flatness modulo their product.
These statements describe the stronger problem for a fixed recurrence.  The
proof above, which varies the group element, does not resolve the Mersenne
radical problem and does not require it.

\item \emph{Valuative results.}
The proper-model descent theorem constructs a central simple algebra and an
order from invariant lattices.  The entropy estimate at a divisorial
valuation excludes every unbounded divisor for which the local translation
length exceeds the explicit cancellation rate.  The
recurrence argument proves the more general local expansion theorem without
having to establish the reverse inequality in the remaining case.
\end{enumerate}

The logical dependence of the principal ingredients is:
\[
\begin{gathered}
\substack{\text{finite-type divisor bounds}\\
          +\text{quadratic norm identity}}
   \longrightarrow
   \text{uniform moving two-trace fibers},\\[2pt]
\substack{\text{reduced-divisor Roth}\\
          +\text{expansion modulo rational }p^2}
   \longrightarrow
   \substack{\text{moving recurrence ideals}\\
             \text{with bounded exponents}},\\[2pt]
\substack{\text{algebraic Fuchsian specialization}\\
          +\text{persistent trace equalities}\\
          +\text{bounded-power expansion}\\
          +\text{congruence-class count}\\
          +\text{linear trace growth}}
   \longrightarrow
   \mathcal J_{k_y}(\Tr_yb)\le\ell_x(b)/2,\\[2pt]
\substack{\text{nearby algebraic specialization}\\
          +\text{the local height inequality}}
   \longrightarrow
   \ell_y(\gamma)\le\ell_x(\gamma)\quad(\gamma\ne1),\\[2pt]
\substack{\ell_y\le\ell_x\\
          +\text{Thurston's asymmetric metric}}
   \longrightarrow
   \text{the original character is algebraic},\\[2pt]
\substack{\text{local height inequality at }x\\
          +\text{distinguished real contribution}\\
          +\text{nonnegative local decomposition}\\
          +\text{Takeuchi and commensurability}}
   \longrightarrow
   \text{compact Schmutz and compact Sarnak},\\[4pt]
\text{independently:}\qquad
\substack{\text{one-term radical lower rate}\\
          +\text{proved one-term upper rate}}
   \longrightarrow\text{the same arithmeticity conclusion},\\[2pt]
\substack{\text{number-field \(abc\)}\\
          +\text{superapproximation modulo rational squares}}
   \longrightarrow\text{the stronger pointwise radical equality}.
\end{gathered}
\tag{16.1}
\]

\section{Comparison with known results}

The following table places the results beside the closest results in the
literature.
\begin{center}
\begin{tabular}{p{0.44\textwidth}p{0.47\textwidth}}
\toprule
Statement & Known result \\
\midrule
Positive trace gap \(\Rightarrow\) derived, Fuchsian/Kleinian
& Bogachev's manuscript; earlier nonuniform work established related
arithmeticity and growth implications, but not this full statement in this
form.\\
Bounded clustering \(\Rightarrow\) arithmetic, nonuniform Fuchsian
& Geninska--Leuzinger.\\
Linear trace growth \(\Rightarrow\) arithmetic, nonuniform Fuchsian
& Hao, in the stronger \(o(T\log\log\log T)\) form.\\
Bounded clustering or linear growth, cocompact Fuchsian
& Proved here in
Theorem~\ref{thm:unconditional-compact-trace-sparsity}.  The more general
Theorem~\ref{thm:persistent-local-expansion-bound} gives the required local
height inequality at every algebraic specialization preserving the trace
equalities at the original character.\\
Same compact conjectures under the one-term Chebyshev radical hypothesis
& Also proved here in
Theorem~\ref{thm:chebyshev-radical-implies-compact}.  The hypothesis has
an exact Mersenne specialization and remains open as a stronger pointwise
statement.\\
Same compact conjectures under number-field \(abc\)
& Theorem~\ref{thm:abc-implies-compact} is retained as an independent
conditional implication through the stronger pointwise radical equality.\\
Same compact conjectures under arbitrary-ideal superapproximation
& Theorem~\ref{thm:superapproximation-implies-compact} is an independent
conditional implication.  The general expansion theorem remains open but is not
needed for the unconditional result.\\
Every closed hyperbolic surface
& Corollary~\ref{cor:universal-trace-growth} gives the sharp universal
exponent \(N_H^*(L)\ge e^{(1/2-o(1))L}\), and arithmetic surfaces show that
\(1/2\) is optimal.  The preceding word-length comparison and removal of
proper powers turn \cite[Theorem~1.2]{BCLM} into exponential primitive-length
growth with an unspecified surface-dependent positive exponent.  We do not
know an earlier result with this universal exponent.\\
Generic compact surfaces of genus at least six
& Hao proves distinct primitive-length exponent \(>1/2\) outside a countable
union of positive-codimension algebraic loci~\cite{HaoLengthSets}.\\
Semi-arithmetic surfaces with modular embeddings
& Weak clustering and exponential mean multiplicity are known in arithmetic
dimension at most two; higher-dimensional results currently require a strong
contraction hypothesis~\cite{BCDPmultiplicities,Zuevsky}.\\
\bottomrule
\end{tabular}
\end{center}

Three natural analytic approaches do not apply under the present
hypotheses.  First, the nonarchimedean Anosov correlation theory does not
apply to a closed surface group.  Sambarino observes that a
\(\vartheta\)-Anosov representation over a nonarchimedean local field would
give continuous transverse boundary maps, forcing the Gromov boundary to be
a Cantor set and the group to be virtually free
\cite[Remark~5.2.6]{SambarinoDichotomy}.

Second, replacing a nondistinguished representation by its Bruhat--Tits translation
length loses essential information.  Cantrell--Reyes prove that the orbit
pseudometric of an action on an \(\R\)-tree without a global fixed point lies in their
Manhattan boundary exactly when the action has bounded backtracking
\cite[Proposition~6.7]{CantrellReyes}.  No such bounded-backtracking theorem
is known for the actions needed here.  Chen--He instead show that the
nonarchimedean harmonic map factors through a folding
\(T_q\to T_F\), with \(q\) Jenkins--Strebel; over a locally compact field,
an unbounded reductive surface-group representation has a non-small
Bruhat--Tits action
\cite[Theorems~1.2--1.3 and Proposition~2.15]{ChenHe}.  Thus the resulting tree
is generally a folded quotient, not simply the dual tree of a measured
foliation.  Moreover its integral translation length has only \(O(R)\)
possible values on an \(R\)-ball for the Fuchsian metric, so at least one
exact length level may retain the full exponential growth rate.  A useful
estimate must distinguish traces modulo high powers, rather than recording
only their valuations.

Third, ordinary congruence equidistribution does not provide enough
independent congruence conditions.
A single trace condition in \(\SL_2(\mathbb F_q)\) has density
\(\asymp q^{-1}\), whereas prescribing
\((\operatorname{tr}A,\operatorname{tr}(AB))\) has generic density
\(\asymp q^{-2}\).  An exact real trace fiber does not prescribe the second
coordinate, so summing over it loses that extra factor.  Existing uniform
congruence counting, such as Magee--Oh--Winter's theorem for Schottky
semigroups in \(\SL_2(\Z)\), carries a \(q^C\) loss and is not a
closed-surface periodic-orbit theorem at the growing-modulus scale required
here~\cite{MageeOhWinter}.  Their vector-valued transfer operator on the
regular representation of the congruence quotient is nevertheless the
relevant model for a possible vector-valued argument.  Such an argument
would require matrix-coefficient bounds against trace-fiber indicators;
the trace values themselves do not form a small right-translation module.

The classification results in Sections~5--7 use the positive trace-gap
theorem, while the divisor estimates and the universal exponent
in Section~12 do not assume arithmeticity or algebraicity of the surface
group.  Theorem~\ref{thm:unconditional-compact-trace-sparsity} proves the
compact bounded-clustering and linear-growth implications.  The conditional
and pointwise results concern stronger arithmetic statements and explain
which parts of an argument based on one fixed recurrence remain open.

\section{Conclusion}

The square-root formula is best understood as a special case of division on
the regular centralizer, whose branches are governed by Kummer theory and
Galois monodromy.  Type \(A_1\) is exceptional: projectivization turns this
multivalued correspondence into the rational map \([1+g]\), and the
Cayley--Hamilton identity expresses every trace \(\Tr(g^nh)\) in terms of two
initial trace coordinates.  In higher type \(A\), polynomial secants remain
rational even when roots do not; their failure to descend to a projective
unitary group distinguishes outer forms from inner forms.

For compact Fuchsian groups, the divisor bounds over finitely generated
fields yield the sharp universal lower exponents for distinct traces and
primitive lengths.  The same estimates hold separately in the components
indexed by a finite quotient, and linear trace growth gives mixed additive
energy \(X^{3-o(1)}\) between components.  Theorem~\ref{thm:same-trace-sum-fiber} also proves that,
apart from the forced inverse sum, applying the injective coordinate map
\(A\mapsto(\Tr(AB_i))_{i=1}^3\) to an exact trace fiber in a word ball of
radius \(n\) gives additive representation multiplicities \(\exp(o(n))\).
These facts explain why an
argument using only additive energy or small doubling cannot prove
arithmeticity.

The arithmeticity proof fixes the recurrence index \(N\) and varies the
group element.  Roth's theorem gives a lower bound for the reduced divisor of
\(q_{N-1}(\Tr a_m)q_N(\Tr a_m)\).  A random-walk choice
\(a_m=b^mh_m\), using expansion only modulo rational \(p^2\), keeps the
global recurrence height close to that of \(b^m\) while ensuring that every
prime occurs to uniformly bounded order in this divisor.  The explicit
quadratic norm identity (13.10mov16) gives a fiber estimate uniform in
\(a_m\):
\[
 g\longmapsto(\Tr_y(a_mg),\Tr_yg).
\]
Known superapproximation for fixed powers of square-free moduli then bounds the number of traces in
each residue class modulo the recurrence ideal.  Lemma
\ref{lem:congruence-wedge} converts this bound into the required upper
estimate.  Comparing the lower coefficient \(2N-4-\epsilon\) with the upper
coefficient \((2N-1)/2\), and then letting \(N\to\infty\), gives the factor
\(1/2\) in the height inequality below.

The result is the inequality
\[
 \mathcal J_{k_y}(\Tr_yb)\le\frac{\ell_x(b)}2
\tag{17.1}
\]
for every algebraic faithful Fuchsian character \(y\) in the rational
closure of a linear-growth character \(x\).  This is the general theorem
from which the compact conjectures follow.  When \(y=x\) is algebraic, the
contribution of the distinguished real place gives equality, and Takeuchi's theorem
\cite{Takeuchi} gives
arithmeticity.  If \(x\) were transcendental, a distinct nearby algebraic
Fuchsian specialization \(y\) in the same component would preserve all of
its trace collisions.  Inequality (17.1) would then give
\(\ell_y(\gamma)\le\ell_x(\gamma)\) for every nontrivial marked closed
curve.  Thurston's asymmetric metric \cite{ThurstonStretch} forces
\(y=x\), a contradiction.  Thus
\[
 \substack{\text{Roth reduced divisors}\\
           +\text{selection using expansion modulo }p^2\\
           +\text{bounded-power congruence counting}}
 \Longrightarrow (17.1)
 \Longrightarrow
 \text{compact Schmutz and compact Sarnak}
\]
is unconditional.

The arbitrary-ideal and one-term radical theorems give stronger statements,
but they are not prerequisites for this conclusion.  In
particular,
\[
 \limsup_n\frac1n
 \log N\operatorname{rad}_S(q_n(\Tr a))
 \le\frac{\ell(a)}2
\]
still gives a pointwise restriction on every fixed recurrence.  Its
matching universal lower bound follows from number-field \(abc\), while
the specialization \(s=3/\sqrt2\) identifies that stronger lower bound
with full exponential radical growth of \(2^n-1\).  By varying the group
element, the proof avoids this Mersenne and non-Wieferich problem; it does
not solve that stronger pointwise problem.

The Lucas--Wieferich localization theorem, the Chebyshev
Chinese-remainder theorem on algebraic curves, proper-model descent, and
the entropy theorem at a divisorial valuation are additional results.  They describe
the limitations of arguments based on one fixed recurrence or on a single
valuation.  Inequality (17.1), proved from the varying family \(a_m\), is
what completes the arithmeticity argument.

\bibliographystyle{siam}
\bibliography{biblio}

\end{document}
